\documentclass[12pt,a4paper]{report}

% Font và ngôn ngữ
\usepackage[utf8]{vietnam}
\usepackage{fontenc}

% Hình ảnh và bảng biểu
\usepackage{graphicx}
\usepackage{float}
\usepackage{array}
\usepackage{longtable}
\usepackage{booktabs}
\usepackage{multirow}
\usepackage{caption}
\usepackage{subcaption}

% Toán học
\usepackage{amsmath}
\usepackage{amssymb}
\usepackage{amsfonts}
\usepackage{amsthm}

% Màu sắc và định dạng
\usepackage{color}
\usepackage{xcolor}
\usepackage{hyperref}
\hypersetup{
    colorlinks=true,
    linkcolor=blue,
    urlcolor=blue,
    citecolor=blue
}

% Định dạng trang
\usepackage{geometry}
\geometry{
    a4paper,
    top=2.5cm,
    bottom=2.5cm,
    left=3cm,
    right=2.5cm
}

% Header và Footer
\usepackage{fancyhdr}
\pagestyle{fancy}
\fancyhf{}
\fancyhead[L]{\leftmark}
\fancyhead[R]{\thepage}
\renewcommand{\headrulewidth}{0.4pt}

% Các package hữu ích khác
\usepackage{enumitem}
\usepackage{tocloft}
\usepackage{indentfirst}
\usepackage{url}
\usepackage{listings}
\usepackage{verbatim}

% Thiết lập cho listings (code)
\lstset{
    basicstyle=\ttfamily\small,
    breaklines=true,
    frame=single,
    numbers=left,
    numberstyle=\tiny,
    tabsize=2
}

%%%%%%%%%%%%%%%%%%%%%%%%%%%%%%%%%%%%%%%%%%%%%%%%%%%%%%%%%%%%%%
% BẮT ĐẦU TÀI LIỆU
%%%%%%%%%%%%%%%%%%%%%%%%%%%%%%%%%%%%%%%%%%%%%%%%%%%%%%%%%%%%%%

\begin{document}

% Trang bìa
\title{\textbf{Nền tảng Quản lý Dự án Phần mềm Hướng AI}}
\author{Phạm Minh Hiển}
\date{07/2023}
\maketitle

% Mục lục
\tableofcontents
\newpage

% Danh sách hình
\listoffigures
\newpage

% Danh sách bảng
\listoftables
\newpage

%%%%%%%%%%%%%%%%%%%%%%%%%%%%%%%%%%%%%%%%%%%%%%%%%%%%%%%%%%%%%%
% CHƯƠNG 1: GIỚI THIỆU
%%%%%%%%%%%%%%%%%%%%%%%%%%%%%%%%%%%%%%%%%%%%%%%%%%%%%%%%%%%%%%
\chapter{GIỚI THIỆU}
\label{chap:introduction}

\section{Bối cảnh}
Trong những năm gần đây, ngành công nghiệp phần mềm đã trở thành một trong những lĩnh vực phát triển nhanh nhất của nền kinh tế số. Các doanh nghiệp ngày càng phụ thuộc vào các hệ thống phần mềm nhằm hỗ trợ hoạt động kinh doanh, quản lý dữ liệu và cung cấp dịch vụ cho khách hàng. Sự phát triển mạnh mẽ của điện toán đám mây (Cloud Computing), Internet vạn vật (IoT), dữ liệu lớn (Big Data) và trí tuệ nhân tạo (Artificial Intelligence -- AI) đã làm thay đổi đáng kể phương thức phát triển cũng như quản lý các dự án phần mềm.

Song song với sự gia tăng về quy mô và độ phức tạp của hệ thống phần mềm, yêu cầu đối với hoạt động quản lý dự án cũng ngày càng cao. Một dự án phần mềm hiện đại thường bao gồm nhiều nhóm phát triển, nhiều bên liên quan, hàng trăm User Story và hàng nghìn Task được thực hiện đồng thời. Ngoài ra, yêu cầu của khách hàng có xu hướng thay đổi liên tục trong suốt vòng đời dự án, làm tăng nguy cơ chậm tiến độ, vượt ngân sách và giảm chất lượng sản phẩm.

Để giải quyết các thách thức trên, phương pháp phát triển phần mềm linh hoạt (Agile Software Development) đã trở thành xu hướng chủ đạo trong ngành công nghiệp phần mềm. Trong đó, Scrum là framework được áp dụng rộng rãi nhất nhờ khả năng chia nhỏ công việc thành các Sprint ngắn, cho phép phản hồi nhanh với sự thay đổi của yêu cầu khách hàng.

Tuy nhiên, việc áp dụng Agile cũng tạo ra lượng lớn dữ liệu quản lý như User Story, Sprint Backlog, Velocity, Burndown Chart và lịch sử thực hiện Task. Khối lượng dữ liệu này vượt quá khả năng phân tích thủ công của Product Owner, Scrum Master và Project Manager.

Các công cụ quản lý dự án phổ biến hiện nay như Jira, Trello, Azure DevOps và GitHub Projects chủ yếu tập trung vào việc lưu trữ và theo dõi công việc. Mặc dù hỗ trợ tốt quá trình cộng tác, các hệ thống này vẫn phụ thuộc đáng kể vào kinh nghiệm của người quản lý trong các hoạt động như:

\begin{itemize}
\item Phân tích và làm rõ yêu cầu;
\item Ước lượng Story Point;
\item Dự báo tiến độ hoàn thành dự án;
\item Phát hiện rủi ro tiềm ẩn;
\item Phân bổ nguồn lực phát triển.
\end{itemize}

Sự phát triển của các mô hình ngôn ngữ lớn (Large Language Models -- LLM) như GPT đã mở ra cơ hội mới trong việc ứng dụng trí tuệ nhân tạo vào quản lý dự án phần mềm. AI không chỉ có khả năng xử lý ngôn ngữ tự nhiên mà còn hỗ trợ phân tích yêu cầu, đề xuất Task, ước lượng khối lượng công việc, phát hiện rủi ro và đưa ra các khuyến nghị hỗ trợ ra quyết định.

Theo nghiên cứu tổng quan của Bachari và cộng sự \cite{bachari2025}, các kỹ thuật trí tuệ nhân tạo đã được ứng dụng hiệu quả trong các lĩnh vực quản lý tiến độ, quản lý chi phí, quản lý nguồn lực và quản lý rủi ro dự án. Tuy nhiên, phần lớn các nghiên cứu hiện nay vẫn tập trung vào từng bài toán riêng lẻ thay vì xây dựng một nền tảng tích hợp AI hỗ trợ toàn bộ quy trình quản lý dự án Agile.

Từ những phân tích trên, đề tài \textbf{``Nền tảng Quản lý Dự án Phần mềm Hướng AI''} được đề xuất nhằm xây dựng một hệ thống quản lý dự án thế hệ mới, trong đó AI đóng vai trò như một trợ lý thông minh hỗ trợ Product Owner, Scrum Master, Project Manager và Developer trong toàn bộ vòng đời phát triển phần mềm.

\section{Lý do chọn đề tài}
Có nhiều nguyên nhân thúc đẩy việc nghiên cứu và phát triển một nền tảng quản lý dự án phần mềm tích hợp trí tuệ nhân tạo.

Thứ nhất, các dự án phần mềm hiện đại ngày càng phức tạp cả về quy mô lẫn số lượng thành viên tham gia. Việc quản lý thủ công các User Story, Task, Sprint và nguồn lực dễ dẫn đến sai sót trong quá trình lập kế hoạch và theo dõi tiến độ.

Thứ hai, mặc dù các công cụ quản lý dự án hiện nay cung cấp nhiều tính năng hỗ trợ cộng tác, phần lớn các quyết định quan trọng vẫn dựa trên kinh nghiệm chủ quan của người quản lý. Điều này làm giảm khả năng mở rộng và gây khó khăn cho các nhóm phát triển mới hoặc thiếu kinh nghiệm.

Thứ ba, sự phát triển mạnh mẽ của trí tuệ nhân tạo tạo điều kiện để tự động hóa nhiều hoạt động vốn tiêu tốn thời gian như phân tích yêu cầu, phân rã Task, ước lượng Story Point hay phát hiện rủi ro.

Thứ tư, xu hướng Human-Centric AI đang được xem là hướng tiếp cận phù hợp trong các hệ thống hỗ trợ ra quyết định. Thay vì thay thế hoàn toàn con người, AI được sử dụng như một trợ lý nhằm tăng cường năng lực của người quản lý và nhóm phát triển.

Những lý do trên cho thấy việc xây dựng một nền tảng quản lý dự án tích hợp AI không chỉ mang ý nghĩa nghiên cứu mà còn có giá trị thực tiễn cao trong bối cảnh chuyển đổi số hiện nay.


\section{Mục tiêu nghiên cứu}
Mục tiêu tổng quát của đề tài là nghiên cứu và phát triển một nền tảng quản lý dự án phần mềm hướng AI hỗ trợ quy trình Agile và Scrum.

Các mục tiêu cụ thể bao gồm:

\begin{enumerate}
\item Xây dựng hệ thống quản lý Project, Sprint, User Story và Task.
\item Tích hợp AI hỗ trợ phân tích và làm rõ yêu cầu.
\item Tự động phân rã User Story thành các Task kỹ thuật.
\item Hỗ trợ ước lượng Story Point dựa trên dữ liệu lịch sử.
\item Dự báo tiến độ hoàn thành dự án.
\item Phát hiện và cảnh báo rủi ro trong quá trình thực hiện Sprint.
\item Đề xuất phương án phân công nguồn lực tối ưu.
\item Xây dựng Dashboard trực quan hỗ trợ theo dõi tiến độ thời gian thực.
\item Áp dụng mô hình Human-in-the-Loop, Human-on-the-Loop và Human-out-of-the-Loop.
\end{enumerate}


\section{Đối tượng nghiên cứu}
Đối tượng nghiên cứu của đề tài bao gồm:

\begin{itemize}
\item Các phương pháp quản lý dự án Agile và Scrum;
\item Các mô hình AI hỗ trợ Software Engineering;
\item Large Language Models (LLM);
\item Multi-Agent AI Systems;
\item Human-Centered Artificial Intelligence;
\item Các kỹ thuật dự báo và quản lý rủi ro trong dự án phần mềm.
\end{itemize}

\section{Phạm vi nghiên cứu}
Đề tài tập trung vào việc xây dựng phiên bản MVP (Minimum Viable Product) của hệ thống quản lý dự án phần mềm hướng AI.

Phạm vi chức năng bao gồm:

\begin{itemize}
\item Quản lý Project;
\item Quản lý Sprint;
\item Quản lý User Story;
\item Quản lý Task;
\item Phân tích yêu cầu bằng AI;
\item Phân rã User Story thành Task;
\item Ước lượng Story Point;
\item Dự báo tiến độ;
\item Phát hiện rủi ro;
\item Đề xuất phân công nguồn lực;
\item Dashboard và báo cáo.
\end{itemize}

Các chức năng nằm ngoài phạm vi đề tài bao gồm:

\begin{itemize}
\item Huấn luyện mô hình AI từ đầu;
\item Tích hợp đầy đủ với Jira hoặc Azure DevOps;
\item Quản lý tài chính doanh nghiệp;
\item Quản lý nhân sự cấp tổ chức.
\end{itemize}

\section{Phương pháp nghiên cứu}
Đề tài sử dụng kết hợp nhiều phương pháp nghiên cứu khác nhau:

\begin{enumerate}
\item Nghiên cứu tài liệu về Agile, Scrum, PMBOK và AI.
\item Khảo sát các công cụ quản lý dự án hiện có.
\item Phân tích yêu cầu nghiệp vụ.
\item Thiết kế hệ thống theo hướng đối tượng.
\item Thiết kế kiến trúc AI Multi-Agent.
\item Xây dựng nguyên mẫu hệ thống.
\item Thực nghiệm và đánh giá hiệu quả.
\end{enumerate}

\section{Đóng góp của đề tài}

Đề tài có các đóng góp chính như sau:

\begin{itemize}
\item Đề xuất mô hình quản lý dự án phần mềm tích hợp AI theo hướng Human-Centric.
\item Xây dựng kiến trúc Multi-Agent AI phục vụ các nghiệp vụ quản lý dự án.
\item Thiết kế quy trình cộng tác Human--AI thông qua HITL, HOTL và HOOTL.
\item Xây dựng nguyên mẫu hệ thống quản lý dự án hỗ trợ Agile Software Development.
\item Đánh giá khả năng ứng dụng AI trong hoạt động quản lý dự án phần mềm.
\end{itemize}

\section{Cấu trúc báo cáo}
Báo cáo được tổ chức thành mười lăm chương.

\begin{itemize}
\item Chương 1: Giới thiệu.
\item Chương 2: Cơ sở lý thuyết.
\item Chương 3: Tổng quan nghiên cứu.
\item Chương 4: Phân tích bài toán.
\item Chương 5: Đặc tả yêu cầu hệ thống.
\item Chương 6: Phân tích nghiệp vụ.
\item Chương 7: Thiết kế hệ thống.
\item Chương 8: Kiến trúc AI và Multi-Agent.
\item Chương 9: Thiết kế cơ sở dữ liệu.
\item Chương 10: Thiết kế Backend.
\item Chương 11: Thiết kế Frontend.
\item Chương 12: Triển khai hệ thống.
\item Chương 13: Thực nghiệm.
\item Chương 14: Đánh giá kết quả.
\item Chương 15: Kết luận và hướng phát triển.
\end{itemize}


%%%%%%%%%%%%%%%%%%%%%%%%%%%%%%%%%%%%%%%%%%%%%%%%%%%%%%%%%%%%%%
% CHƯƠNG 2: CƠ SỞ LÝ THUYẾT
%%%%%%%%%%%%%%%%%%%%%%%%%%%%%%%%%%%%%%%%%%%%%%%%%%%%%%%%%%%%%%
\chapter{CƠ SỞ LÝ THUYẾT}
\label{chap:background}

\section{Giới thiệu chương}
Quản lý dự án phần mềm là một lĩnh vực kết hợp giữa khoa học quản lý và kỹ thuật phần mềm nhằm bảo đảm dự án được triển khai đúng phạm vi, đúng tiến độ, đúng chi phí và đạt chất lượng mong muốn. Trong bối cảnh chuyển đổi số và sự phát triển nhanh chóng của trí tuệ nhân tạo (Artificial Intelligence -- AI), phương pháp quản lý dự án truyền thống đang dần được bổ sung bởi các công cụ hỗ trợ phân tích dữ liệu, dự báo và tự động hóa quá trình ra quyết định.

Để xây dựng một nền tảng quản lý dự án phần mềm hướng AI, việc nghiên cứu các cơ sở lý thuyết liên quan là yêu cầu cần thiết. Chương này trình bày các nền tảng lý thuyết được sử dụng trong luận văn, bao gồm quản lý dự án phần mềm, phương pháp phát triển Agile, Scrum Framework, PMBOK, trí tuệ nhân tạo, mô hình ngôn ngữ lớn (Large Language Models), hệ thống Multi-Agent và mô hình Human-Centric AI. Những nội dung này là cơ sở cho việc phân tích, thiết kế và triển khai hệ thống được trình bày trong các chương tiếp theo.

\section{Tổng quan về quản lý dự án phần mềm}

\subsection{Khái niệm dự án phần mềm}
Theo Viện Quản lý Dự án Hoa Kỳ (Project Management Institute -- PMI), dự án là một nỗ lực mang tính tạm thời được thực hiện nhằm tạo ra một sản phẩm, dịch vụ hoặc kết quả duy nhất \cite{pmbok7}.

Trong lĩnh vực kỹ thuật phần mềm, dự án phần mềm là tập hợp các hoạt động có kế hoạch nhằm phát triển, triển khai hoặc bảo trì một hệ thống phần mềm đáp ứng các yêu cầu nghiệp vụ của tổ chức hoặc khách hàng. Khác với các dự án sản xuất truyền thống, dự án phần mềm có mức độ không chắc chắn cao do yêu cầu thường xuyên thay đổi, sản phẩm mang tính vô hình và phụ thuộc nhiều vào tri thức của đội ngũ phát triển.

Một dự án phần mềm thông thường trải qua các giai đoạn chính như Hình \ref{fig:software_lifecycle}.

\begin{figure}[H]
\centering
\fbox{
\begin{minipage}{0.85\textwidth}
\centering
Requirements $\rightarrow$ Analysis $\rightarrow$ Design $\rightarrow$ Implementation $\rightarrow$ Testing $\rightarrow$ Deployment $\rightarrow$ Maintenance
\end{minipage}
}
\caption{Vòng đời phát triển phần mềm tổng quát}
\label{fig:software_lifecycle}
\end{figure}

Trong mỗi giai đoạn, người quản lý dự án phải thực hiện nhiều hoạt động khác nhau như phân tích yêu cầu, lập kế hoạch, phân công nguồn lực, theo dõi tiến độ, quản lý rủi ro và đánh giá chất lượng.

\subsection{Đặc điểm của dự án phần mềm}
So với các loại dự án khác, dự án phần mềm có một số đặc điểm nổi bật.

\begin{itemize}

\item
Sản phẩm mang tính vô hình, khó đánh giá mức độ hoàn thành thông qua quan sát trực tiếp.

\item
Yêu cầu nghiệp vụ thay đổi thường xuyên trong quá trình phát triển.

\item
Phụ thuộc lớn vào trình độ chuyên môn và khả năng phối hợp của đội ngũ phát triển.

\item
Khó ước lượng chính xác khối lượng công việc cũng như thời gian thực hiện.

\item
Tiềm ẩn nhiều rủi ro về kỹ thuật, nhân sự và thay đổi yêu cầu.

\item
Đòi hỏi quá trình kiểm thử và bảo trì liên tục ngay cả sau khi sản phẩm được triển khai.

\end{itemize}

Những đặc điểm trên làm cho việc quản lý dự án phần mềm trở nên phức tạp hơn nhiều so với các lĩnh vực quản lý truyền thống.

\subsection{Các thành phần của quản lý dự án phần mềm}
Theo PMBOK phiên bản thứ bảy, quản lý dự án bao gồm nhiều lĩnh vực tri thức khác nhau nhằm bảo đảm dự án đạt được mục tiêu đề ra \cite{pmbok7}. Trong phạm vi phát triển phần mềm, các thành phần quan trọng bao gồm:

\begin{enumerate}

\item Quản lý phạm vi (Scope Management)

Định nghĩa và kiểm soát phạm vi của dự án nhằm bảo đảm chỉ thực hiện những công việc cần thiết.

\item Quản lý tiến độ (Schedule Management)

Lập kế hoạch và theo dõi tiến độ thực hiện các công việc.

\item Quản lý chi phí (Cost Management)

Ước lượng, kiểm soát và tối ưu hóa ngân sách dự án.

\item Quản lý chất lượng (Quality Management)

Đảm bảo phần mềm đáp ứng các tiêu chuẩn chất lượng đã đặt ra.

\item Quản lý nguồn lực (Resource Management)

Phân công nhân sự, quản lý năng lực và tối ưu hóa hiệu suất của nhóm phát triển.

\item Quản lý truyền thông (Communication Management)

Đảm bảo luồng thông tin giữa các bên liên quan được truyền đạt đầy đủ và chính xác.

\item Quản lý rủi ro (Risk Management)

Nhận diện, đánh giá và xây dựng phương án xử lý các rủi ro có thể ảnh hưởng đến dự án.

\item Quản lý các bên liên quan (Stakeholder Management)

Phối hợp giữa khách hàng, Product Owner, Scrum Master, Project Manager và nhóm phát triển nhằm đạt được mục tiêu chung.

\end{enumerate}

Trong các dự án hiện đại, các lĩnh vực trên không hoạt động độc lập mà có mối quan hệ chặt chẽ với nhau. Ví dụ, việc thay đổi phạm vi sẽ ảnh hưởng trực tiếp đến tiến độ, chi phí và rủi ro của dự án.

\subsection{Các thách thức trong quản lý dự án phần mềm}
Nhiều nghiên cứu chỉ ra rằng tỷ lệ thất bại của các dự án phần mềm vẫn ở mức cao mặc dù đã có sự hỗ trợ của nhiều công cụ quản lý hiện đại. Các nguyên nhân phổ biến bao gồm:

\begin{itemize}

\item Yêu cầu không đầy đủ hoặc thay đổi liên tục.

\item Ước lượng thời gian và chi phí thiếu chính xác.

\item Thiếu khả năng phát hiện rủi ro từ sớm.

\item Phân bổ nguồn lực chưa hợp lý.

\item Thiếu dữ liệu hỗ trợ quá trình ra quyết định.

\item Giao tiếp không hiệu quả giữa các thành viên dự án.

\end{itemize}

Những hạn chế trên là động lực thúc đẩy việc ứng dụng trí tuệ nhân tạo vào quản lý dự án nhằm hỗ trợ phân tích dữ liệu và nâng cao chất lượng ra quyết định.

\subsection{Xu hướng ứng dụng trí tuệ nhân tạo trong quản lý dự án}
Sự phát triển mạnh mẽ của Machine Learning và Large Language Models đã mở ra hướng tiếp cận mới trong quản lý dự án phần mềm. AI có khả năng xử lý lượng lớn dữ liệu lịch sử, phát hiện các mẫu (patterns) và hỗ trợ người quản lý trong nhiều hoạt động như:

\begin{itemize}

\item Phân tích yêu cầu phần mềm.

\item Sinh User Story và Task.

\item Ước lượng Story Point.

\item Dự báo tiến độ hoàn thành Sprint.

\item Phát hiện nguy cơ chậm tiến độ.

\item Đề xuất phân công nhân sự phù hợp.

\item Sinh báo cáo và tài liệu dự án.

\end{itemize}

Theo Bachari và cộng sự \cite{bachari2025}, AI đặc biệt phát huy hiệu quả trong bốn lĩnh vực chính gồm quản lý tiến độ, quản lý nguồn lực, quản lý chi phí và quản lý rủi ro. Tuy nhiên, phần lớn các nghiên cứu hiện nay vẫn chỉ giải quyết từng bài toán riêng biệt thay vì xây dựng một nền tảng tích hợp toàn diện.

Điều này tạo động lực cho việc nghiên cứu và phát triển một hệ thống quản lý dự án phần mềm hướng AI, trong đó nhiều AI Agent phối hợp với con người để hỗ trợ toàn bộ quy trình phát triển phần mềm theo phương pháp Agile.

\section{Phương pháp phát triển phần mềm Agile}
\label{sec:agile}
\subsection{Lịch sử hình thành Agile}

Trong nhiều năm, mô hình Waterfall là phương pháp phát triển phần mềm phổ biến nhất. Waterfall chia quá trình phát triển thành các giai đoạn tuần tự như phân tích yêu cầu, thiết kế, lập trình, kiểm thử và triển khai. Mỗi giai đoạn chỉ được bắt đầu sau khi giai đoạn trước hoàn thành.

Mặc dù phù hợp với các dự án có yêu cầu ổn định, Waterfall bộc lộ nhiều hạn chế khi áp dụng cho các dự án phần mềm hiện đại, đặc biệt trong môi trường kinh doanh thay đổi nhanh. Những khó khăn thường gặp bao gồm:

\begin{itemize}
\item Khó thích nghi với sự thay đổi yêu cầu.
\item Chu kỳ phát triển dài.
\item Khách hàng chỉ được tiếp cận sản phẩm ở giai đoạn cuối.
\item Chi phí sửa đổi yêu cầu rất lớn.
\item Nguy cơ thất bại cao nếu phân tích yêu cầu ban đầu chưa đầy đủ.
\end{itemize}

Nhằm khắc phục các hạn chế này, tháng 2 năm 2001, mười bảy chuyên gia phát triển phần mềm đã họp tại Snowbird (Utah, Hoa Kỳ) và công bố \textit{Agile Manifesto}. Từ đây, Agile trở thành một triết lý phát triển phần mềm mới, nhấn mạnh tính linh hoạt, khả năng thích ứng và sự hợp tác giữa các thành viên dự án \cite{agilemanifesto}.

Agile không phải là một quy trình cụ thể mà là tập hợp các nguyên tắc và giá trị định hướng việc phát triển phần mềm. Nhiều framework như Scrum, Extreme Programming (XP), Kanban và Lean Software Development đều được xây dựng dựa trên các nguyên tắc của Agile.

\subsection{Giá trị cốt lõi của Agile}

Agile Manifesto xác định bốn giá trị cốt lõi làm nền tảng cho toàn bộ quá trình phát triển phần mềm.

\begin{enumerate}

\item
Cá nhân và sự tương tác quan trọng hơn quy trình và công cụ.

\item
Phần mềm hoạt động được quan trọng hơn tài liệu đầy đủ.

\item
Hợp tác với khách hàng quan trọng hơn đàm phán hợp đồng.

\item
Phản hồi với sự thay đổi quan trọng hơn việc tuân thủ kế hoạch ban đầu.

\end{enumerate}

Các giá trị này không phủ nhận vai trò của quy trình, tài liệu hay kế hoạch, mà nhấn mạnh rằng trong môi trường phát triển phần mềm hiện đại, yếu tố con người và khả năng thích ứng có vai trò quyết định đối với sự thành công của dự án.

\subsection{Mười hai nguyên lý của Agile}

Bên cạnh bốn giá trị cốt lõi, Agile còn được xây dựng trên mười hai nguyên lý cơ bản.

\begin{enumerate}

\item Đáp ứng nhu cầu khách hàng thông qua việc cung cấp phần mềm sớm và liên tục.

\item Chấp nhận sự thay đổi yêu cầu ngay cả ở giai đoạn cuối của dự án.

\item Phát hành phần mềm hoạt động trong thời gian ngắn, thường từ hai đến bốn tuần.

\item Nhà phát triển và khách hàng phải cộng tác thường xuyên.

\item Xây dựng dự án xoay quanh những cá nhân có động lực.

\item Giao tiếp trực tiếp là hình thức truyền đạt thông tin hiệu quả nhất.

\item Phần mềm hoạt động là thước đo chính của tiến độ.

\item Duy trì tốc độ phát triển bền vững.

\item Liên tục chú trọng chất lượng kỹ thuật và thiết kế tốt.

\item Đơn giản hóa tối đa công việc không cần thiết.

\item Các nhóm tự tổ chức sẽ tạo ra thiết kế tốt hơn.

\item Thường xuyên đánh giá và cải tiến quy trình làm việc.

\end{enumerate}

Các nguyên lý này giúp Agile có khả năng thích nghi với môi trường phát triển thay đổi liên tục và phù hợp với các dự án phần mềm có yêu cầu chưa ổn định.

\subsection{Vòng đời phát triển Agile}

Khác với Waterfall, Agile chia toàn bộ dự án thành nhiều chu kỳ phát triển ngắn gọi là Sprint.

Mỗi Sprint thường kéo dài từ hai đến bốn tuần và tạo ra một phiên bản phần mềm có thể sử dụng được.

\begin{figure}[H]
\centering
\fbox{
\begin{minipage}{0.95\textwidth}
\centering
Planning $\rightarrow$ Design $\rightarrow$ Development $\rightarrow$ Testing $\rightarrow$ Review $\rightarrow$ Retrospective $\rightarrow$ Next Sprint
\end{minipage}
}
\caption{Chu trình phát triển Agile}
\label{fig:agile_cycle}
\end{figure}

Sau mỗi Sprint, nhóm phát triển sẽ tổ chức Sprint Review và Sprint Retrospective nhằm đánh giá kết quả đạt được, thu thập phản hồi từ khách hàng và cải tiến quy trình cho Sprint tiếp theo.

Việc chia nhỏ dự án thành nhiều vòng lặp ngắn giúp giảm thiểu rủi ro, tăng khả năng kiểm soát tiến độ và cho phép khách hàng tham gia vào quá trình phát triển ngay từ những giai đoạn đầu.

\subsection{Các vai trò trong Agile}

Trong Agile nói chung và Scrum nói riêng, dự án thường bao gồm các vai trò sau:

\begin{itemize}

\item Product Owner

Đại diện khách hàng, chịu trách nhiệm quản lý Product Backlog và xác định giá trị của sản phẩm.

\item Scrum Master

Đảm bảo nhóm phát triển tuân thủ Scrum Framework, loại bỏ các trở ngại và hỗ trợ quá trình cải tiến liên tục.

\item Development Team

Nhóm kỹ sư phần mềm chịu trách nhiệm thiết kế, lập trình, kiểm thử và triển khai sản phẩm.

\item Stakeholders

Bao gồm khách hàng, người sử dụng cuối cùng và các bên liên quan khác.

\end{itemize}

Trong đề tài này, ngoài các vai trò truyền thống, hệ thống còn bổ sung nhiều AI Agent đóng vai trò trợ lý thông minh cho từng nhóm nghiệp vụ.

\subsection{Ưu điểm của Agile}

Agile mang lại nhiều lợi ích so với các phương pháp truyền thống.

\begin{itemize}

\item Thích nghi nhanh với thay đổi.

\item Rút ngắn chu kỳ phát triển.

\item Tăng mức độ hài lòng của khách hàng.

\item Phát hiện lỗi sớm.

\item Giảm rủi ro dự án.

\item Tăng khả năng cộng tác giữa các thành viên.

\item Nâng cao chất lượng sản phẩm.

\end{itemize}

Đặc biệt, việc chia nhỏ yêu cầu thành User Story và Task giúp AI có thể xử lý từng phần công việc một cách hiệu quả hơn.

\subsection{Hạn chế của Agile}

Mặc dù có nhiều ưu điểm, Agile vẫn tồn tại một số hạn chế.

\begin{itemize}

\item Khó áp dụng đối với dự án có quy mô rất lớn.

\item Phụ thuộc nhiều vào kinh nghiệm của Product Owner.

\item Story Point mang tính chủ quan.

\item Khó dự báo tiến độ dài hạn.

\item Việc phân công nguồn lực chủ yếu dựa trên kinh nghiệm.

\item Khó phát hiện sớm các rủi ro tiềm ẩn.

\end{itemize}

Những hạn chế này chính là cơ hội để ứng dụng trí tuệ nhân tạo nhằm hỗ trợ quá trình quản lý dự án.

\subsection{So sánh Agile và Waterfall}

\begin{table}[H]
\centering
\caption{So sánh Agile và Waterfall}
\begin{tabular}{|p{4cm}|p{5cm}|p{5cm}|}
\hline
Tiêu chí & Waterfall & Agile \\
\hline
Quy trình & Tuần tự & Lặp và tăng trưởng \\
\hline
Thay đổi yêu cầu & Khó & Linh hoạt \\
\hline
Khách hàng & Ít tham gia & Tham gia liên tục \\
\hline
Chu kỳ phát triển & Dài & Ngắn \\
\hline
Quản lý rủi ro & Cuối dự án & Liên tục \\
\hline
Khả năng thích nghi & Thấp & Cao \\
\hline
Phù hợp & Yêu cầu ổn định & Yêu cầu thay đổi \\
\hline
\end{tabular}
\label{tab:agile_vs_waterfall}
\end{table}

Bảng \ref{tab:agile_vs_waterfall} cho thấy Agile có nhiều ưu điểm trong môi trường phát triển phần mềm hiện đại, đặc biệt đối với các dự án có yêu cầu thay đổi liên tục.

\subsection{Lý do lựa chọn Agile cho đề tài}

Hệ thống được đề xuất trong luận văn hướng tới việc hỗ trợ toàn bộ vòng đời phát triển phần mềm theo Agile. Các chức năng như phân tích User Story, sinh Task, ước lượng Story Point, dự báo Sprint, phát hiện rủi ro và hỗ trợ phân công đều được thiết kế xoay quanh quy trình Agile.

Việc lựa chọn Agile thay vì Waterfall xuất phát từ các lý do sau:

\begin{enumerate}

\item Agile phù hợp với các dự án có yêu cầu thay đổi liên tục.

\item User Story là đầu vào tự nhiên cho các mô hình ngôn ngữ lớn.

\item Sprint cung cấp dữ liệu lịch sử phục vụ AI Forecasting.

\item Burndown Chart và Velocity hỗ trợ AI đánh giá tiến độ.

\item AI có thể tham gia hỗ trợ Product Owner và Scrum Master trong mọi Sprint.

\end{enumerate}

Do đó, Agile là nền tảng phù hợp nhất để xây dựng hệ thống quản lý dự án phần mềm hướng AI được đề xuất trong luận văn.

\section{Scrum Framework}
\label{sec:scrum}
\subsection{Giới thiệu Scrum}

Scrum là một framework phát triển phần mềm theo phương pháp Agile được đề xuất bởi Ken Schwaber và Jeff Sutherland nhằm hỗ trợ các nhóm phát triển giải quyết những bài toán có mức độ phức tạp cao thông qua quá trình phát triển lặp và tăng trưởng (Iterative and Incremental Development) \cite{scrumguide2020}. Không giống các mô hình quy trình truyền thống, Scrum không quy định chi tiết cách thức xây dựng sản phẩm mà chỉ cung cấp một khung quản lý giúp các nhóm phát triển có thể tự tổ chức (Self-management), thích nghi nhanh với sự thay đổi và liên tục cải tiến quy trình làm việc.

Trong Scrum, toàn bộ dự án được chia thành nhiều chu kỳ ngắn gọi là Sprint. Mỗi Sprint tạo ra một phiên bản phần mềm có thể sử dụng được (Increment), đồng thời cung cấp phản hồi giúp nhóm phát triển điều chỉnh kế hoạch cho các Sprint tiếp theo. Nhờ đó, Scrum giúp giảm thiểu rủi ro, tăng khả năng thích ứng và cải thiện chất lượng sản phẩm.

Đối với đề tài này, Scrum được lựa chọn làm nền tảng quản lý dự án vì phù hợp với việc tích hợp trí tuệ nhân tạo trong từng giai đoạn của vòng đời phát triển phần mềm.

%%%%%%%%%%%%%%%%%%%%%%%%%%%%%%%%%%%%%%%%%%%%%%%%%%%%%%%%%%%%%%

\subsection{Giá trị cốt lõi của Scrum}

Theo Scrum Guide 2020, Scrum được xây dựng dựa trên năm giá trị cốt lõi gồm:

\begin{enumerate}

\item Commitment

Nhóm phát triển cam kết hoàn thành mục tiêu Sprint và tạo ra giá trị cho khách hàng.

\item Focus

Toàn bộ nhóm tập trung vào Sprint Goal và các công việc ưu tiên cao nhất.

\item Openness

Các thành viên sẵn sàng chia sẻ thông tin, tiến độ và các khó khăn gặp phải.

\item Respect

Các thành viên tôn trọng lẫn nhau, tôn trọng ý kiến và chuyên môn của từng cá nhân.

\item Courage

Nhóm phát triển sẵn sàng đối mặt với khó khăn, thử nghiệm các giải pháp mới và chấp nhận thay đổi.

\end{enumerate}

Năm giá trị này đóng vai trò quan trọng trong việc xây dựng môi trường làm việc minh bạch và cộng tác hiệu quả.

%%%%%%%%%%%%%%%%%%%%%%%%%%%%%%%%%%%%%%%%%%%%%%%%%%%%%%%%%%%%%%

\subsection{Scrum Team}

Scrum Team là đơn vị trung tâm của toàn bộ framework. Theo Scrum Guide 2020, Scrum Team bao gồm ba vai trò chính như Hình \ref{fig:scrumteam}.

\begin{figure}[H]
\centering
\fbox{
\begin{minipage}{0.9\textwidth}
\centering
Product Owner $\rightarrow$ Scrum Team $\rightarrow$ Developers $\rightarrow$ Increment
\end{minipage}
}
\caption{Cấu trúc Scrum Team}
\label{fig:scrumteam}
\end{figure}

\subsubsection{Product Owner}

Product Owner chịu trách nhiệm tối đa hóa giá trị của sản phẩm thông qua việc quản lý Product Backlog. Các nhiệm vụ chính bao gồm:

\begin{itemize}

\item Thu thập yêu cầu.

\item Viết User Story.

\item Xác định mức độ ưu tiên.

\item Phê duyệt Story Point.

\item Quản lý Product Backlog.

\end{itemize}

Trong hệ thống đề xuất, Product Owner còn sử dụng AI Requirement Agent để phân tích yêu cầu và AI Planning Agent để sinh Task.

%%%%%%%%%%%%%%%%%%%%%%%%%%%%%%%%%%%%%%%%%%%%%%%%%%%%%%%%%%%%%%

\subsubsection{Scrum Master}

Scrum Master đảm bảo nhóm phát triển tuân thủ đúng Scrum Framework.

Các nhiệm vụ gồm:

\begin{itemize}

\item Điều phối Sprint.

\item Hỗ trợ nhóm phát triển.

\item Loại bỏ trở ngại.

\item Theo dõi tiến độ.

\item Tổ chức Sprint Review.

\item Tổ chức Sprint Retrospective.

\end{itemize}

Trong đề tài, Scrum Master còn nhận cảnh báo từ AI Risk Agent và AI Forecast Agent nhằm hỗ trợ việc ra quyết định.

%%%%%%%%%%%%%%%%%%%%%%%%%%%%%%%%%%%%%%%%%%%%%%%%%%%%%%%%%%%%%%

\subsubsection{Developers}

Developers là nhóm trực tiếp xây dựng sản phẩm.

Các nhiệm vụ bao gồm:

\begin{itemize}

\item Thiết kế.

\item Lập trình.

\item Kiểm thử.

\item Viết tài liệu.

\item Review mã nguồn.

\item Triển khai hệ thống.

\end{itemize}

AI Task Agent đóng vai trò hỗ trợ Developers trong việc sinh Task, sinh tài liệu và đánh giá tiến độ.

%%%%%%%%%%%%%%%%%%%%%%%%%%%%%%%%%%%%%%%%%%%%%%%%%%%%%%%%%%%%%%

\subsection{Scrum Artifacts}

Scrum định nghĩa ba Artifact quan trọng.

\subsubsection{Product Backlog}

Product Backlog là danh sách toàn bộ các yêu cầu của sản phẩm.

Mỗi mục trong Product Backlog thường bao gồm:

\begin{itemize}

\item User Story

\item Priority

\item Story Point

\item Business Value

\item Acceptance Criteria

\end{itemize}

Trong hệ thống đề xuất, Product Backlog được AI phân tích để phát hiện yêu cầu chưa rõ ràng trước khi chuyển sang Sprint Planning.

%%%%%%%%%%%%%%%%%%%%%%%%%%%%%%%%%%%%%%%%%%%%%%%%%%%%%%%%%%%%%%

\subsubsection{Sprint Backlog}

Sprint Backlog bao gồm các User Story được lựa chọn để thực hiện trong Sprint hiện tại.

Ngoài danh sách User Story, Sprint Backlog còn bao gồm:

\begin{itemize}

\item Task

\item Assignee

\item Remaining Time

\item Status

\item Burndown Data

\end{itemize}

Sprint Backlog là nguồn dữ liệu đầu vào của AI Forecast Agent.

%%%%%%%%%%%%%%%%%%%%%%%%%%%%%%%%%%%%%%%%%%%%%%%%%%%%%%%%%%%%%%

\subsubsection{Increment}

Increment là phiên bản phần mềm hoàn chỉnh được tạo ra sau mỗi Sprint.

Increment phải đáp ứng Definition of Done (DoD) và có thể triển khai ngay nếu cần.

%%%%%%%%%%%%%%%%%%%%%%%%%%%%%%%%%%%%%%%%%%%%%%%%%%%%%%%%%%%%%%

\subsection{Các sự kiện trong Scrum}

Scrum quy định năm sự kiện chính.

\subsubsection{Sprint Planning}

Sprint Planning xác định:

\begin{itemize}

\item Sprint Goal.

\item User Story được chọn.

\item Story Point.

\item Task Breakdown.

\item Resource Allocation.

\end{itemize}

Trong hệ thống đề xuất, AI hỗ trợ Product Owner phân rã User Story thành Task và đề xuất Story Point.

%%%%%%%%%%%%%%%%%%%%%%%%%%%%%%%%%%%%%%%%%%%%%%%%%%%%%%%%%%%%%%

\subsubsection{Daily Scrum}

Daily Scrum là cuộc họp ngắn diễn ra hàng ngày.

Mỗi thành viên trả lời ba câu hỏi:

\begin{enumerate}

\item Hôm qua đã làm gì?

\item Hôm nay sẽ làm gì?

\item Có khó khăn gì?

\end{enumerate}

AI Dashboard có thể tổng hợp tiến độ tự động từ trạng thái Task nhằm giảm thời gian họp.

%%%%%%%%%%%%%%%%%%%%%%%%%%%%%%%%%%%%%%%%%%%%%%%%%%%%%%%%%%%%%%

\subsubsection{Sprint Review}

Sprint Review được tổ chức cuối Sprint nhằm:

\begin{itemize}

\item Demo sản phẩm.

\item Nhận phản hồi.

\item Cập nhật Product Backlog.

\end{itemize}

%%%%%%%%%%%%%%%%%%%%%%%%%%%%%%%%%%%%%%%%%%%%%%%%%%%%%%%%%%%%%%

\subsubsection{Sprint Retrospective}

Sprint Retrospective giúp nhóm phát triển:

\begin{itemize}

\item Đánh giá quy trình.

\item Xác định điểm mạnh.

\item Phát hiện vấn đề.

\item Đề xuất cải tiến.

\end{itemize}

Trong hệ thống đề xuất, AI Documentation Agent hỗ trợ tổng hợp báo cáo Retrospective tự động.

\subsection{Story Point}

Story Point là đơn vị đo lường tương đối dùng để ước lượng mức độ khó của User Story.

Khác với thời gian thực hiện, Story Point phản ánh tổng hợp nhiều yếu tố:

\begin{itemize}

\item Độ phức tạp.

\item Khối lượng công việc.

\item Rủi ro.

\item Mức độ không chắc chắn.

\end{itemize}

Hầu hết các nhóm Scrum sử dụng dãy Fibonacci:

\[
1,\; 2,\; 3,\; 5,\; 8,\; 13,\; 21,\; 34
\]

để biểu diễn Story Point.

Trong hệ thống đề xuất, AI Story Point Agent sử dụng User Story, lịch sử Sprint và dữ liệu Velocity để đề xuất Story Point kèm theo giải thích.

%%%%%%%%%%%%%%%%%%%%%%%%%%%%%%%%%%%%%%%%%%%%%%%%%%%%%%%%%%%%%%

\subsection{Velocity}

Velocity là tổng Story Point hoàn thành trong một Sprint.

Giả sử Sprint có \(n\) User Story, Velocity được tính như sau:

\[
Velocity=\sum_{i=1}^{n}SP_i
\]

trong đó \(SP_i\) là Story Point của User Story thứ \(i\).

Velocity được AI Forecast Agent sử dụng để dự báo thời gian hoàn thành dự án.

%%%%%%%%%%%%%%%%%%%%%%%%%%%%%%%%%%%%%%%%%%%%%%%%%%%%%%%%%%%%%%

\subsection{Burndown Chart}

Burndown Chart biểu diễn lượng công việc còn lại theo thời gian.

Giả sử \(W_t\) là khối lượng công việc còn lại tại ngày \(t\), khi đó

\[
W_t=W_0-\sum_{i=1}^{t}Completed_i
\]

Nếu đường Burndown nằm trên đường lý tưởng, Sprint có nguy cơ chậm tiến độ.

AI Risk Agent liên tục phân tích Burndown Chart để phát hiện các bất thường và gửi cảnh báo đến Scrum Master.

%%%%%%%%%%%%%%%%%%%%%%%%%%%%%%%%%%%%%%%%%%%%%%%%%%%%%%%%%%%%%%

\subsection{Vai trò của Scrum trong hệ thống đề xuất}

Toàn bộ kiến trúc của hệ thống được xây dựng dựa trên quy trình Scrum.

Quy trình xử lý nghiệp vụ tổng quát được thể hiện như sau:

\begin{center}
Requirement $\rightarrow$ User Story $\rightarrow$ AI Requirement Analysis $\rightarrow$ Task Breakdown $\rightarrow$ Story Point Estimation $\rightarrow$ Sprint Planning $\rightarrow$ Task Assignment $\rightarrow$ Sprint Execution $\rightarrow$ AI Risk Monitoring $\rightarrow$ Forecast $\rightarrow$ Sprint Review $\rightarrow$ Retrospective $\rightarrow$ Release
\end{center}

Có thể thấy AI không thay thế Scrum mà đóng vai trò trợ lý thông minh trong từng bước của quy trình. Điều này phù hợp với định hướng Human-Centric AI và các cơ chế Human-in-the-Loop, Human-on-the-Loop và Human-out-of-the-Loop được trình bày trong Product Vision của hệ thống.

\section{PMBOK và Quản lý dự án hiện đại}
\label{sec:pmbok}
\subsection{Giới thiệu PMBOK}

Project Management Body of Knowledge (PMBOK) là bộ tiêu chuẩn quản lý dự án do Viện Quản lý Dự án Hoa Kỳ (Project Management Institute -- PMI) xây dựng và phát triển. PMBOK được xem là một trong những tiêu chuẩn phổ biến nhất trong lĩnh vực quản lý dự án và được áp dụng trong nhiều ngành nghề khác nhau như công nghệ thông tin, xây dựng, sản xuất, tài chính và y tế \cite{pmbok7}.

Khác với Scrum hay Agile, PMBOK không phải là một quy trình phát triển phần mềm mà là tập hợp các nguyên tắc, quy trình và thực tiễn tốt nhất (Best Practices) giúp các nhà quản lý lập kế hoạch, tổ chức, giám sát và kiểm soát dự án.

Phiên bản PMBOK lần thứ bảy (PMBOK 7th Edition) chuyển từ cách tiếp cận dựa trên quy trình (Process-based) sang cách tiếp cận dựa trên giá trị (Value Delivery System), trong đó nhấn mạnh khả năng thích ứng, quản trị dự án và tạo giá trị liên tục cho các bên liên quan.

%%%%%%%%%%%%%%%%%%%%%%%%%%%%%%%%%%%%%%%%%%%%%%%%%%%%%%%%%%%%%%

\subsection{Các nguyên tắc của PMBOK 7}

PMBOK phiên bản thứ bảy xác định mười hai nguyên tắc cốt lõi nhằm định hướng hoạt động quản lý dự án.

\begin{enumerate}

\item Stewardship

Quản lý dự án phải được thực hiện một cách minh bạch, có trách nhiệm và hướng tới lợi ích chung của tổ chức.

\item Team

Xây dựng môi trường làm việc cộng tác, khuyến khích chia sẻ tri thức và phát triển năng lực của nhóm.

\item Stakeholders

Liên tục tương tác với các bên liên quan nhằm bảo đảm các yêu cầu được hiểu đúng và đáp ứng đầy đủ.

\item Value

Mọi hoạt động của dự án đều hướng tới việc tạo ra giá trị cho khách hàng và tổ chức.

\item Systems Thinking

Xem dự án như một hệ thống gồm nhiều thành phần có mối quan hệ phụ thuộc lẫn nhau.

\item Leadership

Nhà quản lý cần tạo động lực và định hướng cho nhóm phát triển.

\item Tailoring

Lựa chọn phương pháp quản lý phù hợp với từng dự án thay vì áp dụng cứng nhắc.

\item Quality

Đảm bảo chất lượng trong toàn bộ vòng đời phát triển.

\item Complexity

Nhận diện và quản lý các yếu tố phức tạp có thể ảnh hưởng đến dự án.

\item Risk

Chủ động phát hiện và xử lý các rủi ro.

\item Adaptability

Khả năng thích ứng với sự thay đổi của môi trường.

\item Change

Quản lý thay đổi nhằm bảo đảm dự án luôn tạo ra giá trị.

\end{enumerate}

Có thể nhận thấy nhiều nguyên tắc của PMBOK phù hợp với triết lý Agile và tạo điều kiện thuận lợi cho việc tích hợp các công nghệ AI vào quá trình quản lý dự án.

%%%%%%%%%%%%%%%%%%%%%%%%%%%%%%%%%%%%%%%%%%%%%%%%%%%%%%%%%%%%%%

\subsection{Các lĩnh vực quản lý dự án}

Trong PMBOK truyền thống, quản lý dự án bao gồm nhiều lĩnh vực tri thức (Knowledge Areas). Các lĩnh vực này vẫn giữ vai trò quan trọng trong PMBOK 7 và được thể hiện trong Bảng \ref{tab:pmbok_knowledge}.

\begin{table}[H]
\centering
\caption{Các lĩnh vực quản lý dự án theo PMBOK}
\begin{tabular}{|p{5cm}|p{8cm}|}
\hline
Lĩnh vực & Mục tiêu \\
\hline
Scope Management & Quản lý phạm vi dự án \\
\hline
Schedule Management & Quản lý tiến độ \\
\hline
Cost Management & Quản lý chi phí \\
\hline
Quality Management & Quản lý chất lượng \\
\hline
Resource Management & Quản lý nguồn lực \\
\hline
Communication Management & Quản lý truyền thông \\
\hline
Risk Management & Quản lý rủi ro \\
\hline
Procurement Management & Quản lý mua sắm \\
\hline
Stakeholder Management & Quản lý các bên liên quan \\
\hline
Integration Management & Điều phối toàn bộ dự án \\
\hline
\end{tabular}
\label{tab:pmbok_knowledge}
\end{table}

Các lĩnh vực trên có mối liên hệ chặt chẽ với nhau. Ví dụ, việc mở rộng phạm vi dự án sẽ ảnh hưởng trực tiếp đến tiến độ, chi phí và rủi ro.

%%%%%%%%%%%%%%%%%%%%%%%%%%%%%%%%%%%%%%%%%%%%%%%%%%%%%%%%%%%%%%

\subsection{Vòng đời quản lý dự án}

PMBOK chia vòng đời dự án thành năm nhóm quy trình chính.

\begin{figure}[H]
\centering
\fbox{
\begin{minipage}{0.9\textwidth}
\centering
Initiating $\rightarrow$ Planning $\rightarrow$ Executing $\rightarrow$ Monitoring \& Controlling $\rightarrow$ Closing
\end{minipage}
}
\caption{Vòng đời quản lý dự án theo PMBOK}
\label{fig:pmbok_process}
\end{figure}

Trong từng giai đoạn, nhà quản lý cần thực hiện nhiều hoạt động khác nhau nhằm đảm bảo dự án đạt được mục tiêu về phạm vi, tiến độ và chất lượng.

%%%%%%%%%%%%%%%%%%%%%%%%%%%%%%%%%%%%%%%%%%%%%%%%%%%%%%%%%%%%%%

\subsection{PMBOK và Agile}

Trong nhiều năm, PMBOK thường bị xem là phương pháp quản lý truyền thống, trong khi Agile đại diện cho phương pháp phát triển linh hoạt.

Tuy nhiên, PMBOK phiên bản mới đã hỗ trợ nhiều mô hình phát triển khác nhau, bao gồm:

\begin{itemize}

\item Predictive

\item Iterative

\item Incremental

\item Agile

\item Hybrid

\end{itemize}

Điều này cho phép kết hợp các nguyên tắc quản lý của PMBOK với quy trình phát triển Scrum để tận dụng ưu điểm của cả hai phương pháp.

Trong đề tài này, Scrum được sử dụng để tổ chức hoạt động phát triển phần mềm, trong khi PMBOK cung cấp các nguyên tắc quản lý dự án ở mức chiến lược.

%%%%%%%%%%%%%%%%%%%%%%%%%%%%%%%%%%%%%%%%%%%%%%%%%%%%%%%%%%%%%%

\subsection{Vai trò của AI trong các lĩnh vực PMBOK}

Theo Bachari và cộng sự \cite{bachari2025}, AI có thể hỗ trợ nhiều lĩnh vực quản lý dự án khác nhau.

Bảng \ref{tab:ai_pmbok} trình bày mối quan hệ giữa các AI Agent trong hệ thống và các lĩnh vực quản lý của PMBOK.

\begin{table}[H]
\centering
\caption{Ánh xạ AI Agent với các lĩnh vực PMBOK}
\begin{tabular}{|p{4cm}|p{6cm}|p{4cm}|}
\hline
AI Agent & Chức năng & PMBOK \\
\hline
Requirement Agent & Phân tích yêu cầu & Scope \\
\hline
Task Planning Agent & Phân rã User Story & Schedule \\
\hline
Story Point Agent & Ước lượng khối lượng công việc & Schedule \\
\hline
Forecast Agent & Dự báo tiến độ & Schedule \\
\hline
Risk Agent & Phát hiện rủi ro & Risk \\
\hline
Assignment Agent & Phân công nhân sự & Resource \\
\hline
Documentation Agent & Sinh báo cáo & Communication \\
\hline
Dashboard Agent & Theo dõi KPI & Integration \\
\hline
\end{tabular}
\label{tab:ai_pmbok}
\end{table}

Bảng trên cho thấy hệ thống được đề xuất không chỉ hỗ trợ một hoạt động đơn lẻ mà bao phủ nhiều lĩnh vực quản lý dự án khác nhau theo PMBOK.

%%%%%%%%%%%%%%%%%%%%%%%%%%%%%%%%%%%%%%%%%%%%%%%%%%%%%%%%%%%%%%

\subsection{Ý nghĩa của PMBOK đối với đề tài}

Việc áp dụng PMBOK trong đề tài mang lại nhiều lợi ích.

Thứ nhất, PMBOK cung cấp một khung quản lý dự án toàn diện giúp hệ thống không chỉ phục vụ nhóm phát triển phần mềm mà còn hỗ trợ các nhà quản lý ở cấp chiến lược.

Thứ hai, các lĩnh vực quản lý của PMBOK tạo cơ sở để xác định các AI Agent cần thiết trong hệ thống.

Thứ ba, việc kết hợp PMBOK với Scrum và trí tuệ nhân tạo tạo nên một mô hình quản lý dự án hiện đại, vừa đảm bảo tính linh hoạt của Agile vừa giữ được khả năng kiểm soát theo chuẩn quản lý dự án quốc tế.

Trong các chương tiếp theo, các nguyên tắc của PMBOK sẽ được sử dụng để thiết kế kiến trúc hệ thống, xây dựng các AI Agent và đánh giá khả năng hỗ trợ quản lý dự án của nền tảng được đề xuất.

\section{Trí tuệ nhân tạo trong quản lý dự án phần mềm}
\label{sec:ai_pm}
\subsection{Giới thiệu}

Trong những năm gần đây, sự phát triển mạnh mẽ của trí tuệ nhân tạo (Artificial Intelligence -- AI) đã tạo ra những thay đổi đáng kể trong nhiều lĩnh vực như y tế, tài chính, giáo dục, sản xuất và công nghệ thông tin. Đặc biệt, sự xuất hiện của các mô hình ngôn ngữ lớn (Large Language Models -- LLMs) đã mở rộng phạm vi ứng dụng của AI sang các bài toán đòi hỏi khả năng suy luận, xử lý ngôn ngữ tự nhiên và hỗ trợ ra quyết định.

Trong quản lý dự án phần mềm, AI không còn chỉ được sử dụng để dự báo tiến độ hoặc phân tích dữ liệu lịch sử mà đã trở thành một thành phần hỗ trợ xuyên suốt vòng đời phát triển phần mềm. AI có thể tham gia vào quá trình phân tích yêu cầu, sinh tài liệu, phân rã User Story, ước lượng Story Point, phát hiện rủi ro, dự báo Sprint, hỗ trợ phân công nhân sự và sinh báo cáo quản lý.

Theo Bachari và cộng sự \cite{bachari2025}, AI hiện được xem là một trong những công nghệ có khả năng thay đổi phương thức quản lý dự án trong thập kỷ tới. Thay vì thay thế Project Manager, AI đóng vai trò là trợ lý thông minh (Intelligent Assistant), hỗ trợ con người xử lý khối lượng dữ liệu lớn và đưa ra các khuyến nghị dựa trên dữ liệu.

%%%%%%%%%%%%%%%%%%%%%%%%%%%%%%%%%%%%%%%%%%%%%%%%%%%%%%%%%%%%%%

\subsection{Khái niệm trí tuệ nhân tạo}

Trí tuệ nhân tạo là lĩnh vực của khoa học máy tính nghiên cứu việc xây dựng các hệ thống có khả năng thực hiện những nhiệm vụ vốn đòi hỏi trí thông minh của con người.

Theo Russell và Norvig \cite{russell2021}, AI được định nghĩa là:

\begin{quote}
``Artificial Intelligence is the study of intelligent agents that perceive their environment and take actions that maximize their chances of achieving their goals.''
\end{quote}

Một hệ thống AI điển hình bao gồm các thành phần:

\begin{itemize}
\item Thu thập dữ liệu;
\item Tiền xử lý dữ liệu;
\item Học mô hình;
\item Suy luận;
\item Đưa ra quyết định;
\item Học từ phản hồi.
\end{itemize}

Trong đề tài này, AI không trực tiếp điều khiển dự án mà hoạt động như một hệ thống hỗ trợ quyết định (Decision Support System).

%%%%%%%%%%%%%%%%%%%%%%%%%%%%%%%%%%%%%%%%%%%%%%%%%%%%%%%%%%%%%%

\subsection{Phân loại AI}

Theo mức độ năng lực, AI có thể được chia thành ba nhóm chính.

\subsubsection{Artificial Narrow Intelligence (ANI)}

ANI là dạng AI phổ biến hiện nay.

Đây là các hệ thống được thiết kế để giải quyết một nhiệm vụ cụ thể.

Ví dụ:

\begin{itemize}
\item ChatGPT
\item GitHub Copilot
\item Google Translate
\item AlphaFold
\end{itemize}

Hệ thống trong luận văn thuộc nhóm ANI.

%%%%%%%%%%%%%%%%%%%%%%%%%%%%%%%%%%%%%%%%%%%%%%%%%%%%%%%%%%%%%%

\subsubsection{Artificial General Intelligence (AGI)}

AGI là AI có khả năng thực hiện mọi nhiệm vụ trí tuệ tương tự con người.

Đến thời điểm hiện tại AGI vẫn đang là mục tiêu nghiên cứu.

%%%%%%%%%%%%%%%%%%%%%%%%%%%%%%%%%%%%%%%%%%%%%%%%%%%%%%%%%%%%%%

\subsubsection{Artificial Super Intelligence (ASI)}

ASI là AI vượt qua năng lực nhận thức của con người trong hầu hết các lĩnh vực.

Hiện nay chưa tồn tại hệ thống ASI thực tế.

%%%%%%%%%%%%%%%%%%%%%%%%%%%%%%%%%%%%%%%%%%%%%%%%%%%%%%%%%%%%%%

\subsection{Machine Learning}

Machine Learning là nhánh của AI cho phép máy tính học từ dữ liệu thay vì lập trình bằng các luật cố định.

Giả sử tập dữ liệu huấn luyện

\[
D=\{(x_i,y_i)\}_{i=1}^{n}
\]

Machine Learning tìm hàm \(f(x)\) sao cho \(f(x_i)\approx y_i\) với sai số nhỏ nhất.

Các phương pháp Machine Learning phổ biến gồm:

\begin{itemize}

\item Supervised Learning

\item Unsupervised Learning

\item Semi-supervised Learning

\item Reinforcement Learning

\end{itemize}

Trong quản lý dự án, Machine Learning được sử dụng để:

\begin{itemize}

\item dự báo thời gian hoàn thành;

\item phân loại mức độ rủi ro;

\item dự đoán chi phí;

\item phát hiện bất thường.

\end{itemize}

%%%%%%%%%%%%%%%%%%%%%%%%%%%%%%%%%%%%%%%%%%%%%%%%%%%%%%%%%%%%%%

\subsection{Deep Learning}

Deep Learning là tập con của Machine Learning sử dụng mạng nơ-ron nhiều lớp để học biểu diễn dữ liệu.

Một neuron được mô tả bởi

\[
y=f(Wx+b)
\]

trong đó \(W\) là ma trận trọng số, \(b\) là bias, \(f\) là hàm kích hoạt.

Deep Learning giúp AI học được các đặc trưng phức tạp mà không cần thiết kế thủ công.

Ngày nay, hầu hết các LLM đều được xây dựng dựa trên Deep Learning.

%%%%%%%%%%%%%%%%%%%%%%%%%%%%%%%%%%%%%%%%%%%%%%%%%%%%%%%%%%%%%%

\subsection{Natural Language Processing}

Natural Language Processing (NLP) là lĩnh vực nghiên cứu việc xử lý ngôn ngữ tự nhiên bằng máy tính.

Trong đề tài này, NLP đóng vai trò rất quan trọng vì đầu vào chủ yếu của hệ thống là:

\begin{itemize}

\item Requirement;

\item Product Vision;

\item User Story;

\item Acceptance Criteria;

\item Sprint Report;

\item Meeting Note.

\end{itemize}

AI cần hiểu được ý nghĩa của các tài liệu này để:

\begin{itemize}

\item phát hiện yêu cầu chưa rõ;

\item sinh Task;

\item sinh tài liệu;

\item đề xuất Story Point.

\end{itemize}

%%%%%%%%%%%%%%%%%%%%%%%%%%%%%%%%%%%%%%%%%%%%%%%%%%%%%%%%%%%%%%

\subsection{Sự phát triển của Large Language Models}

Sự ra đời của Transformer \cite{vaswani2017attention} đánh dấu bước ngoặt quan trọng trong NLP.

Khác với RNN hay LSTM, Transformer sử dụng cơ chế Self-Attention cho phép xử lý song song và học được quan hệ dài hạn giữa các token.

Kiến trúc Transformer bao gồm:

\begin{itemize}

\item Input Embedding;

\item Positional Encoding;

\item Multi-head Self Attention;

\item Feed Forward Network;

\item Layer Normalization;

\item Output Layer.

\end{itemize}

Từ Transformer, nhiều mô hình LLM nổi tiếng đã được phát triển như:

\begin{itemize}

\item GPT Series;

\item BERT;

\item T5;

\item LLaMA;

\item Claude;

\item Gemini;

\item DeepSeek.

\end{itemize}

Các mô hình này có khả năng hiểu và sinh văn bản với chất lượng gần tương đương con người.

%%%%%%%%%%%%%%%%%%%%%%%%%%%%%%%%%%%%%%%%%%%%%%%%%%%%%%%%%%%%%%

\subsection{Khả năng của LLM trong Software Engineering}

Các nghiên cứu gần đây cho thấy LLM có thể hỗ trợ nhiều hoạt động trong Software Engineering.

\begin{table}[H]
\centering
\caption{Ứng dụng của LLM trong Software Engineering}
\begin{tabular}{|p{6cm}|p{8cm}|}
\hline
Hoạt động & Khả năng hỗ trợ \\
\hline
Requirement Analysis & Làm rõ yêu cầu \\
\hline
User Story & Sinh User Story \\
\hline
Task Planning & Phân rã Task \\
\hline
Coding & Sinh mã nguồn \\
\hline
Testing & Sinh Test Case \\
\hline
Documentation & Sinh tài liệu \\
\hline
Code Review & Phân tích chất lượng mã nguồn \\
\hline
Bug Detection & Phát hiện lỗi \\
\hline
\end{tabular}
\end{table}

Có thể thấy LLM không chỉ hỗ trợ lập trình mà còn tham gia vào nhiều hoạt động quản lý dự án.

\subsection{Ứng dụng AI trong quản lý dự án phần mềm}
%%%%%%%%%%%%%%%%%%%%%%%%%%%%%%%%%%%%%%%%%%%%%%%%%%%%%%%%%%%%%%

Trong nhiều năm, các hoạt động quản lý dự án chủ yếu dựa vào kinh nghiệm của Project Manager và dữ liệu lịch sử của tổ chức. Mặc dù các hệ thống như Jira, Trello và Azure DevOps hỗ trợ lưu trữ thông tin và theo dõi tiến độ, việc ra quyết định vẫn mang tính thủ công và phụ thuộc nhiều vào năng lực của người quản lý.

Sự phát triển của Artificial Intelligence đã tạo ra một hướng tiếp cận mới. Thay vì chỉ đóng vai trò là công cụ quản lý dữ liệu, AI có thể phân tích khối lượng lớn dữ liệu lịch sử, phát hiện các mẫu (Pattern), dự báo xu hướng và hỗ trợ người quản lý đưa ra quyết định chính xác hơn.

Theo nghiên cứu tổng quan của Bachari và cộng sự \cite{bachari2025}, AI đang được ứng dụng trong hầu hết các lĩnh vực của quản lý dự án hiện đại, đặc biệt là:

\begin{itemize}
\item Quản lý phạm vi dự án;
\item Quản lý tiến độ;
\item Quản lý chi phí;
\item Quản lý nguồn lực;
\item Quản lý rủi ro;
\item Quản lý chất lượng;
\item Quản lý tri thức.
\end{itemize}

AI không thay thế Project Manager mà đóng vai trò như một hệ thống hỗ trợ quyết định (Decision Support System), giúp giảm thời gian phân tích dữ liệu và nâng cao chất lượng của các quyết định quản lý.

%%%%%%%%%%%%%%%%%%%%%%%%%%%%%%%%%%%%%%%%%%%%%%%%%%%%%%%%%%%%%%

\subsection{AI trong quản lý phạm vi dự án}

Quản lý phạm vi (Scope Management) là quá trình xác định và kiểm soát các yêu cầu của dự án nhằm bảo đảm nhóm phát triển chỉ thực hiện những công việc cần thiết.

Trong các dự án phần mềm, tài liệu yêu cầu thường được viết dưới dạng ngôn ngữ tự nhiên nên rất dễ xuất hiện các vấn đề như:

\begin{itemize}
\item yêu cầu mơ hồ;
\item yêu cầu trùng lặp;
\item yêu cầu mâu thuẫn;
\item yêu cầu thiếu thông tin;
\item yêu cầu không khả thi.
\end{itemize}

Các mô hình ngôn ngữ lớn có khả năng xử lý ngôn ngữ tự nhiên nên có thể hỗ trợ Product Owner trong việc:

\begin{enumerate}
\item phát hiện yêu cầu chưa rõ ràng;
\item đề xuất câu hỏi làm rõ;
\item phân loại Requirement;
\item sinh User Story;
\item sinh Acceptance Criteria;
\item xác định Business Value.
\end{enumerate}

Trong hệ thống được đề xuất, Requirement Agent chịu trách nhiệm thực hiện các chức năng này trước khi Product Backlog được tạo.

%%%%%%%%%%%%%%%%%%%%%%%%%%%%%%%%%%%%%%%%%%%%%%%%%%%%%%%%%%%%%%

\subsection{AI trong quản lý tiến độ}

Schedule Management là lĩnh vực được nghiên cứu nhiều nhất trong ứng dụng AI vào quản lý dự án.

Thông thường, Project Manager phải dự báo:

\begin{itemize}
\item thời gian hoàn thành Sprint;
\item thời gian hoàn thành Release;
\item khả năng chậm tiến độ;
\item Velocity của nhóm.
\end{itemize}

Các phương pháp truyền thống chủ yếu dựa trên kinh nghiệm và dữ liệu thống kê đơn giản.

Trong khi đó, AI có thể học từ:

\begin{itemize}
\item lịch sử Sprint;
\item Burndown Chart;
\item Story Point;
\item Velocity;
\item thời gian thực hiện Task;
\item số lượng Bug.
\end{itemize}

Giả sử \(Velocity_k\) là Velocity của Sprint thứ \(k\).

Velocity trung bình được xác định theo:

\[
\overline{V}
=
\frac{1}{n}
\sum_{i=1}^{n}
Velocity_i
\]

Từ Velocity trung bình, AI Forecast Agent có thể dự đoán số Sprint còn lại của dự án:

\[
Sprint_{remain}
=
\frac{StoryPoint_{remain}}
{\overline{V}}
\]

Mặc dù đây là công thức đơn giản, hệ thống còn kết hợp dữ liệu lịch sử và các yếu tố rủi ro để cải thiện độ chính xác.

%%%%%%%%%%%%%%%%%%%%%%%%%%%%%%%%%%%%%%%%%%%%%%%%%%%%%%%%%%%%%%

\subsection{AI trong quản lý rủi ro}

Theo PMBOK, Risk Management bao gồm:

\begin{enumerate}
\item Nhận diện rủi ro.
\item Phân tích rủi ro.
\item Đánh giá mức độ ảnh hưởng.
\item Lập kế hoạch ứng phó.
\end{enumerate}

AI giúp tự động hóa ba bước đầu tiên.

Ví dụ:

\begin{itemize}
\item Sprint liên tục chậm tiến độ;
\item số Bug tăng nhanh;
\item Developer quá tải;
\item Velocity giảm;
\item Burn-down Chart lệch khỏi đường lý tưởng.
\end{itemize}

AI Risk Agent liên tục phân tích dữ liệu thời gian thực và tính điểm rủi ro:

\[
RiskScore
=
\sum_{i=1}^{m}
w_i
x_i
\]

trong đó \(x_i\) là đặc trưng của dự án, \(w_i\) là trọng số học được.

Nếu \(RiskScore>T\), AI sẽ gửi cảnh báo đến Scrum Master.

%%%%%%%%%%%%%%%%%%%%%%%%%%%%%%%%%%%%%%%%%%%%%%%%%%%%%%%%%%%%%%

\subsection{AI trong quản lý nguồn lực}

Resource Management là bài toán phân công Developer vào các Task phù hợp.

Các yếu tố cần xem xét gồm:

\begin{itemize}
\item kỹ năng;
\item kinh nghiệm;
\item khối lượng công việc hiện tại;
\item lịch sử thực hiện;
\item mức độ ưu tiên.
\end{itemize}

Giả sử \(Skill(d,t)\) biểu diễn mức độ phù hợp giữa Developer \(d\) và Task \(t\).

AI Assignment Agent lựa chọn:

\[
d^*
=
\arg\max
Skill(d,t)
\]

đồng thời cân bằng tải công việc nhằm tránh tình trạng quá tải cho một thành viên.

%%%%%%%%%%%%%%%%%%%%%%%%%%%%%%%%%%%%%%%%%%%%%%%%%%%%%%%%%%%%%%

\subsection{AI trong quản lý chi phí}

Mặc dù đề tài không tập trung vào Cost Management, AI vẫn có khả năng hỗ trợ dự báo chi phí thông qua:

\begin{itemize}
\item Story Point;
\item Velocity;
\item số lượng nhân sự;
\item thời gian phát triển;
\item dữ liệu lịch sử.
\end{itemize}

Thông tin này giúp Project Manager đánh giá khả năng vượt ngân sách từ sớm.

%%%%%%%%%%%%%%%%%%%%%%%%%%%%%%%%%%%%%%%%%%%%%%%%%%%%%%%%%%%%%%

\subsection{AI trong quản lý tri thức}

Một trong những ưu điểm lớn của LLM là khả năng xử lý lượng lớn tài liệu.

AI Documentation Agent có thể:

\begin{itemize}
\item sinh tài liệu SRS;
\item sinh Meeting Minutes;
\item sinh Sprint Report;
\item sinh Retrospective Report;
\item sinh Release Note;
\item sinh API Documentation.
\end{itemize}

Điều này giúp giảm đáng kể khối lượng công việc thủ công của nhóm phát triển.

%%%%%%%%%%%%%%%%%%%%%%%%%%%%%%%%%%%%%%%%%%%%%%%%%%%%%%%%%%%%%%

\subsection{Lợi ích của AI trong quản lý dự án}

Việc tích hợp AI mang lại nhiều lợi ích.

\begin{table}[H]
\centering
\caption{Lợi ích của AI trong quản lý dự án}
\begin{tabular}{|p{6cm}|p{8cm}|}
\hline
Hoạt động & Lợi ích \\
\hline
Requirement Analysis & Giảm yêu cầu mơ hồ \\
\hline
Task Breakdown & Tiết kiệm thời gian lập kế hoạch \\
\hline
Story Point Estimation & Tăng độ chính xác ước lượng \\
\hline
Forecasting & Dự báo tiến độ tốt hơn \\
\hline
Risk Detection & Phát hiện rủi ro sớm \\
\hline
Resource Allocation & Tối ưu phân công nhân sự \\
\hline
Documentation & Tự động sinh tài liệu \\
\hline
Dashboard & Hỗ trợ ra quyết định thời gian thực \\
\hline
\end{tabular}
\label{tab:ai_benefits}
\end{table}

%%%%%%%%%%%%%%%%%%%%%%%%%%%%%%%%%%%%%%%%%%%%%%%%%%%%%%%%%%%%%%

\subsection{Hạn chế của AI}

Bên cạnh những ưu điểm, AI vẫn tồn tại nhiều hạn chế cần được xem xét.

\begin{itemize}
\item Có thể sinh thông tin không chính xác (Hallucination).
\item Phụ thuộc vào chất lượng dữ liệu đầu vào.
\item Khó giải thích hoàn toàn quá trình suy luận.
\item Không thay thế được kinh nghiệm của Project Manager.
\item Có nguy cơ tạo ra các quyết định thiên lệch nếu dữ liệu huấn luyện không cân bằng.
\end{itemize}

Do đó, trong đề tài này AI được thiết kế theo định hướng Human-Centric AI, nghĩa là mọi quyết định quan trọng đều có sự giám sát hoặc phê duyệt của con người.

\subsection{Large Language Models trong quản lý dự án phần mềm}
\label{sec:llm_project_management}
%%%%%%%%%%%%%%%%%%%%%%%%%%%%%%%%%%%%%%%%%%%%%%%%%%%%%%%%%%%%%%

Sự xuất hiện của Large Language Models (LLMs) đã tạo ra bước ngoặt quan trọng trong lĩnh vực Software Engineering. Khác với các hệ thống Machine Learning truyền thống vốn chỉ giải quyết các bài toán dự đoán hoặc phân loại, LLM có khả năng hiểu ngữ cảnh, suy luận trên văn bản tự nhiên và sinh nội dung với chất lượng gần tương đương con người.

Trong các dự án phần mềm, phần lớn dữ liệu đều tồn tại dưới dạng văn bản phi cấu trúc như:

\begin{itemize}
\item Software Requirement Specification (SRS);
\item Product Vision;
\item User Story;
\item Use Case;
\item Meeting Minutes;
\item Sprint Report;
\item Source Code;
\item Documentation.
\end{itemize}

Đây chính là loại dữ liệu mà LLM được thiết kế để xử lý hiệu quả.

%%%%%%%%%%%%%%%%%%%%%%%%%%%%%%%%%%%%%%%%%%%%%%%%%%%%%%%%%%%%%%

\subsubsection{Khả năng hiểu yêu cầu}

Một trong những khó khăn lớn nhất của Product Owner là chuyển đổi yêu cầu của khách hàng thành các User Story rõ ràng và đầy đủ.

Ví dụ, khách hàng có thể đưa ra yêu cầu:

\begin{quote}
"Hệ thống cần quản lý dự án và hỗ trợ AI."
\end{quote}

Thông tin trên chưa trả lời được nhiều câu hỏi quan trọng như:

\begin{itemize}
\item Quản lý những loại dự án nào?
\item Có bao nhiêu vai trò người dùng?
\item AI hỗ trợ ở những giai đoạn nào?
\item Có cần Dashboard hay không?
\item Có hỗ trợ Sprint không?
\end{itemize}

LLM có khả năng:

\begin{enumerate}
\item phát hiện yêu cầu còn thiếu;
\item đặt câu hỏi làm rõ;
\item đề xuất User Story;
\item sinh Acceptance Criteria;
\item phát hiện mâu thuẫn giữa các Requirement.
\end{enumerate}

%%%%%%%%%%%%%%%%%%%%%%%%%%%%%%%%%%%%%%%%%%%%%%%%%%%%%%%%%%%%%%

\subsubsection{Sinh User Story}

Sau khi Requirement được xác nhận, AI có thể chuyển đổi sang User Story theo mẫu chuẩn:

\begin{quote}
As a \textit{Role}

I want \textit{Goal}

So that \textit{Business Value}
\end{quote}

Ví dụ

\begin{quote}
As a Product Owner,

I want AI to decompose User Story into Tasks,

So that project planning becomes faster.
\end{quote}

Điều này giúp chuẩn hóa toàn bộ Product Backlog.

%%%%%%%%%%%%%%%%%%%%%%%%%%%%%%%%%%%%%%%%%%%%%%%%%%%%%%%%%%%%%%

\subsubsection{Sinh Acceptance Criteria}

Từ User Story, AI tiếp tục sinh Acceptance Criteria theo chuẩn Gherkin.

Ví dụ:

\begin{verbatim}
Given Product Owner nhập User Story
When AI phân tích
Then sinh danh sách Task
And Product Owner có thể chỉnh sửa
\end{verbatim}

Việc tự động sinh Acceptance Criteria giúp tăng tính đầy đủ của tài liệu SRS và giảm công sức của Product Owner.

%%%%%%%%%%%%%%%%%%%%%%%%%%%%%%%%%%%%%%%%%%%%%%%%%%%%%%%%%%%%%%

\subsubsection{Sinh Task}

Sau khi User Story được tạo, AI tiếp tục phân rã thành các Task kỹ thuật.

Ví dụ

\begin{table}[H]
\centering
\caption{Ví dụ phân rã User Story}
\begin{tabular}{|p{7cm}|p{7cm}|}
\hline
User Story & Task \\
\hline
AI phân tích Requirement & Thiết kế Prompt \\
& Gọi LLM API \\
& Phân tích JSON \\
& Validate Output \\
& Lưu Database \\
& Hiển thị UI \\
\hline
\end{tabular}
\end{table}

Khả năng này giúp giảm đáng kể thời gian Sprint Planning.

%%%%%%%%%%%%%%%%%%%%%%%%%%%%%%%%%%%%%%%%%%%%%%%%%%%%%%%%%%%%%%

\subsection{Prompt Engineering}

Prompt Engineering là quá trình thiết kế Prompt nhằm hướng dẫn LLM tạo ra kết quả mong muốn.

Một Prompt hiệu quả thường gồm năm thành phần.

\begin{enumerate}

\item Role

Xác định vai trò của AI.

Ví dụ:

\begin{quote}
You are an experienced Scrum Master.
\end{quote}

\item Context

Cung cấp ngữ cảnh.

\item Task

Mô tả nhiệm vụ.

\item Constraints

Các ràng buộc.

\item Output Format

Định dạng đầu ra.

\end{enumerate}

Trong hệ thống đề xuất, mọi AI Agent đều sử dụng Prompt Engineering thay vì Prompt ngẫu nhiên nhằm đảm bảo đầu ra ổn định.

%%%%%%%%%%%%%%%%%%%%%%%%%%%%%%%%%%%%%%%%%%%%%%%%%%%%%%%%%%%%%%

\subsection{Structured Output}

Một yêu cầu quan trọng khi tích hợp AI vào hệ thống phần mềm là đầu ra phải có cấu trúc.

Thay vì trả về văn bản tự do, AI được yêu cầu sinh JSON.

Ví dụ:

\begin{verbatim}
{
  "storyPoint":8,
  "reason":"Complex authentication",
  "risk":"Medium"
}
\end{verbatim}

Structured Output giúp Backend dễ dàng:

\begin{itemize}
\item parse dữ liệu;
\item validate;
\item lưu Database;
\item hiển thị Dashboard.
\end{itemize}

%%%%%%%%%%%%%%%%%%%%%%%%%%%%%%%%%%%%%%%%%%%%%%%%%%%%%%%%%%%%%%

\subsection{Function Calling}

Các LLM hiện đại hỗ trợ Function Calling.

Thay vì trả lời trực tiếp, AI có thể yêu cầu hệ thống thực hiện một hàm.

Ví dụ

\begin{verbatim}
estimateStoryPoint()
createTask()
searchProject()
forecastSprint()
\end{verbatim}

Điều này cho phép AI tương tác trực tiếp với Backend thay vì chỉ đóng vai trò Chatbot.

%%%%%%%%%%%%%%%%%%%%%%%%%%%%%%%%%%%%%%%%%%%%%%%%%%%%%%%%%%%%%%

\subsection{Retrieval-Augmented Generation}

Nếu chỉ dựa trên dữ liệu đã học, LLM có thể tạo ra thông tin không chính xác.

Để khắc phục vấn đề này, nhiều hệ thống hiện đại áp dụng Retrieval-Augmented Generation (RAG).

Quy trình tổng quát gồm:

\begin{center}
User Query $\rightarrow$ Retriever $\rightarrow$ Knowledge Base $\rightarrow$ Relevant Documents $\rightarrow$ LLM $\rightarrow$ Answer
\end{center}

Trong hệ thống được đề xuất, nguồn dữ liệu truy xuất có thể bao gồm:

\begin{itemize}
\item Product Vision;
\item SRS;
\item User Story;
\item Sprint History;
\item Project Database;
\item Internal Documentation.
\end{itemize}

Nhờ đó AI luôn trả lời dựa trên dữ liệu của dự án thay vì chỉ dựa trên tri thức chung của mô hình.

%%%%%%%%%%%%%%%%%%%%%%%%%%%%%%%%%%%%%%%%%%%%%%%%%%%%%%%%%%%%%%

\subsection{Hallucination}

Hallucination là hiện tượng AI sinh ra thông tin không tồn tại hoặc không chính xác.

Ví dụ:

\begin{itemize}
\item tạo Requirement không có;
\item sinh API sai;
\item sinh Story Point vô lý;
\item viện dẫn tài liệu không tồn tại.
\end{itemize}

Để giảm Hallucination, hệ thống áp dụng:

\begin{enumerate}
\item Structured Prompt.
\item RAG.
\item JSON Validation.
\item Human Review.
\item Confidence Score.
\end{enumerate}

Đây cũng là lý do đề tài lựa chọn mô hình Human-Centered AI thay vì để AI tự động đưa ra mọi quyết định.

%%%%%%%%%%%%%%%%%%%%%%%%%%%%%%%%%%%%%%%%%%%%%%%%%%%%%%%%%%%%%%

\subsection{Liên hệ với hệ thống đề xuất}

Trong luận văn này, LLM là thành phần trung tâm của toàn bộ nền tảng AI Project Management.

Thay vì xây dựng một Chatbot đơn lẻ, hệ thống sử dụng nhiều AI Agent chuyên biệt, mỗi Agent đảm nhận một nghiệp vụ riêng:

\begin{table}[H]
\centering
\caption{Các AI Agent trong hệ thống}
\begin{tabular}{|p{5cm}|p{9cm}|}
\hline
AI Agent & Nhiệm vụ \\
\hline
Requirement Agent & Phân tích Requirement \\
\hline
Planning Agent & Sinh User Story và Task \\
\hline
Story Point Agent & Ước lượng Story Point \\
\hline
Forecast Agent & Dự báo tiến độ \\
\hline
Risk Agent & Phát hiện rủi ro \\
\hline
Assignment Agent & Đề xuất phân công \\
\hline
Documentation Agent & Sinh tài liệu dự án \\
\hline
Dashboard Agent & Tổng hợp KPI \\
\hline
\end{tabular}
\label{tab:ai_agents}
\end{table}

Việc chia nhỏ thành nhiều AI Agent giúp tăng khả năng mở rộng, giảm độ phức tạp của Prompt và cho phép từng Agent được tối ưu hóa cho một nghiệp vụ cụ thể. Đây cũng là nền tảng để Chương~\ref{chap:system_design} và Chương~\ref{chap:ai_architecture} trình bày kiến trúc Multi-Agent AI của hệ thống.

\section{Hệ thống Multi-Agent AI trong quản lý dự án}
\label{sec:multiagent}
\subsection{Giới thiệu Multi-Agent System}

Trong giai đoạn đầu của các ứng dụng Generative AI, phần lớn hệ thống chỉ sử dụng một Large Language Model (LLM) duy nhất để xử lý toàn bộ yêu cầu của người dùng. Mặc dù cách tiếp cận này đơn giản trong triển khai, nhưng khi bài toán ngày càng phức tạp, một mô hình đơn lẻ khó có thể đạt hiệu quả cao ở tất cả các nghiệp vụ.

Đối với quản lý dự án phần mềm, hệ thống cần đồng thời thực hiện nhiều nhiệm vụ khác nhau như:

\begin{itemize}
\item Phân tích Requirement;
\item Sinh User Story;
\item Phân rã Task;
\item Ước lượng Story Point;
\item Dự báo Sprint;
\item Phân tích rủi ro;
\item Sinh tài liệu;
\item Hỗ trợ Dashboard.
\end{itemize}

Mỗi nhiệm vụ có dữ liệu đầu vào, mục tiêu và phương pháp xử lý khác nhau. Vì vậy, xu hướng hiện nay là xây dựng các hệ thống Multi-Agent AI thay vì chỉ sử dụng một mô hình duy nhất.

Multi-Agent System (MAS) là hệ thống gồm nhiều tác tử thông minh (Agent), mỗi Agent chịu trách nhiệm giải quyết một nhóm nhiệm vụ chuyên biệt và phối hợp với các Agent khác nhằm hoàn thành mục tiêu chung.

Trong đề tài này, kiến trúc Multi-Agent được lựa chọn nhằm tăng khả năng mở rộng, giảm độ phức tạp của Prompt và cải thiện chất lượng kết quả sinh bởi AI.

%%%%%%%%%%%%%%%%%%%%%%%%%%%%%%%%%%%%%%%%%%%%%%%%%%%%%%%%%%%%%%

\subsection{Đặc điểm của AI Agent}

Một AI Agent hiện đại thường bao gồm năm thành phần chính như Hình \ref{fig:agent_architecture}.

\begin{figure}[H]
\centering
\fbox{
\begin{minipage}{0.9\textwidth}
\centering
Input $\rightarrow$ Memory $\rightarrow$ Reasoning $\rightarrow$ Tool Calling $\rightarrow$ Output
\end{minipage}
}
\caption{Kiến trúc tổng quát của AI Agent}
\label{fig:agent_architecture}
\end{figure}

Trong đó:

\begin{itemize}

\item \textbf{Input}: tiếp nhận yêu cầu từ người dùng hoặc Agent khác.

\item \textbf{Memory}: lưu trữ ngữ cảnh, lịch sử hội thoại và dữ liệu dự án.

\item \textbf{Reasoning}: sử dụng LLM để suy luận và lập kế hoạch.

\item \textbf{Tool Calling}: gọi API hoặc truy cập cơ sở dữ liệu khi cần thiết.

\item \textbf{Output}: trả về kết quả dưới dạng văn bản hoặc dữ liệu có cấu trúc.

\end{itemize}

Kiến trúc này cho phép Agent không chỉ trả lời câu hỏi mà còn tương tác trực tiếp với các thành phần khác của hệ thống.

%%%%%%%%%%%%%%%%%%%%%%%%%%%%%%%%%%%%%%%%%%%%%%%%%%%%%%%%%%%%%%

\subsection{Kiến trúc Multi-Agent của hệ thống}

Dựa trên Product Vision và SRS, hệ thống AI Project Management được thiết kế gồm nhiều Agent độc lập, mỗi Agent đảm nhiệm một vai trò riêng trong quy trình quản lý dự án.

\begin{table}[H]
\centering
\caption{Các AI Agent trong hệ thống đề xuất}
\begin{tabular}{|p{4.5cm}|p{10cm}|}
\hline
Agent & Chức năng chính \\
\hline
Requirement Agent & Phân tích yêu cầu, phát hiện mâu thuẫn, đề xuất chỉnh sửa Requirement. \\
\hline
Planning Agent & Chuyển Requirement thành User Story, Epic và Task. \\
\hline
Story Point Agent & Ước lượng Story Point dựa trên độ phức tạp và dữ liệu lịch sử. \\
\hline
Assignment Agent & Đề xuất phân công thành viên phù hợp cho từng Task. \\
\hline
Forecast Agent & Dự báo tiến độ Sprint và thời gian hoàn thành dự án. \\
\hline
Risk Agent & Phân tích dữ liệu để phát hiện nguy cơ chậm tiến độ hoặc quá tải nguồn lực. \\
\hline
Documentation Agent & Tự động sinh tài liệu SRS, Sprint Report, Meeting Minutes và Release Note. \\
\hline
Dashboard Agent & Tổng hợp KPI và hiển thị thông tin trực quan trên Dashboard. \\
\hline
\end{tabular}
\label{tab:agent_roles}
\end{table}

Việc phân tách chức năng thành nhiều Agent giúp mỗi Agent chỉ tập trung vào một bài toán cụ thể, từ đó giảm kích thước Prompt và cải thiện độ chính xác của kết quả.

%%%%%%%%%%%%%%%%%%%%%%%%%%%%%%%%%%%%%%%%%%%%%%%%%%%%%%%%%%%%%%

\subsection{Luồng phối hợp giữa các Agent}

Các Agent không hoạt động độc lập mà trao đổi dữ liệu với nhau theo một quy trình thống nhất.

\begin{figure}[H]
\centering
\fbox{
\begin{minipage}{0.95\textwidth}
\centering
Requirement $\rightarrow$ Requirement Agent $\rightarrow$ Planning Agent $\rightarrow$ Story Point Agent $\rightarrow$ Assignment Agent $\rightarrow$ Forecast Agent $\rightarrow$ Risk Agent $\rightarrow$ Documentation Agent $\rightarrow$ Dashboard
\end{minipage}
}
\caption{Luồng xử lý của hệ thống Multi-Agent}
\label{fig:multi_agent_flow}
\end{figure}

Sau khi Product Owner nhập Requirement, Requirement Agent sẽ kiểm tra tính đầy đủ và nhất quán của yêu cầu. Planning Agent tiếp tục sinh User Story và phân rã thành các Task kỹ thuật. Story Point Agent ước lượng khối lượng công việc, sau đó Assignment Agent đề xuất phân công nhân sự. Trong quá trình Sprint diễn ra, Forecast Agent và Risk Agent liên tục theo dõi dữ liệu để dự báo tiến độ và phát hiện rủi ro. Cuối cùng, Documentation Agent tổng hợp báo cáo và Dashboard Agent cập nhật các chỉ số quản lý theo thời gian thực.

%%%%%%%%%%%%%%%%%%%%%%%%%%%%%%%%%%%%%%%%%%%%%%%%%%%%%%%%%%%%%%

\subsection{Bộ nhớ của AI Agent}

Để duy trì ngữ cảnh trong suốt vòng đời dự án, mỗi Agent cần sử dụng cơ chế bộ nhớ (Memory). Bộ nhớ có thể được chia thành ba loại:

\begin{itemize}
\item \textbf{Short-term Memory}: lưu trữ thông tin của phiên làm việc hiện tại.
\item \textbf{Long-term Memory}: lưu trữ dữ liệu lịch sử của dự án như Sprint, Velocity và các quyết định trước đây.
\item \textbf{Knowledge Base}: chứa Product Vision, SRS, Use Case, UML, README và các tài liệu nội bộ.
\end{itemize}

Thông qua cơ chế Retrieval-Augmented Generation (RAG), Agent có thể truy xuất các tài liệu liên quan trước khi đưa ra câu trả lời, từ đó giảm nguy cơ Hallucination và tăng tính nhất quán.

%%%%%%%%%%%%%%%%%%%%%%%%%%%%%%%%%%%%%%%%%%%%%%%%%%%%%%%%%%%%%%

\subsection{Ưu điểm của kiến trúc Multi-Agent}

So với việc sử dụng một LLM duy nhất, kiến trúc Multi-Agent mang lại nhiều lợi ích:

\begin{itemize}
\item Chuyên môn hóa từng Agent cho một nghiệp vụ cụ thể.
\item Dễ dàng mở rộng và bảo trì hệ thống.
\item Giảm kích thước Prompt cho từng Agent.
\item Tăng khả năng tái sử dụng.
\item Giảm nguy cơ Hallucination nhờ phân chia nhiệm vụ rõ ràng.
\item Cho phép tích hợp nhiều mô hình AI khác nhau trong cùng một hệ thống.
\end{itemize}

Những ưu điểm này đặc biệt phù hợp với các nền tảng quản lý dự án quy mô lớn, nơi nhiều nghiệp vụ cần được xử lý đồng thời.

%%%%%%%%%%%%%%%%%%%%%%%%%%%%%%%%%%%%%%%%%%%%%%%%%%%%%%%%%%%%%%

\subsection{Liên hệ với hệ thống AI Project Management}

Kiến trúc Multi-Agent là nền tảng cốt lõi của hệ thống được đề xuất trong luận văn. Thay vì xây dựng một chatbot tổng quát, hệ thống tổ chức các Agent theo đúng quy trình Scrum và PMBOK. Mỗi Agent đảm nhiệm một bước trong vòng đời quản lý dự án, đồng thời phối hợp thông qua các API nội bộ và cơ sở dữ liệu chung.

Cách tiếp cận này không chỉ giúp nâng cao chất lượng kết quả mà còn tạo điều kiện thuận lợi cho việc mở rộng hệ thống trong tương lai, ví dụ bổ sung AI Code Review Agent, AI Test Generation Agent hoặc AI DevOps Agent mà không cần thay đổi kiến trúc tổng thể.

\section{Human-Centered AI trong quản lý dự án phần mềm}
\label{sec:human_centered_ai}
\subsection{Khái niệm Human-Centered AI}

Sự phát triển nhanh chóng của các hệ thống trí tuệ nhân tạo đã mang lại khả năng tự động hóa nhiều hoạt động trong quản lý dự án. Tuy nhiên, AI vẫn tồn tại nhiều hạn chế như tạo ra thông tin không chính xác (Hallucination), thiếu khả năng giải thích, khó đánh giá các yếu tố mang tính xã hội và không thể thay thế hoàn toàn kinh nghiệm của nhà quản lý dự án.

Vì vậy, xu hướng hiện nay không hướng tới việc thay thế con người mà tập trung xây dựng các hệ thống \textit{Human-Centered AI} (HCAI). Theo quan điểm này, AI được xem là một công cụ hỗ trợ giúp con người nâng cao năng suất và chất lượng ra quyết định thay vì tự động thực hiện toàn bộ quy trình.

Human-Centered AI nhấn mạnh ba mục tiêu chính:

\begin{itemize}
\item AI phải hỗ trợ con người.
\item Con người luôn giữ quyền kiểm soát cuối cùng.
\item Mọi quyết định quan trọng đều có thể được giải thích và kiểm chứng.
\end{itemize}

Trong lĩnh vực quản lý dự án phần mềm, cách tiếp cận này đặc biệt quan trọng vì nhiều quyết định liên quan đến phạm vi dự án, chi phí, tiến độ và nguồn lực có ảnh hưởng lớn tới doanh nghiệp.

%%%%%%%%%%%%%%%%%%%%%%%%%%%%%%%%%%%%%%%%%%%%%%%%%%%%%%%%%%%%%%

\subsection{Các mức độ tự động hóa của AI}

Theo nhiều nghiên cứu về Human-AI Collaboration, mức độ tham gia của AI có thể chia thành ba mô hình chính.

\subsubsection{Human-in-the-Loop}

Ở mô hình Human-in-the-Loop (HITL), AI chỉ đưa ra gợi ý hoặc dự đoán, trong khi quyết định cuối cùng vẫn do con người thực hiện.

Ví dụ:

\begin{itemize}
\item AI đề xuất Story Point.
\item Product Owner xác nhận hoặc chỉnh sửa.
\end{itemize}

Đây là mô hình được áp dụng cho các nghiệp vụ có mức độ ảnh hưởng cao.

%%%%%%%%%%%%%%%%%%%%%%%%%%%%%%%%%%%%%%%%%%%%%%%%%%%%%%%%%%%%%%

\subsubsection{Human-on-the-Loop}

Trong Human-on-the-Loop (HOTL), AI có thể tự động thực hiện các tác vụ thông thường nhưng con người vẫn theo dõi và có thể can thiệp khi cần thiết.

Ví dụ:

\begin{itemize}
\item AI tự động sinh Sprint Report.
\item Scrum Master kiểm tra trước khi phát hành.
\end{itemize}

%%%%%%%%%%%%%%%%%%%%%%%%%%%%%%%%%%%%%%%%%%%%%%%%%%%%%%%%%%%%%%

\subsubsection{Human-out-of-the-Loop}

Human-out-of-the-Loop (HOOTL) là mô hình AI hoàn toàn tự động mà không cần sự giám sát của con người.

Mặc dù phù hợp với một số hệ thống tự động hóa, mô hình này không được lựa chọn trong đề tài do các quyết định quản lý dự án thường liên quan đến nhiều yếu tố nghiệp vụ và rủi ro.

%%%%%%%%%%%%%%%%%%%%%%%%%%%%%%%%%%%%%%%%%%%%%%%%%%%%%%%%%%%%%%

\subsection{Giải thích quyết định của AI}

Một yêu cầu quan trọng đối với AI trong quản lý dự án là khả năng giải thích (Explainability).

Thay vì chỉ trả về kết quả:

\begin{quote}
Story Point = 13
\end{quote}

AI cần giải thích lý do, ví dụ:

\begin{itemize}
\item User Story có nhiều chức năng.
\item Liên quan đến xác thực và phân quyền.
\item Yêu cầu tích hợp hệ thống ngoài.
\item Có mức độ không chắc chắn cao.
\end{itemize}

Khả năng giải thích giúp Product Owner và Scrum Master dễ dàng đánh giá tính hợp lý của các đề xuất và tăng mức độ tin cậy đối với hệ thống.

%%%%%%%%%%%%%%%%%%%%%%%%%%%%%%%%%%%%%%%%%%%%%%%%%%%%%%%%%%%%%%

\subsection{Vai trò của con người trong hệ thống đề xuất}

Dựa trên Product Vision và các yêu cầu chức năng, hệ thống AI Project Management được thiết kế theo nguyên tắc ``AI đề xuất -- Con người quyết định''. Mọi AI Agent đều hoạt động dưới sự giám sát của người dùng.

Bảng \ref{tab:human_ai_roles} trình bày mối quan hệ giữa AI và người dùng trong từng nghiệp vụ.

\begin{table}[H]
\centering
\caption{Vai trò của AI và con người trong hệ thống}
\begin{tabular}{|p{5cm}|p{4cm}|p{5cm}|}
\hline
Nghiệp vụ & AI & Con người \\
\hline
Phân tích Requirement & Đề xuất & Phê duyệt \\
\hline
Sinh User Story & Đề xuất & Chỉnh sửa \\
\hline
Sinh Task & Tự động & Kiểm tra \\
\hline
Story Point & Đề xuất & Xác nhận \\
\hline
Forecast Sprint & Dự báo & Ra quyết định \\
\hline
Risk Analysis & Cảnh báo & Lập kế hoạch ứng phó \\
\hline
Documentation & Sinh tự động & Kiểm duyệt \\
\hline
Assignment & Gợi ý & Quyết định cuối cùng \\
\hline
\end{tabular}
\label{tab:human_ai_roles}
\end{table}

%%%%%%%%%%%%%%%%%%%%%%%%%%%%%%%%%%%%%%%%%%%%%%%%%%%%%%%%%%%%%%

\subsection{Khoảng trống nghiên cứu}

Qua khảo sát các nghiên cứu hiện có và các nền tảng quản lý dự án phổ biến như Jira, Azure DevOps và ClickUp, có thể nhận thấy phần lớn các hệ thống AI hiện nay chỉ tập trung vào một số tác vụ riêng lẻ như sinh mã nguồn, hỗ trợ viết tài liệu hoặc trả lời câu hỏi.

Các hạn chế chính bao gồm:

\begin{itemize}
\item Chưa hỗ trợ toàn bộ vòng đời quản lý dự án.
\item Thiếu sự phối hợp giữa nhiều AI Agent.
\item Chưa tích hợp đồng thời Scrum và PMBOK.
\item Ít sử dụng dữ liệu lịch sử của chính dự án để hỗ trợ ra quyết định.
\item Chưa áp dụng đầy đủ mô hình Human-Centered AI.
\end{itemize}

Đây chính là khoảng trống mà luận văn hướng tới giải quyết.

%%%%%%%%%%%%%%%%%%%%%%%%%%%%%%%%%%%%%%%%%%%%%%%%%%%%%%%%%%%%%%

\subsection{Định hướng của hệ thống đề xuất}

Dựa trên các phân tích ở trên, luận văn đề xuất một nền tảng AI Project Management có các đặc điểm sau:

\begin{itemize}
\item Kiến trúc Multi-Agent AI.
\item Tích hợp Large Language Models.
\item Hỗ trợ toàn bộ quy trình Scrum.
\item Áp dụng các nguyên tắc của PMBOK.
\item Sử dụng Retrieval-Augmented Generation để khai thác tri thức nội bộ.
\item Thiết kế theo mô hình Human-Centered AI.
\item Cho phép người dùng kiểm soát và phê duyệt các quyết định quan trọng.
\end{itemize}

Hệ thống không hướng tới việc thay thế Project Manager mà đóng vai trò là một nền tảng hỗ trợ thông minh, giúp giảm thời gian xử lý nghiệp vụ, nâng cao chất lượng lập kế hoạch và cải thiện hiệu quả quản lý dự án phần mềm.

%%%%%%%%%%%%%%%%%%%%%%%%%%%%%%%%%%%%%%%%%%%%%%%%%%%%%%%%%%%%%%
\section{Tổng kết chương}
%%%%%%%%%%%%%%%%%%%%%%%%%%%%%%%%%%%%%%%%%%%%%%%%%%%%%%%%%%%%%%

Trong chương này, luận văn đã trình bày các cơ sở lý thuyết làm nền tảng cho việc xây dựng hệ thống AI Project Management. Nội dung bao gồm tổng quan về quản lý dự án phần mềm, phương pháp Agile, Scrum Framework, PMBOK, trí tuệ nhân tạo, Large Language Models, Prompt Engineering, Retrieval-Augmented Generation, kiến trúc Multi-Agent và mô hình Human-Centered AI.

Các phân tích cho thấy sự kết hợp giữa Scrum, PMBOK và các mô hình AI hiện đại có khả năng hỗ trợ hiệu quả nhiều hoạt động trong quản lý dự án như phân tích yêu cầu, lập kế hoạch Sprint, ước lượng Story Point, dự báo tiến độ, phát hiện rủi ro và sinh tài liệu.

Những cơ sở lý thuyết này sẽ được sử dụng trong Chương~\ref{chap:system_analysis} để phân tích yêu cầu nghiệp vụ và trong Chương~\ref{chap:system_design} để thiết kế kiến trúc của hệ thống AI Project Management được đề xuất.

%%%%%%%%%%%%%%%%%%%%%%%%%%%%%%%%%%%%%%%%%%%%%%%%%%%%%%%%%%%%%%
% CHƯƠNG 3: TỔNG QUAN NGHIÊN CỨU
%%%%%%%%%%%%%%%%%%%%%%%%%%%%%%%%%%%%%%%%%%%%%%%%%%%%%%%%%%%%%%
\chapter{TỔNG QUAN NGHIÊN CỨU}
\label{chap:related_work}

\section{Giới thiệu chương}

Chương này trình bày tổng quan các nghiên cứu liên quan đến việc ứng dụng trí tuệ nhân tạo trong quản lý dự án phần mềm, đặc biệt là các phương pháp phân tích yêu cầu, ước lượng công việc, phát hiện rủi ro và tự động hóa quy trình Agile.

\section{Các nghiên cứu về AI trong phân tích yêu cầu}

\subsection{Phương pháp phân tích yêu cầu tự động}

Nhiều nghiên cứu đã chỉ ra rằng các kỹ thuật xử lý ngôn ngữ tự nhiên (NLP) có thể hỗ trợ hiệu quả trong việc phân tích và làm rõ yêu cầu phần mềm \cite{bachari2025}. Các mô hình ngôn ngữ lớn hiện đại có khả năng phát hiện các yêu cầu mơ hồ, trùng lặp hoặc mâu thuẫn từ tài liệu đặc tả.

\subsection{Sinh User Story tự động}

Một số nghiên cứu gần đây đã chứng minh khả năng của LLM trong việc sinh User Story từ mô tả yêu cầu, giúp giảm thời gian và công sức của Product Owner \cite{vaswani2017attention}.

\section{Các nghiên cứu về ước lượng phần mềm}

\subsection{Mô hình học máy trong ước lượng}

Các nghiên cứu đã áp dụng Random Forest, Support Vector Machine và các mô hình học sâu để ước lượng khối lượng công việc dựa trên dữ liệu lịch sử của dự án \cite{russell2021}. Kết quả cho thấy các mô hình học máy có thể đạt độ chính xác tương đương hoặc cao hơn so với ước lượng thủ công.

\section{Các nghiên cứu về Multi-Agent AI}

\subsection{Kiến trúc Multi-Agent trong Software Engineering}

Xu hướng xây dựng các hệ thống Multi-Agent AI đang được áp dụng rộng rãi trong nhiều lĩnh vực, bao gồm cả kỹ thuật phần mềm. Mỗi Agent đảm nhiệm một nhiệm vụ chuyên biệt và phối hợp với các Agent khác để giải quyết các bài toán phức tạp \cite{russell2021}.

\section{Khoảng trống nghiên cứu}

Qua khảo sát các nghiên cứu hiện có, có thể nhận thấy rằng mặc dù AI đã được ứng dụng trong nhiều lĩnh vực của quản lý dự án phần mềm, phần lớn các nghiên cứu vẫn tập trung vào từng bài toán riêng lẻ. Chưa có một nền tảng tích hợp toàn diện hỗ trợ toàn bộ vòng đời phát triển phần mềm theo Agile, từ phân tích yêu cầu đến đánh giá Sprint. Đây chính là khoảng trống mà luận văn hướng tới giải quyết.

%%%%%%%%%%%%%%%%%%%%%%%%%%%%%%%%%%%%%%%%%%%%%%%%%%%%%%%%%%%%%%
%%%%%%%%%%%%%%%%%%%%%%%%%%%%%%%%%%%%%%%%%%%%%%%%%%%%%%%%%%%%%%
\chapter{PHÂN TÍCH BÀI TOÁN}
\label{chap:problem_analysis}
%%%%%%%%%%%%%%%%%%%%%%%%%%%%%%%%%%%%%%%%%%%%%%%%%%%%%%%%%%%%%%

\section{Giới thiệu chương}

Sau khi trình bày các cơ sở lý thuyết về quản lý dự án phần mềm, Agile, Scrum và trí tuệ nhân tạo ở Chương 2 và tổng quan các nghiên cứu liên quan ở Chương 3, chương này tập trung phân tích chi tiết bài toán thực tế mà đề tài hướng tới giải quyết.

Mục tiêu của chương bao gồm:

\begin{itemize}
\item Phân tích hiện trạng quản lý dự án phần mềm trong các doanh nghiệp.
\item Xác định các vấn đề và thách thức trong quy trình hiện tại.
\item Phân tích nguyên nhân gốc rễ của các vấn đề.
\item Xác định các yêu cầu và ràng buộc đối với giải pháp.
\item Đề xuất định hướng giải pháp tổng thể.
\end{itemize}

Nội dung của chương sẽ là cơ sở để xác định các yêu cầu chức năng và phi chức năng cho hệ thống ở Chương 5.

%%%%%%%%%%%%%%%%%%%%%%%%%%%%%%%%%%%%%%%%%%%%%%%%%%%%%%%%%%%%%%
\section{Hiện trạng quản lý dự án phần mềm}
%%%%%%%%%%%%%%%%%%%%%%%%%%%%%%%%%%%%%%%%%%%%%%%%%%%%%%%%%%%%%%

\subsection{Tổng quan về quy trình quản lý dự án hiện tại}

Trong bối cảnh phát triển phần mềm hiện đại, hầu hết các doanh nghiệp công nghệ đều áp dụng phương pháp Agile, đặc biệt là Scrum, để quản lý các dự án phần mềm. Quy trình quản lý dự án điển hình bao gồm các giai đoạn sau:

\begin{figure}[H]
\centering
\fbox{
\begin{minipage}{0.95\textwidth}
\centering
Thu thập yêu cầu $\rightarrow$ Phân tích yêu cầu $\rightarrow$ Lập kế hoạch Sprint $\rightarrow$ Phân công Task $\rightarrow$ Theo dõi tiến độ $\rightarrow$ Đánh giá Sprint $\rightarrow$ Phát hành
\end{minipage}
}
\caption{Quy trình quản lý dự án Agile điển hình}
\label{fig:current_process}
\end{figure}

Mặc dù quy trình này đã được chuẩn hóa, việc thực hiện vẫn phụ thuộc rất nhiều vào kinh nghiệm và năng lực của các thành viên trong nhóm, đặc biệt là Product Owner, Scrum Master và Project Manager.

\subsection{Các công cụ quản lý dự án phổ biến hiện nay}

Hiện nay, có nhiều công cụ hỗ trợ quản lý dự án phần mềm được sử dụng rộng rãi. Bảng \ref{tab:current_tools} so sánh các công cụ phổ biến:

\begin{table}[H]
\centering
\caption{So sánh các công cụ quản lý dự án phổ biến}
\begin{tabular}{|p{3.5cm}|p{3cm}|p{3cm}|p{3cm}|}
\hline
\textbf{Tính năng} & \textbf{Jira} & \textbf{Trello} & \textbf{Azure DevOps} \\
\hline
Quản lý Sprint & \checkmark & $\triangle$ & \checkmark \\
\hline
Quản lý Backlog & \checkmark & $\triangle$ & \checkmark \\
\hline
Kanban Board & \checkmark & \checkmark & \checkmark \\
\hline
Báo cáo & \checkmark & $\triangle$ & \checkmark \\
\hline
Tích hợp AI & $\triangle$ & $\times$ & $\triangle$ \\
\hline
Phân tích yêu cầu AI & $\times$ & $\times$ & $\times$ \\
\hline
Tự động sinh Task & $\times$ & $\times$ & $\times$ \\
\hline
Dự báo tiến độ AI & $\triangle$ & $\times$ & $\triangle$ \\
\hline
\end{tabular}
\label{tab:current_tools}
\end{table}

Qua bảng so sánh, có thể thấy rằng các công cụ hiện nay chủ yếu tập trung vào việc lưu trữ và theo dõi công việc, trong khi khả năng hỗ trợ ra quyết định thông minh vẫn còn rất hạn chế.

%%%%%%%%%%%%%%%%%%%%%%%%%%%%%%%%%%%%%%%%%%%%%%%%%%%%%%%%%%%%%%
\section{Các vấn đề và thách thức}
%%%%%%%%%%%%%%%%%%%%%%%%%%%%%%%%%%%%%%%%%%%%%%%%%%%%%%%%%%%%%%

\subsection{Vấn đề 1: Phân tích yêu cầu thiếu hiệu quả}

\subsubsection{Mô tả vấn đề}

Một trong những vấn đề lớn nhất trong quản lý dự án phần mềm là việc phân tích và làm rõ yêu cầu từ khách hàng. Các báo cáo yêu cầu thường được viết dưới dạng văn bản tự nhiên, dẫn đến nhiều vấn đề:

\begin{itemize}
\item Yêu cầu mơ hồ, thiếu thông tin chi tiết.
\item Yêu cầu trùng lặp hoặc mâu thuẫn.
\item Yêu cầu không khả thi về mặt kỹ thuật.
\item Thiếu các tiêu chí chấp nhận (Acceptance Criteria) rõ ràng.
\end{itemize}

\subsubsection{Tác động}

\begin{table}[H]
\centering
\caption{Tác động của việc phân tích yêu cầu kém hiệu quả}
\begin{tabular}{|p{5cm}|p{8cm}|}
\hline
\textbf{Hệ quả} & \textbf{Mô tả} \\
\hline
Chậm tiến độ & Thời gian làm rõ yêu cầu kéo dài, ảnh hưởng đến lịch trình \\
\hline
Chi phí phát sinh & Phải sửa đổi khi đã phát triển, chi phí cao gấp nhiều lần \\
\hline
Chất lượng kém & Sản phẩm không đáp ứng đúng nhu cầu khách hàng \\
\hline
Xung đột nội bộ & Mâu thuẫn giữa nhóm phát triển và khách hàng \\
\hline
\end{tabular}
\label{tab:requirement_impact}
\end{table}

\subsection{Vấn đề 2: Ước lượng Story Point thiếu chính xác}

\subsubsection{Mô tả vấn đề}

Trong Scrum, việc ước lượng Story Point đóng vai trò quan trọng trong việc lập kế hoạch Sprint. Tuy nhiên, quá trình này thường gặp nhiều khó khăn:

\begin{itemize}
\item Phụ thuộc vào kinh nghiệm của từng thành viên.
\item Thiếu dữ liệu lịch sử để tham khảo.
\item Sự khác biệt trong cách hiểu và đánh giá giữa các thành viên.
\item Áp lực từ quản lý dẫn đến ước lượng thiếu khách quan.
\end{itemize}

\subsubsection{Hệ quả}

\begin{figure}[H]
\centering
\fbox{
\begin{minipage}{0.9\textwidth}
\centering
Ước lượng sai $\rightarrow$ Kế hoạch Sprint sai $\rightarrow$ Tiến độ chậm $\rightarrow$ Chi phí tăng $\rightarrow$ Khách hàng không hài lòng
\end{minipage}
}
\caption{Chuỗi hệ quả của ước lượng sai}
\label{fig:estimation_consequences}
\end{figure}

\subsection{Vấn đề 3: Phát hiện rủi ro muộn}

\subsubsection{Mô tả vấn đề}

Trong quá trình thực hiện Sprint, có nhiều rủi ro có thể phát sinh nhưng thường được phát hiện muộn:

\begin{itemize}
\item Chậm tiến độ do Task phức tạp hơn dự kiến.
\item Lỗi kỹ thuật phát sinh trong quá trình phát triển.
\item Sự thay đổi yêu cầu đột ngột từ khách hàng.
\item Nguồn lực bị giảm do ốm đau hoặc nghỉ việc.
\end{itemize}

\subsubsection{Thống kê thực tế}

Theo các nghiên cứu về quản lý dự án phần mềm, tỷ lệ dự án gặp rủi ro về tiến độ và chi phí ở mức cao. Bảng \ref{tab:risk_statistics} thống kê các nguyên nhân chính:

\begin{table}[H]
\centering
\caption{Các nguyên nhân gây rủi ro dự án}
\begin{tabular}{|p{6cm}|c|}
\hline
\textbf{Nguyên nhân} & \textbf{Tỷ lệ gặp phải} \\
\hline
Yêu cầu thay đổi & 65\% \\
\hline
Ước lượng sai & 58\% \\
\hline
Thiếu nguồn lực & 42\% \\
\hline
Lỗi kỹ thuật & 38\% \\
\hline
Quản lý rủi ro kém & 35\% \\
\hline
Giao tiếp không hiệu quả & 30\% \\
\hline
\end{tabular}
\label{tab:risk_statistics}
\end{table}

\subsection{Vấn đề 4: Tài liệu dự án không đồng nhất}

\subsubsection{Mô tả vấn đề}

Việc duy trì tài liệu dự án là một thách thức lớn đối với các nhóm phát triển:

\begin{itemize}
\item Tài liệu không được cập nhật thường xuyên.
\item Nội dung tài liệu không đồng nhất giữa các phiên bản.
\item Thiếu tài liệu cho các quyết định quan trọng.
\item Khó tìm kiếm và truy xuất thông tin.
\end{itemize}

\subsection{Vấn đề 5: Phân công nguồn lực chưa tối ưu}

\subsubsection{Mô tả vấn đề}

Việc phân công Developer vào các Task thường dựa trên cảm tính của Project Manager:

\begin{itemize}
\item Không tận dụng được thế mạnh của từng thành viên.
\item Mất cân đối khối lượng công việc.
\item Task bị chuyển qua lại giữa các Developer.
\item Giảm năng suất làm việc của nhóm.
\end{itemize}

%%%%%%%%%%%%%%%%%%%%%%%%%%%%%%%%%%%%%%%%%%%%%%%%%%%%%%%%%%%%%%
\section{Phân tích nguyên nhân gốc rễ}
%%%%%%%%%%%%%%%%%%%%%%%%%%%%%%%%%%%%%%%%%%%%%%%%%%%%%%%%%%%%%%

\subsection{Sơ đồ Ishikawa (Fishbone Diagram)}

Để xác định nguyên nhân gốc rễ của các vấn đề, luận văn sử dụng sơ đồ Ishikawa:

\begin{figure}[H]
\centering
\fbox{
\begin{minipage}{0.95\textwidth}
\centering
\begin{tabular}{c}
\textbf{Nguyên nhân chính} \\
Con người $\rightarrow$ Quy trình $\rightarrow$ Công cụ $\rightarrow$ Dữ liệu $\rightarrow$ Môi trường \\
\downarrow \\
\textbf{Vấn đề: Quản lý dự án kém hiệu quả}
\end{tabular}
\end{minipage}
}
\caption{Sơ đồ Ishikawa phân tích nguyên nhân}
\label{fig:ishikawa}
\end{figure}

\subsection{Phân tích chi tiết}

\subsubsection{Nguyên nhân từ con người}

\begin{itemize}
\item Thiếu kỹ năng phân tích yêu cầu.
\item Kinh nghiệm ước lượng không đồng đều.
\item Khó khăn trong việc đánh giá rủi ro.
\item Sự khác biệt trong cách hiểu và truyền đạt thông tin.
\end{itemize}

\subsubsection{Nguyên nhân từ quy trình}

\begin{itemize}
\item Quy trình phân tích yêu cầu chưa được chuẩn hóa.
\item Thiếu quy trình đánh giá và cải tiến ước lượng.
\item Quy trình phát hiện rủi ro chưa được thiết lập.
\item Quy trình quản lý tài liệu chưa hiệu quả.
\end{itemize}

\subsubsection{Nguyên nhân từ công cụ}

\begin{itemize}
\item Công cụ hiện tại chưa hỗ trợ phân tích AI.
\item Thiếu khả năng dự báo và cảnh báo sớm.
\item Chưa có tính năng tự động sinh tài liệu.
\item Khó khai thác dữ liệu lịch sử.
\end{itemize}

\subsubsection{Nguyên nhân từ dữ liệu}

\begin{itemize}
\item Thiếu dữ liệu lịch sử chất lượng cao.
\item Dữ liệu không được tổ chức có cấu trúc.
\item Không có hệ thống lưu trữ tri thức.
\item Dữ liệu phân tán ở nhiều hệ thống khác nhau.
\end{itemize}

%%%%%%%%%%%%%%%%%%%%%%%%%%%%%%%%%%%%%%%%%%%%%%%%%%%%%%%%%%%%%%
\section{Xác định yêu cầu giải pháp}
%%%%%%%%%%%%%%%%%%%%%%%%%%%%%%%%%%%%%%%%%%%%%%%%%%%%%%%%%%%%%%

\subsection{Yêu cầu chức năng cốt lõi}

Dựa trên phân tích vấn đề, hệ thống cần đáp ứng các yêu cầu chức năng cốt lõi sau:

\begin{enumerate}
\item \textbf{Phân tích yêu cầu thông minh}
\begin{itemize}
\item Tự động phát hiện yêu cầu mơ hồ, trùng lặp và mâu thuẫn.
\item Đề xuất câu hỏi làm rõ yêu cầu.
\item Tự động sinh User Story và Acceptance Criteria.
\end{itemize}

\item \textbf{Hỗ trợ ước lượng Story Point}
\begin{itemize}
\item Đề xuất Story Point dựa trên dữ liệu lịch sử.
\item Cung cấp giải thích và độ tin cậy cho đề xuất.
\item Hỗ trợ quá trình Planning Poker.
\end{itemize}

\item \textbf{Dự báo và cảnh báo rủi ro}
\begin{itemize}
\item Theo dõi tiến độ Sprint theo thời gian thực.
\item Phát hiện sớm các dấu hiệu chậm tiến độ.
\item Gửi cảnh báo khi có rủi ro phát sinh.
\end{itemize}

\item \textbf{Tự động sinh tài liệu}
\begin{itemize}
\item Sinh Sprint Report tự động.
\item Tạo Meeting Minutes và Retrospective Report.
\item Xuất tài liệu dưới dạng PDF/Markdown.
\end{itemize}

\item \textbf{Hỗ trợ phân công nguồn lực}
\begin{itemize}
\item Phân tích kỹ năng và kinh nghiệm của Developer.
\item Đề xuất phân công tối ưu cho từng Task.
\item Cân bằng tải công việc trong nhóm.
\end{itemize}
\end{enumerate}

\subsection{Yêu cầu phi chức năng}

Ngoài các yêu cầu chức năng, hệ thống cần đáp ứng các yêu cầu phi chức năng sau:

\begin{table}[H]
\centering
\caption{Yêu cầu phi chức năng của hệ thống}
\begin{tabular}{|p{5cm}|p{8cm}|}
\hline
\textbf{Yêu cầu} & \textbf{Mô tả} \\
\hline
Hiệu năng & Thời gian phản hồi API < 500 ms, AI < 15 giây \\
\hline
Bảo mật & JWT Authentication, HTTPS, RBAC \\
\hline
Khả năng mở rộng & Kiến trúc Microservices, Multi-Agent \\
\hline
Khả năng bảo trì & Clean Architecture, SOLID Principles \\
\hline
Khả năng sử dụng & Giao diện trực quan, hỗ trợ đa thiết bị \\
\hline
\end{tabular}
\label{tab:non_functional_requirements}
\end{table}

%%%%%%%%%%%%%%%%%%%%%%%%%%%%%%%%%%%%%%%%%%%%%%%%%%%%%%%%%%%%%%
\section{Định hướng giải pháp}
%%%%%%%%%%%%%%%%%%%%%%%%%%%%%%%%%%%%%%%%%%%%%%%%%%%%%%%%%%%%%%

\subsection{Mô hình tổng quan}

Dựa trên phân tích vấn đề và yêu cầu, luận văn đề xuất mô hình giải pháp tổng quan như sau:

\begin{figure}[H]
\centering
\fbox{
\begin{minipage}{0.9\textwidth}
\centering
\begin{tabular}{c}
\textbf{Người dùng (Product Owner, Scrum Master, Developer)} \\
$\downarrow$ \\
\textbf{Giao diện (Frontend React)} \\
$\downarrow$ \\
\textbf{API Gateway (Spring Boot)} \\
$\downarrow$ \\
\textbf{AI Orchestrator} \\
$\downarrow$ \\
\textbf{AI Agents (Multi-Agent System)} \\
$\downarrow$ \\
\textbf{Large Language Models + RAG} \\
$\downarrow$ \\
\textbf{Cơ sở dữ liệu (PostgreSQL)}
\end{tabular}
\end{minipage}
}
\caption{Mô hình giải pháp tổng quan}
\label{fig:solution_overview}
\end{figure}

\subsection{Các thành phần chính}

\subsubsection{Tầng giao diện (Frontend)}

Cung cấp giao diện trực quan cho người dùng thực hiện các thao tác quản lý dự án và tương tác với AI.

\subsubsection{Tầng API Gateway (Backend)}

Tiếp nhận và điều phối các yêu cầu từ Frontend, xác thực người dùng và gọi các service nghiệp vụ.

\subsubsection{Tầng AI Orchestrator}

Điều phối các AI Agent, quản lý Prompt và tích hợp với mô hình ngôn ngữ lớn.

\subsubsection{Tầng Multi-Agent AI}

Bao gồm các Agent chuyên biệt: Requirement Agent, Planning Agent, Story Point Agent, Forecast Agent, Risk Agent, Assignment Agent, Documentation Agent và Dashboard Agent.

\subsubsection{Tầng dữ liệu}

Lưu trữ dữ liệu dự án, lịch sử Sprint và kết quả phân tích AI.

%%%%%%%%%%%%%%%%%%%%%%%%%%%%%%%%%%%%%%%%%%%%%%%%%%%%%%%%%%%%%%
\section{Yêu cầu về công nghệ và hạ tầng}
%%%%%%%%%%%%%%%%%%%%%%%%%%%%%%%%%%%%%%%%%%%%%%%%%%%%%%%%%%%%%%

\subsection{Công nghệ sử dụng}

\begin{table}[H]
\centering
\caption{Công nghệ dự kiến cho hệ thống}
\begin{tabular}{|p{4cm}|p{5cm}|p{5cm}|}
\hline
\textbf{Thành phần} & \textbf{Công nghệ} & \textbf{Vai trò} \\
\hline
Frontend & React + TypeScript & Giao diện người dùng \\
\hline
Backend & Spring Boot & API và nghiệp vụ \\
\hline
AI Service & FastAPI (Python) & Suy luận AI \\
\hline
Database & PostgreSQL & Lưu trữ dữ liệu \\
\hline
Vector DB & pgvector / Milvus & Lưu embedding \\
\hline
LLM & OpenAI API / Ollama & Xử lý ngôn ngữ tự nhiên \\
\hline
Container & Docker & Đóng gói ứng dụng \\
\hline
CI/CD & GitHub Actions & Tự động hóa \\
\hline
\end{tabular}
\label{tab:technology_stack}
\end{table}

\subsection{Yêu cầu hạ tầng}

\begin{itemize}
\item Máy chủ có RAM tối thiểu 16GB để chạy các AI Agent.
\item Ổ cứng SSD để đảm bảo hiệu năng truy xuất dữ liệu.
\item Kết nối Internet ổn định để gọi API LLM (nếu sử dụng dịch vụ cloud).
\item Hỗ trợ Docker và Docker Compose để triển khai.
\end{itemize}

%%%%%%%%%%%%%%%%%%%%%%%%%%%%%%%%%%%%%%%%%%%%%%%%%%%%%%%%%%%%%%
\section{Tổng kết chương}
%%%%%%%%%%%%%%%%%%%%%%%%%%%%%%%%%%%%%%%%%%%%%%%%%%%%%%%%%%%%%%

Trong chương này, luận văn đã phân tích chi tiết hiện trạng quản lý dự án phần mềm trong các doanh nghiệp hiện nay, xác định các vấn đề và thách thức trong quy trình hiện tại.

Các vấn đề chính đã được xác định bao gồm:

\begin{enumerate}
\item Phân tích yêu cầu thiếu hiệu quả.
\item Ước lượng Story Point thiếu chính xác.
\item Phát hiện rủi ro muộn.
\item Tài liệu dự án không đồng nhất.
\item Phân công nguồn lực chưa tối ưu.
\end{enumerate}

Từ đó, luận văn đã xác định các yêu cầu chức năng và phi chức năng của hệ thống, đồng thời đề xuất định hướng giải pháp tổng thể với kiến trúc Multi-Agent AI, tích hợp Large Language Models và cơ chế Retrieval-Augmented Generation (RAG).

Đây là cơ sở để Chương 5 tiến hành đặc tả yêu cầu chi tiết và Chương 6 phân tích nghiệp vụ của hệ thống.

%%%%%%%%%%%%%%%%%%%%%%%%%%%%%%%%%%%%%%%%%%%%%%%%%%%%%%%%%%%%%%%%%%%%%%%%%%%%%%%%%%%%%%%%%%%%%%%%%%%%%%%%%%%%%%%%%%%%%%%%%%%%
%%%%%%%%%%%%%%%%%%%%%%%%%%%%%%%%%%%%%%%%%%%%%%%%%%%%%%%%%%%%%%
\chapter{ĐẶC TẢ YÊU CẦU HỆ THỐNG}
\label{chap:requirements_specification}
%%%%%%%%%%%%%%%%%%%%%%%%%%%%%%%%%%%%%%%%%%%%%%%%%%%%%%%%%%%%%%

\section{Giới thiệu chương}

Sau khi phân tích bài toán và xác định các vấn đề cần giải quyết ở Chương 4, chương này trình bày đặc tả yêu cầu chi tiết của hệ thống AI Project Management. Nội dung của chương dựa trên các tài liệu phân tích nghiệp vụ, Product Vision và Software Requirement Specification (SRS) đã được xây dựng.

Mục tiêu của chương bao gồm:

\begin{itemize}
\item Xác định chi tiết các yêu cầu chức năng của hệ thống.
\item Xác định các yêu cầu phi chức năng cần đáp ứng.
\item Xây dựng Product Backlog cho hệ thống.
\item Xác định các ràng buộc và giả định của hệ thống.
\end{itemize}

Đây là cơ sở để Chương 6 tiến hành phân tích nghiệp vụ và Chương 7 thiết kế hệ thống.

%%%%%%%%%%%%%%%%%%%%%%%%%%%%%%%%%%%%%%%%%%%%%%%%%%%%%%%%%%%%%%
\section{Yêu cầu chức năng}
%%%%%%%%%%%%%%%%%%%%%%%%%%%%%%%%%%%%%%%%%%%%%%%%%%%%%%%%%%%%%%

Yêu cầu chức năng của hệ thống được chia thành chín nhóm chính, tương ứng với các phân hệ chức năng đã xác định trong Product Vision.

\subsection{Quản lý người dùng (User Management)}

Hệ thống cần cung cấp các chức năng quản lý người dùng cơ bản để đảm bảo bảo mật và phân quyền truy cập.

\begin{table}[H]
\centering
\caption{Yêu cầu chức năng - Quản lý người dùng}
\begin{tabular}{|p{2.5cm}|p{3cm}|p{8cm}|}
\hline
\textbf{Mã} & \textbf{Tên} & \textbf{Mô tả} \\
\hline
UC-01 & Đăng ký tài khoản & Cho phép người dùng mới đăng ký tài khoản với email và mật khẩu. \\
\hline
UC-02 & Đăng nhập & Xác thực người dùng bằng email và mật khẩu, cấp JWT Token. \\
\hline
UC-03 & Đăng xuất & Hủy phiên đăng nhập hiện tại của người dùng. \\
\hline
UC-04 & Quên mật khẩu & Cho phép người dùng đặt lại mật khẩu qua email. \\
\hline
UC-05 & Đổi mật khẩu & Cho phép người dùng thay đổi mật khẩu khi đã đăng nhập. \\
\hline
UC-06 & Quản lý hồ sơ & Cho phép người dùng cập nhật thông tin cá nhân và avatar. \\
\hline
UC-07 & Xác thực JWT & Xác thực mọi API request thông qua JWT Access Token. \\
\hline
UC-08 & Refresh Token & Làm mới Access Token khi hết hạn. \\
\hline
\end{tabular}
\label{tab:user_management_requirements}
\end{table}

\textbf{Chi tiết yêu cầu UC-02 (Đăng nhập):}

\begin{itemize}
\item Người dùng nhập email và mật khẩu.
\item Hệ thống kiểm tra email tồn tại và mật khẩu khớp (sử dụng BCrypt).
\item Nếu hợp lệ, hệ thống sinh JWT Access Token (thời hạn 15 phút) và Refresh Token (thời hạn 7 ngày).
\item Trả về token và thông tin người dùng.
\item Nếu sai, trả về lỗi 401 Unauthorized.
\end{itemize}

\subsection{Quản lý Workspace}

Workspace là đơn vị tổ chức cao nhất của hệ thống, đại diện cho một tổ chức hoặc nhóm phát triển.

\begin{table}[H]
\centering
\caption{Yêu cầu chức năng - Quản lý Workspace}
\begin{tabular}{|p{2.5cm}|p{3.5cm}|p{8cm}|}
\hline
\textbf{Mã} & \textbf{Tên} & \textbf{Mô tả} \\
\hline
WS-01 & Tạo Workspace & Cho phép người dùng tạo Workspace mới. \\
\hline
WS-02 & Cập nhật Workspace & Cho phép Workspace Owner cập nhật thông tin. \\
\hline
WS-03 & Xóa Workspace & Cho phép Workspace Owner xóa Workspace. \\
\hline
WS-04 & Mời thành viên & Gửi lời mời tham gia Workspace qua email. \\
\hline
WS-05 & Phân quyền & Gán vai trò (Owner, Admin, Member, Viewer) cho thành viên. \\
\hline
WS-06 & Xóa thành viên & Xóa thành viên khỏi Workspace. \\
\hline
WS-07 & Theo dõi hoạt động & Xem lịch sử hoạt động trong Workspace. \\
\hline
\end{tabular}
\label{tab:workspace_management_requirements}
\end{table}

\textbf{Quy tắc phân quyền trong Workspace:}

\begin{itemize}
\item \textbf{Owner}: Toàn quyền quản lý Workspace, bao gồm xóa Workspace.
\item \textbf{Admin}: Có thể mời/xóa thành viên, quản lý Project.
\item \textbf{Member}: Có thể tạo và tham gia Project.
\item \textbf{Viewer}: Chỉ có quyền xem dữ liệu.
\end{itemize}

\subsection{Quản lý Project}

Project là đơn vị quản lý dự án cụ thể, thuộc về một Workspace.

\begin{table}[H]
\centering
\caption{Yêu cầu chức năng - Quản lý Project}
\begin{tabular}{|p{2.5cm}|p{3.5cm}|p{8cm}|}
\hline
\textbf{Mã} & \textbf{Tên} & \textbf{Mô tả} \\
\hline
PR-01 & Tạo Project & Tạo dự án mới với tên, mô tả, ngày bắt đầu và kết thúc. \\
\hline
PR-02 & Cập nhật Project & Cập nhật thông tin dự án. \\
\hline
PR-03 & Xóa Project & Xóa dự án (yêu cầu quyền Project Manager). \\
\hline
PR-04 & Lưu trữ Project & Lưu trữ dự án đã hoàn thành. \\
\hline
PR-05 & Quản lý thành viên & Thêm/xóa thành viên vào Project. \\
\hline
PR-06 & Xem Dashboard & Xem tổng quan tiến độ và KPI của dự án. \\
\hline
\end{tabular}
\label{tab:project_management_requirements}
\end{table}

\textbf{Thông tin Project cần lưu trữ:}

\begin{itemize}
\item Project Name (bắt buộc)
\item Description
\item Status (Planning, Active, Done, Archived)
\item Priority (Low, Medium, High)
\item Start Date
\item End Date
\item Visibility (Public, Private)
\item Members
\item Labels
\end{itemize}

\subsection{Quản lý Sprint}

Sprint là đơn vị lập kế hoạch cơ bản trong Scrum.

\begin{table}[H]
\centering
\caption{Yêu cầu chức năng - Quản lý Sprint}
\begin{tabular}{|p{2.5cm}|p{3.5cm}|p{8cm}|}
\hline
\textbf{Mã} & \textbf{Tên} & \textbf{Mô tả} \\
\hline
SP-01 & Tạo Sprint & Tạo Sprint mới với tên, mục tiêu và thời gian. \\
\hline
SP-02 & Cập nhật Sprint & Cập nhật thông tin Sprint. \\
\hline
SP-03 & Xóa Sprint & Xóa Sprint (chỉ khi chưa bắt đầu). \\
\hline
SP-04 & Bắt đầu Sprint & Khởi động Sprint, không thể thay đổi Backlog. \\
\hline
SP-05 & Kết thúc Sprint & Kết thúc Sprint, tính toán Velocity. \\
\hline
SP-06 & Xem Burndown Chart & Xem biểu đồ tiến độ Sprint. \\
\hline
SP-07 & Xem Velocity & Xem vận tốc của Sprint và lịch sử. \\
\hline
\end{tabular}
\label{tab:sprint_management_requirements}
\end{table}

\textbf{Quy tắc Sprint:}

\begin{itemize}
\item Sprint không thể bắt đầu nếu chưa có ít nhất một User Story.
\item User Story chỉ có thể được thêm vào Sprint trước khi Sprint bắt đầu.
\item Sprint Goal là bắt buộc và phải được xác định trước khi bắt đầu.
\item Thời gian Sprint mặc định là 2 tuần (có thể cấu hình).
\end{itemize}

\subsection{Quản lý User Story}

User Story là đơn vị mô tả yêu cầu từ góc độ người dùng.

\begin{table}[H]
\centering
\caption{Yêu cầu chức năng - Quản lý User Story}
\begin{tabular}{|p{2.5cm}|p{3.5cm}|p{8cm}|}
\hline
\textbf{Mã} & \textbf{Tên} & \textbf{Mô tả} \\
\hline
ST-01 & Tạo User Story & Tạo User Story với title, description, acceptance criteria. \\
\hline
ST-02 & Cập nhật User Story & Cập nhật thông tin User Story. \\
\hline
ST-03 & Xóa User Story & Xóa User Story (nếu chưa được gán vào Sprint). \\
\hline
ST-04 & Gán Epic & Liên kết User Story với Epic. \\
\hline
ST-05 & Gán Sprint & Gán User Story vào Sprint. \\
\hline
ST-06 & Ước lượng Story Point & Gán Story Point cho User Story. \\
\hline
ST-07 & Cập nhật trạng thái & Cập nhật trạng thái (TODO, DOING, DONE). \\
\hline
\end{tabular}
\label{tab:story_management_requirements}
\end{table}

\textbf{Cấu trúc User Story:}

\begin{verbatim}
As a [Role]
I want [Goal]
So that [Business Value]

Acceptance Criteria:
- [Criterion 1]
- [Criterion 2]
- [Criterion 3]
\end{verbatim}

\subsection{Quản lý Task}

Task là đơn vị công việc kỹ thuật nhỏ nhất, được phân rã từ User Story.

\begin{table}[H]
\centering
\caption{Yêu cầu chức năng - Quản lý Task}
\begin{tabular}{|p{2.5cm}|p{3.5cm}|p{8cm}|}
\hline
\textbf{Mã} & \textbf{Tên} & \textbf{Mô tả} \\
\hline
TS-01 & Tạo Task & Tạo Task kỹ thuật từ User Story. \\
\hline
TS-02 & Cập nhật Task & Cập nhật thông tin Task. \\
\hline
TS-03 & Xóa Task & Xóa Task. \\
\hline
TS-04 & Phân công & Gán Task cho Developer. \\
\hline
TS-05 & Cập nhật trạng thái & Cập nhật trạng thái (TODO, DOING, DONE). \\
\hline
TS-06 & Ước lượng thời gian & Ước lượng giờ công cho Task. \\
\hline
TS-07 & Thêm Comment & Thêm bình luận vào Task. \\
\hline
TS-08 & Đính kèm File & Đính kèm tài liệu vào Task. \\
\hline
\end{tabular}
\label{tab:task_management_requirements}
\end{table}

\textbf{Hiển thị Task:}

Hệ thống hỗ trợ hiển thị Task dưới nhiều chế độ:

\begin{itemize}
\item \textbf{Kanban View}: Hiển thị Task theo cột trạng thái (TODO, DOING, DONE).
\item \textbf{List View}: Hiển thị Task dưới dạng danh sách.
\item \textbf{Calendar View}: Hiển thị Task trên lịch theo ngày đến hạn.
\item \textbf{Timeline View}: Hiển thị Task trên trục thời gian.
\end{itemize}

\subsection{AI Requirement Analysis}

Đây là chức năng cốt lõi của hệ thống, sử dụng AI để phân tích và làm rõ yêu cầu.

\begin{table}[H]
\centering
\caption{Yêu cầu chức năng - AI Requirement Analysis}
\begin{tabular}{|p{2.5cm}|p{3.5cm}|p{8cm}|}
\hline
\textbf{Mã} & \textbf{Tên} & \textbf{Mô tả} \\
\hline
AI-01 & Phân tích yêu cầu & AI đọc và phân tích văn bản yêu cầu từ người dùng. \\
\hline
AI-02 & Phát hiện yêu cầu mơ hồ & Phát hiện và đánh dấu các yêu cầu thiếu rõ ràng. \\
\hline
AI-03 & Phát hiện trùng lặp & Phát hiện các yêu cầu trùng lặp. \\
\hline
AI-04 & Phát hiện mâu thuẫn & Phát hiện các yêu cầu mâu thuẫn nhau. \\
\hline
AI-05 & Sinh User Story & Tự động sinh User Story theo mẫu chuẩn. \\
\hline
AI-06 & Sinh Acceptance Criteria & Tự động sinh Acceptance Criteria cho User Story. \\
\hline
AI-07 & Đề xuất câu hỏi làm rõ & Đề xuất câu hỏi để làm rõ yêu cầu chưa đầy đủ. \\
\hline
AI-08 & Tạo Requirement Summary & Tổng hợp và tóm tắt các yêu cầu. \\
\hline
\end{tabular}
\label{tab:ai_requirement_requirements}
\end{table}

\textbf{Quy trình phân tích yêu cầu:}

\begin{enumerate}
\item Người dùng nhập yêu cầu dưới dạng văn bản tự nhiên.
\item AI Agent tiền xử lý văn bản (tokenization, lemmatization).
\item AI phân tích ngữ nghĩa và ngữ cảnh của yêu cầu.
\item AI phát hiện các vấn đề: mơ hồ, trùng lặp, mâu thuẫn.
\item AI đề xuất User Story và Acceptance Criteria.
\item Người dùng xem xét và phê duyệt / chỉnh sửa.
\item Kết quả được lưu vào cơ sở dữ liệu.
\end{enumerate}

\subsection{AI Story Point Estimation}

\begin{table}[H]
\centering
\caption{Yêu cầu chức năng - AI Story Point Estimation}
\begin{tabular}{|p{2.5cm}|p{3.5cm}|p{8cm}|}
\hline
\textbf{Mã} & \textbf{Tên} & \textbf{Mô tả} \\
\hline
AI-09 & Đề xuất Story Point & AI đề xuất Story Point dựa trên phân tích User Story. \\
\hline
AI-10 & Đánh giá độ phức tạp & Đánh giá các yếu tố độ phức tạp của User Story. \\
\hline
AI-11 & Sử dụng dữ liệu lịch sử & Tham khảo lịch sử Sprint và Velocity của nhóm. \\
\hline
AI-12 & Giải thích lý do & Cung cấp giải thích cho đề xuất Story Point. \\
\hline
AI-13 & Độ tin cậy & Hiển thị Confidence Score cho đề xuất. \\
\hline
\end{tabular}
\label{tab:ai_storypoint_requirements}
\end{table}

\textbf{Các tiêu chí đánh giá Story Point:}

\begin{itemize}
\item \textbf{Complexity}: Độ phức tạp về mặt kỹ thuật.
\item \textbf{Business Logic}: Mức độ phức tạp của logic nghiệp vụ.
\item \textbf{Integration}: Mức độ tích hợp với các hệ thống khác.
\item \textbf{UI/UX}: Mức độ phức tạp của giao diện người dùng.
\item \textbf{Testing}: Mức độ phức tạp của kiểm thử.
\item \textbf{Uncertainty}: Mức độ không chắc chắn của yêu cầu.
\item \textbf{Dependency}: Sự phụ thuộc vào các Task khác.
\end{itemize}

\subsection{AI Sprint Forecast}

\begin{table}[H]
\centering
\caption{Yêu cầu chức năng - AI Sprint Forecast}
\begin{tabular}{|p{2.5cm}|p{3.5cm}|p{8cm}|}
\hline
\textbf{Mã} & \textbf{Tên} & \textbf{Mô tả} \\
\hline
AI-14 & Dự báo Sprint & Dự báo tiến độ hoàn thành Sprint hiện tại. \\
\hline
AI-15 & Dự báo thời gian hoàn thành & Dự báo thời điểm hoàn thành dự án. \\
\hline
AI-16 & Dự báo Velocity & Dự báo Velocity dựa trên lịch sử và xu hướng. \\
\hline
AI-17 & Tính xác suất hoàn thành & Tính xác suất hoàn thành Sprint đúng hạn. \\
\hline
\end{tabular}
\label{tab:ai_forecast_requirements}
\end{table}

\textbf{Dữ liệu đầu vào cho Forecast:}

\begin{itemize}
\item Lịch sử Velocity của các Sprint trước.
\item Dữ liệu Burndown Chart của Sprint hiện tại.
\item Story Point còn lại của Sprint.
\item Khả năng của nhóm (capacity).
\item Các yếu tố rủi ro đã được phát hiện.
\end{itemize}

\subsection{AI Risk Detection}

\begin{table}[H]
\centering
\caption{Yêu cầu chức năng - AI Risk Detection}
\begin{tabular}{|p{2.5cm}|p{3.5cm}|p{8cm}|}
\hline
\textbf{Mã} & \textbf{Tên} & \textbf{Mô tả} \\
\hline
AI-18 & Phát hiện rủi ro tiến độ & Phát hiện nguy cơ chậm tiến độ Sprint. \\
\hline
AI-19 & Phát hiện quá tải & Phát hiện Developer bị quá tải công việc. \\
\hline
AI-20 & Phân tích Burndown & Phân tích xu hướng Burndown Chart. \\
\hline
AI-21 & Phân tích Velocity & Phân tích sự thay đổi Velocity qua các Sprint. \\
\hline
AI-22 & Cảnh báo thời gian thực & Gửi cảnh báo khi phát hiện rủi ro. \\
\hline
AI-23 & Đề xuất giải pháp & Đề xuất giải pháp xử lý rủi ro. \\
\hline
\end{tabular}
\label{tab:ai_risk_requirements}
\end{table}

\textbf{Các loại rủi ro được phát hiện:}

\begin{enumerate}
\item \textbf{Sprint Delay}: Sprint có nguy cơ không hoàn thành đúng hạn.
\item \textbf{Workload Imbalance}: Developer quá tải hoặc phân công không đều.
\item \textbf{Velocity Drop}: Velocity giảm đột ngột qua các Sprint.
\item \textbf{Bug Surge}: Số lượng Bug tăng đột biến.
\item \textbf{Requirement Change}: Yêu cầu thay đổi đột ngột.
\item \textbf{Resource Loss}: Mất thành viên trong nhóm.
\end{enumerate}

\subsection{AI Documentation}

\begin{table}[H]
\centering
\caption{Yêu cầu chức năng - AI Documentation}
\begin{tabular}{|p{2.5cm}|p{3.5cm}|p{8cm}|}
\hline
\textbf{Mã} & \textbf{Tên} & \textbf{Mô tả} \\
\hline
AI-24 & Sinh Sprint Report & Tự động sinh báo cáo Sprint. \\
\hline
AI-25 & Sinh Meeting Minutes & Tự động sinh biên bản họp. \\
\hline
AI-26 & Sinh Retrospective Report & Tự động sinh báo cáo đánh giá Sprint. \\
\hline
AI-27 & Sinh SRS Summary & Tổng hợp tài liệu đặc tả yêu cầu. \\
\hline
AI-28 & Sinh Release Note & Tự động sinh ghi chú phát hành. \\
\hline
AI-29 & Sinh API Documentation & Sinh tài liệu API từ code. \\
\hline
AI-30 & Xuất PDF/Markdown & Xuất tài liệu dưới dạng PDF hoặc Markdown. \\
\hline
\end{tabular}
\label{tab:ai_documentation_requirements}
\end{table}

\subsection{AI Assignment}

\begin{table}[H]
\centering
\caption{Yêu cầu chức năng - AI Assignment}
\begin{tabular}{|p{2.5cm}|p{3.5cm}|p{8cm}|}
\hline
\textbf{Mã} & \textbf{Tên} & \textbf{Mô tả} \\
\hline
AI-31 & Đề xuất phân công & Đề xuất Developer phù hợp nhất cho Task. \\
\hline
AI-32 & Phân tích kỹ năng & Phân tích kỹ năng của từng Developer. \\
\hline
AI-33 & Cân bằng tải & Phân công đảm bảo cân bằng công việc. \\
\hline
AI-34 & Gợi ý dựa trên kinh nghiệm & Đề xuất dựa trên kinh nghiệm xử lý lỗi tương tự. \\
\hline
\end{tabular}
\label{tab:ai_assignment_requirements}
\end{table}

\textbf{Thuật toán phân công:}

Điểm phù hợp giữa Developer \(d\) và Task \(t\) được tính theo công thức:

\[
Score(d,t) = \alpha \cdot SkillMatch(d,t) + \beta \cdot Experience(d,t) - \gamma \cdot Workload(d)
\]

Trong đó:
\begin{itemize}
\item \(SkillMatch(d,t)\): Mức độ khớp giữa kỹ năng của Developer và yêu cầu Task.
\item \(Experience(d,t)\): Kinh nghiệm xử lý Task tương tự.
\item \(Workload(d)\): Khối lượng công việc hiện tại của Developer.
\item \(\alpha, \beta, \gamma\): Trọng số điều chỉnh.
\end{itemize}

\subsection{Dashboard and Analytics}

\begin{table}[H]
\centering
\caption{Yêu cầu chức năng - Dashboard}
\begin{tabular}{|p{2.5cm}|p{3.5cm}|p{8cm}|}
\hline
\textbf{Mã} & \textbf{Tên} & \textbf{Mô tả} \\
\hline
DB-01 & Sprint Burndown & Biểu đồ tiến độ Sprint. \\
\hline
DB-02 & Sprint Velocity & Biểu đồ Velocity qua các Sprint. \\
\hline
DB-03 & Task Distribution & Phân bổ Task theo trạng thái và thành viên. \\
\hline
DB-04 & Team Workload & Biểu đồ tải công việc của từng thành viên. \\
\hline
DB-05 & Story Point Progress & Tiến độ hoàn thành Story Point. \\
\hline
DB-06 & Project Timeline & Tổng quan tiến độ dự án. \\
\hline
DB-07 & Risk Level & Mức độ rủi ro hiện tại. \\
\hline
DB-08 & AI Insight & Phân tích và đề xuất từ AI. \\
\hline
\end{tabular}
\label{tab:dashboard_requirements}
\end{table}

%%%%%%%%%%%%%%%%%%%%%%%%%%%%%%%%%%%%%%%%%%%%%%%%%%%%%%%%%%%%%%
\section{Yêu cầu phi chức năng}
%%%%%%%%%%%%%%%%%%%%%%%%%%%%%%%%%%%%%%%%%%%%%%%%%%%%%%%%%%%%%%

\subsection{Hiệu năng (Performance)}

\begin{table}[H]
\centering
\caption{Yêu cầu về hiệu năng}
\begin{tabular}{|p{5cm}|p{8cm}|}
\hline
\textbf{Yêu cầu} & \textbf{Chỉ tiêu} \\
\hline
Thời gian phản hồi API (CRUD) & < 500 ms \\
\hline
Thời gian phản hồi AI & < 15 giây \\
\hline
Số người dùng đồng thời & Tối thiểu 500 \\
\hline
Thời gian tải Dashboard & < 2 giây \\
\hline
Thời gian tìm kiếm & < 1 giây \\
\hline
\end{tabular}
\label{tab:performance_requirements}
\end{table}

\subsection{Bảo mật (Security)}

\begin{itemize}
\item \textbf{JWT Authentication}: Mọi API request đều phải được xác thực bằng JWT.
\item \textbf{RBAC}: Phân quyền truy cập theo vai trò (Role-Based Access Control).
\item \textbf{HTTPS}: Mọi giao tiếp đều sử dụng HTTPS.
\item \textbf{Password Hashing}: Mật khẩu được mã hóa bằng BCrypt.
\item \textbf{Audit Log}: Ghi nhận mọi thao tác quan trọng của người dùng và AI.
\item \textbf{Input Validation}: Kiểm tra dữ liệu đầu vào để tránh tấn công.
\item \textbf{Rate Limiting}: Giới hạn số lượng request để tránh tấn công DDoS.
\end{itemize}

\subsection{Khả năng mở rộng (Scalability)}

\begin{itemize}
\item \textbf{Microservices Architecture}: Hệ thống được chia thành các service độc lập.
\item \textbf{Multi-Agent AI}: Các AI Agent có thể được mở rộng độc lập.
\item \textbf{Horizontal Scaling}: Có thể thêm instance để xử lý tải tăng.
\item \textbf{Database Scaling}: Hỗ trợ read replicas và sharding.
\item \textbf{Caching}: Sử dụng Redis để cache dữ liệu thường xuyên truy cập.
\end{itemize}

\subsection{Khả năng bảo trì (Maintainability)}

\begin{itemize}
\item \textbf{Clean Architecture}: Phân tách rõ ràng các tầng (Presentation, Application, Domain, Infrastructure).
\item \textbf{SOLID Principles}: Tuân thủ các nguyên tắc thiết kế SOLID.
\item \textbf{Domain-Driven Design (DDD)}: Thiết kế dựa trên nghiệp vụ.
\item \textbf{Code Coverage}: Unit test đạt tối thiểu 70\%.
\item \textbf{Docs}: Tài liệu kỹ thuật đầy đủ.
\end{itemize}

\subsection{Khả năng sử dụng (Usability)}

\begin{itemize}
\item \textbf{Responsive Design}: Hỗ trợ Desktop, Tablet và Mobile.
\item \textbf{Intuitive UI}: Giao diện trực quan, dễ sử dụng.
\item \textbf{Contextual AI Suggestions}: AI cung cấp gợi ý theo ngữ cảnh.
\item \textbf{User Onboarding}: Hướng dẫn sử dụng cho người mới.
\item \textbf{Multi-language}: Hỗ trợ tiếng Việt và tiếng Anh (trong tương lai).
\end{itemize}

\subsection{Độ tin cậy (Reliability)}

\begin{itemize}
\item \textbf{High Availability}: Hệ thống sẵn sàng 99.9\% thời gian.
\item \textbf{Disaster Recovery}: Sao lưu dữ liệu định kỳ.
\item \textbf{Fault Tolerance}: Khả năng phục hồi khi gặp sự cố.
\item \textbf{Health Checks}: Giám sát trạng thái của các service.
\item \textbf{Alerts}: Cảnh báo khi hệ thống gặp vấn đề.
\end{itemize}

%%%%%%%%%%%%%%%%%%%%%%%%%%%%%%%%%%%%%%%%%%%%%%%%%%%%%%%%%%%%%%
\section{Product Backlog}
%%%%%%%%%%%%%%%%%%%%%%%%%%%%%%%%%%%%%%%%%%%%%%%%%%%%%%%%%%%%%%

Dựa trên các yêu cầu chức năng đã xác định, Product Backlog của hệ thống được xây dựng với các ưu tiên như sau:

\begin{table}[H]
\centering
\caption{Product Backlog - Ưu tiên 1 (Must Have)}
\begin{tabular}{|p{2cm}|p{4cm}|p{3cm}|p{3cm}|}
\hline
\textbf{Mã} & \textbf{User Story} & \textbf{Ưu tiên} & \textbf{Story Point} \\
\hline
US-01 & Đăng nhập/Đăng ký & P1 & 5 \\
\hline
US-02 & Tạo Workspace & P1 & 3 \\
\hline
US-03 & Tạo Project & P1 & 5 \\
\hline
US-04 & Tạo Sprint & P1 & 3 \\
\hline
US-05 & Quản lý User Story & P1 & 8 \\
\hline
US-06 & Quản lý Task & P1 & 8 \\
\hline
US-07 & AI Phân tích yêu cầu & P1 & 13 \\
\hline
US-08 & AI Sinh User Story & P1 & 8 \\
\hline
\end{tabular}
\label{tab:product_backlog_p1}
\end{table}

\begin{table}[H]
\centering
\caption{Product Backlog - Ưu tiên 2 (Should Have)}
\begin{tabular}{|p{2cm}|p{4cm}|p{3cm}|p{3cm}|}
\hline
\textbf{Mã} & \textbf{User Story} & \textbf{Ưu tiên} & \textbf{Story Point} \\
\hline
US-09 & AI Ước lượng Story Point & P2 & 8 \\
\hline
US-10 & AI Dự báo Sprint & P2 & 8 \\
\hline
US-11 & AI Phát hiện rủi ro & P2 & 13 \\
\hline
US-12 & Dashboard & P2 & 8 \\
\hline
US-13 & AI Phân công Task & P2 & 8 \\
\hline
\end{tabular}
\label{tab:product_backlog_p2}
\end{table}

\begin{table}[H]
\centering
\caption{Product Backlog - Ưu tiên 3 (Could Have)}
\begin{tabular}{|p{2cm}|p{4cm}|p{3cm}|p{3cm}|}
\hline
\textbf{Mã} & \textbf{User Story} & \textbf{Ưu tiên} & \textbf{Story Point} \\
\hline
US-14 & AI Sinh tài liệu & P3 & 13 \\
\hline
US-15 & Audit Log & P3 & 5 \\
\hline
US-16 & Tìm kiếm nâng cao & P3 & 5 \\
\hline
US-17 & Báo cáo PDF & P3 & 8 \\
\hline
\end{tabular}
\label{tab:product_backlog_p3}
\end{table}

%%%%%%%%%%%%%%%%%%%%%%%%%%%%%%%%%%%%%%%%%%%%%%%%%%%%%%%%%%%%%%
\section{User Story Mapping}
%%%%%%%%%%%%%%%%%%%%%%%%%%%%%%%%%%%%%%%%%%%%%%%%%%%%%%%%%%%%%%

\begin{figure}[H]
\centering
\fbox{
\begin{minipage}{0.95\textwidth}
\centering
\begin{tabular}{|c|c|c|c|c|c|}
\hline
\textbf{Epic} & \multicolumn{5}{c|}{\textbf{User Stories}} \\
\hline
Authentication & Đăng ký & Đăng nhập & Quên mật khẩu & Đổi mật khẩu & Đăng xuất \\
\hline
Workspace & Tạo WS & Cập nhật WS & Xóa WS & Mời thành viên & Phân quyền \\
\hline
Project & Tạo Project & Cập nhật & Xóa Project & Quản lý thành viên & Xem Dashboard \\
\hline
Sprint & Tạo Sprint & Cập nhật & Bắt đầu & Kết thúc & Xem Burndown \\
\hline
Backlog & Tạo User Story & Cập nhật & Tạo Task & Cập nhật Task & Ước lượng SP \\
\hline
AI Assistant & Phân tích yêu cầu & Sinh User Story & Ước lượng SP & Dự báo Sprint & Phát hiện rủi ro \\
\hline
Documentation & Sprint Report & Meeting Minutes & Release Note & API Docs & Export PDF \\
\hline
Dashboard & Burndown Chart & Velocity Chart & Task Distribution & Workload & AI Insight \\
\hline
\end{tabular}
\end{minipage}
}
\caption{User Story Mapping của hệ thống}
\label{fig:user_story_mapping}
\end{figure}

%%%%%%%%%%%%%%%%%%%%%%%%%%%%%%%%%%%%%%%%%%%%%%%%%%%%%%%%%%%%%%
\section{Giả định và ràng buộc}
%%%%%%%%%%%%%%%%%%%%%%%%%%%%%%%%%%%%%%%%%%%%%%%%%%%%%%%%%%%%%%

\subsection{Giả định (Assumptions)}

\begin{enumerate}
\item Người dùng có kết nối Internet ổn định để sử dụng hệ thống.
\item Dữ liệu đầu vào là văn bản tiếng Anh hoặc tiếng Việt.
\item LLM API (OpenAI hoặc tương tự) có sẵn và hoạt động.
\item Người dùng có trình độ cơ bản về Agile/Scrum.
\item Dữ liệu lịch sử của dự án có chất lượng tốt.
\end{enumerate}

\subsection{Ràng buộc (Constraints)}

\begin{enumerate}
\item Hệ thống chỉ hoạt động trên môi trường Linux/Windows.
\item Yêu cầu RAM tối thiểu 16GB để chạy các AI Agent.
\item Thời gian phát triển MVP: 6 tháng.
\item Không tích hợp với Jira/Azure DevOps trong giai đoạn MVP.
\item Chỉ hỗ trợ văn bản, chưa hỗ trợ hình ảnh và video.
\end{enumerate}

\subsection{Rủi ro (Risks)}

\begin{enumerate}
\item \textbf{Nguy cơ}: LLM API có thể bị downtime hoặc thay đổi chính sách.
\item \textbf{Giảm thiểu}: Sử dụng multiple LLM providers, fallback mechanism.
\item \textbf{Nguy cơ}: Dữ liệu không đủ chất lượng để huấn luyện AI.
\item \textbf{Giảm thiểu}: Sử dụng RAG thay vì fine-tuning, khai thác tri thức nội bộ.
\item \textbf{Nguy cơ}: Hiệu năng AI không đáp ứng yêu cầu thời gian thực.
\item \textbf{Giảm thiểu}: Tối ưu Prompt, cache kết quả, sử dụng mô hình nhỏ hơn.
\end{enumerate}

%%%%%%%%%%%%%%%%%%%%%%%%%%%%%%%%%%%%%%%%%%%%%%%%%%%%%%%%%%%%%%
\section{Tổng kết chương}
%%%%%%%%%%%%%%%%%%%%%%%%%%%%%%%%%%%%%%%%%%%%%%%%%%%%%%%%%%%%%%

Trong chương này, luận văn đã trình bày chi tiết các yêu cầu của hệ thống AI Project Management dựa trên phân tích bài toán ở Chương 4.

Các nội dung chính đã được thực hiện:

\begin{enumerate}
\item Xác định 9 nhóm yêu cầu chức năng chính với hơn 50 yêu cầu chi tiết.
\item Xác định 6 nhóm yêu cầu phi chức năng (hiệu năng, bảo mật, mở rộng, v.v.).
\item Xây dựng Product Backlog với ưu tiên rõ ràng cho từng chức năng.
\item Xây dựng User Story Mapping để tổng quan toàn bộ hệ thống.
\item Xác định các giả định, ràng buộc và rủi ro của dự án.
\end{enumerate}

Đây là cơ sở để Chương 6 tiến hành phân tích nghiệp vụ chi tiết, và Chương 7 thiết kế kiến trúc hệ thống.

%%%%%%%%%%%%%%%%%%%%%%%%%%%%%%%%%%%%%%%%%%%%%%%%%%%%%%%%%%%%%%

%%%%%%%%%%%%%%%%%%%%%%%%%%%%%%%%%%%%%%%%%%%%%%%%%%%%%%%%%%%%%%
%%%%%%%%%%%%%%%%%%%%%%%%%%%%%%%%%%%%%%%%%%%%%%%%%%%%%%%%%%%%%%
\chapter{PHÂN TÍCH NGHIỆP VỤ}
\label{chap:business_analysis}
%%%%%%%%%%%%%%%%%%%%%%%%%%%%%%%%%%%%%%%%%%%%%%%%%%%%%%%%%%%%%%

\section{Giới thiệu chương}

Sau khi xác định các yêu cầu chức năng và phi chức năng của hệ thống ở Chương 5, chương này tập trung phân tích nghiệp vụ chi tiết của hệ thống AI Project Management. Phân tích nghiệp vụ là bước quan trọng để hiểu rõ cách thức hệ thống tương tác với người dùng và xử lý các quy trình nghiệp vụ.

Mục tiêu của chương bao gồm:

\begin{itemize}
\item Xác định các tác nhân (Actor) của hệ thống.
\item Xây dựng sơ đồ Use Case tổng quan và chi tiết.
\item Phân tích luồng nghiệp vụ chính của hệ thống.
\item Xây dựng sơ đồ hoạt động (Activity Diagram).
\item Xây dựng sơ đồ tuần tự (Sequence Diagram).
\item Phân tích mối quan hệ giữa các thành phần nghiệp vụ.
\end{itemize}

Đây là cơ sở để Chương 7 tiến hành thiết kế kiến trúc hệ thống.

%%%%%%%%%%%%%%%%%%%%%%%%%%%%%%%%%%%%%%%%%%%%%%%%%%%%%%%%%%%%%%
\section{Các tác nhân của hệ thống}
%%%%%%%%%%%%%%%%%%%%%%%%%%%%%%%%%%%%%%%%%%%%%%%%%%%%%%%%%%%%%%

\subsection{Danh sách tác nhân}

Hệ thống có năm nhóm tác nhân chính, mỗi nhóm có vai trò và quyền hạn khác nhau trong hệ thống.

\begin{table}[H]
\centering
\caption{Các tác nhân của hệ thống}
\begin{tabular}{|p{3cm}|p{3.5cm}|p{8cm}|}
\hline
\textbf{Tác nhân} & \textbf{Vai trò} & \textbf{Quyền hạn chính} \\
\hline
Guest & Người dùng chưa đăng nhập & Xem trang giới thiệu, đăng ký tài khoản \\
\hline
Member & Thành viên đã đăng nhập & Xem Workspace, tạo Project, quản lý Task \\
\hline
Project Manager & Quản lý dự án & Tạo/chỉnh sửa/xóa Project, quản lý Sprint, phê duyệt AI \\
\hline
Workspace Owner & Chủ sở hữu Workspace & Quản lý Workspace, phân quyền thành viên \\
\hline
AI Agent & Trợ lý AI & Tự động phân tích, đề xuất, cảnh báo \\
\hline
\end{tabular}
\label{tab:actors}
\end{table}

\subsection{Phân tích chi tiết tác nhân}

\subsubsection{Guest}

Guest là người dùng chưa đăng nhập vào hệ thống. Đây là tác nhân có quyền hạn thấp nhất.

\begin{itemize}
\item \textbf{Mục tiêu}: Đăng ký tài khoản để trở thành Member.
\item \textbf{Chức năng}: Xem trang chủ, đăng ký, đăng nhập.
\item \textbf{Hạn chế}: Không thể truy cập vào Workspace hay Project.
\end{itemize}

\subsubsection{Member}

Member là người dùng đã đăng nhập và có thể tham gia vào Workspace.

\begin{itemize}
\item \textbf{Mục tiêu}: Tham gia vào các dự án và thực hiện công việc được giao.
\item \textbf{Chức năng}: Xem Workspace, tạo Project, quản lý Task.
\item \textbf{Hạn chế}: Không thể xóa Workspace hay Project do người khác tạo.
\end{itemize}

\subsubsection{Project Manager}

Project Manager là người chịu trách nhiệm quản lý một dự án cụ thể.

\begin{itemize}
\item \textbf{Mục tiêu}: Đảm bảo dự án được thực hiện đúng tiến độ và chất lượng.
\item \textbf{Chức năng}: Tạo/chỉnh sửa Project, quản lý Sprint, phê duyệt đề xuất AI.
\item \textbf{Quyền đặc biệt}: Có thể xóa Project, phân công Task, đánh giá Sprint.
\end{itemize}

\subsubsection{Workspace Owner}

Workspace Owner là người sở hữu và quản lý toàn bộ Workspace.

\begin{itemize}
\item \textbf{Mục tiêu}: Quản lý tổ chức và phân quyền cho các thành viên.
\item \textbf{Chức năng}: Tạo/xóa Workspace, mời thành viên, phân quyền.
\item \textbf{Quyền đặc biệt}: Cao nhất trong Workspace, có thể xóa Workspace.
\end{itemize}

\subsubsection{AI Agent}

AI Agent là tác nhân đặc biệt, hoạt động tự động dưới sự điều phối của AI Orchestrator.

\begin{itemize}
\item \textbf{Mục tiêu}: Hỗ trợ người dùng trong các tác vụ quản lý dự án.
\item \textbf{Chức năng}: Phân tích yêu cầu, ước lượng Story Point, dự báo, phát hiện rủi ro.
\item \textbf{Đặc điểm}: Không có quyền thay đổi dữ liệu mà không có sự phê duyệt của con người.
\end{itemize}

%%%%%%%%%%%%%%%%%%%%%%%%%%%%%%%%%%%%%%%%%%%%%%%%%%%%%%%%%%%%%%
\section{Sơ đồ Use Case tổng quan}
%%%%%%%%%%%%%%%%%%%%%%%%%%%%%%%%%%%%%%%%%%%%%%%%%%%%%%%%%%%%%%

\subsection{Use Case Diagram tổng thể}

\begin{figure}[H]
\centering
\fbox{
\begin{minipage}{0.95\textwidth}
\centering
\begin{tabular}{c}
\textbf{Hệ thống Quản lý Dự án Phần mềm Hướng AI} \\
\\
\begin{tabular}{|c|c|}
\hline
\multicolumn{2}{|c|}{\textbf{Guest}} \\
\hline
Đăng ký & Đăng nhập \\
\hline
\multicolumn{2}{|c|}{\textbf{Member}} \\
\hline
Xem Workspace & Tạo Project \\
Quản lý Task & Xem Dashboard \\
\hline
\multicolumn{2}{|c|}{\textbf{Project Manager}} \\
\hline
Quản lý Sprint & Phân công Task \\
Phê duyệt AI & Xóa Project \\
\hline
\multicolumn{2}{|c|}{\textbf{Workspace Owner}} \\
\hline
Quản lý Workspace & Phân quyền \\
Xóa Workspace & Mời thành viên \\
\hline
\multicolumn{2}{|c|}{\textbf{AI Agent}} \\
\hline
Phân tích yêu cầu & Ước lượng SP \\
Dự báo Sprint & Phát hiện rủi ro \\
\hline
\end{tabular}
\end{tabular}
\end{minipage}
}
\caption{Sơ đồ Use Case tổng quan của hệ thống}
\label{fig:usecase_overview}
\end{figure}

\subsection{Phân nhóm Use Case theo tác nhân}

\begin{table}[H]
\centering
\caption{Phân nhóm Use Case theo tác nhân}
\begin{tabular}{|p{3.5cm}|p{9cm}|}
\hline
\textbf{Tác nhân} & \textbf{Use Case} \\
\hline
Guest & Đăng ký, Đăng nhập, Quên mật khẩu \\
\hline
Member & Xem Workspace, Tạo Project, Xem Project, Quản lý Task, Xem Dashboard, Tìm kiếm \\
\hline
Project Manager & Tạo/chỉnh sửa Project, Quản lý Sprint, Phân công Task, Phê duyệt AI, Xóa Project, Báo cáo \\
\hline
Workspace Owner & Quản lý Workspace, Mời thành viên, Phân quyền, Xóa Workspace \\
\hline
AI Agent & Phân tích yêu cầu, Sinh User Story, Ước lượng SP, Dự báo Sprint, Phát hiện rủi ro, Sinh tài liệu, Gợi ý phân công \\
\hline
\end{tabular}
\label{tab:usecase_by_actor}
\end{table}

%%%%%%%%%%%%%%%%%%%%%%%%%%%%%%%%%%%%%%%%%%%%%%%%%%%%%%%%%%%%%%
\section{Đặc tả chi tiết Use Case}
%%%%%%%%%%%%%%%%%%%%%%%%%%%%%%%%%%%%%%%%%%%%%%%%%%%%%%%%%%%%%%

\subsection{UC01 -- Đăng nhập}

\begin{table}[H]
\centering
\caption{Đặc tả UC01 - Đăng nhập}
\begin{tabular}{|p{4cm}|p{11cm}|}
\hline
\textbf{ID} & UC01 \\
\hline
\textbf{Tên} & Đăng nhập \\
\hline
\textbf{Actor} & Guest \\
\hline
\textbf{Mục tiêu} & Cho phép người dùng xác thực để truy cập hệ thống \\
\hline
\textbf{Tiền điều kiện} & Người dùng đã có tài khoản hợp lệ \\
\hline
\textbf{Hậu điều kiện} & JWT Token được cấp, người dùng chuyển đến Dashboard \\
\hline
\textbf{Độ ưu tiên} & Cao \\
\hline
\textbf{Tần suất} & Thường xuyên \\
\hline
\end{tabular}
\end{table}

\textbf{Luồng chính:}

\begin{enumerate}
\item Người dùng mở trang Login.
\item Nhập Email và Password.
\item Hệ thống kiểm tra dữ liệu đầu vào.
\item Backend xác thực thông tin với cơ sở dữ liệu.
\item Nếu hợp lệ, sinh JWT Access Token và Refresh Token.
\item Trả token về Frontend.
\item Frontend lưu token và chuyển đến Dashboard.
\end{enumerate}

\textbf{Luồng thay thế:}

\begin{itemize}
\item \textbf{A1 - Sai Email}: Hiển thị thông báo "Email không tồn tại".
\item \textbf{A2 - Sai Password}: Hiển thị thông báo "Mật khẩu không đúng".
\item \textbf{A3 - Tài khoản bị khóa}: Hiển thị thông báo "Tài khoản đã bị khóa".
\item \textbf{A4 - Email chưa xác thực}: Hiển thị thông báo "Vui lòng xác thực email".
\end{itemize}

\subsection{UC02 -- Tạo Workspace}

\begin{table}[H]
\centering
\caption{Đặc tả UC02 - Tạo Workspace}
\begin{tabular}{|p{4cm}|p{11cm}|}
\hline
\textbf{ID} & UC02 \\
\hline
\textbf{Tên} & Tạo Workspace \\
\hline
\textbf{Actor} & Member \\
\hline
\textbf{Mục tiêu} & Tạo một Workspace mới \\
\hline
\textbf{Tiền điều kiện} & Người dùng đã đăng nhập \\
\hline
\textbf{Hậu điều kiện} & Workspace mới được tạo, người dùng trở thành Owner \\
\hline
\textbf{Độ ưu tiên} & Cao \\
\hline
\end{tabular}
\end{table}

\textbf{Luồng chính:}

\begin{enumerate}
\item Người dùng chọn "Create Workspace".
\item Nhập tên Workspace (bắt buộc).
\item Chọn Logo (tùy chọn).
\item Chọn Visibility (Public/Private).
\item Nhấn nút "Create".
\item Backend kiểm tra tên Workspace không trùng.
\item Lưu Workspace vào cơ sở dữ liệu.
\item Gán người dùng là Owner.
\item Hiển thị Workspace mới tạo.
\end{enumerate}

\subsection{UC03 -- Tạo Project}

\begin{table}[H]
\centering
\caption{Đặc tả UC03 - Tạo Project}
\begin{tabular}{|p{4cm}|p{11cm}|}
\hline
\textbf{ID} & UC03 \\
\hline
\textbf{Tên} & Tạo Project \\
\hline
\textbf{Actor} & Project Manager \\
\hline
\textbf{Mục tiêu} & Tạo một dự án mới trong Workspace \\
\hline
\textbf{Tiền điều kiện} & Workspace tồn tại, người dùng có quyền \\
\hline
\textbf{Hậu điều kiện} & Project mới được tạo \\
\hline
\textbf{Độ ưu tiên} & Cao \\
\hline
\end{tabular}
\end{table}

\textbf{Thông tin Project cần nhập:}

\begin{itemize}
\item Project Name (bắt buộc)
\item Project Key (mã định danh, bắt buộc)
\item Description (tùy chọn)
\item Start Date (tùy chọn)
\item End Date (tùy chọn)
\item Priority (Low/Medium/High)
\item Visibility (Public/Private)
\end{itemize}

\textbf{Luồng chính:}

\begin{enumerate}
\item Người dùng chọn "Create Project" trong Workspace.
\item Nhập các thông tin cần thiết.
\item Nhấn "Create".
\item Backend kiểm tra dữ liệu và quyền.
\item Lưu Project vào cơ sở dữ liệu.
\item Hiển thị Project vừa tạo.
\end{enumerate}

\subsection{UC04 -- AI Requirement Analysis}

Đây là Use Case quan trọng nhất của hệ thống, thể hiện chức năng cốt lõi của AI.

\begin{table}[H]
\centering
\caption{Đặc tả UC04 - AI Requirement Analysis}
\begin{tabular}{|p{4cm}|p{11cm}|}
\hline
\textbf{ID} & UC04 \\
\hline
\textbf{Tên} & AI Requirement Analysis \\
\hline
\textbf{Actor} & Product Owner \\
\hline
\textbf{Supporting Actor} & Requirement Agent \\
\hline
\textbf{Mục tiêu} & AI phân tích và làm rõ yêu cầu \\
\hline
\textbf{Tiền điều kiện} & Yêu cầu đã được nhập dưới dạng văn bản \\
\hline
\textbf{Hậu điều kiện} & User Story và Acceptance Criteria được đề xuất \\
\hline
\textbf{Độ ưu tiên} & Cao \\
\hline
\end{tabular}
\end{table}

\textbf{Luồng chính:}

\begin{enumerate}
\item Product Owner nhập yêu cầu dưới dạng văn bản tự nhiên.
\item Frontend gửi yêu cầu đến Backend.
\item Backend xác thực JWT và kiểm tra quyền.
\item AI Orchestrator nhận yêu cầu và xác định loại nghiệp vụ.
\item Prompt Builder xây dựng Prompt với System Prompt, Business Context, Project Context.
\item RAG Service truy xuất tài liệu liên quan từ Vector Database.
\item Requirement Agent kết hợp Context và gửi đến LLM.
\item LLM phân tích và sinh kết quả.
\item JSON Validator kiểm tra cấu trúc dữ liệu.
\item Business Validator kiểm tra các quy tắc Scrum.
\item Kết quả được lưu vào cơ sở dữ liệu.
\item Trả về kết quả cho Frontend.
\item Product Owner xem xét và phê duyệt/chỉnh sửa.
\end{enumerate}

\textbf{Luồng thay thế:}

\begin{itemize}
\item \textbf{A1 - LLM không khả dụng}: Backend trả về lỗi 503, thông báo cho người dùng.
\item \textbf{A2 - Dữ liệu không hợp lệ}: JSON Validator báo lỗi, yêu cầu kiểm tra lại.
\end{itemize}

\subsection{UC05 -- AI Story Point Estimation}

\begin{table}[H]
\centering
\caption{Đặc tả UC05 - AI Story Point Estimation}
\begin{tabular}{|p{4cm}|p{11cm}|}
\hline
\textbf{ID} & UC05 \\
\hline
\textbf{Tên} & AI Story Point Estimation \\
\hline
\textbf{Actor} & Scrum Master \\
\hline
\textbf{Supporting Actor} & Story Point Agent \\
\hline
\textbf{Mục tiêu} & AI đề xuất Story Point cho User Story \\
\hline
\textbf{Tiền điều kiện} & User Story đã được tạo \\
\hline
\textbf{Hậu điều kiện} & Story Point được đề xuất và lưu \\
\hline
\textbf{Độ ưu tiên} & Cao \\
\hline
\end{tabular}
\end{table}

\textbf{Các tiêu chí đánh giá:}

\begin{itemize}
\item Complexity: Độ phức tạp của logic
\item Business Logic: Độ phức tạp nghiệp vụ
\item Integration: Mức độ tích hợp
\item UI/UX: Độ phức tạp giao diện
\item Testing: Độ phức tạp kiểm thử
\item Uncertainty: Mức độ không chắc chắn
\end{itemize}

\textbf{Luồng chính:}

\begin{enumerate}
\item Scrum Master chọn User Story cần ước lượng.
\item Frontend gửi yêu cầu đến Backend.
\item AI Orchestrator gọi Story Point Agent.
\item Agent truy xuất dữ liệu lịch sử Sprint và Velocity.
\item Agent đánh giá các tiêu chí và gửi Prompt đến LLM.
\item LLM sinh Story Point kèm giải thích.
\item Kết quả được trả về và hiển thị.
\item Scrum Master xem xét và xác nhận/điều chỉnh.
\end{enumerate}

\subsection{UC06 -- AI Sprint Forecast}

\begin{table}[H]
\centering
\caption{Đặc tả UC06 - AI Sprint Forecast}
\begin{tabular}{|p{4cm}|p{11cm}|}
\hline
\textbf{ID} & UC06 \\
\hline
\textbf{Tên} & AI Sprint Forecast \\
\hline
\textbf{Actor} & Scrum Master / Project Manager \\
\hline
\textbf{Supporting Actor} & Forecast Agent \\
\hline
\textbf{Mục tiêu} & Dự báo tiến độ hoàn thành Sprint \\
\hline
\textbf{Tiền điều kiện} & Sprint đang hoạt động, có dữ liệu Burndown \\
\hline
\textbf{Hậu điều kiện} & Kết quả dự báo được hiển thị \\
\hline
\textbf{Độ ưu tiên} & Trung bình \\
\hline
\end{tabular}
\end{table}

\textbf{Dữ liệu đầu vào:}

\begin{itemize}
\item Lịch sử Velocity
\item Dữ liệu Burndown Chart
\item Story Point còn lại
\item Team Capacity
\item Số ngày còn lại của Sprint
\end{itemize}

\textbf{Dự báo đầu ra:}

\begin{itemize}
\item Estimated Finish Date: Ngày dự kiến hoàn thành
\item Completion Probability: Xác suất hoàn thành đúng hạn (\%)
\item Remaining Sprint: Số Sprint còn lại dự kiến
\item Delay Risk: Mức độ rủi ro chậm tiến độ (Low/Medium/High)
\end{itemize}

\subsection{UC07 -- AI Risk Detection}

\begin{table}[H]
\centering
\caption{Đặc tả UC07 - AI Risk Detection}
\begin{tabular}{|p{4cm}|p{11cm}|}
\hline
\textbf{ID} & UC07 \\
\hline
\textbf{Tên} & AI Risk Detection \\
\hline
\textbf{Actor} & Project Manager \\
\hline
\textbf{Supporting Actor} & Risk Agent \\
\hline
\textbf{Mục tiêu} & Phát hiện và cảnh báo rủi ro trong Sprint \\
\hline
\textbf{Tiền điều kiện} & Sprint đang hoạt động \\
\hline
\textbf{Hậu điều kiện} & Cảnh báo rủi ro được gửi (nếu có) \\
\hline
\textbf{Độ ưu tiên} & Trung bình \\
\hline
\end{tabular}
\end{table}

\textbf{Các chỉ số rủi ro:}

\[
RiskScore = \sum_{i=1}^{n} w_i \cdot x_i
\]

Trong đó:
\begin{itemize}
\item \(x_1\): Task Delay (số Task bị chậm)
\item \(x_2\): Velocity Drop (mức giảm Velocity)
\item \(x_3\): Bug Count (số lượng Bug)
\item \(x_4\): Workload Level (mức độ quá tải)
\item \(x_5\): Requirement Change (mức độ thay đổi yêu cầu)
\item \(w_i\): Trọng số của từng chỉ số
\end{itemize}

\textbf{Phân loại rủi ro:}

\begin{table}[H]
\centering
\caption{Phân loại mức độ rủi ro}
\begin{tabular}{|c|c|c|}
\hline
\textbf{Risk Score} & \textbf{Mức độ} & \textbf{Hành động} \\
\hline
$< 0.3$ & Low & Tiếp tục theo dõi \\
\hline
$0.3 - 0.6$ & Medium & Cần xem xét và có kế hoạch dự phòng \\
\hline
$> 0.6$ & High & Cần hành động ngay lập tức \\
\hline
\end{tabular}
\label{tab:risk_levels}
\end{table}

\subsection{UC08 -- AI Documentation}

\begin{table}[H]
\centering
\caption{Đặc tả UC08 - AI Documentation}
\begin{tabular}{|p{4cm}|p{11cm}|}
\hline
\textbf{ID} & UC08 \\
\hline
\textbf{Tên} & AI Documentation \\
\hline
\textbf{Actor} & Project Manager \\
\hline
\textbf{Supporting Actor} & Documentation Agent \\
\hline
\textbf{Mục tiêu} & Tự động sinh tài liệu dự án \\
\hline
\textbf{Tiền điều kiện} & Dữ liệu dự án đã có \\
\hline
\textbf{Hậu điều kiện} & Tài liệu được sinh và lưu \\
\hline
\textbf{Độ ưu tiên} & Trung bình \\
\hline
\end{tabular}
\end{table}

\textbf{Các loại tài liệu hỗ trợ:}

\begin{itemize}
\item Sprint Report
\item Meeting Minutes
\item Retrospective Report
\item Product Backlog
\item Release Note
\item API Documentation
\item User Manual
\item SRS Summary
\end{itemize}

%%%%%%%%%%%%%%%%%%%%%%%%%%%%%%%%%%%%%%%%%%%%%%%%%%%%%%%%%%%%%%
\section{Luồng nghiệp vụ chính}
%%%%%%%%%%%%%%%%%%%%%%%%%%%%%%%%%%%%%%%%%%%%%%%%%%%%%%%%%%%%%%

\subsection{Luồng báo cáo và phân tích yêu cầu}

Quy trình này mô tả luồng từ khi Product Owner nhập yêu cầu đến khi User Story được tạo và phê duyệt.

\begin{figure}[H]
\centering
\fbox{
\begin{minipage}{0.95\textwidth}
\centering
\begin{tabular}{c}
\textbf{Product Owner} \\
$\downarrow$ \\
\textbf{Nhập yêu cầu} \\
$\downarrow$ \\
\textbf{AI Requirement Analysis} \\
$\downarrow$ \\
\textbf{Phân tích và làm rõ yêu cầu} \\
$\downarrow$ \\
\textbf{Sinh User Story \& Acceptance Criteria} \\
$\downarrow$ \\
\textbf{Product Owner xem xét và phê duyệt} \\
$\downarrow$ \\
\textbf{Lưu vào Product Backlog}
\end{tabular}
\end{minipage}
}
\caption{Luồng báo cáo và phân tích yêu cầu}
\label{fig:requirement_flow}
\end{figure}

\subsection{Luồng lập kế hoạch Sprint}

Quy trình này mô tả luồng từ khi Sprint Planning đến khi Sprint được bắt đầu.

\begin{figure}[H]
\centering
\fbox{
\begin{minipage}{0.95\textwidth}
\centering
\begin{tabular}{c}
\textbf{Scrum Master} \\
$\downarrow$ \\
\textbf{Chọn User Story cho Sprint} \\
$\downarrow$ \\
\textbf{AI Story Point Estimation} \\
$\downarrow$ \\
\textbf{AI Task Breakdown} \\
$\downarrow$ \\
\textbf{AI Assignment} \\
$\downarrow$ \\
\textbf{Scrum Master xem xét và phê duyệt} \\
$\downarrow$ \\
\textbf{Bắt đầu Sprint}
\end{tabular}
\end{minipage}
}
\caption{Luồng lập kế hoạch Sprint}
\label{fig:sprint_planning_flow}
\end{figure}

\subsection{Luồng thực hiện và giám sát Sprint}

Quy trình này mô tả luồng trong quá trình Sprint đang diễn ra.

\begin{figure}[H]
\centering
\fbox{
\begin{minipage}{0.95\textwidth}
\centering
\begin{tabular}{c}
\textbf{Developer cập nhật Task} \\
$\downarrow$ \\
\textbf{Cập nhật Burndown Chart} \\
$\downarrow$ \\
\textbf{AI Risk Detection (Liên tục)} \\
$\downarrow$ \\
\textbf{Phát hiện rủi ro?} \\
$\swarrow$ \hspace{3cm} $\searrow$ \\
\textbf{Có} \hspace{3cm} \textbf{Không} \\
$\downarrow$ \hspace{3cm} $\downarrow$ \\
\textbf{Gửi cảnh báo} \hspace{2.5cm} \textbf{Tiếp tục theo dõi} \\
$\downarrow$ \hspace{3cm} \\
\textbf{Project Manager xử lý}
\end{tabular}
\end{minipage}
}
\caption{Luồng thực hiện và giám sát Sprint}
\label{fig:sprint_execution_flow}
\end{figure}

\subsection{Luồng đánh giá và kết thúc Sprint}

Quy trình này mô tả luồng khi Sprint kết thúc và đánh giá.

\begin{figure}[H]
\centering
\fbox{
\begin{minipage}{0.95\textwidth}
\centering
\begin{tabular}{c}
\textbf{Kết thúc Sprint} \\
$\downarrow$ \\
\textbf{Tính Velocity} \\
$\downarrow$ \\
\textbf{AI Documentation (Sinh Sprint Report)} \\
$\downarrow$ \\
\textbf{Tổ chức Sprint Review} \\
$\downarrow$ \\
\textbf{Tổ chức Sprint Retrospective} \\
$\downarrow$ \\
\textbf{AI Documentation (Sinh Retrospective Report)} \\
$\downarrow$ \\
\textbf{Cập nhật Product Backlog}
\end{tabular}
\end{minipage}
}
\caption{Luồng đánh giá và kết thúc Sprint}
\label{fig:sprint_review_flow}
\end{figure}

%%%%%%%%%%%%%%%%%%%%%%%%%%%%%%%%%%%%%%%%%%%%%%%%%%%%%%%%%%%%%%
\section{Sơ đồ hoạt động (Activity Diagram)}
%%%%%%%%%%%%%%%%%%%%%%%%%%%%%%%%%%%%%%%%%%%%%%%%%%%%%%%%%%%%%%

\subsection{Activity Diagram cho AI Requirement Analysis}

\begin{figure}[H]
\centering
\fbox{
\begin{minipage}{0.95\textwidth}
\centering
\begin{tabular}{c}
\textbf{Bắt đầu} \\
$\downarrow$ \\
\textbf{Product Owner nhập yêu cầu} \\
$\downarrow$ \\
\textbf{Frontend gửi yêu cầu} \\
$\downarrow$ \\
\textbf{Backend xác thực} \\
$\downarrow$ \\
\textbf{AI Orchestrator tiếp nhận} \\
$\downarrow$ \\
\textbf{Requirement Agent phân tích} \\
$\downarrow$ \\
\textbf{Phát hiện vấn đề?} \\
$\swarrow$ \hspace{2.5cm} $\searrow$ \\
\textbf{Có} \hspace{2.5cm} \textbf{Không} \\
$\downarrow$ \hspace{2.5cm} $\downarrow$ \\
\textbf{Đề xuất câu hỏi làm rõ} \hspace{1cm} \textbf{Sinh User Story} \\
$\downarrow$ \hspace{2.5cm} $\downarrow$ \\
\textbf{Chờ Product Owner phản hồi} \hspace{0.5cm} \textbf{Sinh Acceptance Criteria} \\
$\downarrow$ \hspace{2.5cm} $\downarrow$ \\
\textbf{Quay lại phân tích} \hspace{0.5cm} \textbf{Product Owner phê duyệt} \\
\hspace{2.5cm} $\downarrow$ \\
\hspace{2.5cm} \textbf{Lưu vào Product Backlog} \\
\hspace{2.5cm} $\downarrow$ \\
\hspace{2.5cm} \textbf{Kết thúc}
\end{tabular}
\end{minipage}
}
\caption{Activity Diagram - AI Requirement Analysis}
\label{fig:activity_requirement}
\end{figure}

\subsection{Activity Diagram cho AI Story Point Estimation}

\begin{figure}[H]
\centering
\fbox{
\begin{minipage}{0.95\textwidth}
\centering
\begin{tabular}{c}
\textbf{Bắt đầu} \\
$\downarrow$ \\
\textbf{Scrum Master chọn User Story} \\
$\downarrow$ \\
\textbf{Truy xuất dữ liệu lịch sử} \\
$\downarrow$ \\
\textbf{Phân tích các tiêu chí} \\
$\downarrow$ \\
\textbf{Đánh giá Complexity} \\
$\downarrow$ \\
\textbf{Đánh giá Business Logic} \\
$\downarrow$ \\
\textbf{Đánh giá Integration} \\
$\downarrow$ \\
\textbf{Đánh giá UI/UX} \\
$\downarrow$ \\
\textbf{Đánh giá Testing} \\
$\downarrow$ \\
\textbf{Đánh giá Uncertainty} \\
$\downarrow$ \\
\textbf{Tính điểm Story Point} \\
$\downarrow$ \\
\textbf{Sinh giải thích} \\
$\downarrow$ \\
\textbf{Hiển thị kết quả} \\
$\downarrow$ \\
\textbf{Scrum Master xác nhận} \\
$\downarrow$ \\
\textbf{Kết thúc}
\end{tabular}
\end{minipage}
}
\caption{Activity Diagram - AI Story Point Estimation}
\label{fig:activity_storypoint}
\end{figure}

%%%%%%%%%%%%%%%%%%%%%%%%%%%%%%%%%%%%%%%%%%%%%%%%%%%%%%%%%%%%%%
\section{Sơ đồ tuần tự (Sequence Diagram)}
%%%%%%%%%%%%%%%%%%%%%%%%%%%%%%%%%%%%%%%%%%%%%%%%%%%%%%%%%%%%%%

\subsection{Sequence Diagram - Đăng nhập}

\begin{figure}[H]
\centering
\fbox{
\begin{minipage}{0.95\textwidth}
\centering
\begin{tabular}{c}
\textbf{User} $\rightarrow$ \textbf{Frontend} $\rightarrow$ \textbf{Backend} $\rightarrow$ \textbf{Database} \\
\\
1. Nhập email/password \\
2. Gửi request \\
3. Kiểm tra thông tin \\
4. Truy vấn user \\
5. Trả về user \\
6. So sánh password \\
7. Tạo JWT \\
8. Trả token \\
9. Lưu token, chuyển trang
\end{tabular}
\end{minipage}
}
\caption{Sequence Diagram - Đăng nhập}
\label{fig:sequence_login}
\end{figure}

\subsection{Sequence Diagram - AI Requirement Analysis}

\begin{figure}[H]
\centering
\fbox{
\begin{minipage}{0.95\textwidth}
\centering
\begin{tabular}{c}
\textbf{PO} $\rightarrow$ \textbf{Frontend} $\rightarrow$ \textbf{Backend} $\rightarrow$ \textbf{AI Orchestrator} $\rightarrow$ \textbf{Agent} $\rightarrow$ \textbf{LLM} \\
\\
1. Nhập yêu cầu \\
2. Gửi request \\
3. Xác thực \\
4. Điều phối \\
5. Build Prompt \\
6. Gọi API \\
7. Phân tích \\
8. Trả kết quả \\
9. Validate \\
10. Lưu DB \\
11. Trả kết quả \\
12. Hiển thị
\end{tabular}
\end{minipage}
}
\caption{Sequence Diagram - AI Requirement Analysis}
\label{fig:sequence_ai_requirement}
\end{figure}

\subsection{Sequence Diagram - AI Sprint Forecast}

\begin{figure}[H]
\centering
\fbox{
\begin{minipage}{0.95\textwidth}
\centering
\begin{tabular}{c}
\textbf{SM} $\rightarrow$ \textbf{Frontend} $\rightarrow$ \textbf{Backend} $\rightarrow$ \textbf{AI Orchestrator} $\rightarrow$ \textbf{Forecast Agent} \\
\\
1. Yêu cầu dự báo \\
2. Gửi request \\
3. Xác thực \\
4. Điều phối \\
5. Lấy dữ liệu Sprint \\
6. Phân tích \\
7. Tính toán dự báo \\
8. Trả kết quả \\
9. Hiển thị
\end{tabular}
\end{minipage}
}
\caption{Sequence Diagram - AI Sprint Forecast}
\label{fig:sequence_forecast}
\end{figure}

%%%%%%%%%%%%%%%%%%%%%%%%%%%%%%%%%%%%%%%%%%%%%%%%%%%%%%%%%%%%%%
\section{Mối quan hệ giữa các thành phần nghiệp vụ}
%%%%%%%%%%%%%%%%%%%%%%%%%%%%%%%%%%%%%%%%%%%%%%%%%%%%%%%%%%%%%%

\subsection{Sơ đồ quan hệ thực thể - liên kết (ERD)}

\begin{figure}[H]
\centering
\fbox{
\begin{minipage}{0.95\textwidth}
\centering
\begin{tabular}{c}
\textbf{User} \\
$\downarrow$ 1--N \\
\textbf{Workspace} \\
$\downarrow$ 1--N \\
\textbf{Project} \\
$\swarrow$ 1--N \hspace{3cm} $\searrow$ 1--N \\
\textbf{Sprint} \hspace{3.5cm} \textbf{Epic} \\
$\downarrow$ 1--N \hspace{3.5cm} $\downarrow$ 1--N \\
\textbf{User Story} \\
$\downarrow$ 1--N \\
\textbf{Task} \\
$\downarrow$ N--1 \\
\textbf{User (Assignee)}
\end{tabular}
\end{minipage}
}
\caption{Quan hệ giữa các thực thể nghiệp vụ}
\label{fig:entity_relationships}
\end{figure}

\subsection{Bảng ánh xạ yêu cầu chức năng với AI Agent}

\begin{table}[H]
\centering
\caption{Ánh xạ yêu cầu chức năng với AI Agent}
\begin{tabular}{|p{4.5cm}|p{4.5cm}|p{4.5cm}|}
\hline
\textbf{Yêu cầu} & \textbf{AI Agent} & \textbf{Công nghệ} \\
\hline
Phân tích yêu cầu & Requirement Agent & LLM + NLP \\
\hline
Sinh User Story & Planning Agent & LLM + Prompt Eng \\
\hline
Ước lượng Story Point & Story Point Agent & LLM + Dữ liệu lịch sử \\
\hline
Dự báo Sprint & Forecast Agent & Machine Learning \\
\hline
Phát hiện rủi ro & Risk Agent & Rule-based + ML \\
\hline
Phân công Task & Assignment Agent & Optimization \\
\hline
Sinh tài liệu & Documentation Agent & LLM + RAG \\
\hline
Dashboard & Dashboard Agent & Data Visualization \\
\hline
\end{tabular}
\label{tab:requirement_agent_mapping}
\end{table}

%%%%%%%%%%%%%%%%%%%%%%%%%%%%%%%%%%%%%%%%%%%%%%%%%%%%%%%%%%%%%%
\section{Quy tắc nghiệp vụ (Business Rules)}
%%%%%%%%%%%%%%%%%%%%%%%%%%%%%%%%%%%%%%%%%%%%%%%%%%%%%%%%%%%%%%

Hệ thống tuân thủ các quy tắc nghiệp vụ sau:

\begin{enumerate}
\item \textbf{BR-01}: Một Workspace có thể chứa nhiều Project.
\item \textbf{BR-02}: Một Project thuộc duy nhất một Workspace.
\item \textbf{BR-03}: Một Sprint thuộc duy nhất một Project.
\item \textbf{BR-04}: Một User Story chỉ thuộc một Sprint tại một thời điểm.
\item \textbf{BR-05}: Một Task chỉ thuộc một User Story.
\item \textbf{BR-06}: Story Point chỉ được chỉnh sửa trước khi Sprint bắt đầu.
\item \textbf{BR-07}: AI chỉ đưa ra đề xuất, không tự động thay đổi dữ liệu nếu chưa có xác nhận của người dùng.
\item \textbf{BR-08}: Mọi thao tác của AI đều được ghi lại trong Audit Log.
\item \textbf{BR-09}: Chỉ Project Manager hoặc Workspace Owner mới có quyền xác nhận các đề xuất quan trọng của AI.
\item \textbf{BR-10}: Tất cả các phản hồi từ AI phải được kiểm tra định dạng (JSON Schema Validation) trước khi lưu vào cơ sở dữ liệu.
\item \textbf{BR-11}: Mỗi Sprint phải có Sprint Goal rõ ràng.
\item \textbf{BR-12}: User Story chỉ có thể được gán vào Sprint nếu đã có Story Point.
\end{enumerate}

%%%%%%%%%%%%%%%%%%%%%%%%%%%%%%%%%%%%%%%%%%%%%%%%%%%%%%%%%%%%%%
\section{Tổng kết chương}
%%%%%%%%%%%%%%%%%%%%%%%%%%%%%%%%%%%%%%%%%%%%%%%%%%%%%%%%%%%%%%

Trong chương này, luận văn đã trình bày phân tích nghiệp vụ chi tiết của hệ thống AI Project Management.

Các nội dung chính đã được thực hiện:

\begin{enumerate}
\item Xác định 5 nhóm tác nhân với vai trò và quyền hạn rõ ràng.
\item Xây dựng sơ đồ Use Case tổng quan và đặc tả chi tiết 8 Use Case quan trọng nhất.
\item Phân tích 4 luồng nghiệp vụ chính: báo cáo yêu cầu, lập kế hoạch Sprint, thực hiện Sprint, đánh giá Sprint.
\item Xây dựng sơ đồ hoạt động (Activity Diagram) cho các nghiệp vụ quan trọng.
\item Xây dựng sơ đồ tuần tự (Sequence Diagram) thể hiện tương tác giữa các thành phần.
\item Phân tích mối quan hệ giữa các thực thể nghiệp vụ.
\item Đề xuất 12 quy tắc nghiệp vụ cho hệ thống.
\end{enumerate}

Đây là cơ sở để Chương 7 tiến hành thiết kế hệ thống chi tiết.

%%%%%%%%%%%%%%%%%%%%%%%%%%%%%%%%%%%%%%%%%%%%%%%%%%%%%%%%%%%%%%

%%%%%%%%%%%%%%%%%%%%%%%%%%%%%%%%%%%%%%%%%%%%%%%%%%%%%%%%%%%%%%
%%%%%%%%%%%%%%%%%%%%%%%%%%%%%%%%%%%%%%%%%%%%%%%%%%%%%%%%%%%%%%
\chapter{THIẾT KẾ HỆ THỐNG}
\label{chap:system_design}
%%%%%%%%%%%%%%%%%%%%%%%%%%%%%%%%%%%%%%%%%%%%%%%%%%%%%%%%%%%%%%

\section{Giới thiệu chương}

Sau khi hoàn thành quá trình phân tích yêu cầu (Chương 5) và phân tích nghiệp vụ (Chương 6), bước tiếp theo là thiết kế hệ thống. Mục tiêu của chương này là chuyển đổi các yêu cầu đã xác định thành một kiến trúc phần mềm hoàn chỉnh, có khả năng triển khai thực tế.

Khác với các hệ thống quản lý dự án truyền thống, nền tảng được đề xuất trong luận văn tích hợp trí tuệ nhân tạo vào hầu hết các quy trình nghiệp vụ. Vì vậy, kiến trúc hệ thống không chỉ bao gồm các thành phần Backend, Frontend và Database mà còn bổ sung tầng AI Service chịu trách nhiệm điều phối các AI Agent, quản lý Prompt, truy xuất tri thức và tương tác với Large Language Models.

Nội dung chương này trình bày:

\begin{itemize}
\item Kiến trúc tổng thể của hệ thống.
\item Kiến trúc phân lớp (Layered Architecture).
\item Kiến trúc Microservices.
\item Kiến trúc AI Service và AI Orchestrator.
\item Thiết kế cơ sở dữ liệu.
\item Thiết kế API RESTful.
\item Thiết kế AI Pipeline.
\item Thiết kế Dashboard.
\item Các nguyên tắc thiết kế áp dụng.
\end{itemize}

%%%%%%%%%%%%%%%%%%%%%%%%%%%%%%%%%%%%%%%%%%%%%%%%%%%%%%%%%%%%%%
\section{Kiến trúc tổng thể}
%%%%%%%%%%%%%%%%%%%%%%%%%%%%%%%%%%%%%%%%%%%%%%%%%%%%%%%%%%%%%%

\subsection{Mô hình kiến trúc}

Hệ thống được thiết kế theo mô hình Client-Server kết hợp với kiến trúc Microservices nhằm đảm bảo tính mở rộng và khả năng tích hợp AI.

\begin{figure}[H]
\centering
\fbox{
\begin{minipage}{0.95\textwidth}
\centering
\begin{tabular}{c}
\textbf{Trình duyệt (Browser)} \\
$\downarrow$ \\
\textbf{Frontend (React + TypeScript)} \\
$\downarrow$ \\
\textbf{API Gateway (Spring Boot)} \\
$\swarrow$ $\downarrow$ $\searrow$ \\
\textbf{Authentication Service} \hspace{0.5cm} \textbf{Project Service} \hspace{0.5cm} \textbf{AI Service} \\
$\downarrow$ $\downarrow$ $\downarrow$ \\
\textbf{PostgreSQL} \hspace{2cm} \textbf{Vector DB (pgvector)} \\
$\downarrow$ \\
\textbf{LLM (OpenAI / Ollama)}
\end{tabular}
\end{minipage}
}
\caption{Kiến trúc tổng thể của hệ thống}
\label{fig:overall_architecture}
\end{figure}

\subsection{Các thành phần chính}

\begin{table}[H]
\centering
\caption{Các thành phần chính của hệ thống}
\begin{tabular}{|p{4cm}|p{4cm}|p{7cm}|}
\hline
\textbf{Thành phần} & \textbf{Công nghệ} & \textbf{Vai trò} \\
\hline
Frontend & React + TypeScript & Giao diện người dùng, tương tác với API \\
\hline
API Gateway & Spring Boot & Điều phối request, xác thực, routing \\
\hline
Authentication Service & Spring Security + JWT & Xác thực và phân quyền \\
\hline
Project Service & Spring Boot & Quản lý Project, Sprint, Task \\
\hline
AI Service & FastAPI (Python) & Điều phối AI Agent, suy luận \\
\hline
Database & PostgreSQL 16 & Lưu trữ dữ liệu nghiệp vụ \\
\hline
Vector DB & pgvector & Lưu embedding cho RAG \\
\hline
LLM & OpenAI API / Ollama & Xử lý ngôn ngữ tự nhiên \\
\hline
Cache & Redis & Cache dữ liệu, session \\
\hline
Message Queue & RabbitMQ & Xử lý bất đồng bộ \\
\hline
\end{tabular}
\label{tab:components}
\end{table}

%%%%%%%%%%%%%%%%%%%%%%%%%%%%%%%%%%%%%%%%%%%%%%%%%%%%%%%%%%%%%%
\section{Kiến trúc phân lớp (Layered Architecture)}
%%%%%%%%%%%%%%%%%%%%%%%%%%%%%%%%%%%%%%%%%%%%%%%%%%%%%%%%%%%%%%

Backend được xây dựng theo mô hình nhiều lớp (Layered Architecture) kết hợp với Clean Architecture để đảm bảo tính độc lập và dễ bảo trì.

\begin{figure}[H]
\centering
\fbox{
\begin{minipage}{0.8\textwidth}
\centering
\begin{tabular}{c}
\textbf{Presentation Layer (Controller)} \\
$\downarrow$ \\
\textbf{Application Layer (Service)} \\
$\downarrow$ \\
\textbf{Domain Layer (Entity, Value Object)} \\
$\downarrow$ \\
\textbf{Infrastructure Layer (Repository, External)} \\
$\downarrow$ \\
\textbf{Database}
\end{tabular}
\end{minipage}
}
\caption{Kiến trúc phân lớp}
\label{fig:layer_architecture}
\end{figure}

\subsection{Presentation Layer (Tầng trình bày)}

Lớp này chịu trách nhiệm tiếp nhận và trả lời các yêu cầu từ phía Client.

\begin{itemize}
\item \textbf{REST Controller}: Tiếp nhận HTTP Request, trả về HTTP Response.
\item \textbf{JWT Filter}: Xác thực token cho mọi request.
\item \textbf{Exception Handler}: Xử lý ngoại lệ tập trung.
\item \textbf{Request Validation}: Kiểm tra dữ liệu đầu vào.
\end{itemize}

\textbf{Nguyên tắc}: Controller chỉ thực hiện xác thực và điều phối, không chứa logic nghiệp vụ.

\begin{verbatim}
@RestController
@RequestMapping("/api/v1/projects")
public class ProjectController {
    
    @Autowired
    private ProjectService projectService;
    
    @PostMapping
    public ResponseEntity<ProjectDTO> createProject(@RequestBody CreateProjectRequest request) {
        return ResponseEntity.ok(projectService.createProject(request));
    }
}
\end{verbatim}

\subsection{Application Layer (Tầng ứng dụng)}

Lớp này hiện thực các Use Case của hệ thống và điều phối các nghiệp vụ.

\begin{itemize}
\item \textbf{Service Classes}: Hiện thực logic nghiệp vụ.
\item \textbf{Use Case Implementations}: Mỗi Use Case là một phương thức trong Service.
\item \textbf{Transaction Management}: Quản lý giao dịch cơ sở dữ liệu.
\end{itemize}

\textbf{Các Service chính:}

\begin{itemize}
\item AuthenticationService
\item WorkspaceService
\item ProjectService
\item SprintService
\item BacklogService (Epic, User Story)
\item TaskService
\item AIService
\item DashboardService
\item ReportService
\item NotificationService
\end{itemize}

\subsection{Domain Layer (Tầng miền)}

Đây là tầng quan trọng nhất, chứa các thực thể nghiệp vụ và quy tắc kinh doanh.

\begin{itemize}
\item \textbf{Entity}: Các đối tượng nghiệp vụ (User, Project, Sprint, v.v.)
\item \textbf{Value Object}: Các đối tượng bất biến (Email, StoryPoint, v.v.)
\item \textbf{Aggregate Root}: Điểm vào cho các thực thể (Project, Sprint)
\item \textbf{Domain Service}: Các logic nghiệp vụ không thuộc về một Entity cụ thể.
\item \textbf{Repository Interface}: Định nghĩa truy xuất dữ liệu.
\item \textbf{Business Rules}: Các quy tắc nghiệp vụ (Business Rules).
\end{itemize}

\textbf{Ví dụ Domain Entity:}

\begin{verbatim}
@Entity
@Table(name = "projects")
public class Project {
    @Id
    private UUID id;
    private String name;
    private String description;
    @Enumerated(EnumType.STRING)
    private ProjectStatus status;
    private LocalDate startDate;
    private LocalDate endDate;
    
    // Business method
    public void startProject() {
        if (this.status != ProjectStatus.PLANNING) {
            throw new IllegalStateException("Project must be in PLANNING status to start");
        }
        this.status = ProjectStatus.ACTIVE;
        this.startDate = LocalDate.now();
    }
}
\end{verbatim}

\subsection{Infrastructure Layer (Tầng hạ tầng)}

Lớp này chịu trách nhiệm kết nối với các thành phần bên ngoài.

\begin{itemize}
\item \textbf{Repository Implementation}: Hiện thực truy xuất database.
\item \textbf{External API Clients}: Gọi API bên ngoài (AI Service, v.v.)
\item \textbf{Message Queue}: Gửi/nhận thông điệp.
\item \textbf{Cache}: Quản lý cache (Redis).
\item \textbf{File Storage}: Lưu trữ file.
\item \textbf{Email Service}: Gửi email thông báo.
\end{itemize}

%%%%%%%%%%%%%%%%%%%%%%%%%%%%%%%%%%%%%%%%%%%%%%%%%%%%%%%%%%%%%%
\section{Kiến trúc Microservices}
%%%%%%%%%%%%%%%%%%%%%%%%%%%%%%%%%%%%%%%%%%%%%%%%%%%%%%%%%%%%%%

\subsection{Phân chia dịch vụ}

Do số lượng chức năng lớn và có nhiều AI Agent, hệ thống được chia thành các dịch vụ độc lập.

\begin{table}[H]
\centering
\caption{Các Microservice trong hệ thống}
\begin{tabular}{|p{4.5cm}|p{9cm}|}
\hline
\textbf{Service} & \textbf{Chức năng} \\
\hline
Authentication Service & Xác thực, phân quyền, JWT management \\
\hline
Workspace Service & Quản lý Workspace và thành viên \\
\hline
Project Service & Quản lý Project và Sprint \\
\hline
Backlog Service & Quản lý Epic, User Story và Task \\
\hline
AI Service & Điều phối AI Agent và suy luận \\
\hline
Dashboard Service & Tổng hợp dữ liệu Dashboard \\
\hline
Report Service & Sinh báo cáo và tài liệu \\
\hline
Notification Service & Gửi email và thông báo \\
\hline
Analytics Service & Phân tích dữ liệu và KPI \\
\hline
\end{tabular}
\label{tab:microservices}
\end{table}

\subsection{Sơ đồ giao tiếp giữa các Service}

\begin{figure}[H]
\centering
\fbox{
\begin{minipage}{0.95\textwidth}
\centering
\begin{tabular}{c}
\textbf{Frontend} \\
$\downarrow$ \\
\textbf{API Gateway (Spring Cloud Gateway)} \\
$\swarrow$ $\downarrow$ $\searrow$ \hspace{1cm} $\searrow$ \\
\textbf{Auth Service} \hspace{0.3cm} \textbf{Project Service} \hspace{0.3cm} \textbf{Backlog Service} \hspace{0.3cm} \textbf{AI Service} \\
$\downarrow$ \hspace{0.8cm} $\downarrow$ \hspace{0.8cm} $\downarrow$ \hspace{0.8cm} $\downarrow$ \\
\textbf{PostgreSQL} \hspace{0.8cm} \textbf{PostgreSQL} \hspace{0.8cm} \textbf{PostgreSQL} \hspace{0.3cm} \textbf{Vector DB}
\end{tabular}
\end{minipage}
}
\caption{Sơ đồ giao tiếp giữa các Microservice}
\label{fig:microservices_communication}
\end{figure}

\subsection{Lợi ích của kiến trúc Microservices}

\begin{itemize}
\item \textbf{Độc lập}: Mỗi service có thể được phát triển, triển khai và mở rộng độc lập.
\item \textbf{Chịu lỗi}: Một service bị lỗi không ảnh hưởng đến toàn bộ hệ thống.
\item \textbf{Mở rộng}: Có thể mở rộng từng service dựa trên nhu cầu.
\item \textbf{Công nghệ đa dạng}: Mỗi service có thể sử dụng công nghệ phù hợp.
\item \textbf{Tổ chức theo nghiệp vụ}: Phân chia dựa trên domain (DDD).
\end{itemize}

%%%%%%%%%%%%%%%%%%%%%%%%%%%%%%%%%%%%%%%%%%%%%%%%%%%%%%%%%%%%%%
\section{Kiến trúc AI Service và AI Orchestrator}
%%%%%%%%%%%%%%%%%%%%%%%%%%%%%%%%%%%%%%%%%%%%%%%%%%%%%%%%%%%%%%

\subsection{Kiến trúc AI Service}

AI Service là thành phần trung tâm của hệ thống, được xây dựng bằng FastAPI (Python).

\begin{figure}[H]
\centering
\fbox{
\begin{minipage}{0.95\textwidth}
\centering
\begin{tabular}{c}
\textbf{REST API (FastAPI)} \\
$\downarrow$ \\
\textbf{AI Orchestrator} \\
$\swarrow$ $\downarrow$ $\searrow$ \hspace{1cm} $\searrow$ \\
\textbf{Requirement Agent} \hspace{0.2cm} \textbf{Story Point Agent} \hspace{0.2cm} \textbf{Forecast Agent} \hspace{0.2cm} \textbf{Risk Agent} \\
$\downarrow$ \hspace{0.8cm} $\downarrow$ \hspace{0.8cm} $\downarrow$ \hspace{0.8cm} $\downarrow$ \\
\textbf{Prompt Builder} \hspace{0.8cm} \textbf{Prompt Builder} \hspace{0.8cm} \textbf{Prompt Builder} \hspace{0.8cm} \textbf{Prompt Builder} \\
$\downarrow$ \hspace{0.8cm} $\downarrow$ \hspace{0.8cm} $\downarrow$ \hspace{0.8cm} $\downarrow$ \\
\textbf{RAG Service} \\
$\downarrow$ \\
\textbf{LLM (OpenAI / Ollama)}
\end{tabular}
\end{minipage}
}
\caption{Kiến trúc AI Service}
\label{fig:ai_service_architecture}
\end{figure}

\subsection{AI Orchestrator}

AI Orchestrator là thành phần điều phối toàn bộ các AI Agent.

\textbf{Nhiệm vụ chính:}

\begin{enumerate}
\item Nhận yêu cầu từ Backend.
\item Xác định loại nghiệp vụ (Intent Detection).
\item Lựa chọn Agent phù hợp.
\item Xây dựng Prompt với Context.
\item Truy xuất dữ liệu liên quan (RAG).
\item Gọi LLM và nhận kết quả.
\item Kiểm tra và validate kết quả.
\item Lưu kết quả vào cơ sở dữ liệu.
\item Trả kết quả về Backend.
\end{enumerate}

\textbf{Thuật toán điều phối:}

\begin{verbatim}
def orchestrate(request):
    # 1. Detect intent
    intent = detect_intent(request)
    
    # 2. Select agent
    agent = select_agent(intent)
    
    # 3. Retrieve context (RAG)
    context = rag_service.retrieve(request)
    
    # 4. Build prompt
    prompt = prompt_builder.build(agent, request, context)
    
    # 5. Invoke LLM
    response = llm_client.invoke(prompt)
    
    # 6. Validate output
    validated = json_validator.validate(response)
    
    # 7. Apply business rules
    result = business_validator.validate(validated)
    
    # 8. Save to database
    db.save(result)
    
    return result
\end{verbatim}

\subsection{Các AI Agent}

\subsubsection{Requirement Agent}

\begin{itemize}
\item \textbf{Input}: Văn bản yêu cầu, Product Vision, Existing Backlog
\item \textbf{Output}: Requirement Summary, Missing Requirements, Suggested User Stories
\item \textbf{Prompt Strategy}: System Prompt + Business Context + Project Context
\end{itemize}

\subsubsection{Planning Agent}

\begin{itemize}
\item \textbf{Input}: User Story, Sprint, Team Skills
\item \textbf{Output}: Task List, Dependencies, Estimates
\item \textbf{Prompt Strategy}: Phân rã User Story thành Task kỹ thuật
\end{itemize}

\subsubsection{Story Point Agent}

\begin{itemize}
\item \textbf{Input}: User Story, Sprint History, Velocity
\item \textbf{Output}: Story Point, Explanation, Confidence Score
\item \textbf{Evaluation Criteria}: Complexity, Business Logic, Integration, UI, Testing, Uncertainty
\end{itemize}

\subsubsection{Forecast Agent}

\begin{itemize}
\item \textbf{Input}: Sprint History, Velocity, Burndown Data, Remaining SP
\item \textbf{Output}: Finish Date, Completion Probability, Remaining Sprint
\item \textbf{Method}: Kết hợp Machine Learning và dữ liệu lịch sử
\end{itemize}

\subsubsection{Risk Agent}

\begin{itemize}
\item \textbf{Input}: Task Progress, Velocity, Bug Count, Workload
\item \textbf{Output}: Risk Score, Risk Level, Warnings, Recommendations
\item \textbf{Method}: Rule-based + Weighted Scoring
\end{itemize}

\subsubsection{Assignment Agent}

\begin{itemize}
\item \textbf{Input}: Task Requirements, Developer Skills, Workload
\item \textbf{Output}: Recommended Developer, Score, Reasoning
\item \textbf{Method}: Optimization (Skill Match + Experience - Workload)
\end{itemize}

\subsubsection{Documentation Agent}

\begin{itemize}
\item \textbf{Input}: Project Data, Sprint Data, Task Data
\item \textbf{Output}: Sprint Report, Meeting Minutes, Release Note
\item \textbf{Method}: LLM + Template-based Generation
\end{itemize}

\subsubsection{Dashboard Agent}

\begin{itemize}
\item \textbf{Input}: All Project Data
\item \textbf{Output}: KPI Summary, AI Insights, Recommendations
\item \textbf{Method}: Data Aggregation + NLP
\end{itemize}

%%%%%%%%%%%%%%%%%%%%%%%%%%%%%%%%%%%%%%%%%%%%%%%%%%%%%%%%%%%%%%
\section{Thiết kế cơ sở dữ liệu}
%%%%%%%%%%%%%%%%%%%%%%%%%%%%%%%%%%%%%%%%%%%%%%%%%%%%%%%%%%%%%%

\subsection{Mô hình dữ liệu tổng thể}

Cơ sở dữ liệu được thiết kế theo chuẩn 3NF, sử dụng PostgreSQL 16.

\begin{figure}[H]
\centering
\fbox{
\begin{minipage}{0.95\textwidth}
\centering
\begin{tabular}{c}
\textbf{users} \\
$\downarrow$ 1--N \\
\textbf{workspaces} \\
$\downarrow$ 1--N \\
\textbf{projects} \\
$\swarrow$ 1--N \hspace{3cm} $\searrow$ 1--N \\
\textbf{sprints} \hspace{3.5cm} \textbf{epics} \\
$\downarrow$ 1--N \hspace{3.5cm} $\downarrow$ 1--N \\
\textbf{user\_stories} \\
$\downarrow$ 1--N \\
\textbf{tasks} \\
$\uparrow$ \\
\textbf{users (assignee)}
\end{tabular}
\end{minipage}
}
\caption{Mô hình dữ liệu tổng thể}
\label{fig:data_model}
\end{figure}

\subsection{Các bảng dữ liệu chính}

\subsubsection{Bảng Users}

\begin{table}[H]
\centering
\caption{Bảng Users}
\begin{tabular}{|p{4cm}|p{3.5cm}|p{6cm}|}
\hline
\textbf{Trường} & \textbf{Loại dữ liệu} & \textbf{Mô tả} \\
\hline
id & UUID & Khóa chính \\
\hline
email & VARCHAR(255) & Email đăng nhập (unique) \\
\hline
password & VARCHAR(255) & BCrypt hash \\
\hline
fullname & VARCHAR(255) & Họ và tên \\
\hline
avatar & TEXT & URL ảnh đại diện \\
\hline
role & ENUM & ADMIN, PROJECT\_MANAGER, DEVELOPER, TESTER \\
\hline
status & ENUM & ACTIVE, INACTIVE, BLOCKED \\
\hline
created\_at & TIMESTAMP & Thời gian tạo \\
\hline
updated\_at & TIMESTAMP & Thời gian cập nhật \\
\hline
last\_login & TIMESTAMP & Lần đăng nhập cuối \\
\hline
\end{tabular}
\label{tab:users_table}
\end{table}

\subsubsection{Bảng Workspaces}

\begin{table}[H]
\centering
\caption{Bảng Workspaces}
\begin{tabular}{|p{4cm}|p{3.5cm}|p{6cm}|}
\hline
\textbf{Trường} & \textbf{Loại dữ liệu} & \textbf{Mô tả} \\
\hline
id & UUID & Khóa chính \\
\hline
name & VARCHAR(255) & Tên Workspace \\
\hline
owner\_id & UUID & Khóa ngoại đến users \\
\hline
visibility & ENUM & PUBLIC, PRIVATE \\
\hline
logo & TEXT & URL logo \\
\hline
description & TEXT & Mô tả Workspace \\
\hline
created\_at & TIMESTAMP & Thời gian tạo \\
\hline
updated\_at & TIMESTAMP & Thời gian cập nhật \\
\hline
\end{tabular}
\label{tab:workspaces_table}
\end{table}

\subsubsection{Bảng Projects}

\begin{table}[H]
\centering
\caption{Bảng Projects}
\begin{tabular}{|p{4cm}|p{3.5cm}|p{6cm}|}
\hline
\textbf{Trường} & \textbf{Loại dữ liệu} & \textbf{Mô tả} \\
\hline
id & UUID & Khóa chính \\
\hline
workspace\_id & UUID & Khóa ngoại đến workspaces \\
\hline
name & VARCHAR(255) & Tên dự án \\
\hline
key & VARCHAR(50) & Mã định danh (unique trong workspace) \\
\hline
description & TEXT & Mô tả dự án \\
\hline
status & ENUM & PLANNING, ACTIVE, DONE, ARCHIVED \\
\hline
priority & ENUM & LOW, MEDIUM, HIGH \\
\hline
visibility & ENUM & PUBLIC, PRIVATE \\
\hline
start\_date & DATE & Ngày bắt đầu \\
\hline
end\_date & DATE & Ngày kết thúc dự kiến \\
\hline
created\_at & TIMESTAMP & Thời gian tạo \\
\hline
updated\_at & TIMESTAMP & Thời gian cập nhật \\
\hline
\end{tabular}
\label{tab:projects_table}
\end{table}

\subsubsection{Bảng Sprints}

\begin{table}[H]
\centering
\caption{Bảng Sprints}
\begin{tabular}{|p{4cm}|p{3.5cm}|p{6cm}|}
\hline
\textbf{Trường} & \textbf{Loại dữ liệu} & \textbf{Mô tả} \\
\hline
id & UUID & Khóa chính \\
\hline
project\_id & UUID & Khóa ngoại đến projects \\
\hline
name & VARCHAR(255) & Tên Sprint \\
\hline
goal & TEXT & Sprint Goal \\
\hline
start\_date & DATE & Ngày bắt đầu \\
\hline
end\_date & DATE & Ngày kết thúc \\
\hline
velocity & INTEGER & Velocity (SP hoàn thành) \\
\hline
status & ENUM & PLANNING, ACTIVE, COMPLETED \\
\hline
created\_at & TIMESTAMP & Thời gian tạo \\
\hline
updated\_at & TIMESTAMP & Thời gian cập nhật \\
\hline
\end{tabular}
\label{tab:sprints_table}
\end{table}

\subsubsection{Bảng User Stories}

\begin{table}[H]
\centering
\caption{Bảng User Stories}
\begin{tabular}{|p{4cm}|p{3.5cm}|p{6cm}|}
\hline
\textbf{Trường} & \textbf{Loại dữ liệu} & \textbf{Mô tả} \\
\hline
id & UUID & Khóa chính \\
\hline
epic\_id & UUID & Khóa ngoại đến epics (nullable) \\
\hline
sprint\_id & UUID & Khóa ngoại đến sprints (nullable) \\
\hline
title & VARCHAR(255) & Tiêu đề \\
\hline
description & TEXT & Mô tả \\
\hline
acceptance\_criteria & TEXT & Tiêu chí chấp nhận \\
\hline
story\_point & INTEGER & Story Point (theo Fibonacci) \\
\hline
priority & INTEGER & Độ ưu tiên (1-5) \\
\hline
status & ENUM & TODO, DOING, DONE \\
\hline
created\_at & TIMESTAMP & Thời gian tạo \\
\hline
updated\_at & TIMESTAMP & Thời gian cập nhật \\
\hline
\end{tabular}
\label{tab:user_stories_table}
\end{table}

\subsubsection{Bảng Tasks}

\begin{table}[H]
\centering
\caption{Bảng Tasks}
\begin{tabular}{|p{4cm}|p{3.5cm}|p{6cm}|}
\hline
\textbf{Trường} & \textbf{Loại dữ liệu} & \textbf{Mô tả} \\
\hline
id & UUID & Khóa chính \\
\hline
story\_id & UUID & Khóa ngoại đến user\_stories \\
\hline
assignee\_id & UUID & Khóa ngoại đến users (nullable) \\
\hline
title & VARCHAR(255) & Tiêu đề Task \\
\hline
description & TEXT & Mô tả chi tiết \\
\hline
estimate & DOUBLE & Thời gian ước lượng (giờ) \\
\hline
actual\_time & DOUBLE & Thời gian thực tế (giờ) \\
\hline
status & ENUM & TODO, DOING, DONE \\
\hline
priority & ENUM & LOW, MEDIUM, HIGH \\
\hline
due\_date & DATE & Hạn hoàn thành \\
\hline
created\_at & TIMESTAMP & Thời gian tạo \\
\hline
updated\_at & TIMESTAMP & Thời gian cập nhật \\
\hline
\end{tabular}
\label{tab:tasks_table}
\end{table}

\subsubsection{Bảng AI Suggestions}

\begin{table}[H]
\centering
\caption{Bảng AI Suggestions}
\begin{tabular}{|p{4cm}|p{3.5cm}|p{6cm}|}
\hline
\textbf{Trường} & \textbf{Loại dữ liệu} & \textbf{Mô tả} \\
\hline
id & UUID & Khóa chính \\
\hline
project\_id & UUID & Khóa ngoại đến projects \\
\hline
agent & VARCHAR(50) & Tên AI Agent \\
\hline
prompt & TEXT & Prompt đã gửi \\
\hline
response & JSONB & Kết quả từ AI \\
\hline
confidence & DOUBLE & Độ tin cậy (0-1) \\
\hline
approved & BOOLEAN & Đã được phê duyệt \\
\hline
approved\_by & UUID & Khóa ngoại đến users \\
\hline
created\_at & TIMESTAMP & Thời gian tạo \\
\hline
updated\_at & TIMESTAMP & Thời gian cập nhật \\
\hline
\end{tabular}
\label{tab:ai_suggestions_table}
\end{table}

\subsubsection{Bảng Audit Logs}

\begin{table}[H]
\centering
\caption{Bảng Audit Logs}
\begin{tabular}{|p{4cm}|p{3.5cm}|p{6cm}|}
\hline
\textbf{Trường} & \textbf{Loại dữ liệu} & \textbf{Mô tả} \\
\hline
id & UUID & Khóa chính \\
\hline
user\_id & UUID & Khóa ngoại đến users \\
\hline
action & VARCHAR(50) & Hành động \\
\hline
entity & VARCHAR(50) & Thực thể bị tác động \\
\hline
entity\_id & UUID & ID của thực thể \\
\hline
old\_value & JSONB & Giá trị cũ \\
\hline
new\_value & JSONB & Giá trị mới \\
\hline
ip\_address & VARCHAR(45) & Địa chỉ IP \\
\hline
user\_agent & TEXT & User Agent \\
\hline
created\_at & TIMESTAMP & Thời gian tạo \\
\hline
\end{tabular}
\label{tab:audit_logs_table}
\end{table}

\subsubsection{Bảng Document Embeddings (Vector DB)}

\begin{table}[H]
\centering
\caption{Bảng Document Embeddings}
\begin{tabular}{|p{4cm}|p{3.5cm}|p{6cm}|}
\hline
\textbf{Trường} & \textbf{Loại dữ liệu} & \textbf{Mô tả} \\
\hline
id & UUID & Khóa chính \\
\hline
document\_id & VARCHAR(255) & ID của tài liệu gốc \\
\hline
chunk\_index & INTEGER & Chỉ số đoạn văn \\
\hline
content & TEXT & Nội dung đoạn văn \\
\hline
embedding & VECTOR(1536) & Vector embedding \\
\hline
metadata & JSONB & Metadata bổ sung \\
\hline
created\_at & TIMESTAMP & Thời gian tạo \\
\hline
\end{tabular}
\label{tab:embeddings_table}
\end{table}

%%%%%%%%%%%%%%%%%%%%%%%%%%%%%%%%%%%%%%%%%%%%%%%%%%%%%%%%%%%%%%
\section{Thiết kế API RESTful}
%%%%%%%%%%%%%%%%%%%%%%%%%%%%%%%%%%%%%%%%%%%%%%%%%%%%%%%%%%%%%%

\subsection{Nguyên tắc thiết kế API}

\begin{itemize}
\item Tuân thủ chuẩn RESTful.
\item Sử dụng HTTP Method đúng ngữ nghĩa (GET, POST, PUT, DELETE, PATCH).
\item Sử dụng HTTP Status Code chuẩn.
\item API Versioning: `/api/v1/`.
\item Format: JSON cho Request và Response.
\item Authentication: JWT qua Header `Authorization: Bearer <token>`.
\item Validation: Kiểm tra dữ liệu đầu vào.
\end{itemize}

\subsection{Cấu trúc Response chuẩn}

\begin{verbatim}
// Thành công
{
  "success": true,
  "message": "Operation successful",
  "data": { ... }
}

// Lỗi
{
  "success": false,
  "errorCode": "ERROR_CODE",
  "message": "Error description",
  "details": { ... }
}
\end{verbatim}

\subsection{Danh sách API Endpoints}

\subsubsection{Authentication API}

\begin{table}[H]
\centering
\caption{Authentication API}
\begin{tabular}{|p{2.5cm}|p{4cm}|p{8cm}|}
\hline
\textbf{Method} & \textbf{Endpoint} & \textbf{Mô tả} \\
\hline
POST & /auth/register & Đăng ký tài khoản mới \\
\hline
POST & /auth/login & Đăng nhập, nhận JWT \\
\hline
POST & /auth/logout & Đăng xuất \\
\hline
POST & /auth/refresh & Refresh Access Token \\
\hline
POST & /auth/forgot-password & Gửi email đặt lại mật khẩu \\
\hline
POST & /auth/reset-password & Đặt lại mật khẩu \\
\hline
GET & /auth/profile & Lấy thông tin cá nhân \\
\hline
PUT & /auth/profile & Cập nhật thông tin cá nhân \\
\hline
\end{tabular}
\label{tab:auth_api}
\end{table}

\subsubsection{Workspace API}

\begin{table}[H]
\centering
\caption{Workspace API}
\begin{tabular}{|p{2.5cm}|p{4cm}|p{8cm}|}
\hline
\textbf{Method} & \textbf{Endpoint} & \textbf{Mô tả} \\
\hline
POST & /workspaces & Tạo Workspace mới \\
\hline
GET & /workspaces & Lấy danh sách Workspace \\
\hline
GET & /workspaces/{id} & Lấy chi tiết Workspace \\
\hline
PUT & /workspaces/{id} & Cập nhật Workspace \\
\hline
DELETE & /workspaces/{id} & Xóa Workspace \\
\hline
POST & /workspaces/{id}/invite & Mời thành viên \\
\hline
DELETE & /workspaces/{id}/members/{userId} & Xóa thành viên \\
\hline
PUT & /workspaces/{id}/members/{userId}/role & Cập nhật vai trò \\
\hline
\end{tabular}
\label{tab:workspace_api}
\end{table}

\subsubsection{Project API}

\begin{table}[H]
\centering
\caption{Project API}
\begin{tabular}{|p{2.5cm}|p{4cm}|p{8cm}|}
\hline
\textbf{Method} & \textbf{Endpoint} & \textbf{Mô tả} \\
\hline
POST & /projects & Tạo Project mới \\
\hline
GET & /projects & Lấy danh sách Project \\
\hline
GET & /projects/{id} & Lấy chi tiết Project \\
\hline
PUT & /projects/{id} & Cập nhật Project \\
\hline
DELETE & /projects/{id} & Xóa Project \\
\hline
POST & /projects/{id}/archive & Lưu trữ Project \\
\hline
POST & /projects/{id}/members & Thêm thành viên \\
\hline
DELETE & /projects/{id}/members/{userId} & Xóa thành viên \\
\hline
\end{tabular}
\label{tab:project_api}
\end{table}

\subsubsection{Sprint API}

\begin{table}[H]
\centering
\caption{Sprint API}
\begin{tabular}{|p{2.5cm}|p{4cm}|p{8cm}|}
\hline
\textbf{Method} & \textbf{Endpoint} & \textbf{Mô tả} \\
\hline
POST & /sprints & Tạo Sprint mới \\
\hline
GET & /sprints & Lấy danh sách Sprint \\
\hline
GET & /sprints/{id} & Lấy chi tiết Sprint \\
\hline
PUT & /sprints/{id} & Cập nhật Sprint \\
\hline
DELETE & /sprints/{id} & Xóa Sprint \\
\hline
POST & /sprints/{id}/start & Bắt đầu Sprint \\
\hline
POST & /sprints/{id}/complete & Kết thúc Sprint \\
\hline
GET & /sprints/{id}/burndown & Lấy dữ liệu Burndown Chart \\
\hline
GET & /sprints/{id}/velocity & Lấy Velocity \\
\hline
\end{tabular}
\label{tab:sprint_api}
\end{table}

\subsubsection{Backlog API}

\begin{table}[H]
\centering
\caption{Backlog API (Epic, User Story)}
\begin{tabular}{|p{2.5cm}|p{4cm}|p{8cm}|}
\hline
\textbf{Method} & \textbf{Endpoint} & \textbf{Mô tả} \\
\hline
POST & /epics & Tạo Epic mới \\
\hline
GET & /epics & Lấy danh sách Epic \\
\hline
PUT & /epics/{id} & Cập nhật Epic \\
\hline
DELETE & /epics/{id} & Xóa Epic \\
\hline
POST & /stories & Tạo User Story mới \\
\hline
GET & /stories & Lấy danh sách User Story \\
\hline
GET & /stories/{id} & Lấy chi tiết User Story \\
\hline
PUT & /stories/{id} & Cập nhật User Story \\
\hline
DELETE & /stories/{id} & Xóa User Story \\
\hline
PATCH & /stories/{id}/status & Cập nhật trạng thái \\
\hline
PATCH & /stories/{id}/sprint & Gán vào Sprint \\
\hline
PATCH & /stories/{id}/story-point & Cập nhật Story Point \\
\hline
\end{tabular}
\label{tab:backlog_api}
\end{table}

\subsubsection{Task API}

\begin{table}[H]
\centering
\caption{Task API}
\begin{tabular}{|p{2.5cm}|p{4cm}|p{8cm}|}
\hline
\textbf{Method} & \textbf{Endpoint} & \textbf{Mô tả} \\
\hline
POST & /tasks & Tạo Task mới \\
\hline
GET & /tasks & Lấy danh sách Task \\
\hline
GET & /tasks/{id} & Lấy chi tiết Task \\
\hline
PUT & /tasks/{id} & Cập nhật Task \\
\hline
DELETE & /tasks/{id} & Xóa Task \\
\hline
PATCH & /tasks/{id}/status & Cập nhật trạng thái \\
\hline
PATCH & /tasks/{id}/assignee & Phân công Developer \\
\hline
POST & /tasks/{id}/comments & Thêm comment \\
\hline
GET & /tasks/{id}/comments & Lấy danh sách comment \\
\hline
\end{tabular}
\label{tab:task_api}
\end{table}

\subsubsection{AI API}

\begin{table}[H]
\centering
\caption{AI Service API}
\begin{tabular}{|p{2.5cm}|p{4cm}|p{8cm}|}
\hline
\textbf{Method} & \textbf{Endpoint} & \textbf{Mô tả} \\
\hline
POST & /ai/analyze-requirement & Phân tích yêu cầu \\
\hline
POST & /ai/generate-stories & Sinh User Story \\
\hline
POST & /ai/estimate-story-point & Ước lượng Story Point \\
\hline
POST & /ai/forecast-sprint & Dự báo Sprint \\
\hline
POST & /ai/detect-risk & Phát hiện rủi ro \\
\hline
POST & /ai/suggest-assignment & Đề xuất phân công \\
\hline
POST & /ai/generate-document & Sinh tài liệu \\
\hline
POST & /ai/get-insights & Lấy Insights cho Dashboard \\
\hline
\end{tabular}
\label{tab:ai_api}
\end{table}

\subsubsection{Dashboard API}

\begin{table}[H]
\centering
\caption{Dashboard API}
\begin{tabular}{|p{2.5cm}|p{4cm}|p{8cm}|}
\hline
\textbf{Method} & \textbf{Endpoint} & \textbf{Mô tả} \\
\hline
GET & /dashboard/overview & Lấy thống kê tổng quan \\
\hline
GET & /dashboard/burndown/{sprintId} & Lấy dữ liệu Burndown \\
\hline
GET & /dashboard/velocity/{projectId} & Lấy dữ liệu Velocity \\
\hline
GET & /dashboard/workload/{projectId} & Lấy dữ liệu Workload \\
\hline
GET & /dashboard/risk/{projectId} & Lấy đánh giá rủi ro \\
\hline
GET & /dashboard/insights/{projectId} & Lấy Insights từ AI \\
\hline
\end{tabular}
\label{tab:dashboard_api}
\end{table}

%%%%%%%%%%%%%%%%%%%%%%%%%%%%%%%%%%%%%%%%%%%%%%%%%%%%%%%%%%%%%%
\section{Thiết kế AI Pipeline}
%%%%%%%%%%%%%%%%%%%%%%%%%%%%%%%%%%%%%%%%%%%%%%%%%%%%%%%%%%%%%%

\subsection{Pipeline xử lý AI}

\begin{figure}[H]
\centering
\fbox{
\begin{minipage}{0.95\textwidth}
\centering
\begin{tabular}{c}
\textbf{Input: Văn bản yêu cầu} \\
$\downarrow$ \\
\textbf{Bước 1: Tiền xử lý văn bản} \\
- Tokenization \\
- Lowercasing \\
- Stopwords Removal \\
- Lemmatization \\
$\downarrow$ \\
\textbf{Bước 2: Truy xuất tri thức (RAG)} \\
- Tạo embedding từ truy vấn \\
- Tìm kiếm ngữ nghĩa trong Vector DB \\
- Lấy tài liệu liên quan \\
$\downarrow$ \\
\textbf{Bước 3: Xây dựng Prompt} \\
- System Prompt \\
- Business Context \\
- Project Context \\
- User Request \\
- Output Constraint \\
$\downarrow$ \\
\textbf{Bước 4: Suy luận LLM} \\
- Gọi LLM API \\
- Nhận phản hồi \\
$\downarrow$ \\
\textbf{Bước 5: Hậu xử lý} \\
- JSON Schema Validation \\
- Business Rules Validation \\
- Lưu kết quả \\
$\downarrow$ \\
\textbf{Output: Kết quả phân tích}
\end{tabular}
\end{minipage}
}
\caption{AI Pipeline xử lý}
\label{fig:ai_pipeline}
\end{figure}

\subsection{Prompt Engineering Template}

\begin{verbatim}
## System Prompt
Bạn là một chuyên gia phân tích phần mềm dày dạn kinh nghiệm.
Nhiệm vụ của bạn là phân tích yêu cầu phần mềm và sinh User Story.

## Business Rules
- User Story phải tuân theo mẫu: "As a [Role], I want [Goal], So that [Value]"
- Story Point phải tuân theo dãy Fibonacci: 1, 2, 3, 5, 8, 13, 21
- Mỗi User Story phải có Acceptance Criteria rõ ràng

## Project Context
- Project: {project_name}
- Domain: {domain}
- Technologies: {technologies}
- Existing Features: {existing_features}

## User Request
{user_requirement}

## Output Format
Trả về JSON với cấu trúc sau:
{
  "summary": "Tóm tắt yêu cầu",
  "userStories": [
    {
      "title": "...",
      "role": "...",
      "goal": "...",
      "value": "...",
      "acceptanceCriteria": ["...", "..."],
      "storyPoint": 8,
      "priority": 3
    }
  ],
  "missingInfo": ["Câu hỏi làm rõ 1", "Câu hỏi làm rõ 2"],
  "risks": ["Rủi ro 1", "Rủi ro 2"]
}
\end{verbatim}

%%%%%%%%%%%%%%%%%%%%%%%%%%%%%%%%%%%%%%%%%%%%%%%%%%%%%%%%%%%%%%
\section{Nguyên tắc thiết kế}
%%%%%%%%%%%%%%%%%%%%%%%%%%%%%%%%%%%%%%%%%%%%%%%%%%%%%%%%%%%%%%

\subsection{SOLID Principles}

\begin{enumerate}
\item \textbf{S - Single Responsibility}: Mỗi class chỉ có một trách nhiệm duy nhất.
\item \textbf{O - Open/Closed}: Mở rộng nhưng đóng với sửa đổi.
\item \textbf{L - Liskov Substitution}: Các lớp con có thể thay thế lớp cha.
\item \textbf{I - Interface Segregation}: Không ép class implement interface không cần thiết.
\item \textbf{D - Dependency Inversion}: Phụ thuộc vào abstraction, không phải concrete.
\end{enumerate}

\subsection{Clean Architecture}

\begin{itemize}
\item \textbf{Independent of Frameworks}: Kiến trúc không phụ thuộc vào framework cụ thể.
\item \textbf{Testable}: Có thể kiểm thử độc lập với UI, database, v.v.
\item \textbf{Independent of UI}: UI có thể thay đổi mà không ảnh hưởng đến business rules.
\item \textbf{Independent of Database}: Database có thể thay đổi mà không ảnh hưởng đến business rules.
\item \textbf{Independent of External Agencies}: External services có thể thay đổi mà không ảnh hưởng đến business rules.
\end{itemize}

\subsection{Domain-Driven Design (DDD)}

\begin{itemize}
\item \textbf{Ubiquitous Language}: Sử dụng ngôn ngữ chung thống nhất giữa business và technical.
\item \textbf{Bounded Context}: Phân chia rõ ràng các context nghiệp vụ.
\item \textbf{Entity}: Đối tượng có định danh riêng.
\item \textbf{Value Object}: Đối tượng bất biến, không có định danh.
\item \textbf{Aggregate}: Nhóm các Entity và Value Object.
\item \textbf{Repository}: Cung cấp interface truy xuất Aggregate.
\end{itemize}

%%%%%%%%%%%%%%%%%%%%%%%%%%%%%%%%%%%%%%%%%%%%%%%%%%%%%%%%%%%%%%
\section{Tổng kết chương}
%%%%%%%%%%%%%%%%%%%%%%%%%%%%%%%%%%%%%%%%%%%%%%%%%%%%%%%%%%%%%%

Trong chương này, luận văn đã trình bày thiết kế hệ thống chi tiết cho nền tảng AI Project Management.

Các nội dung chính đã được thực hiện:

\begin{enumerate}
\item \textbf{Kiến trúc tổng thể}: Mô hình Client-Server với Microservices.
\item \textbf{Kiến trúc phân lớp}: 4 tầng (Presentation, Application, Domain, Infrastructure).
\item \textbf{Microservices}: 9 service độc lập, giao tiếp qua API Gateway.
\item \textbf{AI Service}: Kiến trúc Multi-Agent với 8 AI Agent và AI Orchestrator.
\item \textbf{Cơ sở dữ liệu}: 12 bảng chính, sử dụng PostgreSQL + pgvector.
\item \textbf{API RESTful}: 8 nhóm API với hơn 60 endpoints.
\item \textbf{AI Pipeline}: 5 bước xử lý từ input đến output.
\item \textbf{Nguyên tắc thiết kế}: SOLID, Clean Architecture, DDD.
\end{enumerate}

Đây là cơ sở để Chương 8 trình bày chi tiết kiến trúc AI và Multi-Agent, và Chương 9 trình bày thiết kế cơ sở dữ liệu chi tiết.

%%%%%%%%%%%%%%%%%%%%%%%%%%%%%%%%%%%%%%%%%%%%%%%%%%%%%%%%%%%%%%%%%%%%%%%%%%%%%%%%%%%%%%%%%%%%%%%%%%%%%%%%%%%%%%%%%%%%%%%%%%%%
% CHƯƠNG 8: KIẾN TRÚC AI VÀ MULTI-AGENT
%%%%%%%%%%%%%%%%%%%%%%%%%%%%%%%%%%%%%%%%%%%%%%%%%%%%%%%%%%%%%%
%%%%%%%%%%%%%%%%%%%%%%%%%%%%%%%%%%%%%%%%%%%%%%%%%%%%%%%%%%%%%%
\chapter{KIẾN TRÚC AI VÀ MULTI-AGENT}
\label{chap:ai_architecture}
%%%%%%%%%%%%%%%%%%%%%%%%%%%%%%%%%%%%%%%%%%%%%%%%%%%%%%%%%%%%%%

\section{Giới thiệu chương}

Sau khi trình bày thiết kế tổng thể hệ thống ở Chương 7, chương này tập trung vào thành phần cốt lõi nhất của luận văn: kiến trúc AI và hệ thống Multi-Agent. Đây là điểm khác biệt lớn nhất giữa nền tảng được đề xuất và các hệ thống quản lý dự án truyền thống.

Mục tiêu của chương bao gồm:

\begin{itemize}
\item Trình bày kiến trúc tổng thể của AI Service.
\item Mô tả chi tiết các AI Agent và chức năng của từng Agent.
\item Thiết kế AI Orchestrator điều phối các Agent.
\item Trình bày cơ chế Prompt Engineering và RAG (Retrieval-Augmented Generation).
\item Mô tả luồng xử lý dữ liệu trong hệ thống AI.
\item Trình bày các kỹ thuật đảm bảo chất lượng và độ tin cậy của AI.
\end{itemize}

Đây là cơ sở để Chương 9 trình bày thiết kế cơ sở dữ liệu chi tiết và Chương 10 trình bày thiết kế Backend.

%%%%%%%%%%%%%%%%%%%%%%%%%%%%%%%%%%%%%%%%%%%%%%%%%%%%%%%%%%%%%%
\section{Tổng quan về kiến trúc AI}
%%%%%%%%%%%%%%%%%%%%%%%%%%%%%%%%%%%%%%%%%%%%%%%%%%%%%%%%%%%%%%

\subsection{Định hướng thiết kế}

Kiến trúc AI của hệ thống được thiết kế dựa trên các nguyên tắc sau:

\begin{enumerate}
\item \textbf{Human-Centric AI}: AI đóng vai trò trợ lý (Assistant), không thay thế con người.
\item \textbf{Multi-Agent}: Chia nhỏ nhiệm vụ phức tạp thành các Agent chuyên biệt.
\item \textbf{Human-in-the-Loop}: Mọi quyết định quan trọng đều cần sự phê duyệt của con người.
\item \textbf{Explainability}: AI phải giải thích được lý do của các đề xuất.
\item \textbf{RAG (Retrieval-Augmented Generation)}: Sử dụng tri thức nội bộ của dự án.
\item \textbf{Modularity}: Các Agent có thể được thay thế hoặc nâng cấp độc lập.
\end{enumerate}

\subsection{Kiến trúc tổng thể AI}

\begin{figure}[H]
\centering
\fbox{
\begin{minipage}{0.95\textwidth}
\centering
\begin{tabular}{c}
\textbf{Tầng 1: API Gateway (REST API)} \\
$\downarrow$ \\
\textbf{Tầng 2: AI Orchestrator} \\
- Intent Detection \\
- Agent Selection \\
- Workflow Orchestration \\
$\downarrow$ \\
\textbf{Tầng 3: Multi-Agent System} \\
\swarrow $\searrow$ \hspace{0.5cm} $\searrow$ \hspace{0.5cm} $\searrow$ \hspace{0.5cm} $\searrow$ \hspace{0.5cm} $\searrow$ \hspace{0.5cm} $\searrow$ \\
Req Agent \hspace{0.2cm} Planning \hspace{0.2cm} StoryPoint \hspace{0.2cm} Forecast \hspace{0.2cm} Risk \hspace{0.2cm} Assignment \\
$\downarrow$ \hspace{0.8cm} $\downarrow$ \hspace{0.8cm} $\downarrow$ \hspace{0.8cm} $\downarrow$ \hspace{0.8cm} $\downarrow$ \hspace{0.8cm} $\downarrow$ \\
\textbf{Tầng 4: RAG \& Knowledge Base} \\
- Vector Database (pgvector) \\
- Document Indexing \\
- Semantic Search \\
$\downarrow$ \\
\textbf{Tầng 5: Large Language Model} \\
- OpenAI GPT-4 / GPT-5 \\
- Ollama (Llama 3) \\
- DeepSeek \\
$\downarrow$ \\
\textbf{Tầng 6: Output Validation} \\
- JSON Schema Validation \\
- Business Rules Validation \\
- Confidence Scoring
\end{tabular}
\end{minipage}
}
\caption{Kiến trúc tổng thể AI}
\label{fig:ai_overall_architecture}
\end{figure}

%%%%%%%%%%%%%%%%%%%%%%%%%%%%%%%%%%%%%%%%%%%%%%%%%%%%%%%%%%%%%%
\section{AI Orchestrator}
%%%%%%%%%%%%%%%%%%%%%%%%%%%%%%%%%%%%%%%%%%%%%%%%%%%%%%%%%%%%%%

\subsection{Giới thiệu AI Orchestrator}

AI Orchestrator là thành phần trung tâm điều phối toàn bộ hoạt động của hệ thống AI. Nó đóng vai trò như một "bộ não" quyết định Agent nào được gọi, theo thứ tự nào, và cách tổng hợp kết quả.

\subsection{Nhiệm vụ chính}

\begin{enumerate}
\item \textbf{Intent Detection}: Phân tích yêu cầu đầu vào để xác định loại nghiệp vụ.
\item \textbf{Agent Selection}: Lựa chọn Agent hoặc chuỗi Agent phù hợp.
\item \textbf{Context Building}: Thu thập và tổng hợp dữ liệu ngữ cảnh.
\item \textbf{Prompt Construction}: Xây dựng Prompt hoàn chỉnh cho LLM.
\item \textbf{Orchestration}: Điều phối thứ tự thực hiện các Agent.
\item \textbf{Result Aggregation}: Tổng hợp kết quả từ nhiều Agent.
\item \textbf{Quality Assurance}: Kiểm tra chất lượng đầu ra.
\end{enumerate}

\subsection{Sơ đồ hoạt động của AI Orchestrator}

\begin{figure}[H]
\centering
\fbox{
\begin{minipage}{0.95\textwidth}
\centering
\begin{tabular}{c}
\textbf{Bắt đầu: Nhận request} \\
$\downarrow$ \\
\textbf{Intent Detection} \\
$\downarrow$ \\
\textbf{Phân loại intent} \\
$\swarrow$ $\downarrow$ $\searrow$ \\
\textbf{Phân tích yêu cầu} \hspace{0.5cm} \textbf{Ước lượng SP} \hspace{0.5cm} \textbf{Dự báo Sprint} \\
$\downarrow$ \hspace{1cm} $\downarrow$ \hspace{1cm} $\downarrow$ \\
\textbf{Requirement Agent} \hspace{0.5cm} \textbf{Story Point Agent} \hspace{0.5cm} \textbf{Forecast Agent} \\
$\downarrow$ \hspace{1cm} $\downarrow$ \hspace{1cm} $\downarrow$ \\
\textbf{Build Prompt} \hspace{0.8cm} \textbf{Build Prompt} \hspace{0.8cm} \textbf{Build Prompt} \\
$\downarrow$ \hspace{1cm} $\downarrow$ \hspace{1cm} $\downarrow$ \\
\textbf{Gọi LLM} \hspace{1cm} \textbf{Gọi LLM} \hspace{1cm} \textbf{Gọi LLM} \\
$\downarrow$ \hspace{1cm} $\downarrow$ \hspace{1cm} $\downarrow$ \\
\textbf{Validate \& Save} \hspace{0.5cm} \textbf{Validate \& Save} \hspace{0.5cm} \textbf{Validate \& Save} \\
$\downarrow$ \\
\textbf{Trả kết quả}
\end{tabular}
\end{minipage}
}
\caption{Sơ đồ hoạt động AI Orchestrator}
\label{fig:orchestrator_flow}
\end{figure}

\subsection{Thuật toán điều phối}

\begin{verbatim}
class AIOrchestrator:
    def orchestrate(self, request: AIRequest) -> AIResponse:
        # Bước 1: Phát hiện intent
        intent = self.detect_intent(request)
        
        # Bước 2: Xác định workflow
        workflow = self.get_workflow(intent)
        
        # Bước 3: Thu thập context
        context = self.build_context(request)
        
        # Bước 4: Thực thi workflow
        result = {}
        for step in workflow.steps:
            agent = self.get_agent(step.agent_type)
            prompt = self.build_prompt(
                agent=agent,
                request=request,
                context=context,
                previous_result=result
            )
            response = self.invoke_llm(prompt)
            validated = self.validate_response(response)
            result[step.name] = validated
            context.update(validated)
        
        # Bước 5: Tổng hợp và trả kết quả
        return self.aggregate_result(result)
\end{verbatim}

\subsection{Phân loại Intent}

\begin{table}[H]
\centering
\caption{Phân loại Intent và Agent tương ứng}
\begin{tabular}{|p{4.5cm}|p{4.5cm}|p{5.5cm}|}
\hline
\textbf{Intent} & \textbf{Agent chính} & \textbf{Agent hỗ trợ} \\
\hline
Phân tích yêu cầu & Requirement Agent & Planning Agent \\
\hline
Sinh User Story & Planning Agent & Requirement Agent \\
\hline
Phân rã Task & Planning Agent & - \\
\hline
Ước lượng Story Point & Story Point Agent & - \\
\hline
Dự báo Sprint & Forecast Agent & Risk Agent \\
\hline
Phát hiện rủi ro & Risk Agent & Forecast Agent \\
\hline
Phân công Task & Assignment Agent & - \\
\hline
Sinh tài liệu & Documentation Agent & - \\
\hline
Dashboard Insights & Dashboard Agent & Risk Agent, Forecast Agent \\
\hline
\end{tabular}
\label{tab:intent_mapping}
\end{table}

%%%%%%%%%%%%%%%%%%%%%%%%%%%%%%%%%%%%%%%%%%%%%%%%%%%%%%%%%%%%%%
\section{Các AI Agent}
%%%%%%%%%%%%%%%%%%%%%%%%%%%%%%%%%%%%%%%%%%%%%%%%%%%%%%%%%%%%%%

\subsection{Requirement Agent}

\subsubsection{Mô tả}

Requirement Agent là Agent đầu tiên trong pipeline, chịu trách nhiệm phân tích và làm rõ yêu cầu từ khách hàng.

\subsubsection{Input}

\begin{itemize}
\item Văn bản yêu cầu thô (Raw Requirement)
\item Product Vision (nếu có)
\item Danh sách các yêu cầu hiện có (để phát hiện trùng lặp)
\item Ngữ cảnh dự án (Project Context)
\end{itemize}

\subsubsection{Output}

\begin{verbatim}
{
  "summary": "Tóm tắt yêu cầu",
  "requirements": [
    {
      "id": "REQ-001",
      "description": "Mô tả yêu cầu",
      "category": "Functional/Non-functional",
      "priority": "High/Medium/Low",
      "status": "Validated/Ambiguous/Duplicate"
    }
  ],
  "ambiguous_requirements": [
    {
      "id": "REQ-002",
      "issue": "Yêu cầu mơ hồ",
      "suggested_questions": ["Câu hỏi làm rõ 1", "Câu hỏi làm rõ 2"]
    }
  ],
  "duplicates": [
    {
      "id": "REQ-003",
      "duplicate_of": "REQ-001",
      "similarity_score": 0.85
    }
  ],
  "missing_requirements": [
    "Yêu cầu còn thiếu 1",
    "Yêu cầu còn thiếu 2"
  ],
  "suggested_user_stories": [
    {
      "title": "As a ...",
      "description": "I want ... So that ...",
      "acceptance_criteria": ["Criterion 1", "Criterion 2"]
    }
  ]
}
\end{verbatim}

\subsubsection{Prompt Template}

\begin{verbatim}
You are a Senior Business Analyst with 10+ years of experience in software 
requirements engineering.

Your task is to analyze the following requirement and:
1. Summarize the requirement
2. Identify ambiguous or missing information
3. Detect duplicates with existing requirements
4. Suggest clarifying questions
5. Generate User Stories with Acceptance Criteria

Business Context:
- Project: {project_name}
- Domain: {domain}
- Technologies: {technologies}

Existing Requirements:
{existing_requirements}

Requirement to Analyze:
{requirement_text}

Output Format: JSON
\end{verbatim}

\subsection{Planning Agent}

\subsubsection{Mô tả}

Planning Agent chịu trách nhiệm chuyển đổi User Story thành các Task kỹ thuật cụ thể.

\subsubsection{Input}

\begin{itemize}
\item User Story đã được phê duyệt
\item Sprint mục tiêu (nếu có)
\item Thông tin về team và kỹ năng
\end{itemize}

\subsubsection{Output}

\begin{verbatim}
{
  "tasks": [
    {
      "id": "TASK-001",
      "title": "Thiết kế Database Schema",
      "description": "Thiết kế bảng users và authentication",
      "estimated_hours": 4,
      "priority": "High",
      "dependencies": [],
      "skills_required": ["SQL", "Database Design"]
    },
    {
      "id": "TASK-002",
      "title": "Implement Login API",
      "description": "Xây dựng API đăng nhập với JWT",
      "estimated_hours": 6,
      "priority": "High",
      "dependencies": ["TASK-001"],
      "skills_required": ["Java", "Spring Boot", "JWT"]
    }
  ],
  "estimated_total_hours": 10,
  "suggested_sprint": "Sprint-3",
  "risks": [
    "JWT integration may be complex",
    "Database migration requires careful planning"
  ]
}
\end{verbatim}

\subsection{Story Point Agent}

\subsubsection{Mô tả}

Story Point Agent hỗ trợ ước lượng khối lượng công việc cho User Story dựa trên nhiều tiêu chí.

\subsubsection{Tiêu chí đánh giá}

\begin{table}[H]
\centering
\caption{Tiêu chí đánh giá Story Point}
\begin{tabular}{|p{4cm}|p{5cm}|p{5cm}|}
\hline
\textbf{Tiêu chí} & \textbf{Mô tả} & \textbf{Trọng số} \\
\hline
Complexity & Độ phức tạp về mặt kỹ thuật & 0.25 \\
\hline
Business Logic & Độ phức tạp của logic nghiệp vụ & 0.20 \\
\hline
Integration & Mức độ tích hợp với hệ thống khác & 0.15 \\
\hline
UI/UX & Độ phức tạp của giao diện & 0.10 \\
\hline
Testing & Độ phức tạp của kiểm thử & 0.15 \\
\hline
Uncertainty & Mức độ không chắc chắn của yêu cầu & 0.15 \\
\hline
\end{tabular}
\label{tab:story_point_criteria}
\end{table}

\subsubsection{Output}

\begin{verbatim}
{
  "story_point": 8,
  "fibonacci_sequence": 8,
  "confidence_score": 0.92,
  "explanation": {
    "complexity": "High - Multiple authentication flows",
    "business_logic": "Medium - Standard OAuth2 flow",
    "integration": "High - Need to integrate with existing user service",
    "ui_ux": "Low - Standard login form",
    "testing": "Medium - Multiple test scenarios needed",
    "uncertainty": "Low - Well-defined requirements",
    "historical_reference": "Similar story last sprint was 5 points"
  },
  "recommended_team_size": 2,
  "risks": [
    "OAuth2 configuration may have security implications",
    "Integration with legacy system may cause delays"
  ]
}
\end{verbatim}

\subsection{Forecast Agent}

\subsubsection{Mô tả}

Forecast Agent dự báo tiến độ Sprint và thời gian hoàn thành dự án dựa trên dữ liệu lịch sử.

\subsubsection{Dữ liệu đầu vào}

\begin{itemize}
\item Lịch sử Velocity của các Sprint trước
\item Dữ liệu Burndown Chart của Sprint hiện tại
\item Story Point còn lại của Sprint
\item Team Capacity (ngày công)
\item Số ngày còn lại của Sprint
\end{itemize}

\subsubsection{Output}

\begin{verbatim}
{
  "sprint_forecast": {
    "estimated_completion_date": "2024-12-20",
    "completion_probability": 0.87,
    "remaining_days": 7,
    "required_velocity": 12.5,
    "current_velocity": 11.2,
    "gap": -1.3,
    "recommendation": "Cần tăng tốc độ hoặc giảm scope"
  },
  "project_forecast": {
    "total_remaining_sp": 45,
    "average_velocity": 10.5,
    "estimated_sprints_remaining": 4.3,
    "estimated_completion_date": "2025-02-15"
  },
  "trend_analysis": {
    "velocity_trend": "Stable",
    "velocity_change": "+0.5%",
    "sprint_health": "Good"
  }
}
\end{verbatim}

\subsection{Risk Agent}

\subsubsection{Mô tả}

Risk Agent liên tục phân tích dữ liệu dự án để phát hiện và cảnh báo rủi ro.

\subsubsection{Chỉ số rủi ro}

\[
RiskScore = \sum_{i=1}^{n} w_i \cdot x_i
\]

Trong đó:
- \(x_1\): Task Delay (\%) - Task bị chậm so với kế hoạch
- \(x_2\): Velocity Drop (\%) - Velocity giảm so với Sprint trước
- \(x_3\): Bug Density (bugs/SP) - Số lượng Bug trên mỗi Story Point
- \(x_4\): Workload Level (\%) - Mức độ quá tải của team
- \(x_5\): Requirement Change Rate (\%) - Tỷ lệ thay đổi yêu cầu

\textbf{Trọng số mặc định:}
- \(w_1 = 0.30\) (Task Delay)
- \(w_2 = 0.25\) (Velocity Drop)
- \(w_3 = 0.20\) (Bug Density)
- \(w_4 = 0.15\) (Workload Level)
- \(w_5 = 0.10\) (Requirement Change)

\subsubsection{Phân loại rủi ro}

\begin{table}[H]
\centering
\caption{Phân loại rủi ro}
\begin{tabular}{|c|c|c|}
\hline
\textbf{Risk Score} & \textbf{Mức độ} & \textbf{Hành động} \\
\hline
$0.0 - 0.3$ & Low & Tiếp tục theo dõi \\
\hline
$0.3 - 0.6$ & Medium & Có kế hoạch dự phòng \\
\hline
$0.6 - 0.8$ & High & Hành động ngay \\
\hline
$0.8 - 1.0$ & Critical & Báo cáo cấp trên \\
\hline
\end{tabular}
\label{tab:risk_levels}
\end{table}

\subsubsection{Output}

\begin{verbatim}
{
  "risk_score": 0.65,
  "risk_level": "High",
  "risk_factors": [
    {
      "factor": "Task Delay",
      "value": 35,
      "threshold": 20,
      "severity": 0.8
    },
    {
      "factor": "Velocity Drop",
      "value": 15,
      "threshold": 10,
      "severity": 0.7
    }
  ],
  "warnings": [
    "Team Backend đang vượt quá capacity 18%",
    "Sprint hiện tại có nguy cơ chậm 3 ngày"
  ],
  "recommendations": [
    "Chuyển 2 Task từ Backend sang API team",
    "Xem xét giảm scope Sprint",
    "Họp khẩn với team để đánh giá tình hình"
  ],
  "affected_sprints": ["Sprint-5"],
  "created_at": "2024-12-15T14:30:00Z"
}
\end{verbatim}

\subsection{Assignment Agent}

\subsubsection{Mô tả}

Assignment Agent đề xuất phân công Developer phù hợp nhất cho từng Task.

\subsubsection{Thuật toán phân công}

Điểm phù hợp giữa Developer \(d\) và Task \(t\):

\[
Score(d,t) = \alpha \cdot SkillMatch(d,t) + \beta \cdot Experience(d,t) - \gamma \cdot Workload(d)
\]

Trong đó:
- \(SkillMatch(d,t)\): Mức độ khớp kỹ năng (0-1)
- \(Experience(d,t)\): Kinh nghiệm với Task tương tự (0-1)
- \(Workload(d)\): Khối lượng công việc hiện tại (0-1)
- \(\alpha = 0.45, \beta = 0.35, \gamma = 0.20\)

\subsubsection{Output}

\begin{verbatim}
{
  "recommendations": [
    {
      "task_id": "TASK-001",
      "developer_id": "DEV-001",
      "developer_name": "Nguyễn Văn A",
      "score": 0.92,
      "skill_match": 0.95,
      "experience": 0.90,
      "current_workload": 0.30,
      "reason": "Có kinh nghiệm với JWT và Spring Security"
    },
    {
      "task_id": "TASK-002",
      "developer_id": "DEV-003",
      "developer_name": "Trần Thị B",
      "score": 0.88,
      "skill_match": 0.90,
      "experience": 0.85,
      "current_workload": 0.25,
      "reason": "Có kinh nghiệm với UI và React"
    }
  ],
  "workload_balance": {
    "dev001": 0.30,
    "dev002": 0.45,
    "dev003": 0.25,
    "dev004": 0.50
  },
  "is_balanced": true
}
\end{verbatim}

\subsection{Documentation Agent}

\subsubsection{Mô tả}

Documentation Agent tự động sinh các tài liệu dự án từ dữ liệu có sẵn.

\subsubsection{Các loại tài liệu}

\begin{table}[H]
\centering
\caption{Các loại tài liệu được hỗ trợ}
\begin{tabular}{|p{4cm}|p{4.5cm}|p{5cm}|}
\hline
\textbf{Loại tài liệu} & \textbf{Dữ liệu đầu vào} & \textbf{Định dạng đầu ra} \\
\hline
Sprint Report & Sprint Data, Task Status, Velocity & PDF, Markdown \\
\hline
Meeting Minutes & Meeting Notes (input) & PDF, Markdown \\
\hline
Retrospective Report & Sprint Data, Team Feedback & PDF, Markdown \\
\hline
Release Note & Completed Features, Bug Fixes & PDF, Markdown \\
\hline
API Documentation & Code Comments, Swagger & HTML, Markdown \\
\hline
User Manual & Feature Descriptions & PDF \\
\hline
SRS Summary & All Requirements & PDF, Markdown \\
\hline
\end{tabular}
\label{tab:document_types}
\end{table}

\subsubsection{Output}

\begin{verbatim}
{
  "title": "Sprint Report - Sprint 5",
  "content": "# Sprint Report\n\n## Sprint Goal\n...",
  "sections": [
    {
      "title": "Sprint Overview",
      "content": "..."
    },
    {
      "title": "Completed Items",
      "content": "..."
    },
    {
      "title": "Blocked Issues",
      "content": "..."
    },
    {
      "title": "Velocity Analysis",
      "content": "..."
    },
    {
      "title": "Retrospective Notes",
      "content": "..."
    }
  ],
  "metadata": {
    "created_at": "2024-12-15T14:30:00Z",
    "generated_by": "Documentation Agent",
    "version": "1.0"
  }
}
\end{verbatim}

\subsection{Dashboard Agent}

\subsubsection{Mô tả}

Dashboard Agent tổng hợp dữ liệu và sinh Insights cho Dashboard của hệ thống.

\subsubsection{Các KPI chính}

\begin{itemize}
\item Sprint Velocity (hiện tại và lịch sử)
\item Sprint Burndown (tiến độ)
\item Story Point Completion (\%)
\item Task Distribution (Theo trạng thái và thành viên)
\item Team Workload (Balance)
\item Risk Level (Tổng quan)
\item Bug Trends (Số lượng Bug theo thời gian)
\item AI Insights (Phân tích và đề xuất)
\end{itemize}

\subsubsection{Output}

\begin{verbatim}
{
  "insights": [
    {
      "type": "positive",
      "message": "Sprint hiện tại đang đi đúng tiến độ với tỷ lệ 87%",
      "confidence": 0.92
    },
    {
      "type": "warning",
      "message": "Team Backend đang vượt quá capacity 18%",
      "confidence": 0.85,
      "recommendation": "Xem xét phân công lại 2 Task"
    },
    {
      "type": "info",
      "message": "Velocity ổn định ở mức 10.5 SP/sprint",
      "confidence": 0.95
    },
    {
      "type": "opportunity",
      "message": "Có thể hoàn thành dự án sớm hơn 1 Sprint nếu duy trì velocity hiện tại",
      "confidence": 0.78
    }
  ],
  "kpis": {
    "velocity": 10.5,
    "completion_rate": 0.87,
    "risk_score": 0.35,
    "workload_balance": 0.82,
    "bug_density": 0.12
  },
  "generated_at": "2024-12-15T14:30:00Z"
}
\end{verbatim}

%%%%%%%%%%%%%%%%%%%%%%%%%%%%%%%%%%%%%%%%%%%%%%%%%%%%%%%%%%%%%%
\section{Prompt Engineering}
%%%%%%%%%%%%%%%%%%%%%%%%%%%%%%%%%%%%%%%%%%%%%%%%%%%%%%%%%%%%%%

\subsection{Giới thiệu}

Prompt Engineering là quá trình thiết kế và tối ưu hóa Prompt để hướng dẫn LLM tạo ra kết quả mong muốn. Đây là yếu tố quan trọng nhất để đảm bảo chất lượng của AI.

\subsection{Cấu trúc Prompt}

Mỗi Prompt trong hệ thống bao gồm 5 thành phần:

\begin{enumerate}
\item \textbf{System Prompt}: Định nghĩa vai trò và hành vi của AI
\item \textbf{Business Context}: Các quy tắc nghiệp vụ và ràng buộc
\item \textbf{Project Context}: Dữ liệu cụ thể của dự án
\item \textbf{User Request}: Yêu cầu cụ thể của người dùng
\item \textbf{Output Constraint}: Định dạng đầu ra (JSON)
\end{enumerate}

\subsection{Ví dụ Prompt Template}

\begin{verbatim}
## System Prompt
You are an experienced Senior Software Architect and Agile Project Manager 
with 15+ years of experience. Your task is to help teams manage software 
projects effectively.

You must:
1. Provide accurate and actionable recommendations
2. Always explain your reasoning
3. Return results in the specified JSON format
4. Never invent information that is not in the context
5. If information is missing, ask clarifying questions

## Business Rules
1. Story Points must follow Fibonacci sequence: 1, 2, 3, 5, 8, 13, 21
2. A Sprint can only be started if it has at least 3 User Stories
3. A User Story must have Acceptance Criteria
4. A Task must be assigned to exactly one Developer
5. AI suggestions require human approval before being applied

## Project Context
- Project Name: {project_name}
- Domain: {domain}
- Technologies: {technologies}
- Team Size: {team_size}
- Current Sprint: {current_sprint}
- Velocity History: {velocity_history}

## User Request
{user_request}

## Output Format
Return a valid JSON object with the following structure:
{output_schema}
\end{verbatim}

\subsection{Chiến lược tối ưu Prompt}

\begin{table}[H]
\centering
\caption{Các chiến lược tối ưu Prompt}
\begin{tabular}{|p{4.5cm}|p{4.5cm}|p{5cm}|}
\hline
\textbf{Chiến lược} & \textbf{Mô tả} & \textbf{Khi sử dụng} \\
\hline
Few-shot Learning & Cung cấp ví dụ mẫu & Đầu tiên, khi cần định hướng \\
\hline
Chain-of-Thought & Yêu cầu giải thích từng bước & Khi cần suy luận phức tạp \\
\hline
Self-Consistency & Yêu cầu nhiều lần, chọn kết quả tốt nhất & Khi cần độ chính xác cao \\
\hline
Structured Output & JSON Schema Validation & Luôn luôn \\
\hline
RAG & Truy xuất tri thức nội bộ & Khi cần thông tin cụ thể \\
\hline
\end{tabular}
\label{tab:prompt_strategies}
\end{table}

%%%%%%%%%%%%%%%%%%%%%%%%%%%%%%%%%%%%%%%%%%%%%%%%%%%%%%%%%%%%%%
\section{Retrieval-Augmented Generation (RAG)}
%%%%%%%%%%%%%%%%%%%%%%%%%%%%%%%%%%%%%%%%%%%%%%%%%%%%%%%%%%%%%%

\subsection{Giới thiệu RAG}

RAG (Retrieval-Augmented Generation) là kỹ thuật kết hợp truy xuất thông tin với sinh văn bản. Thay vì chỉ dựa vào kiến thức đã học của LLM, RAG truy xuất tài liệu liên quan từ kho tri thức nội bộ để cung cấp ngữ cảnh chính xác hơn.

\subsection{Kiến trúc RAG}

\begin{figure}[H]
\centering
\fbox{
\begin{minipage}{0.95\textwidth}
\centering
\begin{tabular}{c}
\textbf{Tài liệu dự án} \\
$\downarrow$ \\
\textbf{Chunking \& Embedding} \\
$\downarrow$ \\
\textbf{Vector Database (pgvector)} \\
$\uparrow$ \hspace{1cm} $\uparrow$ \\
\textbf{Truy vấn} \hspace{1.5cm} \textbf{Semantic Search} \\
$\downarrow$ \hspace{1cm} $\downarrow$ \\
\textbf{Tạo embedding} \hspace{1cm} \textbf{Tài liệu liên quan} \\
$\downarrow$ \hspace{1cm} $\downarrow$ \\
\textbf{Prompt Builder} \\
$\downarrow$ \\
\textbf{LLM} \\
$\downarrow$ \\
\textbf{Kết quả}
\end{tabular}
\end{minipage}
}
\caption{Kiến trúc RAG}
\label{fig:rag_architecture}
\end{figure}

\subsection{Nguồn dữ liệu RAG}

\begin{table}[H]
\centering
\caption{Nguồn dữ liệu cho RAG}
\begin{tabular}{|p{4.5cm}|p{5cm}|p{4.5cm}|}
\hline
\textbf{Nguồn dữ liệu} & \textbf{Loại dữ liệu} & \textbf{Tần suất cập nhật} \\
\hline
Product Vision & Văn bản & Một lần \\
\hline
SRS Document & Văn bản & Theo giai đoạn \\
\hline
User Stories & Văn bản + Metadata & Liên tục \\
\hline
Sprint Reports & Văn bản + Số liệu & Mỗi Sprint \\
\hline
Meeting Minutes & Văn bản & Mỗi cuộc họp \\
\hline
Code Documentation & Văn bản + Mã & Liên tục \\
\hline
API Specifications & JSON/YAML & Theo giai đoạn \\
\hline
UML Diagrams & Văn bản mô tả & Theo giai đoạn \\
\hline
\end{tabular}
\label{tab:rag_sources}
\end{table}

\subsection{Quy trình xử lý RAG}

\begin{enumerate}
\item \textbf{Ingestion}: Thu thập và tiền xử lý tài liệu
\item \textbf{Chunking}: Chia tài liệu thành các đoạn nhỏ (512-1024 tokens)
\item \textbf{Embedding}: Chuyển đổi đoạn văn thành vector (OpenAI ada-002)
\item \textbf{Indexing}: Lưu vector vào Vector Database
\item \textbf{Retrieval}: Tìm kiếm ngữ nghĩa khi có truy vấn
\item \textbf{Augmentation}: Kết hợp kết quả tìm kiếm vào Prompt
\end{enumerate}

\subsection{Vector Database}

Hệ thống sử dụng PostgreSQL với extension pgvector để lưu trữ và truy xuất embedding.

\begin{verbatim}
-- Tạo bảng lưu embedding
CREATE TABLE document_embeddings (
    id UUID PRIMARY KEY DEFAULT gen_random_uuid(),
    document_id VARCHAR(255) NOT NULL,
    chunk_index INTEGER NOT NULL,
    content TEXT NOT NULL,
    embedding VECTOR(1536) NOT NULL,
    metadata JSONB,
    created_at TIMESTAMP DEFAULT CURRENT_TIMESTAMP
);

-- Tạo index cho tìm kiếm ngữ nghĩa
CREATE INDEX idx_embedding ON document_embeddings 
    USING ivfflat (embedding vector_cosine_ops);
\end{verbatim}

%%%%%%%%%%%%%%%%%%%%%%%%%%%%%%%%%%%%%%%%%%%%%%%%%%%%%%%%%%%%%%
\section{Đảm bảo chất lượng AI}
%%%%%%%%%%%%%%%%%%%%%%%%%%%%%%%%%%%%%%%%%%%%%%%%%%%%%%%%%%%%%%

\subsection{JSON Schema Validation}

Mọi đầu ra từ LLM đều được kiểm tra với JSON Schema Validation.

\begin{verbatim}
{
  "$schema": "http://json-schema.org/draft-07/schema#",
  "type": "object",
  "required": ["story_point", "confidence", "reason"],
  "properties": {
    "story_point": {
      "type": "integer",
      "enum": [1, 2, 3, 5, 8, 13, 21]
    },
    "confidence": {
      "type": "number",
      "minimum": 0,
      "maximum": 1
    },
    "reason": {
      "type": "string",
      "minLength": 10
    },
    "risks": {
      "type": "array",
      "items": {"type": "string"}
    }
  }
}
\end{verbatim}

\subsection{Business Rules Validation}

Sau JSON Schema, kết quả còn được kiểm tra với các quy tắc nghiệp vụ.

\begin{verbatim}
class BusinessValidator:
    def validate_story_point(self, data):
        # Kiểm tra Story Point theo Fibonacci
        valid_sp = [1, 2, 3, 5, 8, 13, 21]
        if data['story_point'] not in valid_sp:
            raise ValidationError("Invalid Story Point")
        return True
    
    def validate_acceptance_criteria(self, data):
        # Kiểm tra Acceptance Criteria
        if not data.get('acceptance_criteria'):
            raise ValidationError("Missing Acceptance Criteria")
        if len(data['acceptance_criteria']) < 2:
            raise ValidationError("Need at least 2 criteria")
        return True
\end{verbatim}

\subsection{Confidence Scoring}

Mỗi đề xuất của AI đều kèm theo Confidence Score (0-1).

\begin{table}[H]
\centering
\caption{Ý nghĩa Confidence Score}
\begin{tabular}{|c|c|c|}
\hline
\textbf{Score} & \textbf{Mức độ} & \textbf{Hành động} \\
\hline
$0.9 - 1.0$ & Very High & Có thể tự động áp dụng \\
\hline
$0.7 - 0.9$ & High & Người dùng kiểm tra nhanh \\
\hline
$0.5 - 0.7$ & Medium & Cần người dùng xem xét \\
\hline
$0.0 - 0.5$ & Low & Cần người dùng đánh giá kỹ \\
\hline
\end{tabular}
\label{tab:confidence_levels}
\end{table}

\subsection{Human-in-the-Loop}

Hệ thống áp dụng cơ chế Human-in-the-Loop (HITL) cho các quyết định quan trọng.

\begin{figure}[H]
\centering
\fbox{
\begin{minipage}{0.95\textwidth}
\centering
\begin{tabular}{c}
\textbf{AI đề xuất} \\
$\downarrow$ \\
\textbf{Confidence Score} \\
$\swarrow$ $\searrow$ \\
$\ge 0.9$ \hspace{3cm} $< 0.9$ \\
$\downarrow$ \hspace{3cm} $\downarrow$ \\
\textbf{Tự động áp dụng} \hspace{2cm} \textbf{Cần phê duyệt} \\
\hspace{3.5cm} $\downarrow$ \\
\hspace{3.5cm} \textbf{Người dùng xem xét} \\
\hspace{3.5cm} $\downarrow$ \\
\hspace{3.5cm} \textbf{Phê duyệt/Từ chối/Chỉnh sửa} \\
\hspace{3.5cm} $\downarrow$ \\
\hspace{3.5cm} \textbf{Lưu kết quả}
\end{tabular}
\end{minipage}
}
\caption{Human-in-the-Loop Process}
\label{fig:human_in_loop}
\end{figure}

%%%%%%%%%%%%%%%%%%%%%%%%%%%%%%%%%%%%%%%%%%%%%%%%%%%%%%%%%%%%%%
\section{Luồng xử lý dữ liệu trong AI}
%%%%%%%%%%%%%%%%%%%%%%%%%%%%%%%%%%%%%%%%%%%%%%%%%%%%%%%%%%%%%%

\subsection{Luồng xử lý tổng quát}

\begin{figure}[H]
\centering
\fbox{
\begin{minipage}{0.95\textwidth}
\centering
\begin{tabular}{c}
\textbf{Frontend} \\
$\downarrow$ \\
\textbf{Backend API} \\
$\downarrow$ \\
\textbf{AI Orchestrator} \\
$\downarrow$ \\
\textbf{Intent Detection} \\
$\downarrow$ \\
\textbf{Agent Selection} \\
$\downarrow$ \\
\textbf{RAG Retrieval} \\
$\downarrow$ \\
\textbf{Prompt Builder} \\
$\downarrow$ \\
\textbf{LLM} \\
$\downarrow$ \\
\textbf{JSON Validation} \\
$\downarrow$ \\
\textbf{Business Validation} \\
$\downarrow$ \\
\textbf{Save Database} \\
$\downarrow$ \\
\textbf{Return Result}
\end{tabular}
\end{minipage}
}
\caption{Luồng xử lý dữ liệu AI}
\label{fig:ai_data_flow}
\end{figure}

\subsection{Luồng xử lý chi tiết cho Requirement Analysis}

\begin{enumerate}
\item Product Owner nhập yêu cầu qua Frontend.
\item Frontend gửi request đến Backend API `/ai/analyze-requirement`.
\item Backend gọi AI Orchestrator với request.
\item AI Orchestrator phát hiện intent = REQUIREMENT\_ANALYSIS.
\item Orchestrator chọn Requirement Agent.
\item RAG Service truy xuất tài liệu liên quan từ Vector DB.
\item Prompt Builder xây dựng Prompt với System Prompt, Business Rules, Project Context, User Request.
\item LLM phân tích và sinh kết quả.
\item JSON Validator kiểm tra cấu trúc.
\item Business Validator kiểm tra quy tắc nghiệp vụ.
\item Kết quả được lưu vào bảng AI\_Suggestions.
\item Trả kết quả về Frontend.
\item Product Owner xem xét và phê duyệt.
\end{enumerate}

%%%%%%%%%%%%%%%%%%%%%%%%%%%%%%%%%%%%%%%%%%%%%%%%%%%%%%%%%%%%%%
\section{Tổng kết chương}
%%%%%%%%%%%%%%%%%%%%%%%%%%%%%%%%%%%%%%%%%%%%%%%%%%%%%%%%%%%%%%

Trong chương này, luận văn đã trình bày chi tiết kiến trúc AI và hệ thống Multi-Agent của nền tảng AI Project Management.

Các nội dung chính đã được thực hiện:

\begin{enumerate}
\item \textbf{Kiến trúc AI}: 6 tầng từ API Gateway đến Output Validation.
\item \textbf{AI Orchestrator}: Thành phần trung tâm điều phối các Agent.
\item \textbf{Các AI Agent}: 8 Agent chuyên biệt (Requirement, Planning, Story Point, Forecast, Risk, Assignment, Documentation, Dashboard).
\item \textbf{Prompt Engineering}: 5 thành phần Prompt, chiến lược tối ưu.
\item \textbf{RAG}: Kiến trúc và quy trình truy xuất tri thức nội bộ.
\item \textbf{Đảm bảo chất lượng}: JSON Schema Validation, Business Rules, Confidence Scoring, HITL.
\item \textbf{Luồng xử lý}: Mô tả chi tiết luồng từ Frontend đến kết quả.
\end{enumerate}

Đây là nền tảng để Chương 9 trình bày thiết kế cơ sở dữ liệu chi tiết và Chương 10 trình bày thiết kế Backend.

%%%%%%%%%%%%%%%%%%%%%%%%%%%%%%%%%%%%%%%%%%%%%%%%%%%%%%%%%%%%%%

%%%%%%%%%%%%%%%%%%%%%%%%%%%%%%%%%%%%%%%%%%%%%%%%%%%%%%%%%%%%%%
% CHƯƠNG 9: THIẾT KẾ CƠ SỞ DỮ LIỆU
%%%%%%%%%%%%%%%%%%%%%%%%%%%%%%%%%%%%%%%%%%%%%%%%%%%%%%%%%%%%%%
\chapter{THIẾT KẾ CƠ SỞ DỮ LIỆU}
\label{chap:database_design}
%%%%%%%%%%%%%%%%%%%%%%%%%%%%%%%%%%%%%%%%%%%%%%%%%%%%%%%%%%%%%%
\chapter{THIẾT KẾ CƠ SỞ DỮ LIỆU}
\label{chap:database_design}
%%%%%%%%%%%%%%%%%%%%%%%%%%%%%%%%%%%%%%%%%%%%%%%%%%%%%%%%%%%%%%

\section{Giới thiệu chương}

Sau khi trình bày kiến trúc AI và Multi-Agent ở Chương 8, chương này tập trung vào thiết kế cơ sở dữ liệu cho hệ thống AI Project Management. Cơ sở dữ liệu là thành phần nền tảng lưu trữ toàn bộ dữ liệu nghiệp vụ, lịch sử dự án, kết quả phân tích AI và các thông tin liên quan.

Mục tiêu của chương bao gồm:

\begin{itemize}
\item Trình bày mô hình dữ liệu tổng thể của hệ thống.
\item Thiết kế chi tiết các bảng dữ liệu.
\item Xác định mối quan hệ giữa các bảng.
\item Thiết kế các ràng buộc toàn vẹn (Integrity Constraints).
\item Thiết kế Vector Database cho RAG.
\item Đề xuất các chỉ mục (Indexes) tối ưu hiệu năng.
\item Thiết kế các view và stored procedure hỗ trợ Dashboard.
\end{itemize}

Đây là cơ sở để Chương 10 trình bày thiết kế Backend và Chương 11 trình bày thiết kế Frontend.

%%%%%%%%%%%%%%%%%%%%%%%%%%%%%%%%%%%%%%%%%%%%%%%%%%%%%%%%%%%%%%
\section{Mô hình dữ liệu tổng thể}
%%%%%%%%%%%%%%%%%%%%%%%%%%%%%%%%%%%%%%%%%%%%%%%%%%%%%%%%%%%%%%

\subsection{Sơ đồ thực thể - liên kết (ERD)}

\begin{figure}[H]
\centering
\fbox{
\begin{minipage}{0.95\textwidth}
\centering
\begin{tabular}{c}
\textbf{users} \\
$\downarrow$ 1--N \\
\textbf{workspaces} \\
$\downarrow$ 1--N \\
\textbf{workspace\_members} \\
$\uparrow$ N--1 \\
\textbf{users} \\
\\
\textbf{workspaces} $\rightarrow$ 1--N $\rightarrow$ \textbf{projects} \\
\textbf{projects} $\rightarrow$ 1--N $\rightarrow$ \textbf{sprints} \\
\textbf{projects} $\rightarrow$ 1--N $\rightarrow$ \textbf{epics} \\
\textbf{epics} $\rightarrow$ 1--N $\rightarrow$ \textbf{user\_stories} \\
\textbf{sprints} $\rightarrow$ 1--N $\rightarrow$ \textbf{user\_stories} \\
\textbf{user\_stories} $\rightarrow$ 1--N $\rightarrow$ \textbf{tasks} \\
\textbf{tasks} $\rightarrow$ N--1 $\rightarrow$ \textbf{users (assignee)} \\
\\
\textbf{projects} $\rightarrow$ 1--N $\rightarrow$ \textbf{ai\_suggestions} \\
\textbf{projects} $\rightarrow$ 1--N $\rightarrow$ \textbf{audit\_logs} \\
\textbf{users} $\rightarrow$ 1--N $\rightarrow$ \textbf{notifications} \\
\textbf{projects} $\rightarrow$ 1--N $\rightarrow$ \textbf{documents}
\end{tabular}
\end{minipage}
}
\caption{Sơ đồ thực thể - liên kết (ERD)}
\label{fig:erd}
\end{figure}

\subsection{Các miền dữ liệu chính}

\begin{table}[H]
\centering
\caption{Các miền dữ liệu trong hệ thống}
\begin{tabular}{|p{4cm}|p{4.5cm}|p{5cm}|}
\hline
\textbf{Miền dữ liệu} & \textbf{Bảng chính} & \textbf{Mô tả} \\
\hline
Người dùng và bảo mật & users, roles, permissions & Quản lý tài khoản và phân quyền \\
\hline
Tổ chức & workspaces, workspace\_members & Quản lý Workspace và thành viên \\
\hline
Dự án & projects, project\_members & Quản lý Project và thành viên \\
\hline
Kế hoạch & sprints, epics, user\_stories & Quản lý Sprint, Epic, User Story \\
\hline
Công việc & tasks, task\_comments, task\_attachments & Quản lý Task và chi tiết \\
\hline
AI & ai\_suggestions, ai\_feedbacks & Lưu kết quả và phản hồi AI \\
\hline
Kiểm toán & audit\_logs & Lịch sử thay đổi \\
\hline
Tài liệu & documents, document\_embeddings & Lưu trữ tài liệu và vector \\
\hline
Thông báo & notifications & Hệ thống thông báo \\
\hline
\end{tabular}
\label{tab:data_domains}
\end{table}

%%%%%%%%%%%%%%%%%%%%%%%%%%%%%%%%%%%%%%%%%%%%%%%%%%%%%%%%%%%%%%
\section{Thiết kế chi tiết các bảng}
%%%%%%%%%%%%%%%%%%%%%%%%%%%%%%%%%%%%%%%%%%%%%%%%%%%%%%%%%%%%%%

\subsection{Bảng Users}

Bảng lưu thông tin tài khoản người dùng.

\begin{table}[H]
\centering
\caption{Bảng Users}
\begin{tabular}{|p{4cm}|p{4cm}|p{7cm}|}
\hline
\textbf{Tên trường} & \textbf{Loại dữ liệu} & \textbf{Mô tả} \\
\hline
id & UUID & Khóa chính, sinh tự động \\
\hline
email & VARCHAR(255) & Email đăng nhập, unique \\
\hline
password\_hash & VARCHAR(255) & Mật khẩu mã hóa BCrypt \\
\hline
full\_name & VARCHAR(255) & Họ và tên đầy đủ \\
\hline
username & VARCHAR(100) & Tên hiển thị, unique \\
\hline
avatar\_url & TEXT & URL ảnh đại diện \\
\hline
role & ENUM & ADMIN, PROJECT\_MANAGER, DEVELOPER, TESTER, VIEWER \\
\hline
status & ENUM & ACTIVE, INACTIVE, BLOCKED, PENDING\_VERIFICATION \\
\hline
email\_verified & BOOLEAN & Đã xác thực email \\
\hline
last\_login\_at & TIMESTAMP & Lần đăng nhập cuối \\
\hline
created\_at & TIMESTAMP & Thời gian tạo \\
\hline
updated\_at & TIMESTAMP & Thời gian cập nhật \\
\hline
deleted\_at & TIMESTAMP & Thời gian xóa (soft delete) \\
\hline
\end{tabular}
\label{tab:users_table}
\end{table}

\begin{verbatim}
-- SQL Script
CREATE TABLE users (
    id UUID PRIMARY KEY DEFAULT gen_random_uuid(),
    email VARCHAR(255) NOT NULL UNIQUE,
    password_hash VARCHAR(255) NOT NULL,
    full_name VARCHAR(255) NOT NULL,
    username VARCHAR(100) NOT NULL UNIQUE,
    avatar_url TEXT,
    role VARCHAR(50) NOT NULL DEFAULT 'DEVELOPER',
    status VARCHAR(50) NOT NULL DEFAULT 'ACTIVE',
    email_verified BOOLEAN DEFAULT FALSE,
    last_login_at TIMESTAMP,
    created_at TIMESTAMP NOT NULL DEFAULT CURRENT_TIMESTAMP,
    updated_at TIMESTAMP NOT NULL DEFAULT CURRENT_TIMESTAMP,
    deleted_at TIMESTAMP,
    CONSTRAINT chk_role CHECK (role IN ('ADMIN', 'PROJECT_MANAGER', 'DEVELOPER', 'TESTER', 'VIEWER')),
    CONSTRAINT chk_status CHECK (status IN ('ACTIVE', 'INACTIVE', 'BLOCKED', 'PENDING_VERIFICATION'))
);

CREATE INDEX idx_users_email ON users(email);
CREATE INDEX idx_users_status ON users(status);
CREATE INDEX idx_users_deleted_at ON users(deleted_at) WHERE deleted_at IS NOT NULL;
\end{verbatim}

\subsection{Bảng Workspaces}

Bảng lưu thông tin Workspace.

\begin{table}[H]
\centering
\caption{Bảng Workspaces}
\begin{tabular}{|p{4cm}|p{4cm}|p{7cm}|}
\hline
\textbf{Tên trường} & \textbf{Loại dữ liệu} & \textbf{Mô tả} \\
\hline
id & UUID & Khóa chính, sinh tự động \\
\hline
name & VARCHAR(255) & Tên Workspace \\
\hline
slug & VARCHAR(255) & Định danh URL, unique \\
\hline
description & TEXT & Mô tả Workspace \\
\hline
owner\_id & UUID & Khóa ngoại đến users (Workspace Owner) \\
\hline
logo\_url & TEXT & URL logo \\
\hline
visibility & ENUM & PUBLIC, PRIVATE \\
\hline
settings & JSONB & Cấu hình Workspace (JSON) \\
\hline
created\_at & TIMESTAMP & Thời gian tạo \\
\hline
updated\_at & TIMESTAMP & Thời gian cập nhật \\
\hline
deleted\_at & TIMESTAMP & Thời gian xóa (soft delete) \\
\hline
\end{tabular}
\label{tab:workspaces_table}
\end{table}

\begin{verbatim}
CREATE TABLE workspaces (
    id UUID PRIMARY KEY DEFAULT gen_random_uuid(),
    name VARCHAR(255) NOT NULL,
    slug VARCHAR(255) NOT NULL UNIQUE,
    description TEXT,
    owner_id UUID NOT NULL REFERENCES users(id),
    logo_url TEXT,
    visibility VARCHAR(50) NOT NULL DEFAULT 'PRIVATE',
    settings JSONB DEFAULT '{}'::jsonb,
    created_at TIMESTAMP NOT NULL DEFAULT CURRENT_TIMESTAMP,
    updated_at TIMESTAMP NOT NULL DEFAULT CURRENT_TIMESTAMP,
    deleted_at TIMESTAMP,
    CONSTRAINT chk_visibility CHECK (visibility IN ('PUBLIC', 'PRIVATE'))
);

CREATE INDEX idx_workspaces_owner_id ON workspaces(owner_id);
CREATE INDEX idx_workspaces_slug ON workspaces(slug);
\end{verbatim}

\subsection{Bảng Workspace Members}

Bảng trung gian quản lý thành viên Workspace.

\begin{table}[H]
\centering
\caption{Bảng Workspace Members}
\begin{tabular}{|p{4cm}|p{4cm}|p{7cm}|}
\hline
\textbf{Tên trường} & \textbf{Loại dữ liệu} & \textbf{Mô tả} \\
\hline
id & UUID & Khóa chính, sinh tự động \\
\hline
workspace\_id & UUID & Khóa ngoại đến workspaces \\
\hline
user\_id & UUID & Khóa ngoại đến users \\
\hline
role & ENUM & OWNER, ADMIN, MEMBER, VIEWER \\
\hline
joined\_at & TIMESTAMP & Thời gian tham gia \\
\hline
invited\_by & UUID & Khóa ngoại đến users (người mời) \\
\hline
status & ENUM & ACTIVE, PENDING, BLOCKED \\
\hline
\end{tabular}
\label{tab:workspace_members_table}
\end{table}

\begin{verbatim}
CREATE TABLE workspace_members (
    id UUID PRIMARY KEY DEFAULT gen_random_uuid(),
    workspace_id UUID NOT NULL REFERENCES workspaces(id) ON DELETE CASCADE,
    user_id UUID NOT NULL REFERENCES users(id) ON DELETE CASCADE,
    role VARCHAR(50) NOT NULL DEFAULT 'MEMBER',
    joined_at TIMESTAMP NOT NULL DEFAULT CURRENT_TIMESTAMP,
    invited_by UUID REFERENCES users(id),
    status VARCHAR(50) NOT NULL DEFAULT 'ACTIVE',
    UNIQUE(workspace_id, user_id),
    CONSTRAINT chk_role CHECK (role IN ('OWNER', 'ADMIN', 'MEMBER', 'VIEWER')),
    CONSTRAINT chk_status CHECK (status IN ('ACTIVE', 'PENDING', 'BLOCKED'))
);

CREATE INDEX idx_workspace_members_workspace_id ON workspace_members(workspace_id);
CREATE INDEX idx_workspace_members_user_id ON workspace_members(user_id);
\end{verbatim}

\subsection{Bảng Projects}

Bảng lưu thông tin dự án.

\begin{table}[H]
\centering
\caption{Bảng Projects}
\begin{tabular}{|p{4cm}|p{4cm}|p{7cm}|}
\hline
\textbf{Tên trường} & \textbf{Loại dữ liệu} & \textbf{Mô tả} \\
\hline
id & UUID & Khóa chính, sinh tự động \\
\hline
workspace\_id & UUID & Khóa ngoại đến workspaces \\
\hline
name & VARCHAR(255) & Tên dự án \\
\hline
key & VARCHAR(50) & Mã dự án, unique trong workspace \\
\hline
description & TEXT & Mô tả dự án \\
\hline
status & ENUM & PLANNING, ACTIVE, DONE, ARCHIVED \\
\hline
priority & ENUM & LOW, MEDIUM, HIGH \\
\hline
visibility & ENUM & PUBLIC, PRIVATE \\
\hline
start\_date & DATE & Ngày bắt đầu \\
\hline
end\_date & DATE & Ngày kết thúc dự kiến \\
\hline
actual\_end\_date & DATE & Ngày kết thúc thực tế \\
\hline
settings & JSONB & Cấu hình dự án \\
\hline
created\_at & TIMESTAMP & Thời gian tạo \\
\hline
updated\_at & TIMESTAMP & Thời gian cập nhật \\
\hline
deleted\_at & TIMESTAMP & Thời gian xóa (soft delete) \\
\hline
\end{tabular}
\label{tab:projects_table}
\end{table}

\begin{verbatim}
CREATE TABLE projects (
    id UUID PRIMARY KEY DEFAULT gen_random_uuid(),
    workspace_id UUID NOT NULL REFERENCES workspaces(id) ON DELETE CASCADE,
    name VARCHAR(255) NOT NULL,
    key VARCHAR(50) NOT NULL,
    description TEXT,
    status VARCHAR(50) NOT NULL DEFAULT 'PLANNING',
    priority VARCHAR(50) NOT NULL DEFAULT 'MEDIUM',
    visibility VARCHAR(50) NOT NULL DEFAULT 'PRIVATE',
    start_date DATE,
    end_date DATE,
    actual_end_date DATE,
    settings JSONB DEFAULT '{}'::jsonb,
    created_at TIMESTAMP NOT NULL DEFAULT CURRENT_TIMESTAMP,
    updated_at TIMESTAMP NOT NULL DEFAULT CURRENT_TIMESTAMP,
    deleted_at TIMESTAMP,
    UNIQUE(workspace_id, key),
    CONSTRAINT chk_status CHECK (status IN ('PLANNING', 'ACTIVE', 'DONE', 'ARCHIVED')),
    CONSTRAINT chk_priority CHECK (priority IN ('LOW', 'MEDIUM', 'HIGH')),
    CONSTRAINT chk_visibility CHECK (visibility IN ('PUBLIC', 'PRIVATE'))
);

CREATE INDEX idx_projects_workspace_id ON projects(workspace_id);
CREATE INDEX idx_projects_status ON projects(status);
CREATE INDEX idx_projects_key ON projects(key);
\end{verbatim}

\subsection{Bảng Project Members}

Bảng trung gian quản lý thành viên Project.

\begin{table}[H]
\centering
\caption{Bảng Project Members}
\begin{tabular}{|p{4cm}|p{4cm}|p{7cm}|}
\hline
\textbf{Tên trường} & \textbf{Loại dữ liệu} & \textbf{Mô tả} \\
\hline
id & UUID & Khóa chính, sinh tự động \\
\hline
project\_id & UUID & Khóa ngoại đến projects \\
\hline
user\_id & UUID & Khóa ngoại đến users \\
\hline
role & ENUM & PROJECT\_MANAGER, DEVELOPER, TESTER, VIEWER \\
\hline
joined\_at & TIMESTAMP & Thời gian tham gia \\
\hline
assigned\_by & UUID & Khóa ngoại đến users (người phân công) \\
\hline
\end{tabular}
\label{tab:project_members_table}
\end{table}

\begin{verbatim}
CREATE TABLE project_members (
    id UUID PRIMARY KEY DEFAULT gen_random_uuid(),
    project_id UUID NOT NULL REFERENCES projects(id) ON DELETE CASCADE,
    user_id UUID NOT NULL REFERENCES users(id) ON DELETE CASCADE,
    role VARCHAR(50) NOT NULL DEFAULT 'DEVELOPER',
    joined_at TIMESTAMP NOT NULL DEFAULT CURRENT_TIMESTAMP,
    assigned_by UUID REFERENCES users(id),
    UNIQUE(project_id, user_id),
    CONSTRAINT chk_role CHECK (role IN ('PROJECT_MANAGER', 'DEVELOPER', 'TESTER', 'VIEWER'))
);

CREATE INDEX idx_project_members_project_id ON project_members(project_id);
CREATE INDEX idx_project_members_user_id ON project_members(user_id);
\end{verbatim}

\subsection{Bảng Sprints}

Bảng lưu thông tin Sprint.

\begin{table}[H]
\centering
\caption{Bảng Sprints}
\begin{tabular}{|p{4cm}|p{4cm}|p{7cm}|}
\hline
\textbf{Tên trường} & \textbf{Loại dữ liệu} & \textbf{Mô tả} \\
\hline
id & UUID & Khóa chính, sinh tự động \\
\hline
project\_id & UUID & Khóa ngoại đến projects \\
\hline
name & VARCHAR(255) & Tên Sprint \\
\hline
goal & TEXT & Sprint Goal \\
\hline
start\_date & DATE & Ngày bắt đầu \\
\hline
end\_date & DATE & Ngày kết thúc \\
\hline
actual\_end\_date & DATE & Ngày kết thúc thực tế \\
\hline
velocity & INTEGER & Velocity (SP hoàn thành) \\
\hline
status & ENUM & PLANNING, ACTIVE, COMPLETED, CANCELLED \\
\hline
created\_at & TIMESTAMP & Thời gian tạo \\
\hline
updated\_at & TIMESTAMP & Thời gian cập nhật \\
\hline
\end{tabular}
\label{tab:sprints_table}
\end{table}

\begin{verbatim}
CREATE TABLE sprints (
    id UUID PRIMARY KEY DEFAULT gen_random_uuid(),
    project_id UUID NOT NULL REFERENCES projects(id) ON DELETE CASCADE,
    name VARCHAR(255) NOT NULL,
    goal TEXT,
    start_date DATE,
    end_date DATE,
    actual_end_date DATE,
    velocity INTEGER DEFAULT 0,
    status VARCHAR(50) NOT NULL DEFAULT 'PLANNING',
    created_at TIMESTAMP NOT NULL DEFAULT CURRENT_TIMESTAMP,
    updated_at TIMESTAMP NOT NULL DEFAULT CURRENT_TIMESTAMP,
    CONSTRAINT chk_status CHECK (status IN ('PLANNING', 'ACTIVE', 'COMPLETED', 'CANCELLED'))
);

CREATE INDEX idx_sprints_project_id ON sprints(project_id);
CREATE INDEX idx_sprints_status ON sprints(status);
CREATE INDEX idx_sprints_dates ON sprints(start_date, end_date);
\end{verbatim}

\subsection{Bảng Epics}

Bảng lưu thông tin Epic.

\begin{table}[H]
\centering
\caption{Bảng Epics}
\begin{tabular}{|p{4cm}|p{4cm}|p{7cm}|}
\hline
\textbf{Tên trường} & \textbf{Loại dữ liệu} & \textbf{Mô tả} \\
\hline
id & UUID & Khóa chính, sinh tự động \\
\hline
project\_id & UUID & Khóa ngoại đến projects \\
\hline
name & VARCHAR(255) & Tên Epic \\
\hline
description & TEXT & Mô tả Epic \\
\hline
priority & INTEGER & Độ ưu tiên (1-5) \\
\hline
status & ENUM & TODO, IN\_PROGRESS, DONE \\
\hline
created\_at & TIMESTAMP & Thời gian tạo \\
\hline
updated\_at & TIMESTAMP & Thời gian cập nhật \\
\hline
\end{tabular}
\label{tab:epics_table}
\end{table}

\begin{verbatim}
CREATE TABLE epics (
    id UUID PRIMARY KEY DEFAULT gen_random_uuid(),
    project_id UUID NOT NULL REFERENCES projects(id) ON DELETE CASCADE,
    name VARCHAR(255) NOT NULL,
    description TEXT,
    priority INTEGER DEFAULT 3,
    status VARCHAR(50) NOT NULL DEFAULT 'TODO',
    created_at TIMESTAMP NOT NULL DEFAULT CURRENT_TIMESTAMP,
    updated_at TIMESTAMP NOT NULL DEFAULT CURRENT_TIMESTAMP,
    CONSTRAINT chk_priority CHECK (priority BETWEEN 1 AND 5),
    CONSTRAINT chk_status CHECK (status IN ('TODO', 'IN_PROGRESS', 'DONE'))
);

CREATE INDEX idx_epics_project_id ON epics(project_id);
\end{verbatim}

\subsection{Bảng User Stories}

Bảng lưu thông tin User Story - thực thể quan trọng nhất của Product Backlog.

\begin{table}[H]
\centering
\caption{Bảng User Stories}
\begin{tabular}{|p{4cm}|p{4cm}|p{7cm}|}
\hline
\textbf{Tên trường} & \textbf{Loại dữ liệu} & \textbf{Mô tả} \\
\hline
id & UUID & Khóa chính, sinh tự động \\
\hline
epic\_id & UUID & Khóa ngoại đến epics (nullable) \\
\hline
sprint\_id & UUID & Khóa ngoại đến sprints (nullable) \\
\hline
title & VARCHAR(255) & Tiêu đề User Story \\
\hline
description & TEXT & Mô tả User Story \\
\hline
acceptance\_criteria & TEXT & Tiêu chí chấp nhận \\
\hline
story\_point & INTEGER & Story Point (Fibonacci) \\
\hline
priority & INTEGER & Độ ưu tiên (1-5) \\
\hline
business\_value & INTEGER & Giá trị kinh doanh (1-10) \\
\hline
status & ENUM & TODO, IN\_PROGRESS, DONE, BLOCKED \\
\hline
order\_in\_backlog & INTEGER & Thứ tự trong Product Backlog \\
\hline
created\_at & TIMESTAMP & Thời gian tạo \\
\hline
updated\_at & TIMESTAMP & Thời gian cập nhật \\
\hline
\end{tabular}
\label{tab:user_stories_table}
\end{table}

\begin{verbatim}
CREATE TABLE user_stories (
    id UUID PRIMARY KEY DEFAULT gen_random_uuid(),
    epic_id UUID REFERENCES epics(id) ON DELETE SET NULL,
    sprint_id UUID REFERENCES sprints(id) ON DELETE SET NULL,
    title VARCHAR(255) NOT NULL,
    description TEXT,
    acceptance_criteria TEXT,
    story_point INTEGER,
    priority INTEGER DEFAULT 3,
    business_value INTEGER DEFAULT 5,
    status VARCHAR(50) NOT NULL DEFAULT 'TODO',
    order_in_backlog INTEGER,
    created_at TIMESTAMP NOT NULL DEFAULT CURRENT_TIMESTAMP,
    updated_at TIMESTAMP NOT NULL DEFAULT CURRENT_TIMESTAMP,
    CONSTRAINT chk_story_point CHECK (story_point IN (1, 2, 3, 5, 8, 13, 21)),
    CONSTRAINT chk_priority CHECK (priority BETWEEN 1 AND 5),
    CONSTRAINT chk_business_value CHECK (business_value BETWEEN 1 AND 10),
    CONSTRAINT chk_status CHECK (status IN ('TODO', 'IN_PROGRESS', 'DONE', 'BLOCKED'))
);

CREATE INDEX idx_user_stories_epic_id ON user_stories(epic_id);
CREATE INDEX idx_user_stories_sprint_id ON user_stories(sprint_id);
CREATE INDEX idx_user_stories_status ON user_stories(status);
CREATE INDEX idx_user_stories_order ON user_stories(order_in_backlog);
\end{verbatim}

\subsection{Bảng Tasks}

Bảng lưu thông tin Task - đơn vị công việc nhỏ nhất.

\begin{table}[H]
\centering
\caption{Bảng Tasks}
\begin{tabular}{|p{4cm}|p{4cm}|p{7cm}|}
\hline
\textbf{Tên trường} & \textbf{Loại dữ liệu} & \textbf{Mô tả} \\
\hline
id & UUID & Khóa chính, sinh tự động \\
\hline
story\_id & UUID & Khóa ngoại đến user\_stories \\
\hline
assignee\_id & UUID & Khóa ngoại đến users (nullable) \\
\hline
title & VARCHAR(255) & Tiêu đề Task \\
\hline
description & TEXT & Mô tả chi tiết \\
\hline
estimated\_hours & DECIMAL(5,2) & Thời gian ước lượng (giờ) \\
\hline
actual\_hours & DECIMAL(5,2) & Thời gian thực tế (giờ) \\
\hline
status & ENUM & TODO, IN\_PROGRESS, DONE, BLOCKED \\
\hline
priority & ENUM & LOW, MEDIUM, HIGH, URGENT \\
\hline
due\_date & DATE & Hạn hoàn thành \\
\hline
started\_at & TIMESTAMP & Thời gian bắt đầu \\
\hline
completed\_at & TIMESTAMP & Thời gian hoàn thành \\
\hline
created\_at & TIMESTAMP & Thời gian tạo \\
\hline
updated\_at & TIMESTAMP & Thời gian cập nhật \\
\hline
\end{tabular}
\label{tab:tasks_table}
\end{table}

\begin{verbatim}
CREATE TABLE tasks (
    id UUID PRIMARY KEY DEFAULT gen_random_uuid(),
    story_id UUID NOT NULL REFERENCES user_stories(id) ON DELETE CASCADE,
    assignee_id UUID REFERENCES users(id) ON DELETE SET NULL,
    title VARCHAR(255) NOT NULL,
    description TEXT,
    estimated_hours DECIMAL(5,2),
    actual_hours DECIMAL(5,2),
    status VARCHAR(50) NOT NULL DEFAULT 'TODO',
    priority VARCHAR(50) NOT NULL DEFAULT 'MEDIUM',
    due_date DATE,
    started_at TIMESTAMP,
    completed_at TIMESTAMP,
    created_at TIMESTAMP NOT NULL DEFAULT CURRENT_TIMESTAMP,
    updated_at TIMESTAMP NOT NULL DEFAULT CURRENT_TIMESTAMP,
    CONSTRAINT chk_task_status CHECK (status IN ('TODO', 'IN_PROGRESS', 'DONE', 'BLOCKED')),
    CONSTRAINT chk_task_priority CHECK (priority IN ('LOW', 'MEDIUM', 'HIGH', 'URGENT'))
);

CREATE INDEX idx_tasks_story_id ON tasks(story_id);
CREATE INDEX idx_tasks_assignee_id ON tasks(assignee_id);
CREATE INDEX idx_tasks_status ON tasks(status);
CREATE INDEX idx_tasks_due_date ON tasks(due_date);
\end{verbatim}

\subsection{Bảng Task Comments}

Bảng lưu comment của Task.

\begin{table}[H]
\centering
\caption{Bảng Task Comments}
\begin{tabular}{|p{4cm}|p{4cm}|p{7cm}|}
\hline
\textbf{Tên trường} & \textbf{Loại dữ liệu} & \textbf{Mô tả} \\
\hline
id & UUID & Khóa chính, sinh tự động \\
\hline
task\_id & UUID & Khóa ngoại đến tasks \\
\hline
user\_id & UUID & Khóa ngoại đến users \\
\hline
content & TEXT & Nội dung comment \\
\hline
parent\_id & UUID & Khóa ngoại đến comments (reply) \\
\hline
created\_at & TIMESTAMP & Thời gian tạo \\
\hline
updated\_at & TIMESTAMP & Thời gian cập nhật \\
\hline
\end{tabular}
\label{tab:task_comments_table}
\end{table}

\begin{verbatim}
CREATE TABLE task_comments (
    id UUID PRIMARY KEY DEFAULT gen_random_uuid(),
    task_id UUID NOT NULL REFERENCES tasks(id) ON DELETE CASCADE,
    user_id UUID NOT NULL REFERENCES users(id) ON DELETE CASCADE,
    content TEXT NOT NULL,
    parent_id UUID REFERENCES task_comments(id) ON DELETE CASCADE,
    created_at TIMESTAMP NOT NULL DEFAULT CURRENT_TIMESTAMP,
    updated_at TIMESTAMP NOT NULL DEFAULT CURRENT_TIMESTAMP
);

CREATE INDEX idx_task_comments_task_id ON task_comments(task_id);
CREATE INDEX idx_task_comments_user_id ON task_comments(user_id);
\end{verbatim}

\subsection{Bảng Task Attachments}

Bảng lưu file đính kèm của Task.

\begin{table}[H]
\centering
\caption{Bảng Task Attachments}
\begin{tabular}{|p{4cm}|p{4cm}|p{7cm}|}
\hline
\textbf{Tên trường} & \textbf{Loại dữ liệu} & \textbf{Mô tả} \\
\hline
id & UUID & Khóa chính, sinh tự động \\
\hline
task\_id & UUID & Khóa ngoại đến tasks \\
\hline
user\_id & UUID & Khóa ngoại đến users \\
\hline
file\_name & VARCHAR(255) & Tên file gốc \\
\hline
file\_url & TEXT & URL file trên storage \\
\hline
file\_size & BIGINT & Kích thước file (bytes) \\
\hline
mime\_type & VARCHAR(100) & Loại MIME \\
\hline
created\_at & TIMESTAMP & Thời gian tạo \\
\hline
\end{tabular}
\label{tab:task_attachments_table}
\end{table}

\begin{verbatim}
CREATE TABLE task_attachments (
    id UUID PRIMARY KEY DEFAULT gen_random_uuid(),
    task_id UUID NOT NULL REFERENCES tasks(id) ON DELETE CASCADE,
    user_id UUID NOT NULL REFERENCES users(id) ON DELETE CASCADE,
    file_name VARCHAR(255) NOT NULL,
    file_url TEXT NOT NULL,
    file_size BIGINT,
    mime_type VARCHAR(100),
    created_at TIMESTAMP NOT NULL DEFAULT CURRENT_TIMESTAMP
);

CREATE INDEX idx_task_attachments_task_id ON task_attachments(task_id);
\end{verbatim}

\subsection{Bảng AI Suggestions}

Bảng lưu kết quả phân tích và đề xuất của AI.

\begin{table}[H]
\centering
\caption{Bảng AI Suggestions}
\begin{tabular}{|p{4cm}|p{4cm}|p{7cm}|}
\hline
\textbf{Tên trường} & \textbf{Loại dữ liệu} & \textbf{Mô tả} \\
\hline
id & UUID & Khóa chính, sinh tự động \\
\hline
project\_id & UUID & Khóa ngoại đến projects \\
\hline
agent\_type & VARCHAR(50) & Loại Agent (REQUIREMENT, STORY\_POINT, v.v.) \\
\hline
target\_id & UUID & ID của thực thể liên quan \\
\hline
target\_type & VARCHAR(50) & Loại thực thể (USER\_STORY, SPRINT, v.v.) \\
\hline
prompt & TEXT & Prompt đã gửi \\
\hline
response & JSONB & Kết quả từ AI \\
\hline
confidence & DECIMAL(3,2) & Độ tin cậy (0-1) \\
\hline
approved & BOOLEAN & Đã được phê duyệt \\
\hline
approved\_by & UUID & Khóa ngoại đến users \\
\hline
approved\_at & TIMESTAMP & Thời gian phê duyệt \\
\hline
feedback & TEXT & Phản hồi của người dùng \\
\hline
created\_at & TIMESTAMP & Thời gian tạo \\
\hline
updated\_at & TIMESTAMP & Thời gian cập nhật \\
\hline
\end{tabular}
\label{tab:ai_suggestions_table}
\end{table}

\begin{verbatim}
CREATE TABLE ai_suggestions (
    id UUID PRIMARY KEY DEFAULT gen_random_uuid(),
    project_id UUID NOT NULL REFERENCES projects(id) ON DELETE CASCADE,
    agent_type VARCHAR(50) NOT NULL,
    target_id UUID,
    target_type VARCHAR(50),
    prompt TEXT,
    response JSONB NOT NULL,
    confidence DECIMAL(3,2),
    approved BOOLEAN DEFAULT FALSE,
    approved_by UUID REFERENCES users(id) ON DELETE SET NULL,
    approved_at TIMESTAMP,
    feedback TEXT,
    created_at TIMESTAMP NOT NULL DEFAULT CURRENT_TIMESTAMP,
    updated_at TIMESTAMP NOT NULL DEFAULT CURRENT_TIMESTAMP,
    CONSTRAINT chk_confidence CHECK (confidence >= 0 AND confidence <= 1)
);

CREATE INDEX idx_ai_suggestions_project_id ON ai_suggestions(project_id);
CREATE INDEX idx_ai_suggestions_agent_type ON ai_suggestions(agent_type);
CREATE INDEX idx_ai_suggestions_target ON ai_suggestions(target_type, target_id);
CREATE INDEX idx_ai_suggestions_approved ON ai_suggestions(approved);
\end{verbatim}

\subsection{Bảng Audit Logs}

Bảng lưu lịch sử thay đổi (Audit Log).

\begin{table}[H]
\centering
\caption{Bảng Audit Logs}
\begin{tabular}{|p{4cm}|p{4cm}|p{7cm}|}
\hline
\textbf{Tên trường} & \textbf{Loại dữ liệu} & \textbf{Mô tả} \\
\hline
id & UUID & Khóa chính, sinh tự động \\
\hline
user\_id & UUID & Khóa ngoại đến users \\
\hline
action & VARCHAR(50) & Loại hành động (CREATE, UPDATE, DELETE, v.v.) \\
\hline
entity\_type & VARCHAR(50) & Loại thực thể (PROJECT, SPRINT, TASK, v.v.) \\
\hline
entity\_id & UUID & ID của thực thể \\
\hline
old\_values & JSONB & Giá trị cũ \\
\hline
new\_values & JSONB & Giá trị mới \\
\hline
ip\_address & VARCHAR(45) & Địa chỉ IP \\
\hline
user\_agent & TEXT & User Agent của trình duyệt \\
\hline
created\_at & TIMESTAMP & Thời gian tạo \\
\hline
\end{tabular}
\label{tab:audit_logs_table}
\end{table}

\begin{verbatim}
CREATE TABLE audit_logs (
    id UUID PRIMARY KEY DEFAULT gen_random_uuid(),
    user_id UUID NOT NULL REFERENCES users(id) ON DELETE CASCADE,
    action VARCHAR(50) NOT NULL,
    entity_type VARCHAR(50) NOT NULL,
    entity_id UUID NOT NULL,
    old_values JSONB,
    new_values JSONB,
    ip_address VARCHAR(45),
    user_agent TEXT,
    created_at TIMESTAMP NOT NULL DEFAULT CURRENT_TIMESTAMP
);

CREATE INDEX idx_audit_logs_user_id ON audit_logs(user_id);
CREATE INDEX idx_audit_logs_entity ON audit_logs(entity_type, entity_id);
CREATE INDEX idx_audit_logs_created_at ON audit_logs(created_at);
\end{verbatim}

\subsection{Bảng Notifications}

Bảng lưu thông báo cho người dùng.

\begin{table}[H]
\centering
\caption{Bảng Notifications}
\begin{tabular}{|p{4cm}|p{4cm}|p{7cm}|}
\hline
\textbf{Tên trường} & \textbf{Loại dữ liệu} & \textbf{Mô tả} \\
\hline
id & UUID & Khóa chính, sinh tự động \\
\hline
user\_id & UUID & Khóa ngoại đến users \\
\hline
type & VARCHAR(50) & Loại thông báo (TASK\_ASSIGNED, RISK\_ALERT, v.v.) \\
\hline
title & VARCHAR(255) & Tiêu đề thông báo \\
\hline
content & TEXT & Nội dung thông báo \\
\hline
link & TEXT & Đường dẫn liên quan \\
\hline
is\_read & BOOLEAN & Đã đọc \\
\hline
read\_at & TIMESTAMP & Thời gian đọc \\
\hline
created\_at & TIMESTAMP & Thời gian tạo \\
\hline
\end{tabular}
\label{tab:notifications_table}
\end{table}

\begin{verbatim}
CREATE TABLE notifications (
    id UUID PRIMARY KEY DEFAULT gen_random_uuid(),
    user_id UUID NOT NULL REFERENCES users(id) ON DELETE CASCADE,
    type VARCHAR(50) NOT NULL,
    title VARCHAR(255) NOT NULL,
    content TEXT NOT NULL,
    link TEXT,
    is_read BOOLEAN DEFAULT FALSE,
    read_at TIMESTAMP,
    created_at TIMESTAMP NOT NULL DEFAULT CURRENT_TIMESTAMP
);

CREATE INDEX idx_notifications_user_id ON notifications(user_id);
CREATE INDEX idx_notifications_is_read ON notifications(is_read);
CREATE INDEX idx_notifications_created_at ON notifications(created_at);
\end{verbatim}

%%%%%%%%%%%%%%%%%%%%%%%%%%%%%%%%%%%%%%%%%%%%%%%%%%%%%%%%%%%%%%
\section{Vector Database cho RAG}
%%%%%%%%%%%%%%%%%%%%%%%%%%%%%%%%%%%%%%%%%%%%%%%%%%%%%%%%%%%%%%

\subsection{Tổng quan}

Hệ thống sử dụng PostgreSQL với extension pgvector để lưu trữ và truy xuất vector embedding cho RAG (Retrieval-Augmented Generation).

\subsection{Cài đặt pgvector}

\begin{verbatim}
-- Cài đặt extension
CREATE EXTENSION IF NOT EXISTS vector;

-- Tạo bảng lưu embedding
CREATE TABLE document_embeddings (
    id UUID PRIMARY KEY DEFAULT gen_random_uuid(),
    document_id VARCHAR(255) NOT NULL,
    source_type VARCHAR(50) NOT NULL, -- 'SRS', 'USER_STORY', 'SPRINT_REPORT', etc.
    project_id UUID REFERENCES projects(id) ON DELETE CASCADE,
    chunk_index INTEGER NOT NULL,
    content TEXT NOT NULL,
    embedding VECTOR(1536) NOT NULL,
    metadata JSONB DEFAULT '{}'::jsonb,
    created_at TIMESTAMP NOT NULL DEFAULT CURRENT_TIMESTAMP,
    updated_at TIMESTAMP NOT NULL DEFAULT CURRENT_TIMESTAMP
);

-- Tạo index cho tìm kiếm ngữ nghĩa
CREATE INDEX idx_document_embeddings_embedding 
    ON document_embeddings 
    USING ivfflat (embedding vector_cosine_ops)
    WITH (lists = 100);

-- Tạo index cho truy xuất theo project
CREATE INDEX idx_document_embeddings_project_id 
    ON document_embeddings(project_id);

-- Tạo index cho truy xuất theo source_type
CREATE INDEX idx_document_embeddings_source_type 
    ON document_embeddings(source_type);
\end{verbatim}

\subsection{Hàm tìm kiếm ngữ nghĩa}

\begin{verbatim}
-- Hàm tìm kiếm ngữ nghĩa
CREATE OR REPLACE FUNCTION semantic_search(
    query_embedding VECTOR(1536),
    project_id UUID,
    limit_count INTEGER DEFAULT 10,
    similarity_threshold DECIMAL DEFAULT 0.7
)
RETURNS TABLE(
    id UUID,
    content TEXT,
    source_type VARCHAR(50),
    similarity DECIMAL
) AS $$
BEGIN
    RETURN QUERY
    SELECT 
        de.id,
        de.content,
        de.source_type,
        1 - (de.embedding <=> query_embedding) AS similarity
    FROM document_embeddings de
    WHERE de.project_id = semantic_search.project_id
        AND 1 - (de.embedding <=> query_embedding) >= similarity_threshold
    ORDER BY de.embedding <=> query_embedding
    LIMIT limit_count;
END;
$$ LANGUAGE plpgsql;
\end{verbatim}

\subsection{Quy trình xử lý tài liệu}

\begin{enumerate}
\item \textbf{Ingestion}: Thu thập tài liệu (SRS, User Story, Sprint Report, v.v.)
\item \textbf{Chunking}: Chia tài liệu thành các đoạn nhỏ (512-1024 tokens)
\item \textbf{Embedding}: Tạo vector embedding cho mỗi đoạn (OpenAI ada-002)
\item \textbf{Storage}: Lưu vector vào bảng document\_embeddings
\item \textbf{Retrieval}: Khi có truy vấn, tìm kiếm ngữ nghĩa bằng hàm semantic\_search
\end{enumerate}

%%%%%%%%%%%%%%%%%%%%%%%%%%%%%%%%%%%%%%%%%%%%%%%%%%%%%%%%%%%%%%
\section{Ràng buộc toàn vẹn và Trigger}
%%%%%%%%%%%%%%%%%%%%%%%%%%%%%%%%%%%%%%%%%%%%%%%%%%%%%%%%%%%%%%

\subsection{Trigger cập nhật updated\_at}

\begin{verbatim}
-- Hàm cập nhật updated_at
CREATE OR REPLACE FUNCTION update_updated_at_column()
RETURNS TRIGGER AS $$
BEGIN
    NEW.updated_at = CURRENT_TIMESTAMP;
    RETURN NEW;
END;
$$ LANGUAGE plpgsql;

-- Áp dụng cho các bảng cần thiết
CREATE TRIGGER update_users_updated_at 
    BEFORE UPDATE ON users 
    FOR EACH ROW EXECUTE FUNCTION update_updated_at_column();

CREATE TRIGGER update_workspaces_updated_at 
    BEFORE UPDATE ON workspaces 
    FOR EACH ROW EXECUTE FUNCTION update_updated_at_column();

CREATE TRIGGER update_projects_updated_at 
    BEFORE UPDATE ON projects 
    FOR EACH ROW EXECUTE FUNCTION update_updated_at_column();

CREATE TRIGGER update_sprints_updated_at 
    BEFORE UPDATE ON sprints 
    FOR EACH ROW EXECUTE FUNCTION update_updated_at_column();

CREATE TRIGGER update_user_stories_updated_at 
    BEFORE UPDATE ON user_stories 
    FOR EACH ROW EXECUTE FUNCTION update_updated_at_column();

CREATE TRIGGER update_tasks_updated_at 
    BEFORE UPDATE ON tasks 
    FOR EACH ROW EXECUTE FUNCTION update_updated_at_column();
\end{verbatim}

\subsection{Trigger ghi Audit Log}

\begin{verbatim}
-- Hàm ghi audit log
CREATE OR REPLACE FUNCTION log_audit()
RETURNS TRIGGER AS $$
BEGIN
    INSERT INTO audit_logs (
        user_id,
        action,
        entity_type,
        entity_id,
        old_values,
        new_values
    )
    VALUES (
        current_user_id(), -- Cần hàm lấy user_id hiện tại
        TG_OP,
        TG_TABLE_NAME,
        NEW.id,
        CASE WHEN TG_OP = 'DELETE' THEN row_to_json(OLD) ELSE NULL END,
        CASE WHEN TG_OP <> 'DELETE' THEN row_to_json(NEW) ELSE NULL END
    );
    RETURN NEW;
END;
$$ LANGUAGE plpgsql;
\end{verbatim}

%%%%%%%%%%%%%%%%%%%%%%%%%%%%%%%%%%%%%%%%%%%%%%%%%%%%%%%%%%%%%%
\section{Chỉ mục (Indexes) tối ưu hiệu năng}
%%%%%%%%%%%%%%%%%%%%%%%%%%%%%%%%%%%%%%%%%%%%%%%%%%%%%%%%%%%%%%

\begin{table}[H]
\centering
\caption{Các chỉ mục quan trọng}
\begin{tabular}{|p{4cm}|p{4.5cm}|p{5cm}|}
\hline
\textbf{Bảng} & \textbf{Trường Index} & \textbf{Mục đích} \\
\hline
users & email, username & Tìm kiếm nhanh khi đăng nhập \\
\hline
users & status, deleted\_at & Lọc người dùng hoạt động \\
\hline
workspaces & owner\_id & Lấy Workspace của Owner \\
\hline
workspace\_members & workspace\_id, user\_id & Truy vấn thành viên \\
\hline
projects & workspace\_id, status & Lấy danh sách Project \\
\hline
sprints & project\_id, status & Lấy Sprint theo Project \\
\hline
user\_stories & sprint\_id, status & Lấy Backlog Sprint \\
\hline
tasks & assignee\_id, status & Lấy Task của Developer \\
\hline
tasks & story\_id & Lấy Task của User Story \\
\hline
ai\_suggestions & project\_id, approved & Lấy đề xuất AI \\
\hline
audit\_logs & user\_id, created\_at & Truy vấn lịch sử \\
\hline
notifications & user\_id, is\_read & Lấy thông báo chưa đọc \\
\hline
document\_embeddings & embedding (vector) & Tìm kiếm ngữ nghĩa \\
\hline
\end{tabular}
\label{tab:indexes}
\end{table}

%%%%%%%%%%%%%%%%%%%%%%%%%%%%%%%%%%%%%%%%%%%%%%%%%%%%%%%%%%%%%%
\section{Views cho Dashboard}
%%%%%%%%%%%%%%%%%%%%%%%%%%%%%%%%%%%%%%%%%%%%%%%%%%%%%%%%%%%%%%

\subsection{View Sprint Overview}

\begin{verbatim}
CREATE VIEW sprint_overview AS
SELECT 
    s.id AS sprint_id,
    s.name AS sprint_name,
    s.project_id,
    p.name AS project_name,
    s.goal,
    s.start_date,
    s.end_date,
    s.velocity,
    s.status,
    COUNT(DISTINCT us.id) AS total_stories,
    COUNT(DISTINCT us.id) FILTER (WHERE us.status = 'DONE') AS done_stories,
    SUM(us.story_point) AS total_sp,
    SUM(us.story_point) FILTER (WHERE us.status = 'DONE') AS done_sp,
    COUNT(DISTINCT t.id) AS total_tasks,
    COUNT(DISTINCT t.id) FILTER (WHERE t.status = 'DONE') AS done_tasks,
    (SUM(us.story_point) FILTER (WHERE us.status = 'DONE')::FLOAT / NULLIF(SUM(us.story_point), 0)) * 100 AS completion_rate
FROM sprints s
JOIN projects p ON s.project_id = p.id
LEFT JOIN user_stories us ON us.sprint_id = s.id
LEFT JOIN tasks t ON t.story_id = us.id
GROUP BY s.id, s.name, s.project_id, p.name, s.goal, s.start_date, s.end_date, s.velocity, s.status;
\end{verbatim}

\subsection{View Project Dashboard}

\begin{verbatim}
CREATE VIEW project_dashboard AS
SELECT 
    p.id AS project_id,
    p.name AS project_name,
    p.status AS project_status,
    COUNT(DISTINCT s.id) AS total_sprints,
    COUNT(DISTINCT s.id) FILTER (WHERE s.status = 'ACTIVE') AS active_sprints,
    COUNT(DISTINCT us.id) AS total_stories,
    COUNT(DISTINCT us.id) FILTER (WHERE us.status = 'DONE') AS done_stories,
    SUM(us.story_point) AS total_sp,
    SUM(us.story_point) FILTER (WHERE us.status = 'DONE') AS done_sp,
    COUNT(DISTINCT t.id) AS total_tasks,
    COUNT(DISTINCT t.id) FILTER (WHERE t.status = 'DONE') AS done_tasks,
    AVG(s.velocity) AS avg_velocity,
    COUNT(DISTINCT pm.user_id) AS team_size
FROM projects p
LEFT JOIN sprints s ON s.project_id = p.id
LEFT JOIN user_stories us ON us.sprint_id = s.id
LEFT JOIN tasks t ON t.story_id = us.id
LEFT JOIN project_members pm ON pm.project_id = p.id
GROUP BY p.id, p.name, p.status;
\end{verbatim}

%%%%%%%%%%%%%%%%%%%%%%%%%%%%%%%%%%%%%%%%%%%%%%%%%%%%%%%%%%%%%%
\section{Tổng kết chương}
%%%%%%%%%%%%%%%%%%%%%%%%%%%%%%%%%%%%%%%%%%%%%%%%%%%%%%%%%%%%%%

Trong chương này, luận văn đã trình bày thiết kế cơ sở dữ liệu chi tiết cho hệ thống AI Project Management.

Các nội dung chính đã được thực hiện:

\begin{enumerate}
\item \textbf{Mô hình dữ liệu tổng thể}: Sơ đồ ERD và các miền dữ liệu chính.
\item \textbf{Thiết kế chi tiết các bảng}: 12 bảng chính với đầy đủ trường, kiểu dữ liệu và ràng buộc.
\item \textbf{Vector Database}: Thiết kế bảng document\_embeddings và hàm semantic\_search cho RAG.
\item \textbf{Ràng buộc toàn vẹn}: Triggers tự động cập nhật và ghi audit log.
\item \textbf{Chỉ mục}: Đề xuất các index tối ưu hiệu năng.
\item \textbf{Views}: Các view hỗ trợ Dashboard và báo cáo.
\end{enumerate}

Đây là cơ sở để Chương 10 trình bày thiết kế Backend và các service tương tác với cơ sở dữ liệu.

%%%%%%%%%%%%%%%%%%%%%%%%%%%%%%%%%%%%%%%%%%%%%%%%%%%%%%%%%%%%%%
%%%%%%%%%%%%%%%%%%%%%%%%%%%%%%%%%%%%%%%%%%%%%%%%%%%%%%%%%%%%%%
%%%%%%%%%%%%%%%%%%%%%%%%%%%%%%%%%%%%%%%%%%%%%%%%%%%%%%%%%%%%%%
\chapter{THIẾT KẾ BACKEND}
\label{chap:backend_design}
%%%%%%%%%%%%%%%%%%%%%%%%%%%%%%%%%%%%%%%%%%%%%%%%%%%%%%%%%%%%%%

\section{Giới thiệu chương}

Sau khi trình bày thiết kế cơ sở dữ liệu ở Chương 9, chương này tập trung vào thiết kế chi tiết tầng Backend của hệ thống AI Project Management. Backend là thành phần trung tâm, đóng vai trò kết nối giữa Frontend, Cơ sở dữ liệu và AI Service.

Mục tiêu của chương bao gồm:

\begin{itemize}
\item Trình bày kiến trúc tổng thể của Backend.
\item Thiết kế các Service Layer xử lý nghiệp vụ.
\item Thiết kế Repository Layer truy xuất dữ liệu.
\item Thiết kế hệ thống xác thực và phân quyền (Security).
\item Thiết kế API Controllers và Routing.
\item Thiết kế tích hợp với AI Service.
\item Thiết kế xử lý ngoại lệ và logging.
\item Thiết kế cấu hình và deployment.
\end{itemize}

Đây là cơ sở để Chương 11 trình bày thiết kế Frontend và Chương 12 trình bày triển khai hệ thống.

%%%%%%%%%%%%%%%%%%%%%%%%%%%%%%%%%%%%%%%%%%%%%%%%%%%%%%%%%%%%%%
\section{Kiến trúc Backend tổng thể}
%%%%%%%%%%%%%%%%%%%%%%%%%%%%%%%%%%%%%%%%%%%%%%%%%%%%%%%%%%%%%%

\subsection{Mô hình kiến trúc}

Backend được xây dựng theo kiến trúc Clean Architecture kết hợp với Domain-Driven Design (DDD) và Layered Architecture.

\begin{figure}[H]
\centering
\fbox{
\begin{minipage}{0.95\textwidth}
\centering
\begin{tabular}{c}
\textbf{Tầng Presentation (Controller)} \\
- REST API Controllers \\
- DTOs (Data Transfer Objects) \\
- Request/Response Validation \\
$\downarrow$ \\
\textbf{Tầng Application (Service)} \\
- Use Case Implementations \\
- Business Logic Orchestration \\
- Transaction Management \\
$\downarrow$ \\
\textbf{Tầng Domain} \\
- Entities \\
- Value Objects \\
- Domain Services \\
- Repository Interfaces \\
- Business Rules \\
$\downarrow$ \\
\textbf{Tầng Infrastructure} \\
- Repository Implementations \\
- External API Clients \\
- Database Access (JPA/Hibernate) \\
- Message Queue \\
- Cache (Redis)
\end{tabular}
\end{minipage}
}
\caption{Kiến trúc Backend - Clean Architecture}
\label{fig:backend_architecture}
\end{figure}

\subsection{Công nghệ sử dụng}

\begin{table}[H]
\centering
\caption{Công nghệ Backend}
\begin{tabular}{|p{4cm}|p{5cm}|p{6cm}|}
\hline
\textbf{Thành phần} & \textbf{Công nghệ} & \textbf{Vai trò} \\
\hline
Ngôn ngữ & Java 21 & Ngôn ngữ lập trình chính \\
\hline
Framework & Spring Boot 3.x & Framework ứng dụng \\
\hline
ORM & Spring Data JPA / Hibernate & Truy xuất cơ sở dữ liệu \\
\hline
Security & Spring Security 6.x & Xác thực và phân quyền \\
\hline
JWT & JJWT & Tạo và xác thực JWT Token \\
\hline
Database & PostgreSQL 16 & Cơ sở dữ liệu quan hệ \\
\hline
Migration & Flyway & Quản lý schema database \\
\hline
Cache & Redis & Cache và session \\
\hline
Message Queue & RabbitMQ & Xử lý bất đồng bộ \\
\hline
API Client & RestTemplate / WebClient & Gọi API ngoài \\
\hline
Logging & SLF4J + Logback & Ghi log \\
\hline
Testing & JUnit 5, Mockito & Kiểm thử \\
\hline
Build & Maven / Gradle & Quản lý dependency \\
\hline
\end{tabular}
\label{tab:backend_technologies}
\end{table}

\subsection{Cấu trúc thư mục dự án}

\begin{verbatim}
src/
├── main/
│   ├── java/
│   │   └── com/ai/pms/
│   │       ├── AiProjectManagementApplication.java
│   │       ├── config/               # Cấu hình Spring
│   │       │   ├── SecurityConfig.java
│   │       │   ├── JwtConfig.java
│   │       │   ├── WebConfig.java
│   │       │   └── AppConfig.java
│   │       ├── controller/            # REST Controllers
│   │       │   ├── AuthController.java
│   │       │   ├── WorkspaceController.java
│   │       │   ├── ProjectController.java
│   │       │   ├── SprintController.java
│   │       │   ├── BacklogController.java
│   │       │   ├── TaskController.java
│   │       │   ├── AIController.java
│   │       │   └── DashboardController.java
│   │       ├── service/               # Service Layer
│   │       │   ├── AuthService.java
│   │       │   ├── WorkspaceService.java
│   │       │   ├── ProjectService.java
│   │       │   ├── SprintService.java
│   │       │   ├── BacklogService.java
│   │       │   ├── TaskService.java
│   │       │   ├── AIService.java
│   │       │   └── DashboardService.java
│   │       ├── repository/            # Repository Layer
│   │       │   ├── UserRepository.java
│   │       │   ├── WorkspaceRepository.java
│   │       │   ├── ProjectRepository.java
│   │       │   ├── SprintRepository.java
│   │       │   ├── UserStoryRepository.java
│   │       │   └── TaskRepository.java
│   │       ├── domain/                # Domain Layer
│   │       │   ├── entity/            # Entities
│   │       │   │   ├── User.java
│   │       │   │   ├── Workspace.java
│   │       │   │   ├── Project.java
│   │       │   │   └── ...
│   │       │   ├── valueobject/       # Value Objects
│   │       │   │   ├── Email.java
│   │       │   │   ├── StoryPoint.java
│   │       │   │   └── ...
│   │       │   └── enums/             # Enums
│   │       │       ├── UserRole.java
│   │       │       ├── ProjectStatus.java
│   │       │       └── ...
│   │       ├── dto/                   # Data Transfer Objects
│   │       │   ├── request/
│   │       │   │   ├── LoginRequest.java
│   │       │   │   ├── CreateProjectRequest.java
│   │       │   │   └── ...
│   │       │   └── response/
│   │       │       ├── AuthResponse.java
│   │       │       ├── ProjectResponse.java
│   │       │       └── ...
│   │       ├── security/              # Security
│   │       │   ├── JwtAuthenticationFilter.java
│   │       │   ├── JwtTokenProvider.java
│   │       │   ├── UserDetailsServiceImpl.java
│   │       │   └── SecurityUtils.java
│   │       ├── exception/             # Exception Handling
│   │       │   ├── GlobalExceptionHandler.java
│   │       │   ├── BusinessException.java
│   │       │   └── ResourceNotFoundException.java
│   │       ├── client/                # External Clients
│   │       │   ├── AIServiceClient.java
│   │       │   └── EmailServiceClient.java
│   │       ├── util/                  # Utilities
│   │       │   ├── JwtUtil.java
│   │       │   ├── ValidationUtil.java
│   │       │   └── DateUtil.java
│   │       └── config/
│   │           └── audit/             # Audit Configuration
│   │               ├── AuditAware.java
│   │               └── AuditConfig.java
│   └── resources/
│       ├── application.yml
│       ├── application-dev.yml
│       ├── application-prod.yml
│       └── db/migration/              # Flyway migrations
│           ├── V1__create_users_table.sql
│           ├── V2__create_workspaces_table.sql
│           └── ...
└── test/
    └── java/
        └── com/ai/pms/
            ├── controller/
            ├── service/
            └── repository/
\end{verbatim}

%%%%%%%%%%%%%%%%%%%%%%%%%%%%%%%%%%%%%%%%%%%%%%%%%%%%%%%%%%%%%%
\section{Thiết kế Security (Xác thực và Phân quyền)}
%%%%%%%%%%%%%%%%%%%%%%%%%%%%%%%%%%%%%%%%%%%%%%%%%%%%%%%%%%%%%%

\subsection{JWT Authentication Flow}

\begin{figure}[H]
\centering
\fbox{
\begin{minipage}{0.95\textwidth}
\centering
\begin{tabular}{c}
\textbf{1. Đăng nhập} \\
User $\rightarrow$ POST /auth/login $\rightarrow$ {email, password} \\
$\downarrow$ \\
\textbf{2. Xác thực} \\
Backend kiểm tra email/password \\
$\downarrow$ \\
\textbf{3. Tạo JWT} \\
Sinh Access Token (15 phút) và Refresh Token (7 ngày) \\
$\downarrow$ \\
\textbf{4. Trả token} \\
{ accessToken, refreshToken, userInfo } \\
$\downarrow$ \\
\textbf{5. Gọi API} \\
Client gửi Authorization: Bearer <token> \\
$\downarrow$ \\
\textbf{6. Xác thực token} \\
JwtAuthenticationFilter kiểm tra token \\
$\downarrow$ \\
\textbf{7. Phân quyền} \\
Kiểm tra vai trò và quyền truy cập \\
$\downarrow$ \\
\textbf{8. Xử lý request} \\
Controller thực hiện nghiệp vụ
\end{tabular}
\end{minipage}
}
\caption{Luồng xác thực JWT}
\label{fig:jwt_flow}
\end{figure}

\subsection{Spring Security Configuration}

\begin{verbatim}
@Configuration
@EnableWebSecurity
@EnableMethodSecurity(prePostEnabled = true)
public class SecurityConfig {
    
    @Bean
    public SecurityFilterChain securityFilterChain(HttpSecurity http) throws Exception {
        http
            .csrf(csrf -> csrf.disable())
            .cors(cors -> cors.configurationSource(corsConfigurationSource()))
            .authorizeHttpRequests(auth -> auth
                // Public endpoints
                .requestMatchers("/api/v1/auth/**").permitAll()
                .requestMatchers("/swagger-ui/**", "/v3/api-docs/**").permitAll()
                .requestMatchers("/actuator/health").permitAll()
                // Protected endpoints
                .requestMatchers("/api/v1/workspaces/**").hasAnyRole("ADMIN", "PROJECT_MANAGER")
                .requestMatchers("/api/v1/projects/**").hasAnyRole("ADMIN", "PROJECT_MANAGER", "DEVELOPER")
                .requestMatchers("/api/v1/admin/**").hasRole("ADMIN")
                .anyRequest().authenticated()
            )
            .sessionManagement(session -> session.sessionCreationPolicy(SessionCreationPolicy.STATELESS))
            .authenticationProvider(authenticationProvider())
            .addFilterBefore(jwtAuthenticationFilter(), UsernamePasswordAuthenticationFilter.class)
            .exceptionHandling(ex -> ex
                .authenticationEntryPoint(unauthorizedHandler)
                .accessDeniedHandler(accessDeniedHandler)
            );
        
        return http.build();
    }
    
    @Bean
    public CorsConfigurationSource corsConfigurationSource() {
        CorsConfiguration configuration = new CorsConfiguration();
        configuration.setAllowedOrigins(Arrays.asList("http://localhost:5173", "http://localhost:3000"));
        configuration.setAllowedMethods(Arrays.asList("GET", "POST", "PUT", "DELETE", "PATCH", "OPTIONS"));
        configuration.setAllowedHeaders(Arrays.asList("*"));
        configuration.setAllowCredentials(true);
        configuration.setMaxAge(3600L);
        
        UrlBasedCorsConfigurationSource source = new UrlBasedCorsConfigurationSource();
        source.registerCorsConfiguration("/**", configuration);
        return source;
    }
    
    @Bean
    public AuthenticationManager authenticationManager(AuthenticationConfiguration config) throws Exception {
        return config.getAuthenticationManager();
    }
}
\end{verbatim}

\subsection{JWT Token Provider}

\begin{verbatim}
@Component
public class JwtTokenProvider {
    
    @Value("${jwt.secret}")
    private String jwtSecret;
    
    @Value("${jwt.access-token-expiration}")
    private int accessTokenExpiration;
    
    @Value("${jwt.refresh-token-expiration}")
    private int refreshTokenExpiration;
    
    public String generateAccessToken(UserDetails userDetails) {
        Map<String, Object> claims = new HashMap<>();
        claims.put("roles", userDetails.getAuthorities().stream()
            .map(auth -> auth.getAuthority())
            .collect(Collectors.toList()));
        
        return Jwts.builder()
            .setClaims(claims)
            .setSubject(userDetails.getUsername())
            .setIssuedAt(new Date())
            .setExpiration(new Date(System.currentTimeMillis() + accessTokenExpiration))
            .signWith(getSignKey(), SignatureAlgorithm.HS256)
            .compact();
    }
    
    public String generateRefreshToken(String email) {
        return Jwts.builder()
            .setSubject(email)
            .setIssuedAt(new Date())
            .setExpiration(new Date(System.currentTimeMillis() + refreshTokenExpiration))
            .signWith(getSignKey(), SignatureAlgorithm.HS256)
            .compact();
    }
    
    public boolean validateToken(String token) {
        try {
            Jwts.parserBuilder()
                .setSigningKey(getSignKey())
                .build()
                .parseClaimsJws(token);
            return true;
        } catch (JwtException | IllegalArgumentException e) {
            return false;
        }
    }
    
    public String getEmailFromToken(String token) {
        Claims claims = Jwts.parserBuilder()
            .setSigningKey(getSignKey())
            .build()
            .parseClaimsJws(token)
            .getBody();
        return claims.getSubject();
    }
    
    private Key getSignKey() {
        byte[] keyBytes = Decoders.BASE64.decode(jwtSecret);
        return Keys.hmacShaKeyFor(keyBytes);
    }
}
\end{verbatim}

\subsection{JWT Authentication Filter}

\begin{verbatim}
@Component
public class JwtAuthenticationFilter extends OncePerRequestFilter {
    
    @Autowired
    private JwtTokenProvider tokenProvider;
    
    @Autowired
    private UserDetailsService userDetailsService;
    
    @Override
    protected void doFilterInternal(HttpServletRequest request, 
                                    HttpServletResponse response, 
                                    FilterChain filterChain) 
            throws ServletException, IOException {
        
        try {
            String jwt = getJwtFromRequest(request);
            
            if (StringUtils.hasText(jwt) && tokenProvider.validateToken(jwt)) {
                String email = tokenProvider.getEmailFromToken(jwt);
                UserDetails userDetails = userDetailsService.loadUserByUsername(email);
                UsernamePasswordAuthenticationToken authentication = 
                    new UsernamePasswordAuthenticationToken(
                        userDetails, null, userDetails.getAuthorities()
                    );
                authentication.setDetails(new WebAuthenticationDetailsSource().buildDetails(request));
                SecurityContextHolder.getContext().setAuthentication(authentication);
            }
        } catch (Exception ex) {
            logger.error("Could not set user authentication in security context", ex);
        }
        
        filterChain.doFilter(request, response);
    }
    
    private String getJwtFromRequest(HttpServletRequest request) {
        String bearerToken = request.getHeader("Authorization");
        if (StringUtils.hasText(bearerToken) && bearerToken.startsWith("Bearer ")) {
            return bearerToken.substring(7);
        }
        return null;
    }
}
\end{verbatim}

\subsection{Phân quyền (RBAC)}

\begin{table}[H]
\centering
\caption{Vai trò và quyền hạn}
\begin{tabular}{|p{3.5cm}|p{4.5cm}|p{7cm}|}
\hline
\textbf{Vai trò} & \textbf{Phạm vi} & \textbf{Quyền hạn} \\
\hline
ADMIN & Toàn hệ thống & Toàn quyền quản lý \\
\hline
WORKSPACE\_OWNER & Workspace & Quản lý Workspace, phân quyền \\
\hline
PROJECT\_MANAGER & Project & Quản lý Project, Sprint, phê duyệt AI \\
\hline
DEVELOPER & Project & Xem Project, quản lý Task được giao \\
\hline
TESTER & Project & Báo cáo lỗi, kiểm thử \\
\hline
VIEWER & Project & Chỉ xem dữ liệu \\
\hline
\end{tabular}
\label{tab:roles_permissions}
\end{table}

\begin{verbatim}
@PreAuthorize("hasRole('PROJECT_MANAGER') or hasRole('ADMIN')")
@PostMapping("/projects")
public ResponseEntity<ProjectResponse> createProject(@RequestBody CreateProjectRequest request) {
    // Only PROJECT_MANAGER and ADMIN can create projects
}

@PreAuthorize("hasRole('DEVELOPER') and @securityService.isTaskAssignee(#taskId)")
@PatchMapping("/tasks/{taskId}/status")
public ResponseEntity<TaskResponse> updateTaskStatus(
        @PathVariable UUID taskId,
        @RequestBody UpdateTaskStatusRequest request) {
    // Only the assigned developer can update task status
}

@PreAuthorize("hasRole('PROJECT_MANAGER') or hasRole('WORKSPACE_OWNER')")
@PostMapping("/projects/{projectId}/ai/analyze")
public ResponseEntity<AIAnalysisResponse> analyzeWithAI(@PathVariable UUID projectId) {
    // Only PROJECT_MANAGER or WORKSPACE_OWNER can trigger AI analysis
}
\end{verbatim}

%%%%%%%%%%%%%%%%%%%%%%%%%%%%%%%%%%%%%%%%%%%%%%%%%%%%%%%%%%%%%%
\section{Thiết kế Service Layer}
%%%%%%%%%%%%%%%%%%%%%%%%%%%%%%%%%%%%%%%%%%%%%%%%%%%%%%%%%%%%%%

\subsection{Tổng quan Service Layer}

Service Layer là nơi chứa toàn bộ logic nghiệp vụ của hệ thống. Mỗi service tương ứng với một domain hoặc một nhóm chức năng.

\subsection{AuthService}

\begin{verbatim}
@Service
@Transactional
public class AuthService {
    
    @Autowired
    private UserRepository userRepository;
    
    @Autowired
    private JwtTokenProvider tokenProvider;
    
    @Autowired
    private PasswordEncoder passwordEncoder;
    
    public AuthResponse login(LoginRequest request) {
        // 1. Tìm user theo email
        User user = userRepository.findByEmail(request.getEmail())
            .orElseThrow(() -> new BusinessException("Email không tồn tại"));
        
        // 2. Kiểm tra trạng thái
        if (user.getStatus() == UserStatus.BLOCKED) {
            throw new BusinessException("Tài khoản đã bị khóa");
        }
        
        // 3. Kiểm tra mật khẩu
        if (!passwordEncoder.matches(request.getPassword(), user.getPasswordHash())) {
            throw new BusinessException("Mật khẩu không đúng");
        }
        
        // 4. Tạo token
        UserDetails userDetails = new org.springframework.security.core.userdetails.User(
            user.getEmail(),
            user.getPasswordHash(),
            getAuthorities(user)
        );
        String accessToken = tokenProvider.generateAccessToken(userDetails);
        String refreshToken = tokenProvider.generateRefreshToken(user.getEmail());
        
        // 5. Cập nhật last_login
        user.setLastLoginAt(LocalDateTime.now());
        userRepository.save(user);
        
        return AuthResponse.builder()
            .accessToken(accessToken)
            .refreshToken(refreshToken)
            .tokenType("Bearer")
            .expiresIn(900) // 15 phút
            .user(UserDTO.fromEntity(user))
            .build();
    }
    
    public User register(RegisterRequest request) {
        // Kiểm tra email đã tồn tại
        if (userRepository.existsByEmail(request.getEmail())) {
            throw new BusinessException("Email đã được sử dụng");
        }
        
        // Tạo user mới
        User user = User.builder()
            .email(request.getEmail())
            .passwordHash(passwordEncoder.encode(request.getPassword()))
            .fullName(request.getFullName())
            .username(request.getUsername())
            .role(UserRole.DEVELOPER)
            .status(UserStatus.PENDING_VERIFICATION)
            .build();
        
        return userRepository.save(user);
    }
    
    public AuthResponse refreshToken(RefreshTokenRequest request) {
        if (!tokenProvider.validateToken(request.getRefreshToken())) {
            throw new BusinessException("Refresh token không hợp lệ");
        }
        
        String email = tokenProvider.getEmailFromToken(request.getRefreshToken());
        User user = userRepository.findByEmail(email)
            .orElseThrow(() -> new BusinessException("User không tồn tại"));
        
        UserDetails userDetails = new org.springframework.security.core.userdetails.User(
            user.getEmail(),
            user.getPasswordHash(),
            getAuthorities(user)
        );
        String newAccessToken = tokenProvider.generateAccessToken(userDetails);
        
        return AuthResponse.builder()
            .accessToken(newAccessToken)
            .refreshToken(request.getRefreshToken())
            .tokenType("Bearer")
            .expiresIn(900)
            .build();
    }
    
    private Collection<? extends GrantedAuthority> getAuthorities(User user) {
        return List.of(new SimpleGrantedAuthority("ROLE_" + user.getRole().name()));
    }
}
\end{verbatim}

\subsection{ProjectService}

\begin{verbatim}
@Service
@Transactional
public class ProjectService {
    
    @Autowired
    private ProjectRepository projectRepository;
    
    @Autowired
    private WorkspaceRepository workspaceRepository;
    
    @Autowired
    private ProjectMemberRepository projectMemberRepository;
    
    @Autowired
    private AuditLogService auditLogService;
    
    public Project createProject(CreateProjectRequest request, UUID workspaceId, UUID creatorId) {
        // 1. Kiểm tra Workspace tồn tại
        Workspace workspace = workspaceRepository.findById(workspaceId)
            .orElseThrow(() -> new ResourceNotFoundException("Workspace không tồn tại"));
        
        // 2. Kiểm tra quyền
        if (!workspace.getOwnerId().equals(creatorId)) {
            throw new BusinessException("Không có quyền tạo Project trong Workspace này");
        }
        
        // 3. Kiểm tra key trùng lặp
        if (projectRepository.existsByWorkspaceIdAndKey(workspaceId, request.getKey())) {
            throw new BusinessException("Project Key đã tồn tại trong Workspace");
        }
        
        // 4. Tạo Project
        Project project = Project.builder()
            .workspace(workspace)
            .name(request.getName())
            .key(request.getKey())
            .description(request.getDescription())
            .status(ProjectStatus.PLANNING)
            .priority(request.getPriority())
            .visibility(request.getVisibility())
            .startDate(request.getStartDate())
            .endDate(request.getEndDate())
            .build();
        
        Project savedProject = projectRepository.save(project);
        
        // 5. Thêm creator là Project Manager
        ProjectMember member = ProjectMember.builder()
            .project(savedProject)
            .userId(creatorId)
            .role(ProjectRole.PROJECT_MANAGER)
            .build();
        projectMemberRepository.save(member);
        
        // 6. Ghi audit log
        auditLogService.logAction(creatorId, "CREATE", "PROJECT", savedProject.getId(), null, savedProject);
        
        return savedProject;
    }
    
    public Project updateProject(UUID projectId, UpdateProjectRequest request, UUID userId) {
        Project project = findById(projectId);
        
        // Kiểm tra quyền
        if (!isProjectManager(projectId, userId)) {
            throw new BusinessException("Không có quyền cập nhật Project");
        }
        
        // Lưu giá trị cũ cho audit
        Project oldProject = project.toBuilder().build();
        
        // Cập nhật
        if (request.getName() != null) {
            project.setName(request.getName());
        }
        if (request.getDescription() != null) {
            project.setDescription(request.getDescription());
        }
        if (request.getStatus() != null) {
            project.setStatus(request.getStatus());
        }
        if (request.getPriority() != null) {
            project.setPriority(request.getPriority());
        }
        if (request.getEndDate() != null) {
            project.setEndDate(request.getEndDate());
        }
        
        Project updatedProject = projectRepository.save(project);
        
        // Ghi audit log
        auditLogService.logAction(userId, "UPDATE", "PROJECT", projectId, oldProject, updatedProject);
        
        return updatedProject;
    }
    
    public void deleteProject(UUID projectId, UUID userId) {
        Project project = findById(projectId);
        
        // Chỉ Owner mới có quyền xóa
        if (!isProjectOwner(projectId, userId)) {
            throw new BusinessException("Chỉ Project Owner mới có quyền xóa Project");
        }
        
        project.setStatus(ProjectStatus.ARCHIVED);
        project.setDeletedAt(LocalDateTime.now());
        projectRepository.save(project);
        
        auditLogService.logAction(userId, "DELETE", "PROJECT", projectId, project, null);
    }
    
    public Project findById(UUID projectId) {
        return projectRepository.findById(projectId)
            .orElseThrow(() -> new ResourceNotFoundException("Project không tồn tại"));
    }
    
    public List<Project> getProjectsByWorkspace(UUID workspaceId, UUID userId) {
        // Kiểm tra quyền truy cập workspace
        if (!hasWorkspaceAccess(workspaceId, userId)) {
            throw new BusinessException("Không có quyền truy cập Workspace");
        }
        return projectRepository.findByWorkspaceId(workspaceId);
    }
    
    private boolean isProjectManager(UUID projectId, UUID userId) {
        return projectMemberRepository.existsByProjectIdAndUserIdAndRole(
            projectId, userId, ProjectRole.PROJECT_MANAGER
        );
    }
    
    private boolean isProjectOwner(UUID projectId, UUID userId) {
        Project project = findById(projectId);
        return project.getWorkspace().getOwnerId().equals(userId);
    }
    
    private boolean hasWorkspaceAccess(UUID workspaceId, UUID userId) {
        // Kiểm tra thành viên workspace hoặc là owner
        return workspaceMemberRepository.existsByWorkspaceIdAndUserId(workspaceId, userId)
            || workspaceRepository.existsByIdAndOwnerId(workspaceId, userId);
    }
}
\end{verbatim}

\subsection{SprintService}

\begin{verbatim}
@Service
@Transactional
public class SprintService {
    
    @Autowired
    private SprintRepository sprintRepository;
    
    @Autowired
    private ProjectRepository projectRepository;
    
    @Autowired
    private UserStoryRepository userStoryRepository;
    
    @Autowired
    private AuditLogService auditLogService;
    
    public Sprint createSprint(CreateSprintRequest request, UUID projectId, UUID userId) {
        Project project = projectRepository.findById(projectId)
            .orElseThrow(() -> new ResourceNotFoundException("Project không tồn tại"));
        
        // Kiểm tra quyền
        if (!isProjectManager(projectId, userId)) {
            throw new BusinessException("Không có quyền tạo Sprint");
        }
        
        Sprint sprint = Sprint.builder()
            .project(project)
            .name(request.getName())
            .goal(request.getGoal())
            .startDate(request.getStartDate())
            .endDate(request.getEndDate())
            .status(SprintStatus.PLANNING)
            .build();
        
        Sprint savedSprint = sprintRepository.save(sprint);
        
        auditLogService.logAction(userId, "CREATE", "SPRINT", savedSprint.getId(), null, savedSprint);
        
        return savedSprint;
    }
    
    public Sprint startSprint(UUID sprintId, UUID userId) {
        Sprint sprint = findById(sprintId);
        
        // Kiểm tra quyền
        if (!isProjectManager(sprint.getProject().getId(), userId)) {
            throw new BusinessException("Không có quyền bắt đầu Sprint");
        }
        
        // Kiểm tra trạng thái
        if (sprint.getStatus() != SprintStatus.PLANNING) {
            throw new BusinessException("Sprint đã được bắt đầu hoặc đã kết thúc");
        }
        
        // Kiểm tra có User Story không
        long storyCount = userStoryRepository.countBySprintId(sprintId);
        if (storyCount == 0) {
            throw new BusinessException("Sprint phải có ít nhất 1 User Story");
        }
        
        sprint.setStatus(SprintStatus.ACTIVE);
        sprint.setStartDate(LocalDate.now());
        
        Sprint updatedSprint = sprintRepository.save(sprint);
        
        auditLogService.logAction(userId, "START", "SPRINT", sprintId, null, updatedSprint);
        
        return updatedSprint;
    }
    
    public Sprint completeSprint(UUID sprintId, UUID userId) {
        Sprint sprint = findById(sprintId);
        
        if (!isProjectManager(sprint.getProject().getId(), userId)) {
            throw new BusinessException("Không có quyền kết thúc Sprint");
        }
        
        if (sprint.getStatus() != SprintStatus.ACTIVE) {
            throw new BusinessException("Sprint không ở trạng thái ACTIVE");
        }
        
        // Tính Velocity
        int totalStoryPoints = userStoryRepository.sumStoryPointsBySprintId(sprintId);
        int completedStoryPoints = userStoryRepository.sumCompletedStoryPointsBySprintId(sprintId);
        int velocity = completedStoryPoints;
        
        sprint.setVelocity(velocity);
        sprint.setStatus(SprintStatus.COMPLETED);
        sprint.setActualEndDate(LocalDate.now());
        
        Sprint updatedSprint = sprintRepository.save(sprint);
        
        auditLogService.logAction(userId, "COMPLETE", "SPRINT", sprintId, null, updatedSprint);
        
        return updatedSprint;
    }
    
    public BurndownData getBurndownData(UUID sprintId) {
        Sprint sprint = findById(sprintId);
        
        // Lấy dữ liệu burndown từ bảng daily_log hoặc tính từ tasks
        List<BurndownPoint> points = calculateBurndown(sprint);
        
        return BurndownData.builder()
            .sprintId(sprintId)
            .sprintName(sprint.getName())
            .startDate(sprint.getStartDate())
            .endDate(sprint.getEndDate())
            .totalStoryPoints(userStoryRepository.sumStoryPointsBySprintId(sprintId))
            .points(points)
            .build();
    }
    
    private List<BurndownPoint> calculateBurndown(Sprint sprint) {
        // Logic tính toán burndown từ dữ liệu tasks
        // ...
        return new ArrayList<>();
    }
    
    public Sprint findById(UUID sprintId) {
        return sprintRepository.findById(sprintId)
            .orElseThrow(() -> new ResourceNotFoundException("Sprint không tồn tại"));
    }
    
    private boolean isProjectManager(UUID projectId, UUID userId) {
        // Kiểm tra quyền Project Manager
        return projectMemberRepository.existsByProjectIdAndUserIdAndRole(
            projectId, userId, ProjectRole.PROJECT_MANAGER
        );
    }
}
\end{verbatim}

\subsection{TaskService}

\begin{verbatim}
@Service
@Transactional
public class TaskService {
    
    @Autowired
    private TaskRepository taskRepository;
    
    @Autowired
    private UserStoryRepository userStoryRepository;
    
    @Autowired
    private UserRepository userRepository;
    
    @Autowired
    private NotificationService notificationService;
    
    @Autowired
    private AuditLogService auditLogService;
    
    public Task createTask(CreateTaskRequest request, UUID storyId, UUID userId) {
        UserStory story = userStoryRepository.findById(storyId)
            .orElseThrow(() -> new ResourceNotFoundException("User Story không tồn tại"));
        
        Task task = Task.builder()
            .userStory(story)
            .title(request.getTitle())
            .description(request.getDescription())
            .estimatedHours(request.getEstimatedHours())
            .priority(request.getPriority())
            .dueDate(request.getDueDate())
            .status(TaskStatus.TODO)
            .build();
        
        Task savedTask = taskRepository.save(task);
        
        auditLogService.logAction(userId, "CREATE", "TASK", savedTask.getId(), null, savedTask);
        
        return savedTask;
    }
    
    public Task assignTask(UUID taskId, UUID assigneeId, UUID userId) {
        Task task = findById(taskId);
        
        // Kiểm tra quyền (Project Manager hoặc Workspace Owner)
        if (!canAssignTask(task, userId)) {
            throw new BusinessException("Không có quyền phân công Task");
        }
        
        User assignee = userRepository.findById(assigneeId)
            .orElseThrow(() -> new ResourceNotFoundException("Developer không tồn tại"));
        
        Task oldTask = task.toBuilder().build();
        task.setAssignee(assignee);
        
        Task updatedTask = taskRepository.save(task);
        
        // Gửi thông báo cho Developer
        notificationService.sendTaskAssignmentNotification(assigneeId, updatedTask);
        
        auditLogService.logAction(userId, "ASSIGN", "TASK", taskId, oldTask, updatedTask);
        
        return updatedTask;
    }
    
    public Task updateTaskStatus(UUID taskId, TaskStatus newStatus, UUID userId) {
        Task task = findById(taskId);
        
        // Kiểm tra quyền: chỉ assignee hoặc Project Manager mới được cập nhật
        if (!canUpdateStatus(task, userId)) {
            throw new BusinessException("Không có quyền cập nhật trạng thái Task");
        }
        
        Task oldTask = task.toBuilder().build();
        task.setStatus(newStatus);
        
        if (newStatus == TaskStatus.IN_PROGRESS && task.getStartedAt() == null) {
            task.setStartedAt(LocalDateTime.now());
        }
        
        if (newStatus == TaskStatus.DONE && task.getCompletedAt() == null) {
            task.setCompletedAt(LocalDateTime.now());
        }
        
        Task updatedTask = taskRepository.save(task);
        
        auditLogService.logAction(userId, "UPDATE_STATUS", "TASK", taskId, oldTask, updatedTask);
        
        return updatedTask;
    }
    
    public Task addComment(UUID taskId, AddCommentRequest request, UUID userId) {
        Task task = findById(taskId);
        
        Comment comment = Comment.builder()
            .task(task)
            .userId(userId)
            .content(request.getContent())
            .build();
        
        // Lưu comment
        commentRepository.save(comment);
        
        // Gửi thông báo
        if (task.getAssignee() != null && !task.getAssignee().getId().equals(userId)) {
            notificationService.sendCommentNotification(task.getAssignee().getId(), task, comment);
        }
        
        return task;
    }
    
    public Task findById(UUID taskId) {
        return taskRepository.findById(taskId)
            .orElseThrow(() -> new ResourceNotFoundException("Task không tồn tại"));
    }
    
    public List<Task> getTasksByStory(UUID storyId) {
        return taskRepository.findByUserStoryId(storyId);
    }
    
    public List<Task> getTasksByAssignee(UUID assigneeId) {
        return taskRepository.findByAssigneeId(assigneeId);
    }
    
    private boolean canAssignTask(Task task, UUID userId) {
        // Project Manager hoặc Workspace Owner
        return isProjectManager(task.getUserStory().getSprint().getProject().getId(), userId)
            || isWorkspaceOwner(task.getUserStory().getSprint().getProject().getWorkspace().getId(), userId);
    }
    
    private boolean canUpdateStatus(Task task, UUID userId) {
        // Assignee hoặc Project Manager
        return (task.getAssignee() != null && task.getAssignee().getId().equals(userId))
            || isProjectManager(task.getUserStory().getSprint().getProject().getId(), userId);
    }
}
\end{verbatim}

\subsection{AIService - Tích hợp AI}

\begin{verbatim}
@Service
public class AIService {
    
    @Autowired
    private AIServiceClient aiServiceClient;
    
    @Autowired
    private ProjectRepository projectRepository;
    
    @Autowired
    private UserStoryRepository userStoryRepository;
    
    @Autowired
    private TaskRepository taskRepository;
    
    @Autowired
    private AISuggestionRepository aiSuggestionRepository;
    
    @Autowired
    private AuditLogService auditLogService;
    
    public AIAnalysisResponse analyzeRequirement(UUID projectId, String requirement, UUID userId) {
        // 1. Lấy context dự án
        Project project = projectRepository.findById(projectId)
            .orElseThrow(() -> new ResourceNotFoundException("Project không tồn tại"));
        
        // 2. Lấy dữ liệu liên quan
        List<UserStory> existingStories = userStoryRepository.findByProjectId(projectId);
        List<Task> existingTasks = taskRepository.findByProjectId(projectId);
        
        // 3. Xây dựng request cho AI Service
        AnalyzeRequest request = AnalyzeRequest.builder()
            .projectId(projectId)
            .projectName(project.getName())
            .domain(project.getSettings().get("domain"))
            .technologies(project.getSettings().get("technologies"))
            .requirement(requirement)
            .existingStories(existingStories)
            .existingTasks(existingTasks)
            .build();
        
        // 4. Gọi AI Service
        try {
            AnalysisResult result = aiServiceClient.analyzeRequirement(request);
            
            // 5. Lưu kết quả
            AISuggestion suggestion = AISuggestion.builder()
                .projectId(projectId)
                .agentType(AgentType.REQUIREMENT)
                .targetType("REQUIREMENT")
                .prompt(request.toString())
                .response(result.getRawResponse())
                .confidence(result.getConfidence())
                .build();
            aiSuggestionRepository.save(suggestion);
            
            // 6. Ghi audit log
            auditLogService.logAction(userId, "AI_ANALYZE", "REQUIREMENT", projectId, null, result);
            
            return result;
            
        } catch (Exception e) {
            throw new BusinessException("Lỗi khi gọi AI Service: " + e.getMessage());
        }
    }
    
    public StoryPointResponse estimateStoryPoint(UUID storyId, UUID userId) {
        UserStory story = userStoryRepository.findById(storyId)
            .orElseThrow(() -> new ResourceNotFoundException("User Story không tồn tại"));
        
        // Lấy dữ liệu lịch sử
        List<Sprint> historicalSprints = sprintRepository.findByProjectId(story.getSprint().getProject().getId());
        List<UserStory> historicalStories = userStoryRepository.findByProjectId(story.getSprint().getProject().getId());
        
        EstimateRequest request = EstimateRequest.builder()
            .story(story)
            .historicalSprints(historicalSprints)
            .historicalStories(historicalStories)
            .build();
        
        try {
            StoryPointResponse result = aiServiceClient.estimateStoryPoint(request);
            
            // Lưu kết quả
            AISuggestion suggestion = AISuggestion.builder()
                .projectId(story.getSprint().getProject().getId())
                .agentType(AgentType.STORY_POINT)
                .targetType("USER_STORY")
                .targetId(storyId)
                .response(result.getRawResponse())
                .confidence(result.getConfidence())
                .build();
            aiSuggestionRepository.save(suggestion);
            
            return result;
            
        } catch (Exception e) {
            throw new BusinessException("Lỗi khi ước lượng Story Point: " + e.getMessage());
        }
    }
    
    public ForecastResponse forecastSprint(UUID sprintId, UUID userId) {
        Sprint sprint = sprintRepository.findById(sprintId)
            .orElseThrow(() -> new ResourceNotFoundException("Sprint không tồn tại"));
        
        // Lấy dữ liệu burndown và velocity
        List<Sprint> previousSprints = sprintRepository.findPreviousSprints(sprintId, 5);
        List<BurndownPoint> burndownData = getBurndownData(sprintId);
        
        ForecastRequest request = ForecastRequest.builder()
            .sprint(sprint)
            .previousSprints(previousSprints)
            .burndownData(burndownData)
            .build();
        
        try {
            ForecastResponse result = aiServiceClient.forecastSprint(request);
            
            // Lưu kết quả
            AISuggestion suggestion = AISuggestion.builder()
                .projectId(sprint.getProject().getId())
                .agentType(AgentType.FORECAST)
                .targetType("SPRINT")
                .targetId(sprintId)
                .response(result.getRawResponse())
                .confidence(result.getConfidence())
                .build();
            aiSuggestionRepository.save(suggestion);
            
            return result;
            
        } catch (Exception e) {
            throw new BusinessException("Lỗi khi dự báo Sprint: " + e.getMessage());
        }
    }
    
    public RiskResponse detectRisk(UUID projectId, UUID userId) {
        // Thu thập dữ liệu
        List<Sprint> sprints = sprintRepository.findByProjectId(projectId);
        List<Task> tasks = taskRepository.findByProjectId(projectId);
        
        // Xây dựng request
        RiskRequest request = RiskRequest.builder()
            .projectId(projectId)
            .sprints(sprints)
            .tasks(tasks)
            .build();
        
        try {
            RiskResponse result = aiServiceClient.detectRisk(request);
            
            // Lưu kết quả
            AISuggestion suggestion = AISuggestion.builder()
                .projectId(projectId)
                .agentType(AgentType.RISK)
                .targetType("PROJECT")
                .targetId(projectId)
                .response(result.getRawResponse())
                .confidence(result.getConfidence())
                .build();
            aiSuggestionRepository.save(suggestion);
            
            return result;
            
        } catch (Exception e) {
            throw new BusinessException("Lỗi khi phát hiện rủi ro: " + e.getMessage());
        }
    }
}
\end{verbatim}

\subsection{AIServiceClient - Gọi AI Service}

\begin{verbatim}
@Component
public class AIServiceClient {
    
    @Value("${ai.service.url}")
    private String aiServiceUrl;
    
    @Autowired
    private RestTemplate restTemplate;
    
    public AnalysisResult analyzeRequirement(AnalyzeRequest request) {
        String url = aiServiceUrl + "/api/analyze-requirement";
        
        try {
            ResponseEntity<AnalysisResult> response = restTemplate.postForEntity(
                url,
                request,
                AnalysisResult.class
            );
            
            if (!response.getStatusCode().is2xxSuccessful()) {
                throw new AIServiceException("AI Service returned error: " + response.getStatusCode());
            }
            
            return response.getBody();
            
        } catch (RestClientException e) {
            throw new AIServiceException("Failed to call AI Service: " + e.getMessage());
        }
    }
    
    public StoryPointResponse estimateStoryPoint(EstimateRequest request) {
        String url = aiServiceUrl + "/api/estimate-story-point";
        
        try {
            ResponseEntity<StoryPointResponse> response = restTemplate.postForEntity(
                url,
                request,
                StoryPointResponse.class
            );
            return response.getBody();
        } catch (RestClientException e) {
            throw new AIServiceException("Failed to estimate story point: " + e.getMessage());
        }
    }
    
    public ForecastResponse forecastSprint(ForecastRequest request) {
        String url = aiServiceUrl + "/api/forecast-sprint";
        
        try {
            ResponseEntity<ForecastResponse> response = restTemplate.postForEntity(
                url,
                request,
                ForecastResponse.class
            );
            return response.getBody();
        } catch (RestClientException e) {
            throw new AIServiceException("Failed to forecast sprint: " + e.getMessage());
        }
    }
    
    public RiskResponse detectRisk(RiskRequest request) {
        String url = aiServiceUrl + "/api/detect-risk";
        
        try {
            ResponseEntity<RiskResponse> response = restTemplate.postForEntity(
                url,
                request,
                RiskResponse.class
            );
            return response.getBody();
        } catch (RestClientException e) {
            throw new AIServiceException("Failed to detect risk: " + e.getMessage());
        }
    }
}
\end{verbatim}

%%%%%%%%%%%%%%%%%%%%%%%%%%%%%%%%%%%%%%%%%%%%%%%%%%%%%%%%%%%%%%
\section{Thiết kế Repository Layer}
%%%%%%%%%%%%%%%%%%%%%%%%%%%%%%%%%%%%%%%%%%%%%%%%%%%%%%%%%%%%%%

\subsection{Tổng quan Repository}

Repository Layer sử dụng Spring Data JPA để tương tác với cơ sở dữ liệu.

\subsection{UserRepository}

\begin{verbatim}
@Repository
public interface UserRepository extends JpaRepository<User, UUID> {
    
    Optional<User> findByEmail(String email);
    
    Optional<User> findByUsername(String username);
    
    boolean existsByEmail(String email);
    
    boolean existsByUsername(String username);
    
    @Query("SELECT u FROM User u WHERE u.status = 'ACTIVE' AND u.deletedAt IS NULL")
    List<User> findActiveUsers();
    
    @Query("SELECT u FROM User u WHERE u.role = :role AND u.status = 'ACTIVE'")
    List<User> findByRole(@Param("role") UserRole role);
    
    @Query("SELECT u FROM User u WHERE LOWER(u.fullName) LIKE LOWER(CONCAT('%', :keyword, '%'))")
    List<User> searchByKeyword(@Param("keyword") String keyword);
}
\end{verbatim}

\subsection{ProjectRepository}

\begin{verbatim}
@Repository
public interface ProjectRepository extends JpaRepository<Project, UUID> {
    
    List<Project> findByWorkspaceId(UUID workspaceId);
    
    Optional<Project> findByWorkspaceIdAndKey(UUID workspaceId, String key);
    
    boolean existsByWorkspaceIdAndKey(UUID workspaceId, String key);
    
    @Query("SELECT p FROM Project p WHERE p.workspace.id = :workspaceId AND p.status = :status")
    List<Project> findByWorkspaceIdAndStatus(@Param("workspaceId") UUID workspaceId, 
                                              @Param("status") ProjectStatus status);
    
    @Query("SELECT p FROM Project p WHERE p.workspace.id = :workspaceId AND p.deletedAt IS NULL")
    List<Project> findActiveByWorkspaceId(@Param("workspaceId") UUID workspaceId);
    
    @Query("SELECT p FROM Project p WHERE p.workspace.id = :workspaceId " +
           "AND (LOWER(p.name) LIKE LOWER(CONCAT('%', :keyword, '%')) " +
           "OR LOWER(p.description) LIKE LOWER(CONCAT('%', :keyword, '%')))")
    List<Project> searchByKeyword(@Param("workspaceId") UUID workspaceId, 
                                   @Param("keyword") String keyword);
}
\end{verbatim}

\subsection{UserStoryRepository}

\begin{verbatim}
@Repository
public interface UserStoryRepository extends JpaRepository<UserStory, UUID> {
    
    List<UserStory> findBySprintId(UUID sprintId);
    
    List<UserStory> findByEpicId(UUID epicId);
    
    List<UserStory> findByProjectId(UUID projectId);
    
    @Query("SELECT SUM(us.storyPoint) FROM UserStory us WHERE us.sprint.id = :sprintId")
    Integer sumStoryPointsBySprintId(@Param("sprintId") UUID sprintId);
    
    @Query("SELECT SUM(us.storyPoint) FROM UserStory us WHERE us.sprint.id = :sprintId AND us.status = 'DONE'")
    Integer sumCompletedStoryPointsBySprintId(@Param("sprintId") UUID sprintId);
    
    @Query("SELECT COUNT(us) FROM UserStory us WHERE us.sprint.id = :sprintId AND us.status = 'TODO'")
    Long countTodoBySprintId(@Param("sprintId") UUID sprintId);
    
    @Query("SELECT us FROM UserStory us WHERE us.sprint.id = :sprintId ORDER BY us.orderInBacklog")
    List<UserStory> findBySprintIdOrdered(@Param("sprintId") UUID sprintId);
    
    @Query("SELECT us FROM UserStory us WHERE us.sprint.id IS NULL AND us.status = 'TODO' ORDER BY us.priority")
    List<UserStory> findTodoInBacklog(@Param("projectId") UUID projectId);
}
\end{verbatim}

\subsection{TaskRepository}

\begin{verbatim}
@Repository
public interface TaskRepository extends JpaRepository<Task, UUID> {
    
    List<Task> findByUserStoryId(UUID storyId);
    
    List<Task> findByAssigneeId(UUID assigneeId);
    
    List<Task> findByProjectId(UUID projectId);
    
    @Query("SELECT t FROM Task t WHERE t.assignee.id = :assigneeId AND t.status = 'TODO'")
    List<Task> findTodoByAssignee(@Param("assigneeId") UUID assigneeId);
    
    @Query("SELECT t FROM Task t WHERE t.assignee.id = :assigneeId AND t.status = 'IN_PROGRESS'")
    List<Task> findInProgressByAssignee(@Param("assigneeId") UUID assigneeId);
    
    @Query("SELECT COUNT(t) FROM Task t WHERE t.assignee.id = :assigneeId AND t.status != 'DONE'")
    Long countActiveTasksByAssignee(@Param("assigneeId") UUID assigneeId);
    
    @Query("SELECT t FROM Task t WHERE t.userStory.sprint.id = :sprintId ORDER BY t.priority DESC")
    List<Task> findBySprintId(@Param("sprintId") UUID sprintId);
}
\end{verbatim}

%%%%%%%%%%%%%%%%%%%%%%%%%%%%%%%%%%%%%%%%%%%%%%%%%%%%%%%%%%%%%%
\section{Xử lý ngoại lệ (Exception Handling)}
%%%%%%%%%%%%%%%%%%%%%%%%%%%%%%%%%%%%%%%%%%%%%%%%%%%%%%%%%%%%%%

\subsection{Global Exception Handler}

\begin{verbatim}
@RestControllerAdvice
public class GlobalExceptionHandler {
    
    private static final Logger logger = LoggerFactory.getLogger(GlobalExceptionHandler.class);
    
    @ExceptionHandler(ResourceNotFoundException.class)
    @ResponseStatus(HttpStatus.NOT_FOUND)
    public ResponseEntity<ErrorResponse> handleResourceNotFound(ResourceNotFoundException ex) {
        logger.warn("Resource not found: {}", ex.getMessage());
        return buildErrorResponse(HttpStatus.NOT_FOUND, "RESOURCE_NOT_FOUND", ex.getMessage());
    }
    
    @ExceptionHandler(BusinessException.class)
    @ResponseStatus(HttpStatus.BAD_REQUEST)
    public ResponseEntity<ErrorResponse> handleBusinessException(BusinessException ex) {
        logger.warn("Business error: {}", ex.getMessage());
        return buildErrorResponse(HttpStatus.BAD_REQUEST, "BUSINESS_ERROR", ex.getMessage());
    }
    
    @ExceptionHandler(AIServiceException.class)
    @ResponseStatus(HttpStatus.SERVICE_UNAVAILABLE)
    public ResponseEntity<ErrorResponse> handleAIServiceException(AIServiceException ex) {
        logger.error("AI Service error: {}", ex.getMessage());
        return buildErrorResponse(HttpStatus.SERVICE_UNAVAILABLE, "AI_SERVICE_ERROR", "AI Service temporarily unavailable");
    }
    
    @ExceptionHandler(AccessDeniedException.class)
    @ResponseStatus(HttpStatus.FORBIDDEN)
    public ResponseEntity<ErrorResponse> handleAccessDenied(AccessDeniedException ex) {
        logger.warn("Access denied: {}", ex.getMessage());
        return buildErrorResponse(HttpStatus.FORBIDDEN, "ACCESS_DENIED", "You don't have permission");
    }
    
    @ExceptionHandler(MethodArgumentNotValidException.class)
    @ResponseStatus(HttpStatus.BAD_REQUEST)
    public ResponseEntity<ValidationErrorResponse> handleValidationException(MethodArgumentNotValidException ex) {
        Map<String, String> errors = new HashMap<>();
        ex.getBindingResult().getFieldErrors().forEach(error -> 
            errors.put(error.getField(), error.getDefaultMessage())
        );
        
        ValidationErrorResponse response = ValidationErrorResponse.builder()
            .success(false)
            .errorCode("VALIDATION_ERROR")
            .message("Validation failed")
            .errors(errors)
            .build();
        
        return ResponseEntity.badRequest().body(response);
    }
    
    @ExceptionHandler(Exception.class)
    @ResponseStatus(HttpStatus.INTERNAL_SERVER_ERROR)
    public ResponseEntity<ErrorResponse> handleGenericException(Exception ex) {
        logger.error("Unexpected error occurred", ex);
        return buildErrorResponse(
            HttpStatus.INTERNAL_SERVER_ERROR,
            "INTERNAL_ERROR",
            "An unexpected error occurred"
        );
    }
    
    private ResponseEntity<ErrorResponse> buildErrorResponse(HttpStatus status, String errorCode, String message) {
        ErrorResponse response = ErrorResponse.builder()
            .success(false)
            .errorCode(errorCode)
            .message(message)
            .timestamp(LocalDateTime.now())
            .build();
        return ResponseEntity.status(status).body(response);
    }
}
\end{verbatim}

\subsection{Error Response DTO}

\begin{verbatim}
@Data
@Builder
public class ErrorResponse {
    private boolean success;
    private String errorCode;
    private String message;
    private LocalDateTime timestamp;
    
    @JsonInclude(JsonInclude.Include.NON_NULL)
    private Map<String, String> details;
}

@Data
@Builder
public class ValidationErrorResponse {
    private boolean success;
    private String errorCode;
    private String message;
    private Map<String, String> errors;
    private LocalDateTime timestamp;
}
\end{verbatim}

%%%%%%%%%%%%%%%%%%%%%%%%%%%%%%%%%%%%%%%%%%%%%%%%%%%%%%%%%%%%%%
\section{Logging và Monitoring}
%%%%%%%%%%%%%%%%%%%%%%%%%%%%%%%%%%%%%%%%%%%%%%%%%%%%%%%%%%%%%%

\subsection{Logging Configuration}

\begin{verbatim}
# application.yml
logging:
  level:
    com.ai.pms: DEBUG
    org.springframework.security: INFO
    org.hibernate: INFO
  pattern:
    console: "%d{yyyy-MM-dd HH:mm:ss} - %level - %msg%n"
    file: "%d{yyyy-MM-dd HH:mm:ss} - %level - [%thread] - %class{36} - %msg%n"
  file:
    name: logs/application.log
    max-size: 10MB
    max-history: 30
\end{verbatim}

\subsection{Audit Log Service}

\begin{verbatim}
@Service
@Transactional
public class AuditLogService {
    
    @Autowired
    private AuditLogRepository auditLogRepository;
    
    @Autowired
    private SecurityUtils securityUtils;
    
    public void logAction(UUID userId, String action, String entityType, 
                          UUID entityId, Object oldValue, Object newValue) {
        AuditLog log = AuditLog.builder()
            .userId(userId)
            .action(action)
            .entityType(entityType)
            .entityId(entityId)
            .oldValues(convertToJson(oldValue))
            .newValues(convertToJson(newValue))
            .ipAddress(securityUtils.getClientIp())
            .userAgent(securityUtils.getUserAgent())
            .createdAt(LocalDateTime.now())
            .build();
        
        auditLogRepository.save(log);
    }
    
    public void logAction(String action, String entityType, UUID entityId, 
                          Object oldValue, Object newValue) {
        UUID userId = securityUtils.getCurrentUserId();
        logAction(userId, action, entityType, entityId, oldValue, newValue);
    }
    
    private JsonNode convertToJson(Object value) {
        if (value == null) return null;
        ObjectMapper mapper = new ObjectMapper();
        return mapper.convertValue(value, JsonNode.class);
    }
    
    public List<AuditLog> getAuditLogsForEntity(String entityType, UUID entityId) {
        return auditLogRepository.findByEntityTypeAndEntityIdOrderByCreatedAtDesc(entityType, entityId);
    }
    
    public List<AuditLog> getAuditLogsForUser(UUID userId) {
        return auditLogRepository.findByUserIdOrderByCreatedAtDesc(userId);
    }
}
\end{verbatim}

%%%%%%%%%%%%%%%%%%%%%%%%%%%%%%%%%%%%%%%%%%%%%%%%%%%%%%%%%%%%%%
\section{Tổng kết chương}
%%%%%%%%%%%%%%%%%%%%%%%%%%%%%%%%%%%%%%%%%%%%%%%%%%%%%%%%%%%%%%

Trong chương này, luận văn đã trình bày thiết kế chi tiết tầng Backend của hệ thống AI Project Management.

Các nội dung chính đã được thực hiện:

\begin{enumerate}
\item \textbf{Kiến trúc Backend}: Clean Architecture với 4 tầng (Presentation, Application, Domain, Infrastructure).
\item \textbf{Security}: JWT Authentication, RBAC, Spring Security configuration.
\item \textbf{Service Layer}: Thiết kế chi tiết các Service (AuthService, ProjectService, SprintService, TaskService, AIService).
\item \textbf{Repository Layer}: Sử dụng Spring Data JPA với các repository interfaces.
\item \textbf{Integration với AI}: AIServiceClient gọi FastAPI service.
\item \textbf{Xử lý ngoại lệ}: GlobalExceptionHandler với các error codes.
\item \textbf{Logging và Audit}: Ghi log và audit trail cho mọi thao tác.
\end{enumerate}

Đây là cơ sở để Chương 11 trình bày thiết kế Frontend và Chương 12 trình bày triển khai hệ thống.

%%%%%%%%%%%%%%%%%%%%%%%%%%%%%%%%%%%%%%%%%%%%%%%%%%%%%%%%%%%%%%

%%%%%%%%%%%%%%%%%%%%%%%%%%%%%%%%%%%%%%%%%%%%%%%%%%%%%%%%%%%%%%
%%%%%%%%%%%%%%%%%%%%%%%%%%%%%%%%%%%%%%%%%%%%%%%%%%%%%%%%%%%%%%
\chapter{THIẾT KẾ FRONTEND}
\label{chap:frontend_design}
%%%%%%%%%%%%%%%%%%%%%%%%%%%%%%%%%%%%%%%%%%%%%%%%%%%%%%%%%%%%%%

\section{Giới thiệu chương}

Sau khi trình bày thiết kế Backend ở Chương 10, chương này tập trung vào thiết kế Frontend của hệ thống AI Project Management. Frontend là thành phần giao tiếp trực tiếp với người dùng, cung cấp giao diện trực quan và tương tác với các tính năng của hệ thống.

Mục tiêu của chương bao gồm:

\begin{itemize}
\item Trình bày kiến trúc tổng thể của Frontend.
\item Thiết kế các thành phần giao diện (UI Components).
\item Thiết kế luồng điều hướng (Navigation Flow).
\item Thiết kế trạng thái và quản lý dữ liệu (State Management).
\item Thiết kế tích hợp API với Backend.
\item Thiết kế responsive và trải nghiệm người dùng (UX).
\item Thiết kế các trang chính của hệ thống.
\end{itemize}

Đây là cơ sở để Chương 12 trình bày triển khai hệ thống.

%%%%%%%%%%%%%%%%%%%%%%%%%%%%%%%%%%%%%%%%%%%%%%%%%%%%%%%%%%%%%%
\section{Kiến trúc Frontend tổng thể}
%%%%%%%%%%%%%%%%%%%%%%%%%%%%%%%%%%%%%%%%%%%%%%%%%%%%%%%%%%%%%%

\subsection{Mô hình kiến trúc}

Frontend được xây dựng theo kiến trúc Component-Based với React và TypeScript, tuân theo mô hình MVC (Model-View-Controller) trên client.

\begin{figure}[H]
\centering
\fbox{
\begin{minipage}{0.95\textwidth}
\centering
\begin{tabular}{c}
\textbf{Tầng View (Components)} \\
- Pages \\
- Layouts \\
- UI Components \\
- HOCs (Higher-Order Components) \\
$\downarrow$ \\
\textbf{Tầng Controller (Hooks \& Services)} \\
- Custom Hooks \\
- API Services \\
- Context Providers \\
$\downarrow$ \\
\textbf{Tầng Model (State \& Data)} \\
- React Context \\
- Redux / Zustand \\
- Local Storage \\
- API Data
\end{tabular}
\end{minipage}
}
\caption{Kiến trúc Frontend}
\label{fig:frontend_architecture}
\end{figure}

\subsection{Công nghệ sử dụng}

\begin{table}[H]
\centering
\caption{Công nghệ Frontend}
\begin{tabular}{|p{4cm}|p{5cm}|p{6cm}|}
\hline
\textbf{Thành phần} & \textbf{Công nghệ} & \textbf{Vai trò} \\
\hline
Framework & React 19 & Xây dựng giao diện người dùng \\
\hline
Ngôn ngữ & TypeScript 5.x & Lập trình type-safe \\
\hline
Build Tool & Vite 5.x & Build và dev server \\
\hline
State Management & Zustand / Redux Toolkit & Quản lý trạng thái toàn cục \\
\hline
Routing & React Router 6.x & Điều hướng trang \\
\hline
UI Library & Material-UI / Ant Design & Components UI \\
\hline
Styling & Tailwind CSS & CSS utility-first \\
\hline
HTTP Client & Axios & Gọi API \\
\hline
Form Handling & React Hook Form & Quản lý form \\
\hline
Validation & Zod & Validation schema \\
\hline
Charts & Chart.js / Recharts & Biểu đồ \\
\hline
Testing & Jest + React Testing Library & Kiểm thử \\
\hline
\end{tabular}
\label{tab:frontend_technologies}
\end{table}

\subsection{Cấu trúc thư mục dự án}

\begin{verbatim}
src/
├── main.tsx                    # Entry point
├── App.tsx                     # Root component
├── index.css                   # Global styles
├── api/                        # API Services
│   ├── client.ts               # Axios instance
│   ├── auth.ts                 # Auth API
│   ├── workspace.ts            # Workspace API
│   ├── project.ts              # Project API
│   ├── sprint.ts               # Sprint API
│   ├── backlog.ts              # Backlog API
│   ├── task.ts                 # Task API
│   ├── ai.ts                   # AI API
│   └── dashboard.ts            # Dashboard API
├── assets/                     # Static assets
│   ├── images/
│   ├── fonts/
│   └── icons/
├── components/                 # Reusable Components
│   ├── common/                 # Common components
│   │   ├── Button/
│   │   ├── Input/
│   │   ├── Card/
│   │   ├── Modal/
│   │   ├── Table/
│   │   ├── Loading/
│   │   └── Toast/
│   ├── layout/                 # Layout components
│   │   ├── MainLayout/
│   │   ├── AuthLayout/
│   │   ├── Header/
│   │   ├── Sidebar/
│   │   └── Footer/
│   └── shared/                 # Shared business components
│       ├── KanbanBoard/
│       ├── BurndownChart/
│       ├── VelocityChart/
│       ├── AIPanel/
│       ├── CommentSection/
│       └── IssueCard/
├── pages/                      # Page components
│   ├── auth/
│   │   ├── Login/
│   │   ├── Register/
│   │   └── ForgotPassword/
│   ├── dashboard/
│   │   └── Dashboard/
│   ├── workspace/
│   │   ├── WorkspaceList/
│   │   ├── WorkspaceDetail/
│   │   └── WorkspaceSettings/
│   ├── project/
│   │   ├── ProjectList/
│   │   ├── ProjectDetail/
│   │   └── ProjectSettings/
│   ├── sprint/
│   │   ├── SprintList/
│   │   └── SprintDetail/
│   ├── backlog/
│   │   ├── BacklogList/
│   │   ├── EpicDetail/
│   │   └── StoryDetail/
│   ├── task/
│   │   ├── TaskBoard/
│   │   └── TaskDetail/
│   └── ai/
│       ├── AIAnalysis/
│       └── AIReports/
├── hooks/                      # Custom Hooks
│   ├── useAuth.ts
│   ├── useApi.ts
│   ├── useToast.ts
│   ├── useModal.ts
│   └── useDebounce.ts
├── store/                      # State Management
│   ├── authStore.ts
│   ├── workspaceStore.ts
│   ├── projectStore.ts
│   └── uiStore.ts
├── types/                      # TypeScript Types
│   ├── api.ts
│   ├── auth.ts
│   ├── project.ts
│   ├── sprint.ts
│   ├── task.ts
│   └── ai.ts
├── utils/                      # Utilities
│   ├── constants.ts
│   ├── helpers.ts
│   ├── validators.ts
│   └── formatters.ts
├── routes/                     # Route configuration
│   └── index.tsx
└── config/                     # Configuration
    ├── env.ts
    └── api.ts
\end{verbatim}

%%%%%%%%%%%%%%%%%%%%%%%%%%%%%%%%%%%%%%%%%%%%%%%%%%%%%%%%%%%%%%
\section{Thiết kế Quản lý trạng thái (State Management)}
%%%%%%%%%%%%%%%%%%%%%%%%%%%%%%%%%%%%%%%%%%%%%%%%%%%%%%%%%%%%%%

\subsection{Zustand Store - Auth Store}

\begin{verbatim}
// store/authStore.ts
import { create } from 'zustand';
import { persist } from 'zustand/middleware';
import { AuthService } from '@/api/auth';

interface AuthState {
  user: User | null;
  token: string | null;
  isAuthenticated: boolean;
  isLoading: boolean;
  login: (email: string, password: string) => Promise<void>;
  logout: () => void;
  refreshToken: () => Promise<void>;
  setUser: (user: User) => void;
}

export const useAuthStore = create<AuthState>()(
  persist(
    (set, get) => ({
      user: null,
      token: null,
      isAuthenticated: false,
      isLoading: false,

      login: async (email: string, password: string) => {
        set({ isLoading: true });
        try {
          const response = await AuthService.login({ email, password });
          set({
            user: response.user,
            token: response.accessToken,
            isAuthenticated: true,
            isLoading: false,
          });
        } catch (error) {
          set({ isLoading: false });
          throw error;
        }
      },

      logout: () => {
        set({
          user: null,
          token: null,
          isAuthenticated: false,
        });
        localStorage.removeItem('auth-storage');
      },

      refreshToken: async () => {
        try {
          const response = await AuthService.refreshToken();
          set({ token: response.accessToken });
        } catch (error) {
          get().logout();
        }
      },

      setUser: (user: User) => {
        set({ user });
      },
    }),
    {
      name: 'auth-storage',
      partialize: (state) => ({
        user: state.user,
        token: state.token,
        isAuthenticated: state.isAuthenticated,
      }),
    }
  )
);
\end{verbatim}

\subsection{Zustand Store - Project Store}

\begin{verbatim}
// store/projectStore.ts
import { create } from 'zustand';
import { ProjectService } from '@/api/project';

interface ProjectState {
  projects: Project[];
  currentProject: Project | null;
  isLoading: boolean;
  fetchProjects: (workspaceId: string) => Promise<void>;
  fetchProject: (projectId: string) => Promise<void>;
  createProject: (data: CreateProjectData) => Promise<void>;
  updateProject: (projectId: string, data: UpdateProjectData) => Promise<void>;
  deleteProject: (projectId: string) => Promise<void>;
  setCurrentProject: (project: Project) => void;
  clearProjects: () => void;
}

export const useProjectStore = create<ProjectState>((set, get) => ({
  projects: [],
  currentProject: null,
  isLoading: false,

  fetchProjects: async (workspaceId: string) => {
    set({ isLoading: true });
    try {
      const projects = await ProjectService.getProjectsByWorkspace(workspaceId);
      set({ projects, isLoading: false });
    } catch (error) {
      set({ isLoading: false });
      throw error;
    }
  },

  fetchProject: async (projectId: string) => {
    set({ isLoading: true });
    try {
      const project = await ProjectService.getProject(projectId);
      set({ currentProject: project, isLoading: false });
    } catch (error) {
      set({ isLoading: false });
      throw error;
    }
  },

  createProject: async (data: CreateProjectData) => {
    set({ isLoading: true });
    try {
      const project = await ProjectService.createProject(data);
      set((state) => ({
        projects: [project, ...state.projects],
        isLoading: false,
      }));
    } catch (error) {
      set({ isLoading: false });
      throw error;
    }
  },

  updateProject: async (projectId: string, data: UpdateProjectData) => {
    set({ isLoading: true });
    try {
      const updated = await ProjectService.updateProject(projectId, data);
      set((state) => ({
        projects: state.projects.map((p) =>
          p.id === projectId ? updated : p
        ),
        currentProject: state.currentProject?.id === projectId ? updated : state.currentProject,
        isLoading: false,
      }));
    } catch (error) {
      set({ isLoading: false });
      throw error;
    }
  },

  deleteProject: async (projectId: string) => {
    set({ isLoading: true });
    try {
      await ProjectService.deleteProject(projectId);
      set((state) => ({
        projects: state.projects.filter((p) => p.id !== projectId),
        currentProject: state.currentProject?.id === projectId ? null : state.currentProject,
        isLoading: false,
      }));
    } catch (error) {
      set({ isLoading: false });
      throw error;
    }
  },

  setCurrentProject: (project: Project) => {
    set({ currentProject: project });
  },

  clearProjects: () => {
    set({ projects: [], currentProject: null });
  },
}));
\end{verbatim}

%%%%%%%%%%%%%%%%%%%%%%%%%%%%%%%%%%%%%%%%%%%%%%%%%%%%%%%%%%%%%%
\section{Thiết kế API Integration}
%%%%%%%%%%%%%%%%%%%%%%%%%%%%%%%%%%%%%%%%%%%%%%%%%%%%%%%%%%%%%%

\subsection{Axios Client Configuration}

\begin{verbatim}
// api/client.ts
import axios, { AxiosInstance, AxiosError } from 'axios';
import { useAuthStore } from '@/store/authStore';

const API_BASE_URL = import.meta.env.VITE_API_BASE_URL || 'http://localhost:8080/api/v1';

class ApiClient {
  private instance: AxiosInstance;

  constructor() {
    this.instance = axios.create({
      baseURL: API_BASE_URL,
      timeout: 30000,
      headers: {
        'Content-Type': 'application/json',
      },
    });

    this.setupInterceptors();
  }

  private setupInterceptors() {
    // Request interceptor - add token
    this.instance.interceptors.request.use(
      (config) => {
        const token = useAuthStore.getState().token;
        if (token) {
          config.headers.Authorization = `Bearer ${token}`;
        }
        return config;
      },
      (error) => Promise.reject(error)
    );

    // Response interceptor - handle errors
    this.instance.interceptors.response.use(
      (response) => response.data,
      async (error: AxiosError) => {
        if (error.response?.status === 401) {
          // Token expired - try refresh
          try {
            const { refreshToken } = useAuthStore.getState();
            if (refreshToken) {
              const response = await AuthService.refreshToken();
              useAuthStore.getState().setToken(response.accessToken);
              // Retry original request
              return this.instance.request(error.config!);
            }
          } catch (refreshError) {
            useAuthStore.getState().logout();
            window.location.href = '/login';
          }
        }
        return Promise.reject(error);
      }
    );
  }

  public getInstance(): AxiosInstance {
    return this.instance;
  }
}

export const apiClient = new ApiClient().getInstance();
\end{verbatim}

\subsection{API Services}

\begin{verbatim}
// api/auth.ts
import { apiClient } from './client';

export interface LoginRequest {
  email: string;
  password: string;
}

export interface LoginResponse {
  accessToken: string;
  refreshToken: string;
  tokenType: string;
  expiresIn: number;
  user: User;
}

export interface RegisterRequest {
  email: string;
  password: string;
  fullName: string;
  username: string;
}

export const AuthService = {
  login: (data: LoginRequest): Promise<LoginResponse> =>
    apiClient.post('/auth/login', data),

  register: (data: RegisterRequest): Promise<User> =>
    apiClient.post('/auth/register', data),

  logout: (): Promise<void> =>
    apiClient.post('/auth/logout'),

  refreshToken: (): Promise<{ accessToken: string }> =>
    apiClient.post('/auth/refresh'),

  getProfile: (): Promise<User> =>
    apiClient.get('/auth/profile'),

  updateProfile: (data: Partial<User>): Promise<User> =>
    apiClient.put('/auth/profile', data),

  forgotPassword: (email: string): Promise<void> =>
    apiClient.post('/auth/forgot-password', { email }),

  resetPassword: (token: string, password: string): Promise<void> =>
    apiClient.post('/auth/reset-password', { token, password }),
};

// api/project.ts
export interface CreateProjectRequest {
  workspaceId: string;
  name: string;
  key: string;
  description?: string;
  priority: 'LOW' | 'MEDIUM' | 'HIGH';
  visibility: 'PUBLIC' | 'PRIVATE';
  startDate?: string;
  endDate?: string;
}

export interface UpdateProjectRequest {
  name?: string;
  description?: string;
  status?: 'PLANNING' | 'ACTIVE' | 'DONE' | 'ARCHIVED';
  priority?: 'LOW' | 'MEDIUM' | 'HIGH';
  endDate?: string;
}

export const ProjectService = {
  getProjectsByWorkspace: (workspaceId: string): Promise<Project[]> =>
    apiClient.get(`/workspaces/${workspaceId}/projects`),

  getProject: (projectId: string): Promise<Project> =>
    apiClient.get(`/projects/${projectId}`),

  createProject: (data: CreateProjectRequest): Promise<Project> =>
    apiClient.post('/projects', data),

  updateProject: (projectId: string, data: UpdateProjectRequest): Promise<Project> =>
    apiClient.patch(`/projects/${projectId}`, data),

  deleteProject: (projectId: string): Promise<void> =>
    apiClient.delete(`/projects/${projectId}`),

  archiveProject: (projectId: string): Promise<void> =>
    apiClient.post(`/projects/${projectId}/archive`),

  getProjectMembers: (projectId: string): Promise<ProjectMember[]> =>
    apiClient.get(`/projects/${projectId}/members`),

  addProjectMember: (projectId: string, userId: string, role: string): Promise<void> =>
    apiClient.post(`/projects/${projectId}/members`, { userId, role }),

  removeProjectMember: (projectId: string, userId: string): Promise<void> =>
    apiClient.delete(`/projects/${projectId}/members/${userId}`),
};

// api/sprint.ts
export interface CreateSprintRequest {
  name: string;
  goal: string;
  startDate: string;
  endDate: string;
}

export const SprintService = {
  getSprintsByProject: (projectId: string): Promise<Sprint[]> =>
    apiClient.get(`/projects/${projectId}/sprints`),

  getSprint: (sprintId: string): Promise<Sprint> =>
    apiClient.get(`/sprints/${sprintId}`),

  createSprint: (projectId: string, data: CreateSprintRequest): Promise<Sprint> =>
    apiClient.post(`/projects/${projectId}/sprints`, data),

  updateSprint: (sprintId: string, data: Partial<CreateSprintRequest>): Promise<Sprint> =>
    apiClient.patch(`/sprints/${sprintId}`, data),

  deleteSprint: (sprintId: string): Promise<void> =>
    apiClient.delete(`/sprints/${sprintId}`),

  startSprint: (sprintId: string): Promise<Sprint> =>
    apiClient.post(`/sprints/${sprintId}/start`),

  completeSprint: (sprintId: string): Promise<Sprint> =>
    apiClient.post(`/sprints/${sprintId}/complete`),

  getBurndown: (sprintId: string): Promise<BurndownData> =>
    apiClient.get(`/sprints/${sprintId}/burndown`),

  getVelocity: (projectId: string): Promise<VelocityData> =>
    apiClient.get(`/projects/${projectId}/velocity`),
};

// api/task.ts
export interface CreateTaskRequest {
  storyId: string;
  title: string;
  description?: string;
  estimatedHours?: number;
  priority: 'LOW' | 'MEDIUM' | 'HIGH' | 'URGENT';
  dueDate?: string;
}

export const TaskService = {
  getTasksByStory: (storyId: string): Promise<Task[]> =>
    apiClient.get(`/stories/${storyId}/tasks`),

  getTask: (taskId: string): Promise<Task> =>
    apiClient.get(`/tasks/${taskId}`),

  createTask: (data: CreateTaskRequest): Promise<Task> =>
    apiClient.post('/tasks', data),

  updateTask: (taskId: string, data: Partial<Task>): Promise<Task> =>
    apiClient.patch(`/tasks/${taskId}`, data),

  deleteTask: (taskId: string): Promise<void> =>
    apiClient.delete(`/tasks/${taskId}`),

  assignTask: (taskId: string, assigneeId: string): Promise<Task> =>
    apiClient.post(`/tasks/${taskId}/assign`, { assigneeId }),

  updateTaskStatus: (taskId: string, status: string): Promise<Task> =>
    apiClient.patch(`/tasks/${taskId}/status`, { status }),

  addComment: (taskId: string, content: string): Promise<Comment> =>
    apiClient.post(`/tasks/${taskId}/comments`, { content }),

  getComments: (taskId: string): Promise<Comment[]> =>
    apiClient.get(`/tasks/${taskId}/comments`),

  getMyTasks: (): Promise<Task[]> =>
    apiClient.get('/tasks/my'),
};

// api/ai.ts
export const AIService = {
  analyzeRequirement: (projectId: string, requirement: string): Promise<AIAnalysisResult> =>
    apiClient.post(`/projects/${projectId}/ai/analyze`, { requirement }),

  estimateStoryPoint: (storyId: string): Promise<StoryPointEstimate> =>
    apiClient.post(`/stories/${storyId}/ai/estimate`),

  forecastSprint: (sprintId: string): Promise<ForecastResult> =>
    apiClient.post(`/sprints/${sprintId}/ai/forecast`),

  detectRisk: (projectId: string): Promise<RiskResult> =>
    apiClient.post(`/projects/${projectId}/ai/risk`),

  generateDocument: (projectId: string, type: string): Promise<DocumentResult> =>
    apiClient.post(`/projects/${projectId}/ai/document`, { type }),

  getDashboardInsights: (projectId: string): Promise<InsightResult> =>
    apiClient.get(`/projects/${projectId}/ai/insights`),
};
\end{verbatim}

%%%%%%%%%%%%%%%%%%%%%%%%%%%%%%%%%%%%%%%%%%%%%%%%%%%%%%%%%%%%%%
\section{Thiết kế các trang chính}
%%%%%%%%%%%%%%%%%%%%%%%%%%%%%%%%%%%%%%%%%%%%%%%%%%%%%%%%%%%%%%

\subsection{Trang Login}

\begin{figure}[H]
\centering
\fbox{
\begin{minipage}{0.95\textwidth}
\centering
\begin{tabular}{c}
\textbf{Logo + Tên hệ thống} \\
$\downarrow$ \\
\textbf{Form Đăng nhập} \\
- Email Input \\
- Password Input \\
- Remember Me checkbox \\
$\downarrow$ \\
\textbf{Nút Đăng nhập} \\
$\downarrow$ \\
\textbf{Liên kết: Quên mật khẩu?} \\
\textbf{Liên kết: Đăng ký tài khoản mới}
\end{tabular}
\end{minipage}
}
\caption{Thiết kế trang Login}
\label{fig:login_page}
\end{figure}

\subsection{Trang Dashboard}

\begin{figure}[H]
\centering
\fbox{
\begin{minipage}{0.95\textwidth}
\centering
\begin{tabular}{c}
\textbf{Header: Logo | Search | Notification | User Profile} \\
$\downarrow$ \\
\textbf{Stats Cards: Tổng Project | Active Sprint | Tasks | Members} \\
$\downarrow$ \\
\textbf{Biểu đồ Velocity} \hspace{4cm} \textbf{Biểu đồ Burndown} \\
$\downarrow$ \\
\textbf{AI Insights} \\
$\downarrow$ \\
\textbf{Recent Activities}
\end{tabular}
\end{minipage}
}
\caption{Thiết kế trang Dashboard}
\label{fig:dashboard_page}
\end{figure}

\subsection{Trang Project Detail}

\begin{figure}[H]
\centering
\fbox{
\begin{minipage}{0.95\textwidth}
\centering
\begin{tabular}{c}
\textbf{Project Header: Tên | Key | Status | Priority | Actions} \\
$\downarrow$ \\
\textbf{Project Info: Description | Dates | Members} \\
$\downarrow$ \\
\textbf{Tabs: Overview | Sprints | Backlog | Tasks | Settings} \\
$\downarrow$ \\
\textbf{Nội dung theo tab}
\end{tabular}
\end{minipage}
}
\caption{Thiết kế trang Project Detail}
\label{fig:project_detail_page}
\end{figure}

\subsection{Trang Backlog}

\begin{figure}[H]
\centering
\fbox{
\begin{minipage}{0.95\textwidth}
\centering
\begin{tabular}{c}
\textbf{Header: Backlog | Filters | Search | Create Story} \\
$\downarrow$ \\
\textbf{Epics List} \\
$\downarrow$ \\
\textbf{User Stories (drag \& drop)} \\
- Priority | Title | Story Point | Status | Assignee \\
$\downarrow$ \\
\textbf{Bulk Actions}
\end{tabular}
\end{minipage}
}
\caption{Thiết kế trang Backlog}
\label{fig:backlog_page}
\end{figure}

\subsection{Trang Kanban Board}

\begin{figure}[H]
\centering
\fbox{
\begin{minipage}{0.95\textwidth}
\centering
\begin{tabular}{c}
\textbf{Header: Board Name | Filters | Add Task} \\
$\downarrow$ \\
\textbf{Columns: TODO | IN PROGRESS | REVIEW | DONE} \\
$\downarrow$ \\
\textbf{Task Cards trong mỗi column} \\
- Title | Assignee | Priority | Labels | Due Date
\end{tabular}
\end{minipage}
}
\caption{Thiết kế trang Kanban Board}
\label{fig:kanban_page}
\end{figure}

\subsection{Trang Sprint Detail}

\begin{figure}[H]
\centering
\fbox{
\begin{minipage}{0.95\textwidth}
\centering
\begin{tabular}{c}
\textbf{Sprint Header: Tên | Goal | Dates | Status} \\
$\downarrow$ \\
\textbf{Sprint Metrics: Velocity | Completion | Remaining} \\
$\downarrow$ \\
\textbf{Biểu đồ Burndown} \\
$\downarrow$ \\
\textbf{User Stories trong Sprint} \\
$\downarrow$ \\
\textbf{AI Forecast Panel}
\end{tabular}
\end{minipage}
}
\caption{Thiết kế trang Sprint Detail}
\label{fig:sprint_detail_page}
\end{figure}

\subsection{Trang AI Analysis}

\begin{figure}[H]
\centering
\fbox{
\begin{minipage}{0.95\textwidth}
\centering
\begin{tabular}{c}
\textbf{Header: AI Analysis | Project Selector} \\
$\downarrow$ \\
\textbf{Requirement Input} \\
- Text Area \\
- Analyze Button \\
$\downarrow$ \\
\textbf{Analysis Results} \\
- Summary \\
- User Stories Sinh ra \\
- Acceptance Criteria \\
- Missing Info \\
$\downarrow$ \\
\textbf{Phê duyệt / Chỉnh sửa}
\end{tabular}
\end{minipage}
}
\caption{Thiết kế trang AI Analysis}
\label{fig:ai_analysis_page}
\end{figure}

%%%%%%%%%%%%%%%%%%%%%%%%%%%%%%%%%%%%%%%%%%%%%%%%%%%%%%%%%%%%%%
\section{Thiết kế Component}
%%%%%%%%%%%%%%%%%%%%%%%%%%%%%%%%%%%%%%%%%%%%%%%%%%%%%%%%%%%%%%

\subsection{KanbanBoard Component}

\begin{verbatim}
// components/shared/KanbanBoard/KanbanBoard.tsx
import React, { useState } from 'react';
import { DragDropContext, Droppable, Draggable } from 'react-beautiful-dnd';

interface KanbanBoardProps {
  columns: Column[];
  onTaskMove: (taskId: string, fromColumn: string, toColumn: string) => void;
  onTaskClick: (taskId: string) => void;
  onAddTask: (columnId: string) => void;
}

export const KanbanBoard: React.FC<KanbanBoardProps> = ({
  columns,
  onTaskMove,
  onTaskClick,
  onAddTask,
}) => {
  const [isDragging, setIsDragging] = useState(false);

  const handleDragEnd = (result: any) => {
    setIsDragging(false);
    if (!result.destination) return;

    const { draggableId, source, destination } = result;
    if (source.droppableId === destination.droppableId) return;

    onTaskMove(draggableId, source.droppableId, destination.droppableId);
  };

  return (
    <DragDropContext onDragStart={() => setIsDragging(true)} onDragEnd={handleDragEnd}>
      <div className="kanban-board grid grid-cols-4 gap-4 p-4">
        {columns.map((column) => (
          <div key={column.id} className="kanban-column bg-gray-50 rounded-lg p-3">
            <div className="column-header flex justify-between items-center mb-3">
              <h3 className="font-semibold text-gray-700">
                {column.title}
                <span className="ml-2 text-sm text-gray-400">
                  ({column.tasks.length})
                </span>
              </h3>
              <button
                onClick={() => onAddTask(column.id)}
                className="text-blue-500 hover:text-blue-700"
              >
                +
              </button>
            </div>

            <Droppable droppableId={column.id}>
              {(provided, snapshot) => (
                <div
                  ref={provided.innerRef}
                  {...provided.droppableProps}
                  className={`task-list min-h-[200px] ${
                    snapshot.isDraggingOver ? 'bg-blue-50' : ''
                  }`}
                >
                  {column.tasks.map((task, index) => (
                    <Draggable
                      key={task.id}
                      draggableId={task.id}
                      index={index}
                    >
                      {(provided, snapshot) => (
                        <div
                          ref={provided.innerRef}
                          {...provided.draggableProps}
                          {...provided.dragHandleProps}
                          className={`task-card bg-white rounded-lg shadow-sm p-3 mb-2 cursor-pointer hover:shadow-md transition-shadow ${
                            snapshot.isDragging ? 'shadow-lg' : ''
                          }`}
                          onClick={() => onTaskClick(task.id)}
                        >
                          <div className="task-title text-sm font-medium text-gray-800">
                            {task.title}
                          </div>
                          <div className="task-meta flex justify-between items-center mt-2">
                            <div className="flex items-center gap-2">
                              <span
                                className={`priority-badge text-xs px-2 py-0.5 rounded ${
                                  task.priority === 'HIGH'
                                    ? 'bg-red-100 text-red-700'
                                    : task.priority === 'MEDIUM'
                                    ? 'bg-yellow-100 text-yellow-700'
                                    : 'bg-green-100 text-green-700'
                                }`}
                              >
                                {task.priority}
                              </span>
                              {task.storyPoint && (
                                <span className="story-point text-xs px-2 py-0.5 bg-purple-100 text-purple-700 rounded">
                                  {task.storyPoint}
                                </span>
                              )}
                            </div>
                            {task.assignee && (
                              <div className="assignee-avatar w-6 h-6 rounded-full bg-gray-300 flex items-center justify-center text-xs">
                                {task.assignee.name[0]}
                              </div>
                            )}
                          </div>
                          {task.dueDate && (
                            <div className="task-due-date text-xs text-gray-400 mt-1">
                              Due: {new Date(task.dueDate).toLocaleDateString()}
                            </div>
                          )}
                        </div>
                      )}
                    </Draggable>
                  ))}
                  {provided.placeholder}
                </div>
              )}
            </Droppable>
          </div>
        ))}
      </div>
    </DragDropContext>
  );
};
\end{verbatim}

\subsection{BurndownChart Component}

\begin{verbatim}
// components/shared/BurndownChart/BurndownChart.tsx
import React from 'react';
import {
  LineChart,
  Line,
  XAxis,
  YAxis,
  CartesianGrid,
  Tooltip,
  Legend,
  ResponsiveContainer,
} from 'recharts';

interface BurndownChartProps {
  data: BurndownPoint[];
  totalStoryPoints: number;
  sprintDays: number;
  title?: string;
}

export const BurndownChart: React.FC<BurndownChartProps> = ({
  data,
  totalStoryPoints,
  sprintDays,
  title = 'Burndown Chart',
}) => {
  // Tạo đường lý tưởng
  const idealLine = data.map((point, index) => ({
    day: point.day,
    remaining: totalStoryPoints - (totalStoryPoints / sprintDays) * point.day,
  }));

  return (
    <div className="burndown-chart bg-white rounded-lg shadow p-4">
      <h3 className="text-lg font-semibold mb-4">{title}</h3>
      <ResponsiveContainer width="100%" height={300}>
        <LineChart
          data={data}
          margin={{ top: 5, right: 30, left: 20, bottom: 5 }}
        >
          <CartesianGrid strokeDasharray="3 3" />
          <XAxis dataKey="day" label={{ value: 'Days', position: 'bottom' }} />
          <YAxis
            label={{ value: 'Story Points Remaining', angle: -90, position: 'left' }}
            domain={[0, totalStoryPoints]}
          />
          <Tooltip />
          <Legend />
          <Line
            type="monotone"
            dataKey="remaining"
            stroke="#3B82F6"
            strokeWidth={2}
            name="Actual"
            dot={{ r: 4 }}
          />
          <Line
            type="linear"
            data={idealLine}
            dataKey="remaining"
            stroke="#9CA3AF"
            strokeWidth={2}
            strokeDasharray="5 5"
            name="Ideal"
            dot={false}
          />
        </LineChart>
      </ResponsiveContainer>
      <div className="chart-metrics grid grid-cols-3 gap-4 mt-4 text-center">
        <div>
          <div className="text-sm text-gray-500">Total SP</div>
          <div className="text-lg font-bold">{totalStoryPoints}</div>
        </div>
        <div>
          <div className="text-sm text-gray-500">Remaining</div>
          <div className="text-lg font-bold">
            {data[data.length - 1]?.remaining ?? 0}
          </div>
        </div>
        <div>
          <div className="text-sm text-gray-500">Completion</div>
          <div className="text-lg font-bold">
            {data.length > 0
              ? Math.round(
                  ((totalStoryPoints - data[data.length - 1].remaining) /
                    totalStoryPoints) *
                    100
                )
              : 0}
            %
          </div>
        </div>
      </div>
    </div>
  );
};
\end{verbatim}

\subsection{AIPanel Component}

\begin{verbatim}
// components/shared/AIPanel/AIPanel.tsx
import React, { useState } from 'react';
import { AIService } from '@/api/ai';

interface AIPanelProps {
  projectId: string;
  onSuggestionApply: (suggestion: any) => void;
}

export const AIPanel: React.FC<AIPanelProps> = ({ projectId, onSuggestionApply }) => {
  const [requirement, setRequirement] = useState('');
  const [isAnalyzing, setIsAnalyzing] = useState(false);
  const [result, setResult] = useState<AIAnalysisResult | null>(null);
  const [activeTab, setActiveTab] = useState<'analysis' | 'forecast' | 'risk'>('analysis');

  const handleAnalyze = async () => {
    if (!requirement.trim()) return;

    setIsAnalyzing(true);
    try {
      const response = await AIService.analyzeRequirement(projectId, requirement);
      setResult(response);
    } catch (error) {
      console.error('Analysis failed:', error);
    } finally {
      setIsAnalyzing(false);
    }
  };

  return (
    <div className="ai-panel bg-white rounded-lg shadow">
      <div className="panel-header border-b p-4">
        <h3 className="text-lg font-semibold">AI Assistant</h3>
        <div className="tabs flex gap-2 mt-2">
          <button
            className={`px-3 py-1 rounded ${
              activeTab === 'analysis' ? 'bg-blue-500 text-white' : 'bg-gray-100'
            }`}
            onClick={() => setActiveTab('analysis')}
          >
            Analysis
          </button>
          <button
            className={`px-3 py-1 rounded ${
              activeTab === 'forecast' ? 'bg-blue-500 text-white' : 'bg-gray-100'
            }`}
            onClick={() => setActiveTab('forecast')}
          >
            Forecast
          </button>
          <button
            className={`px-3 py-1 rounded ${
              activeTab === 'risk' ? 'bg-blue-500 text-white' : 'bg-gray-100'
            }`}
            onClick={() => setActiveTab('risk')}
          >
            Risk
          </button>
        </div>
      </div>

      <div className="panel-body p-4">
        {activeTab === 'analysis' && (
          <div className="analysis-section">
            <div className="requirement-input">
              <label className="block text-sm font-medium text-gray-700 mb-1">
                Enter Requirement
              </label>
              <textarea
                value={requirement}
                onChange={(e) => setRequirement(e.target.value)}
                className="w-full border rounded-lg p-2 h-32 resize-none"
                placeholder="Describe the feature or requirement..."
              />
              <button
                onClick={handleAnalyze}
                disabled={isAnalyzing || !requirement.trim()}
                className="mt-2 px-4 py-2 bg-blue-500 text-white rounded hover:bg-blue-600 disabled:opacity-50"
              >
                {isAnalyzing ? 'Analyzing...' : 'Analyze with AI'}
              </button>
            </div>

            {result && (
              <div className="analysis-result mt-4 border-t pt-4">
                <div className="summary bg-blue-50 p-3 rounded-lg">
                  <h4 className="font-semibold">Summary</h4>
                  <p className="text-sm text-gray-700">{result.summary}</p>
                </div>

                {result.userStories.length > 0 && (
                  <div className="user-stories mt-3">
                    <h4 className="font-semibold">Suggested User Stories</h4>
                    {result.userStories.map((story, index) => (
                      <div
                        key={index}
                        className="story-card border rounded-lg p-3 mt-2 hover:bg-gray-50"
                      >
                        <div className="flex justify-between items-start">
                          <div>
                            <p className="font-medium">{story.title}</p>
                            <p className="text-sm text-gray-600">
                              As a {story.role}, I want {story.goal}, So that{' '}
                              {story.value}
                            </p>
                            <div className="mt-1">
                              <span className="text-xs font-medium text-gray-500">
                                Acceptance Criteria:
                              </span>
                              <ul className="text-sm text-gray-600 list-disc list-inside">
                                {story.acceptanceCriteria.map((criterion, i) => (
                                  <li key={i}>{criterion}</li>
                                ))}
                              </ul>
                            </div>
                          </div>
                          <div className="flex items-center gap-2">
                            <span className="px-2 py-1 bg-purple-100 text-purple-700 rounded text-xs">
                              SP: {story.storyPoint}
                            </span>
                            <button
                              onClick={() => onSuggestionApply(story)}
                              className="px-3 py-1 bg-green-500 text-white rounded text-sm hover:bg-green-600"
                            >
                              Apply
                            </button>
                          </div>
                        </div>
                      </div>
                    ))}
                  </div>
                )}

                {result.missingInfo.length > 0 && (
                  <div className="missing-info mt-3 bg-yellow-50 p-3 rounded-lg">
                    <h4 className="font-semibold text-yellow-700">
                      ⚠️ Need Clarification
                    </h4>
                    <ul className="text-sm text-yellow-700 list-disc list-inside">
                      {result.missingInfo.map((info, index) => (
                        <li key={index}>{info}</li>
                      ))}
                    </ul>
                  </div>
                )}
              </div>
            )}
          </div>
        )}

        {activeTab === 'forecast' && (
          <div className="forecast-section text-center text-gray-400">
            <p>AI Sprint Forecast will appear here</p>
          </div>
        )}

        {activeTab === 'risk' && (
          <div className="risk-section text-center text-gray-400">
            <p>AI Risk Detection will appear here</p>
          </div>
        )}
      </div>

      <div className="panel-footer border-t p-3 bg-gray-50 text-xs text-gray-400">
        <span>Powered by AI • Confidence score: {result?.confidence || 'N/A'}</span>
      </div>
    </div>
  );
};
\end{verbatim}

%%%%%%%%%%%%%%%%%%%%%%%%%%%%%%%%%%%%%%%%%%%%%%%%%%%%%%%%%%%%%%
\section{Thiết kế Routing và Navigation}
%%%%%%%%%%%%%%%%%%%%%%%%%%%%%%%%%%%%%%%%%%%%%%%%%%%%%%%%%%%%%%

\subsection{Route Configuration}

\begin{verbatim}
// routes/index.tsx
import { createBrowserRouter } from 'react-router-dom';
import { ProtectedRoute } from '@/components/common/ProtectedRoute';
import { MainLayout } from '@/components/layout/MainLayout';
import { AuthLayout } from '@/components/layout/AuthLayout';

// Pages
import { Login } from '@/pages/auth/Login';
import { Register } from '@/pages/auth/Register';
import { ForgotPassword } from '@/pages/auth/ForgotPassword';
import { Dashboard } from '@/pages/dashboard/Dashboard';
import { WorkspaceList } from '@/pages/workspace/WorkspaceList';
import { WorkspaceDetail } from '@/pages/workspace/WorkspaceDetail';
import { ProjectList } from '@/pages/project/ProjectList';
import { ProjectDetail } from '@/pages/project/ProjectDetail';
import { SprintDetail } from '@/pages/sprint/SprintDetail';
import { BacklogList } from '@/pages/backlog/BacklogList';
import { TaskBoard } from '@/pages/task/TaskBoard';
import { AIAnalysis } from '@/pages/ai/AIAnalysis';

export const router = createBrowserRouter([
  {
    path: '/',
    element: <ProtectedRoute />,
    children: [
      {
        element: <MainLayout />,
        children: [
          { path: '/', element: <Dashboard /> },
          { path: '/workspaces', element: <WorkspaceList /> },
          { path: '/workspaces/:workspaceId', element: <WorkspaceDetail /> },
          { path: '/projects', element: <ProjectList /> },
          { path: '/projects/:projectId', element: <ProjectDetail /> },
          { path: '/sprints/:sprintId', element: <SprintDetail /> },
          { path: '/backlog', element: <BacklogList /> },
          { path: '/tasks', element: <TaskBoard /> },
          { path: '/ai-analysis', element: <AIAnalysis /> },
        ],
      },
    ],
  },
  {
    path: '/auth',
    element: <AuthLayout />,
    children: [
      { path: 'login', element: <Login /> },
      { path: 'register', element: <Register /> },
      { path: 'forgot-password', element: <ForgotPassword /> },
    ],
  },
]);
\end{verbatim}

\subsection{ProtectedRoute Component}

\begin{verbatim}
// components/common/ProtectedRoute.tsx
import { Navigate, Outlet } from 'react-router-dom';
import { useAuthStore } from '@/store/authStore';

export const ProtectedRoute: React.FC = () => {
  const { isAuthenticated, token } = useAuthStore();

  if (!isAuthenticated || !token) {
    return <Navigate to="/auth/login" replace />;
  }

  return <Outlet />;
};
\end{verbatim}

%%%%%%%%%%%%%%%%%%%%%%%%%%%%%%%%%%%%%%%%%%%%%%%%%%%%%%%%%%%%%%
\section{Thiết kế Responsive và UX}
%%%%%%%%%%%%%%%%%%%%%%%%%%%%%%%%%%%%%%%%%%%%%%%%%%%%%%%%%%%%%%

\subsection{Responsive Breakpoints}

\begin{table}[H]
\centering
\caption{Responsive Breakpoints}
\begin{tabular}{|p{3.5cm}|p{4cm}|p{7cm}|}
\hline
\textbf{Breakpoint} & \textbf{Kích thước} & \textbf{Hiển thị} \\
\hline
Mobile & $< 640$px & Sidebar thu gọn, 1 cột \\
\hline
Tablet & $640 - 1024$px & Sidebar thu gọn, 2 cột \\
\hline
Desktop & $> 1024$px & Sidebar đầy đủ, đa cột \\
\hline
\end{tabular}
\label{tab:responsive_breakpoints}
\end{table}

\subsection{UX Design Principles}

\begin{enumerate}
\item \textbf{Simplicity}: Giao diện đơn giản, dễ sử dụng.
\item \textbf{Consistency}: Nhất quán trong thiết kế và tương tác.
\item \textbf{Feedback}: Phản hồi ngay lập tức cho mọi hành động của người dùng.
\item \textbf{Progressive Disclosure}: Hiển thị thông tin theo mức độ cần thiết.
\item \textbf{Error Prevention}: Ngăn chặn lỗi trước khi xảy ra.
\item \textbf{Flexibility}: Cho phép người dùng tùy chỉnh.
\item \textbf{Efficiency}: Tối ưu hóa thao tác cho người dùng thường xuyên.
\end{enumerate}

%%%%%%%%%%%%%%%%%%%%%%%%%%%%%%%%%%%%%%%%%%%%%%%%%%%%%%%%%%%%%%
\section{Tổng kết chương}
%%%%%%%%%%%%%%%%%%%%%%%%%%%%%%%%%%%%%%%%%%%%%%%%%%%%%%%%%%%%%%

Trong chương này, luận văn đã trình bày thiết kế chi tiết Frontend của hệ thống AI Project Management.

Các nội dung chính đã được thực hiện:

\begin{enumerate}
\item \textbf{Kiến trúc Frontend}: Component-Based với React và TypeScript.
\item \textbf{State Management}: Zustand cho quản lý trạng thái toàn cục.
\item \textbf{API Integration}: Axios client với interceptors cho token và error handling.
\item \textbf{Pages}: Thiết kế 7 trang chính (Login, Dashboard, Project, Backlog, Kanban, Sprint, AI Analysis).
\item \textbf{Components}: Thiết kế các component tái sử dụng (KanbanBoard, BurndownChart, AIPanel).
\item \textbf{Routing}: React Router với Protected Routes.
\item \textbf{Responsive \& UX}: Thiết kế đáp ứng và nguyên tắc UX.
\end{enumerate}

Đây là cơ sở để Chương 12 trình bày triển khai hệ thống và Chương 13 trình bày thực nghiệm.

%%%%%%%%%%%%%%%%%%%%%%%%%%%%%%%%%%%%%%%%%%%%%%%%%%%%%%%%%%%%%%
%%%%%%%%%%%%%%%%%%%%%%%%%%%%%%%%%%%%%%%%%%%%%%%%%%%%%%%%%%%%%%
%%%%%%%%%%%%%%%%%%%%%%%%%%%%%%%%%%%%%%%%%%%%%%%%%%%%%%%%%%%%%%
\chapter{TRIỂN KHAI HỆ THỐNG}
\label{chap:deployment}
%%%%%%%%%%%%%%%%%%%%%%%%%%%%%%%%%%%%%%%%%%%%%%%%%%%%%%%%%%%%%%

\section{Giới thiệu chương}

Sau khi trình bày thiết kế Frontend ở Chương 11, chương này tập trung vào quy trình triển khai hệ thống AI Project Management từ môi trường phát triển đến môi trường production. Triển khai đóng vai trò quan trọng trong việc đảm bảo hệ thống hoạt động ổn định, có thể mở rộng và dễ dàng bảo trì.

Mục tiêu của chương bao gồm:

\begin{itemize}
\item Trình bày môi trường phát triển và các công cụ sử dụng.
\item Thiết kế quy trình containerization với Docker.
\item Thiết kế orchestration với Docker Compose.
\item Thiết lập CI/CD Pipeline với GitHub Actions.
\item Hướng dẫn cài đặt và cấu hình từng thành phần.
\item Trình bày các lệnh quản lý và vận hành hệ thống.
\item Đề xuất kiến trúc triển khai production với Kubernetes.
\end{itemize}

Đây là cơ sở để Chương 13 trình bày thực nghiệm và đánh giá hệ thống.

%%%%%%%%%%%%%%%%%%%%%%%%%%%%%%%%%%%%%%%%%%%%%%%%%%%%%%%%%%%%%%
\section{Môi trường phát triển}
%%%%%%%%%%%%%%%%%%%%%%%%%%%%%%%%%%%%%%%%%%%%%%%%%%%%%%%%%%%%%%

\subsection{Yêu cầu phần mềm}

\begin{table}[H]
\centering
\caption{Yêu cầu phần mềm cho môi trường phát triển}
\begin{tabular}{|p{4cm}|p{4cm}|p{6cm}|}
\hline
\textbf{Công cụ} & \textbf{Phiên bản} & \textbf{Mục đích} \\
\hline
Node.js & 20.x trở lên & Chạy Frontend và Backend \\
\hline
Java & JDK 21 & Chạy Backend Spring Boot \\
\hline
Python & 3.11 trở lên & Chạy AI Service \\
\hline
Docker Desktop & 24.x trở lên & Containerization \\
\hline
Docker Compose & 2.x trở lên & Orchestration \\
\hline
PostgreSQL & 16.x & Cơ sở dữ liệu \\
\hline
Redis & 7.x & Cache \\
\hline
Git & 2.x & Quản lý mã nguồn \\
\hline
Google Chrome & Mới nhất & Kiểm thử Extension \\
\hline
\end{tabular}
\label{tab:dev_requirements}
\end{table}

\subsection{Cấu trúc thư mục dự án}

\begin{verbatim}
ai-project-management/
├── docker-compose.yml          # Docker Compose cho toàn bộ hệ thống
├── docker-compose.dev.yml      # Docker Compose cho môi trường dev
├── .env.example                # Mẫu biến môi trường
├── .github/
│   └── workflows/
│       ├── ci.yml              # CI Pipeline
│       └── cd.yml              # CD Pipeline
├── backend/
│   ├── Dockerfile              # Dockerfile cho Backend
│   ├── pom.xml / build.gradle  # Build configuration
│   ├── src/
│   ├── .env                    # Biến môi trường Backend
│   └── scripts/
│       └── entrypoint.sh       # Script khởi động
├── frontend/
│   ├── Dockerfile              # Dockerfile cho Frontend
│   ├── package.json
│   ├── vite.config.ts
│   ├── src/
│   └── nginx.conf              # Cấu hình Nginx (production)
├── ai-service/
│   ├── Dockerfile              # Dockerfile cho AI Service
│   ├── requirements.txt
│   ├── app/
│   └── models/                 # Thư mục chứa model đã huấn luyện
├── chrome-extension/
│   ├── manifest.json
│   ├── popup.html
│   ├── popup.js
│   └── content-script.js
├── database/
│   ├── init.sql                # Script khởi tạo database
│   └── migrations/             # Flyway migrations
├── docs/
│   ├── api/
│   └── deployment/
└── README.md
\end{verbatim}

%%%%%%%%%%%%%%%%%%%%%%%%%%%%%%%%%%%%%%%%%%%%%%%%%%%%%%%%%%%%%%
\section{Containerization với Docker}
%%%%%%%%%%%%%%%%%%%%%%%%%%%%%%%%%%%%%%%%%%%%%%%%%%%%%%%%%%%%%%

\subsection{Dockerfile cho Backend}

\begin{verbatim}
# backend/Dockerfile
FROM openjdk:21-jdk-slim AS build

WORKDIR /app

# Copy và build ứng dụng
COPY mvnw .
COPY .mvn .mvn
COPY pom.xml .
RUN ./mvnw dependency:go-offline -B

COPY src src
RUN ./mvnw package -DskipTests

# Runtime stage
FROM openjdk:21-jdk-slim

WORKDIR /app

# Copy JAR file từ build stage
COPY --from=build /app/target/*.jar app.jar

# Copy entrypoint script
COPY scripts/entrypoint.sh /entrypoint.sh
RUN chmod +x /entrypoint.sh

# Tạo user không root
RUN addgroup --system --gid 1001 appgroup && \
    adduser --system --uid 1001 --gid 1001 appuser

USER appuser

EXPOSE 8080

ENTRYPOINT ["/entrypoint.sh"]
\end{verbatim}

\subsection{Dockerfile cho Frontend (Production)}

\begin{verbatim}
# frontend/Dockerfile
FROM node:20-alpine AS build

WORKDIR /app

# Copy và cài đặt dependencies
COPY package*.json ./
RUN npm ci --only=production

# Copy source và build
COPY . .
RUN npm run build

# Production stage - Nginx
FROM nginx:alpine

# Copy Nginx configuration
COPY nginx.conf /etc/nginx/conf.d/default.conf

# Copy build files
COPY --from=build /app/dist /usr/share/nginx/html

EXPOSE 80

CMD ["nginx", "-g", "daemon off;"]
\end{verbatim}

\subsection{Dockerfile cho AI Service}

\begin{verbatim}
# ai-service/Dockerfile
FROM python:3.11-slim

WORKDIR /app

# Cài đặt dependencies
COPY requirements.txt .
RUN pip install --no-cache-dir -r requirements.txt

# Copy source code và model
COPY app/ app/
COPY models/ models/

# Tạo user không root
RUN addgroup --system --gid 1001 aigroup && \
    adduser --system --uid 1001 --gid 1001 aiuser

USER aiuser

EXPOSE 8000

CMD ["uvicorn", "app.main:app", "--host", "0.0.0.0", "--port", "8000"]
\end{verbatim}

\subsection{Dockerfile cho Database (PostgreSQL + pgvector)}

\begin{verbatim}
# database/Dockerfile
FROM postgres:16-alpine

# Cài đặt pgvector
RUN apk add --no-cache git build-base postgresql-dev && \
    git clone https://github.com/pgvector/pgvector.git /tmp/pgvector && \
    cd /tmp/pgvector && \
    make && make install && \
    apk del git build-base

# Copy init script
COPY init.sql /docker-entrypoint-initdb.d/

EXPOSE 5432
\end{verbatim}

%%%%%%%%%%%%%%%%%%%%%%%%%%%%%%%%%%%%%%%%%%%%%%%%%%%%%%%%%%%%%%
\section{Docker Compose Configuration}
%%%%%%%%%%%%%%%%%%%%%%%%%%%%%%%%%%%%%%%%%%%%%%%%%%%%%%%%%%%%%%

\subsection{docker-compose.yml}

\begin{verbatim}
# docker-compose.yml
version: '3.8'

services:
  # PostgreSQL Database
  postgres:
    build:
      context: ./database
      dockerfile: Dockerfile
    container_name: ai-pms-postgres
    environment:
      POSTGRES_USER: ${POSTGRES_USER:-ai_pms_user}
      POSTGRES_PASSWORD: ${POSTGRES_PASSWORD:-ai_pms_pass}
      POSTGRES_DB: ${POSTGRES_DB:-ai_pms_db}
    ports:
      - "5432:5432"
    volumes:
      - postgres_data:/var/lib/postgresql/data
    healthcheck:
      test: ["CMD-SHELL", "pg_isready -U ${POSTGRES_USER} -d ${POSTGRES_DB}"]
      interval: 10s
      timeout: 5s
      retries: 5
    networks:
      - ai-pms-network

  # Redis Cache
  redis:
    image: redis:7-alpine
    container_name: ai-pms-redis
    ports:
      - "6379:6379"
    volumes:
      - redis_data:/data
    healthcheck:
      test: ["CMD", "redis-cli", "ping"]
      interval: 10s
      timeout: 5s
      retries: 5
    networks:
      - ai-pms-network

  # Backend API
  backend:
    build:
      context: ./backend
      dockerfile: Dockerfile
    container_name: ai-pms-backend
    environment:
      SPRING_PROFILES_ACTIVE: ${SPRING_PROFILES_ACTIVE:-docker}
      DATABASE_URL: ${DATABASE_URL:-jdbc:postgresql://postgres:5432/ai_pms_db}
      DATABASE_USER: ${DATABASE_USER:-ai_pms_user}
      DATABASE_PASSWORD: ${DATABASE_PASSWORD:-ai_pms_pass}
      REDIS_HOST: ${REDIS_HOST:-redis}
      AI_SERVICE_URL: ${AI_SERVICE_URL:-http://ai-service:8000}
      JWT_SECRET: ${JWT_SECRET:-your-secret-key}
      JWT_EXPIRATION: ${JWT_EXPIRATION:-900}
    ports:
      - "8080:8080"
    depends_on:
      postgres:
        condition: service_healthy
      redis:
        condition: service_healthy
    healthcheck:
      test: ["CMD", "curl", "-f", "http://localhost:8080/actuator/health"]
      interval: 30s
      timeout: 10s
      retries: 3
    networks:
      - ai-pms-network

  # AI Service
  ai-service:
    build:
      context: ./ai-service
      dockerfile: Dockerfile
    container_name: ai-pms-ai-service
    ports:
      - "8000:8000"
    environment:
      MODEL_PATH: ${MODEL_PATH:-/app/models}
      LOG_LEVEL: ${LOG_LEVEL:-INFO}
    healthcheck:
      test: ["CMD", "curl", "-f", "http://localhost:8000/health"]
      interval: 30s
      timeout: 10s
      retries: 3
    networks:
      - ai-pms-network

  # Frontend Dashboard
  frontend:
    build:
      context: ./frontend
      dockerfile: Dockerfile
    container_name: ai-pms-frontend
    ports:
      - "5173:80"
    depends_on:
      - backend
    environment:
      VITE_API_BASE_URL: ${VITE_API_BASE_URL:-http://localhost:8080/api/v1}
    networks:
      - ai-pms-network

  # Nginx Reverse Proxy (optional)
  nginx:
    image: nginx:alpine
    container_name: ai-pms-nginx
    ports:
      - "80:80"
      - "443:443"
    volumes:
      - ./nginx/nginx.conf:/etc/nginx/nginx.conf
      - ./nginx/ssl:/etc/nginx/ssl
    depends_on:
      - backend
      - frontend
      - ai-service
    networks:
      - ai-pms-network

networks:
  ai-pms-network:
    driver: bridge

volumes:
  postgres_data:
  redis_data:
\end{verbatim}

\subsection{docker-compose.dev.yml (Phát triển)}

\begin{verbatim}
# docker-compose.dev.yml
version: '3.8'

services:
  postgres:
    image: postgres:16-alpine
    container_name: ai-pms-postgres-dev
    environment:
      POSTGRES_USER: ai_pms_user
      POSTGRES_PASSWORD: ai_pms_pass
      POSTGRES_DB: ai_pms_db
    ports:
      - "5433:5432"
    volumes:
      - postgres_data_dev:/var/lib/postgresql/data

  redis:
    image: redis:7-alpine
    container_name: ai-pms-redis-dev
    ports:
      - "6380:6379"

  # Backend chạy từ source (hot-reload)
  backend-dev:
    image: eclipse-temurin:21-jdk
    container_name: ai-pms-backend-dev
    working_dir: /app
    volumes:
      - ./backend:/app
      - ~/.m2:/root/.m2
    ports:
      - "8080:8080"
      - "5005:5005"
    environment:
      SPRING_PROFILES_ACTIVE: dev
      DATABASE_URL: jdbc:postgresql://postgres:5432/ai_pms_db
      DATABASE_USER: ai_pms_user
      DATABASE_PASSWORD: ai_pms_pass
      REDIS_HOST: redis
      AI_SERVICE_URL: http://ai-service:8000
    command: ./mvnw spring-boot:run -Dspring-boot.run.jvmArguments="-agentlib:jdwp=transport=dt_socket,server=y,suspend=n,address=*:5005"
    depends_on:
      - postgres
      - redis

  # AI Service chạy từ source (hot-reload)
  ai-service-dev:
    image: python:3.11-slim
    container_name: ai-pms-ai-service-dev
    working_dir: /app
    volumes:
      - ./ai-service:/app
    ports:
      - "8000:8000"
    command: uvicorn app.main:app --host 0.0.0.0 --port 8000 --reload
    environment:
      MODEL_PATH: /app/models

  # Frontend chạy từ source (hot-reload)
  frontend-dev:
    image: node:20-alpine
    container_name: ai-pms-frontend-dev
    working_dir: /app
    volumes:
      - ./frontend:/app
      - /app/node_modules
    ports:
      - "5173:5173"
    command: npm run dev -- --host
    environment:
      VITE_API_BASE_URL: http://localhost:8080/api/v1

volumes:
  postgres_data_dev:
\end{verbatim}

%%%%%%%%%%%%%%%%%%%%%%%%%%%%%%%%%%%%%%%%%%%%%%%%%%%%%%%%%%%%%%
\section{CI/CD Pipeline với GitHub Actions}
%%%%%%%%%%%%%%%%%%%%%%%%%%%%%%%%%%%%%%%%%%%%%%%%%%%%%%%%%%%%%%

\subsection{CI Pipeline (.github/workflows/ci.yml)}

\begin{verbatim}
# .github/workflows/ci.yml
name: CI - Continuous Integration

on:
  push:
    branches: [ main, develop ]
  pull_request:
    branches: [ main ]

jobs:
  # Backend CI
  backend-build:
    name: Backend - Build & Test
    runs-on: ubuntu-latest
    
    services:
      postgres:
        image: postgres:16
        env:
          POSTGRES_USER: ai_pms_user
          POSTGRES_PASSWORD: ai_pms_pass
          POSTGRES_DB: ai_pms_db
        ports:
          - 5432:5432
        options: >-
          --health-cmd pg_isready
          --health-interval 10s
          --health-timeout 5s
          --health-retries 5
    
    steps:
      - name: Checkout code
        uses: actions/checkout@v4
      
      - name: Set up JDK 21
        uses: actions/setup-java@v4
        with:
          java-version: '21'
          distribution: 'temurin'
      
      - name: Cache Maven dependencies
        uses: actions/cache@v3
        with:
          path: ~/.m2
          key: ${{ runner.os }}-m2-${{ hashFiles('**/pom.xml') }}
          restore-keys: ${{ runner.os }}-m2
      
      - name: Build and Test Backend
        run: |
          cd backend
          ./mvnw clean compile test package -DskipTests=false
      
      - name: Upload test results
        uses: actions/upload-artifact@v3
        if: always()
        with:
          name: backend-test-results
          path: backend/target/surefire-reports/

  # Frontend CI
  frontend-build:
    name: Frontend - Build & Test
    runs-on: ubuntu-latest
    
    steps:
      - name: Checkout code
        uses: actions/checkout@v4
      
      - name: Set up Node.js
        uses: actions/setup-node@v4
        with:
          node-version: '20'
          cache: 'npm'
          cache-dependency-path: frontend/package-lock.json
      
      - name: Install dependencies
        run: |
          cd frontend
          npm ci
      
      - name: Lint
        run: |
          cd frontend
          npm run lint
      
      - name: Build
        run: |
          cd frontend
          npm run build
      
      - name: Run tests
        run: |
          cd frontend
          npm run test -- --watchAll=false --coverage
      
      - name: Upload build artifacts
        uses: actions/upload-artifact@v3
        with:
          name: frontend-build
          path: frontend/dist/

  # AI Service CI
  ai-service-check:
    name: AI Service - Check
    runs-on: ubuntu-latest
    
    steps:
      - name: Checkout code
        uses: actions/checkout@v4
      
      - name: Set up Python
        uses: actions/setup-python@v5
        with:
          python-version: '3.11'
      
      - name: Install dependencies
        run: |
          cd ai-service
          pip install -r requirements.txt
      
      - name: Check syntax
        run: |
          cd ai-service
          python -m compileall app/
          python -c "import app.main; print('Import OK')"
      
      - name: Run tests
        run: |
          cd ai-service
          pip install pytest
          python -m pytest tests/ -v

  # Chrome Extension CI
  extension-check:
    name: Chrome Extension - Check
    runs-on: ubuntu-latest
    
    steps:
      - name: Checkout code
        uses: actions/checkout@v4
      
      - name: Validate manifest.json
        run: |
          cd chrome-extension
          if command -v jq &> /dev/null; then
            jq empty manifest.json
          else
            echo "jq not installed, skipping JSON validation"
          fi
      
      - name: Check required files
        run: |
          cd chrome-extension
          test -f manifest.json
          test -f popup.html
          test -f popup.js
          echo "All required files exist"

  # Security Scan
  security-scan:
    name: Security Scan
    runs-on: ubuntu-latest
    needs: [backend-build, frontend-build]
    
    steps:
      - name: Checkout code
        uses: actions/checkout@v4
      
      - name: Run Trivy vulnerability scanner
        uses: aquasecurity/trivy-action@master
        with:
          scan-type: 'fs'
          scan-ref: '.'
          format: 'sarif'
          output: 'trivy-results.sarif'
      
      - name: Upload Trivy results to GitHub Security tab
        uses: github/codeql-action/upload-sarif@v2
        with:
          sarif_file: 'trivy-results.sarif'
\end{verbatim}

\subsection{CD Pipeline (.github/workflows/cd.yml)}

\begin{verbatim}
# .github/workflows/cd.yml
name: CD - Continuous Deployment

on:
  push:
    branches: [ main ]
  workflow_dispatch:

jobs:
  deploy:
    name: Deploy to Production
    runs-on: ubuntu-latest
    
    steps:
      - name: Checkout code
        uses: actions/checkout@v4
      
      - name: Set up Docker Buildx
        uses: docker/setup-buildx-action@v3
      
      - name: Log in to Docker Hub
        uses: docker/login-action@v3
        with:
          username: ${{ secrets.DOCKER_USERNAME }}
          password: ${{ secrets.DOCKER_PASSWORD }}
      
      - name: Build and push Backend
        uses: docker/build-push-action@v5
        with:
          context: ./backend
          file: ./backend/Dockerfile
          push: true
          tags: |
            ${{ secrets.DOCKER_USERNAME }}/ai-pms-backend:latest
            ${{ secrets.DOCKER_USERNAME }}/ai-pms-backend:${{ github.sha }}
      
      - name: Build and push Frontend
        uses: docker/build-push-action@v5
        with:
          context: ./frontend
          file: ./frontend/Dockerfile
          push: true
          tags: |
            ${{ secrets.DOCKER_USERNAME }}/ai-pms-frontend:latest
            ${{ secrets.DOCKER_USERNAME }}/ai-pms-frontend:${{ github.sha }}
      
      - name: Build and push AI Service
        uses: docker/build-push-action@v5
        with:
          context: ./ai-service
          file: ./ai-service/Dockerfile
          push: true
          tags: |
            ${{ secrets.DOCKER_USERNAME }}/ai-pms-ai-service:latest
            ${{ secrets.DOCKER_USERNAME }}/ai-pms-ai-service:${{ github.sha }}
      
      - name: Deploy to Server
        uses: appleboy/ssh-action@v1.0.0
        with:
          host: ${{ secrets.DEPLOY_HOST }}
          username: ${{ secrets.DEPLOY_USER }}
          key: ${{ secrets.DEPLOY_SSH_KEY }}
          script: |
            cd /opt/ai-pms
            docker compose pull
            docker compose up -d --remove-orphans
            docker system prune -f
      
      - name: Health Check
        run: |
          echo "Waiting for deployment to be healthy..."
          sleep 30
          curl -f http://${{ secrets.DEPLOY_HOST }}/api/v1/actuator/health || exit 1
          curl -f http://${{ secrets.DEPLOY_HOST }}/ai/health || exit 1
          echo "Deployment successful!"
      
      - name: Notify Deployment Status
        uses: actions/github-script@v6
        with:
          script: |
            const success = process.env.SUCCESS === 'true';
            const message = success ? '✅ Deployment successful!' : '❌ Deployment failed!';
            // Add notification logic here (Slack, Email, etc.)
\end{verbatim}

%%%%%%%%%%%%%%%%%%%%%%%%%%%%%%%%%%%%%%%%%%%%%%%%%%%%%%%%%%%%%%
\section{Cài đặt và Cấu hình từng thành phần}
%%%%%%%%%%%%%%%%%%%%%%%%%%%%%%%%%%%%%%%%%%%%%%%%%%%%%%%%%%%%%%

\subsection{Cấu hình Backend (application.yml)}

\begin{verbatim}
# backend/src/main/resources/application.yml
spring:
  application:
    name: ai-pms-backend
  
  datasource:
    url: ${DATABASE_URL:jdbc:postgresql://localhost:5432/ai_pms_db}
    username: ${DATABASE_USER:ai_pms_user}
    password: ${DATABASE_PASSWORD:ai_pms_pass}
    driver-class-name: org.postgresql.Driver
    hikari:
      maximum-pool-size: 10
      minimum-idle: 5
      connection-timeout: 30000
  
  jpa:
    hibernate:
      ddl-auto: validate
    properties:
      hibernate:
        dialect: org.hibernate.dialect.PostgreSQLDialect
        format_sql: true
        jdbc:
          batch_size: 20
    show-sql: false
  
  flyway:
    enabled: true
    baseline-on-migrate: true
    locations: classpath:db/migration
  
  data:
    redis:
      host: ${REDIS_HOST:localhost}
      port: ${REDIS_PORT:6379}
      timeout: 5000ms

server:
  port: ${PORT:8080}
  servlet:
    context-path: /api/v1
  
  # CORS Configuration
  cors:
    allowed-origins: ${CORS_ORIGINS:http://localhost:5173,http://localhost:3000}
    allowed-methods: GET,POST,PUT,DELETE,PATCH,OPTIONS
    allowed-headers: "*"
    allow-credentials: true

# JWT Configuration
jwt:
  secret: ${JWT_SECRET:your-secret-key-here}
  access-token-expiration: ${JWT_ACCESS_EXPIRATION:900}  # 15 phút
  refresh-token-expiration: ${JWT_REFRESH_EXPIRATION:604800}  # 7 ngày

# AI Service
ai:
  service:
    url: ${AI_SERVICE_URL:http://localhost:8000}
    timeout: 30000

# Logging
logging:
  level:
    com.ai.pms: ${LOG_LEVEL:INFO}
    org.springframework.security: WARN
    org.hibernate: WARN

# Actuator
management:
  endpoints:
    web:
      exposure:
        include: health,info,metrics,prometheus
  endpoint:
    health:
      show-details: always
  metrics:
    export:
      prometheus:
        enabled: true
\end{verbatim}

\subsection{Cấu hình Frontend (.env)}

\begin{verbatim}
# frontend/.env
VITE_API_BASE_URL=http://localhost:8080/api/v1
VITE_AI_SERVICE_URL=http://localhost:8000
VITE_APP_NAME=AI Project Management
VITE_APP_VERSION=1.0.0
VITE_ENVIRONMENT=development
\end{verbatim}

\subsection{Cấu hình AI Service (.env)}

\begin{verbatim}
# ai-service/.env
MODEL_PATH=./models
LOG_LEVEL=INFO
MAX_CONCURRENT_REQUESTS=10
REQUEST_TIMEOUT=30
\end{verbatim}

%%%%%%%%%%%%%%%%%%%%%%%%%%%%%%%%%%%%%%%%%%%%%%%%%%%%%%%%%%%%%%
\section{Hướng dẫn cài đặt}
%%%%%%%%%%%%%%%%%%%%%%%%%%%%%%%%%%%%%%%%%%%%%%%%%%%%%%%%%%%%%%

\subsection{Cài đặt với Docker Compose (Khuyến nghị)}

\begin{verbatim}
# 1. Clone repository
git clone https://github.com/your-username/ai-project-management.git
cd ai-project-management

# 2. Tạo file .env từ .env.example
cp .env.example .env
# Chỉnh sửa các biến môi trường trong .env

# 3. Khởi động toàn bộ hệ thống
docker compose up -d --build

# 4. Chạy database migration và seed data
docker compose exec backend npx prisma migrate deploy
docker compose exec backend npx prisma db seed

# 5. Kiểm tra trạng thái
docker compose ps

# 6. Truy cập hệ thống
# Frontend: http://localhost:5173
# Backend API: http://localhost:8080/api/v1
# AI Service: http://localhost:8000
# Swagger UI: http://localhost:8080/swagger-ui.html
\end{verbatim}

\subsection{Cài đặt thủ công}

\begin{verbatim}
# 1. Cài đặt Database
docker run -d \
  --name postgres \
  -e POSTGRES_USER=ai_pms_user \
  -e POSTGRES_PASSWORD=ai_pms_pass \
  -e POSTGRES_DB=ai_pms_db \
  -p 5432:5432 \
  postgres:16

# 2. Cài đặt Redis
docker run -d \
  --name redis \
  -p 6379:6379 \
  redis:7-alpine

# 3. Cài đặt Backend
cd backend
./mvnw clean install
./mvnw spring-boot:run

# 4. Cài đặt AI Service
cd ai-service
python -m venv .venv
source .venv/bin/activate  # Windows: .venv\Scripts\activate
pip install -r requirements.txt
uvicorn app.main:app --reload --port 8000

# 5. Cài đặt Frontend
cd frontend
npm install
npm run dev

# 6. Cài đặt Chrome Extension
# Mở Chrome, vào chrome://extensions
# Bật Developer mode
# Click "Load unpacked"
# Chọn thư mục chrome-extension/
\end{verbatim}

%%%%%%%%%%%%%%%%%%%%%%%%%%%%%%%%%%%%%%%%%%%%%%%%%%%%%%%%%%%%%%
\section{Quản lý và Vận hành}
%%%%%%%%%%%%%%%%%%%%%%%%%%%%%%%%%%%%%%%%%%%%%%%%%%%%%%%%%%%%%%

\subsection{Các lệnh Docker Compose thường dùng}

\begin{table}[H]
\centering
\caption{Các lệnh Docker Compose hữu ích}
\begin{tabular}{|p{5cm}|p{9cm}|}
\hline
\textbf{Lệnh} & \textbf{Mô tả} \\
\hline
docker compose up -d & Khởi động tất cả services \\
\hline
docker compose down & Dừng và xóa tất cả containers \\
\hline
docker compose logs -f & Xem log của tất cả services \\
\hline
docker compose logs -f backend & Xem log của Backend \\
\hline
docker compose ps & Xem trạng thái các containers \\
\hline
docker compose restart backend & Khởi động lại Backend \\
\hline
docker compose exec backend bash & Truy cập vào container Backend \\
\hline
docker compose exec postgres psql -U ai_pms_user ai_pms_db & Truy cập PostgreSQL \\
\hline
docker compose exec redis redis-cli & Truy cập Redis \\
\hline
docker compose build & Build lại images \\
\hline
docker compose pull & Pull images mới nhất \\
\hline
docker system prune -f & Dọn dẹp resources không sử dụng \\
\hline
\end{tabular}
\label{tab:docker_commands}
\end{table}

\subsection{Health Check và Monitoring}

\begin{verbatim}
# Kiểm tra health của từng service
# Backend
curl http://localhost:8080/actuator/health

# AI Service
curl http://localhost:8000/health

# Database
docker compose exec postgres pg_isready -U ai_pms_user

# Redis
docker compose exec redis redis-cli ping

# Kiểm tra toàn bộ hệ thống
curl http://localhost:8080/actuator/health | jq
\end{verbatim}

\subsection{Backup và Restore Database}

\begin{verbatim}
# Backup
docker compose exec postgres pg_dump -U ai_pms_user ai_pms_db > backup.sql

# Restore
cat backup.sql | docker compose exec -T postgres psql -U ai_pms_user ai_pms_db

# Backup với nén
docker compose exec postgres pg_dump -U ai_pms_user ai_pms_db | gzip > backup.sql.gz
\end{verbatim}

%%%%%%%%%%%%%%%%%%%%%%%%%%%%%%%%%%%%%%%%%%%%%%%%%%%%%%%%%%%%%%
\section{Production Deployment}
%%%%%%%%%%%%%%%%%%%%%%%%%%%%%%%%%%%%%%%%%%%%%%%%%%%%%%%%%%%%%%

\subsection{Kiến trúc Production}

\begin{figure}[H]
\centering
\fbox{
\begin{minipage}{0.95\textwidth}
\centering
\begin{tabular}{c}
\textbf{User} \\
$\downarrow$ \\
\textbf{CDN / CloudFront} \\
$\downarrow$ \\
\textbf{Load Balancer (Nginx)} \\
$\swarrow$ $\searrow$ \\
\textbf{Frontend (React)} \hspace{3cm} \textbf{API Gateway} \\
\hspace{3cm} $\downarrow$ \\
\hspace{3cm} \textbf{Backend Service (Spring Boot)} \\
\hspace{3cm} $\swarrow$ $\searrow$ \\
\hspace{2cm} \textbf{PostgreSQL} \hspace{2cm} \textbf{AI Service (FastAPI)}
\end{tabular}
\end{minipage}
}
\caption{Kiến trúc Production}
\label{fig:production_architecture}
\end{figure}

\subsection{Kubernetes Deployment (K8s)}

\begin{verbatim}
# deployment.yaml (Kubernetes)
apiVersion: apps/v1
kind: Deployment
metadata:
  name: ai-pms-backend
  namespace: production
spec:
  replicas: 3
  selector:
    matchLabels:
      app: ai-pms-backend
  template:
    metadata:
      labels:
        app: ai-pms-backend
    spec:
      containers:
      - name: backend
        image: your-registry/ai-pms-backend:latest
        ports:
        - containerPort: 8080
        env:
        - name: SPRING_PROFILES_ACTIVE
          value: "prod"
        - name: DATABASE_URL
          valueFrom:
            secretKeyRef:
              name: db-secret
              key: url
        resources:
          requests:
            memory: "512Mi"
            cpu: "250m"
          limits:
            memory: "1Gi"
            cpu: "500m"
        livenessProbe:
          httpGet:
            path: /actuator/health/liveness
            port: 8080
          initialDelaySeconds: 30
          periodSeconds: 10
        readinessProbe:
          httpGet:
            path: /actuator/health/readiness
            port: 8080
          initialDelaySeconds: 20
          periodSeconds: 5
---
apiVersion: v1
kind: Service
metadata:
  name: ai-pms-backend-service
  namespace: production
spec:
  selector:
    app: ai-pms-backend
  ports:
  - port: 8080
    targetPort: 8080
  type: ClusterIP
---
apiVersion: autoscaling/v2
kind: HorizontalPodAutoscaler
metadata:
  name: ai-pms-backend-hpa
  namespace: production
spec:
  scaleTargetRef:
    apiVersion: apps/v1
    kind: Deployment
    name: ai-pms-backend
  minReplicas: 2
  maxReplicas: 10
  metrics:
  - type: Resource
    resource:
      name: cpu
      target:
        type: Utilization
        averageUtilization: 70
  - type: Resource
    resource:
      name: memory
      target:
        type: Utilization
        averageUtilization: 80
\end{verbatim}

\subsection{Monitoring Stack}

\begin{table}[H]
\centering
\caption{Monitoring Stack đề xuất}
\begin{tabular}{|p{4cm}|p{4cm}|p{6cm}|}
\hline
\textbf{Thành phần} & \textbf{Công cụ} & \textbf{Mục đích} \\
\hline
Metrics Collection & Prometheus & Thu thập metrics \\
\hline
Visualization & Grafana & Dashboard và biểu đồ \\
\hline
Logging & ELK Stack & Tập trung log \\
\hline
Tracing & Jaeger & Distributed tracing \\
\hline
Alerting & AlertManager & Cảnh báo \\
\hline
APM & New Relic / Datadog & Application Performance \\
\hline
\end{tabular}
\label{tab:monitoring_stack}
\end{table}

%%%%%%%%%%%%%%%%%%%%%%%%%%%%%%%%%%%%%%%%%%%%%%%%%%%%%%%%%%%%%%
\section{Tổng kết chương}
%%%%%%%%%%%%%%%%%%%%%%%%%%%%%%%%%%%%%%%%%%%%%%%%%%%%%%%%%%%%%%

Trong chương này, luận văn đã trình bày chi tiết quy trình triển khai hệ thống AI Project Management từ môi trường phát triển đến production.

Các nội dung chính đã được thực hiện:

\begin{enumerate}
\item \textbf{Môi trường phát triển}: Yêu cầu phần mềm và cấu trúc thư mục.
\item \textbf{Containerization}: Dockerfiles cho Backend, Frontend, AI Service và Database.
\item \textbf{Docker Compose}: Cấu hình orchestration cho development và production.
\item \textbf{CI/CD Pipeline}: GitHub Actions cho CI và CD.
\item \textbf{Cài đặt và Cấu hình}: Hướng dẫn cấu hình từng thành phần.
\item \textbf{Hướng dẫn cài đặt}: Quy trình cài đặt với Docker Compose và thủ công.
\item \textbf{Quản lý và Vận hành}: Các lệnh hữu ích và backup/restore.
\item \textbf{Production Deployment}: Kubernetes và monitoring stack.
\end{enumerate}

Đây là cơ sở để Chương 13 trình bày thực nghiệm và đánh giá hệ thống.

%%%%%%%%%%%%%%%%%%%%%%%%%%%%%%%%%%%%%%%%%%%%%%%%%%%%%%%%%%%%%%

%%%%%%%%%%%%%%%%%%%%%%%%%%%%%%%%%%%%%%%%%%%%%%%%%%%%%%%%%%%%%%
%%%%%%%%%%%%%%%%%%%%%%%%%%%%%%%%%%%%%%%%%%%%%%%%%%%%%%%%%%%%%%
\chapter{THỰC NGHIỆM}
\label{chap:experiments}
%%%%%%%%%%%%%%%%%%%%%%%%%%%%%%%%%%%%%%%%%%%%%%%%%%%%%%%%%%%%%%

\section{Giới thiệu chương}

Sau khi trình bày quy trình triển khai hệ thống ở Chương 12, chương này tập trung vào việc thực nghiệm và đánh giá hệ thống AI Project Management. Thực nghiệm là bước quan trọng để kiểm chứng rằng hệ thống đáp ứng các yêu cầu đã đề ra và hoạt động hiệu quả trong thực tế.

Mục tiêu của chương bao gồm:

\begin{itemize}
\item Trình bày môi trường và kịch bản thực nghiệm.
\item Đánh giá hiệu năng của các thành phần hệ thống.
\item Đánh giá chất lượng của các AI Agent.
\item Đánh giá trải nghiệm người dùng.
\item So sánh với các hệ thống hiện có.
\item Phân tích kết quả và rút ra nhận xét.
\end{itemize}

Đây là cơ sở để Chương 14 trình bày đánh giá kết quả và Chương 15 trình bày kết luận.

%%%%%%%%%%%%%%%%%%%%%%%%%%%%%%%%%%%%%%%%%%%%%%%%%%%%%%%%%%%%%%
\section{Môi trường thực nghiệm}
%%%%%%%%%%%%%%%%%%%%%%%%%%%%%%%%%%%%%%%%%%%%%%%%%%%%%%%%%%%%%%

\subsection{Cấu hình phần cứng}

\begin{table}[H]
\centering
\caption{Cấu hình máy tính thực nghiệm}
\begin{tabular}{|p{5cm}|p{8cm}|}
\hline
\textbf{Thành phần} & \textbf{Thông số} \\
\hline
CPU & Intel Core i7-12700H (14 nhân, 20 luồng, 2.3GHz - 4.7GHz) \\
\hline
RAM & 32 GB DDR5 4800MHz \\
\hline
GPU & NVIDIA GeForce RTX 3060 (6GB VRAM) \\
\hline
Ổ cứng & SSD NVMe 1TB (Đọc: 7000 MB/s, Ghi: 5000 MB/s) \\
\hline
Hệ điều hành & Ubuntu 24.04 LTS (Kernel 6.8.0) \\
\hline
Docker & Docker Desktop 4.28.0 \\
\hline
\end{tabular}
\label{tab:hw_config}
\end{table}

\subsection{Cấu hình phần mềm}

\begin{table}[H]
\centering
\caption{Cấu hình phần mềm thực nghiệm}
\begin{tabular}{|p{5cm}|p{8cm}|}
\hline
\textbf{Thành phần} & \textbf{Phiên bản} \\
\hline
JDK & OpenJDK 21.0.2 \\
\hline
Spring Boot & 3.2.4 \\
\hline
Node.js & 20.12.0 \\
\hline
React & 19.0.0 \\
\hline
Python & 3.11.8 \\
\hline
FastAPI & 0.110.0 \\
\hline
PostgreSQL & 16.2 \\
\hline
Redis & 7.2.4 \\
\hline
LLM & GPT-4o (OpenAI API) \\
\hline
\end{tabular}
\label{tab:sw_config}
\end{table}

\subsection{Các công cụ đo lường}

\begin{itemize}
\item \textbf{Apache JMeter 5.6.2}: Đo lường hiệu năng API.
\item \textbf{Lighthouse}: Đánh giá hiệu năng Frontend.
\item \textbf{Prometheus + Grafana}: Giám sát metrics.
\item \textbf{Python scripts}: Đánh giá độ chính xác AI.
\item \textbf{Postman}: Kiểm thử API.
\end{itemize}

%%%%%%%%%%%%%%%%%%%%%%%%%%%%%%%%%%%%%%%%%%%%%%%%%%%%%%%%%%%%%%
\section{Bộ dữ liệu thực nghiệm}
%%%%%%%%%%%%%%%%%%%%%%%%%%%%%%%%%%%%%%%%%%%%%%%%%%%%%%%%%%%%%%

\subsection{Dữ liệu cho AI Agents}

\subsubsection{Dữ liệu huấn luyện và kiểm thử}

\begin{table}[H]
\centering
\caption{Bộ dữ liệu AI}
\begin{tabular}{|p{4cm}|p{4.5cm}|p{5cm}|}
\hline
\textbf{Loại dữ liệu} & \textbf{Số lượng} & \textbf{Nguồn} \\
\hline
User Stories (có Story Point) & 1,200 & Jira RedHat (ENTMQMAAS) \\
\hline
Sprint Reports & 48 & Dự án thực tế \\
\hline
Bug Reports (đã phân loại) & 2,500 & Jira RedHat \\
\hline
Requirements (văn bản) & 850 & Tổng hợp từ nhiều dự án \\
\hline
Task Breakdowns & 600 & Kho lưu trữ nội bộ \\
\hline
\end{tabular}
\label{tab:ai_dataset}
\end{table}

\subsubsection{Phân bố dữ liệu Story Point}

\begin{table}[H]
\centering
\caption{Phân bố Story Point trong tập dữ liệu}
\begin{tabular}{|c|c|c|}
\hline
\textbf{Story Point} & \textbf{Số lượng} & \textbf{Tỷ lệ (\%)} \\
\hline
1 & 85 & 7.1 \\
\hline
2 & 210 & 17.5 \\
\hline
3 & 340 & 28.3 \\
\hline
5 & 285 & 23.8 \\
\hline
8 & 180 & 15.0 \\
\hline
13 & 75 & 6.3 \\
\hline
21 & 25 & 2.0 \\
\hline
\textbf{Tổng} & \textbf{1,200} & \textbf{100} \\
\hline
\end{tabular}
\label{tab:sp_distribution}
\end{table}

\subsection{Dữ liệu cho hiệu năng hệ thống}

\begin{table}[H]
\centering
\caption{Kịch bản kiểm thử hiệu năng}
\begin{tabular}{|p{4cm}|p{4.5cm}|p{5cm}|}
\hline
\textbf{Kịch bản} & \textbf{Số request} & \textbf{Thời gian chạy} \\
\hline
Đăng nhập & 1,000 & 5 phút \\
\hline
Tạo Project & 500 & 5 phút \\
\hline
Tạo Sprint & 500 & 5 phút \\
\hline
Tạo Task & 1,000 & 5 phút \\
\hline
AI Requirement Analysis & 200 & 10 phút \\
\hline
AI Story Point Estimation & 200 & 10 phút \\
\hline
API tổng hợp (hỗn hợp) & 5,000 & 15 phút \\
\hline
\end{tabular}
\label{tab:perf_scenarios}
\end{table}

%%%%%%%%%%%%%%%%%%%%%%%%%%%%%%%%%%%%%%%%%%%%%%%%%%%%%%%%%%%%%%
\section{Kịch bản thực nghiệm}
%%%%%%%%%%%%%%%%%%%%%%%%%%%%%%%%%%%%%%%%%%%%%%%%%%%%%%%%%%%%%%

\subsection{Kịch bản 1: Đánh giá hiệu năng API}

\textbf{Mục tiêu}: Đo thời gian phản hồi và thông lượng của các API.

\textbf{Các bước thực hiện}:

\begin{enumerate}
\item Cấu hình Apache JMeter với các thread groups.
\item Thực hiện các request đến từng endpoint.
\item Ghi lại thời gian phản hồi trung bình, percentile (p50, p95, p99).
\item Phân tích kết quả.
\end{enumerate}

\subsection{Kịch bản 2: Đánh giá chất lượng AI}

\textbf{Mục tiêu}: Đánh giá độ chính xác của các AI Agent.

\textbf{Các bước thực hiện}:

\begin{enumerate}
\item Chia tập dữ liệu thành training (80\%) và test (20\%).
\item Huấn luyện các mô hình (đối với các thuật toán ML cổ điển).
\item Đánh giá trên tập test với các metrics (Accuracy, Precision, Recall, F1).
\item So sánh với baseline (ước lượng thủ công).
\end{enumerate}

\subsection{Kịch bản 3: Đánh giá trải nghiệm người dùng}

\textbf{Mục tiêu}: Đánh giá mức độ hài lòng của người dùng.

\textbf{Các bước thực hiện}:

\begin{enumerate}
\item Mời 20 người dùng (5 PO, 5 SM, 10 Developer) tham gia.
\item Hướng dẫn sử dụng hệ thống trong 2 tuần.
\item Thu thập phản hồi qua khảo sát.
\item Phân tích kết quả.
\end{enumerate}

%%%%%%%%%%%%%%%%%%%%%%%%%%%%%%%%%%%%%%%%%%%%%%%%%%%%%%%%%%%%%%
\section{Kết quả thực nghiệm}
%%%%%%%%%%%%%%%%%%%%%%%%%%%%%%%%%%%%%%%%%%%%%%%%%%%%%%%%%%%%%%

\subsection{Kết quả đánh giá hiệu năng API}

\subsubsection{Thời gian phản hồi trung bình}

\begin{table}[H]
\centering
\caption{Thời gian phản hồi API trung bình}
\begin{tabular}{|p{5cm}|p{3cm}|p{3cm}|p{3cm}|}
\hline
\textbf{API} & \textbf{Trung bình (ms)} & \textbf{p95 (ms)} & \textbf{p99 (ms)} \\
\hline
POST /auth/login & 82 & 145 & 210 \\
\hline
GET /projects & 96 & 180 & 250 \\
\hline
POST /projects & 108 & 195 & 280 \\
\hline
POST /sprints & 95 & 170 & 240 \\
\hline
POST /tasks & 73 & 140 & 200 \\
\hline
GET /dashboard/overview & 145 & 260 & 380 \\
\hline
POST /ai/analyze-requirement & 2,380 & 4,200 & 5,800 \\
\hline
POST /ai/estimate-story-point & 1,954 & 3,500 & 4,800 \\
\hline
GET /sprints/burndown & 85 & 150 & 220 \\
\hline
POST /ai/risk & 2,145 & 3,800 & 5,200 \\
\hline
\end{tabular}
\label{tab:api_response_time}
\end{table}

\begin{figure}[H]
\centering
\fbox{
\begin{minipage}{0.9\textwidth}
\centering
\begin{tabular}{c}
\textbf{Biểu đồ so sánh thời gian phản hồi API} \\
\\
\begin{tabular}{|c|c|}
\hline
API & Thời gian (ms) \\
\hline
Login & 82 \\
\hline
Get Projects & 96 \\
\hline
Create Project & 108 \\
\hline
Create Task & 73 \\
\hline
Dashboard & 145 \\
\hline
AI Analysis & 2380 \\
\hline
AI Story Point & 1954 \\
\hline
AI Risk & 2145 \\
\hline
\end{tabular}
\end{tabular}
\end{minipage}
}
\caption{So sánh thời gian phản hồi API}
\label{fig:api_response_chart}
\end{figure}

\subsubsection{Thông lượng (Throughput)}

\begin{table}[H]
\centering
\caption{Thông lượng API}
\begin{tabular}{|p{5cm}|p{3cm}|p{3cm}|}
\hline
\textbf{API} & \textbf{Requests/sec} & \textbf{Concurrent Users} \\
\hline
GET /projects & 185 & 100 \\
\hline
POST /projects & 42 & 50 \\
\hline
GET /dashboard/overview & 28 & 50 \\
\hline
POST /ai/analyze-requirement & 5 & 20 \\
\hline
POST /tasks & 156 & 100 \\
\hline
\end{tabular}
\label{tab:api_throughput}
\end{table}

\subsection{Kết quả đánh giá chất lượng AI}

\subsubsection{Requirement Agent}

\begin{table}[H]
\centering
\caption{Đánh giá Requirement Agent}
\begin{tabular}{|p{5cm}|c|}
\hline
\textbf{Tiêu chí} & \textbf{Kết quả (\%)} \\
\hline
Requirement Summary Accuracy & 95.8 \\
\hline
Missing Requirement Detection & 92.6 \\
\hline
Duplicate Detection & 94.3 \\
\hline
Requirement Classification & 96.1 \\
\hline
User Story Generation (Human evaluation) & 87.2 \\
\hline
Acceptance Criteria Generation & 91.5 \\
\hline
\end{tabular}
\label{tab:req_agent_results}
\end{table}

\subsubsection{Story Point Agent}

\begin{table}[H]
\centering
\caption{Đánh giá Story Point Agent}
\begin{tabular}{|p{5cm}|c|}
\hline
\textbf{Tiêu chí} & \textbf{Kết quả} \\
\hline
Accuracy (Đúng hoàn toàn) & 81.5\% \\
\hline
Accuracy (Sai lệch ±1 SP) & 14.8\% \\
\hline
Accuracy (Sai lệch >2 SP) & 3.7\% \\
\hline
Precision & 82.3\% \\
\hline
Recall & 80.1\% \\
\hline
F1-Score & 81.2\% \\
\hline
Mean Absolute Error (MAE) & 1.2 \\
\hline
Root Mean Square Error (RMSE) & 1.8 \\
\hline
\end{tabular}
\label{tab:sp_agent_results}
\end{table}

\begin{figure}[H]
\centering
\fbox{
\begin{minipage}{0.9\textwidth}
\centering
\begin{tabular}{c}
\textbf{Ma trận nhầm lẫn - Story Point Agent} \\
\\
\begin{tabular}{|c|c|c|c|c|c|c|c|}
\hline
\multicolumn{1}{|c|}{} & \multicolumn{7}{c|}{\textbf{Dự đoán}} \\
\cline{2-8}
\multicolumn{1}{|c|}{\textbf{Thực tế}} & 1 & 2 & 3 & 5 & 8 & 13 & 21 \\
\hline
1 & 14 & 2 & 1 & 0 & 0 & 0 & 0 \\
\hline
2 & 3 & 28 & 5 & 2 & 0 & 0 & 0 \\
\hline
3 & 0 & 4 & 50 & 8 & 1 & 0 & 0 \\
\hline
5 & 0 & 1 & 6 & 38 & 5 & 1 & 0 \\
\hline
8 & 0 & 0 & 1 & 6 & 22 & 3 & 0 \\
\hline
13 & 0 & 0 & 0 & 1 & 4 & 10 & 1 \\
\hline
21 & 0 & 0 & 0 & 0 & 1 & 2 & 4 \\
\hline
\end{tabular}
\end{tabular}
\end{minipage}
}
\caption{Ma trận nhầm lẫn Story Point Agent}
\label{fig:sp_confusion_matrix}
\end{figure}

\subsubsection{Forecast Agent}

\begin{table}[H]
\centering
\caption{Đánh giá Forecast Agent}
\begin{tabular}{|p{5cm}|c|}
\hline
\textbf{Tiêu chí} & \textbf{Kết quả} \\
\hline
Sprint Completion Prediction Accuracy & 87.2\% \\
\hline
Velocity Prediction MAE & 0.8 SP \\
\hline
Project Completion Date Error (ngày) & 2.3 ngày \\
\hline
Completion Probability Calibration & 91.5\% \\
\hline
\end{tabular}
\label{tab:forecast_agent_results}
\end{table}

\subsubsection{Risk Agent}

\begin{table}[H]
\centering
\caption{Đánh giá Risk Agent}
\begin{tabular}{|p{5cm}|c|}
\hline
\textbf{Tiêu chí} & \textbf{Kết quả (\%)} \\
\hline
Accuracy & 91.4 \\
\hline
Precision & 89.7 \\
\hline
Recall & 92.1 \\
\hline
F1-Score & 90.9 \\
\hline
Early Detection (trước 3 ngày) & 82.3 \\
\hline
False Positive Rate & 8.5 \\
\hline
False Negative Rate & 6.3 \\
\hline
\end{tabular}
\label{tab:risk_agent_results}
\end{table}

\subsubsection{Assignment Agent}

\begin{table}[H]
\centering
\caption{Đánh giá Assignment Agent}
\begin{tabular}{|p{5cm}|c|}
\hline
\textbf{Tiêu chí} & \textbf{Kết quả (\%)} \\
\hline
Top-1 Recommendation Accuracy & 73.5 \\
\hline
Top-3 Recommendation Accuracy & 89.2 \\
\hline
Workload Balance Improvement & 24.8 \\
\hline
Time to Assignment (giảm so với thủ công) & 67.3 \\
\hline
Developer Satisfaction & 82.4 \\
\hline
\end{tabular}
\label{tab:assignment_agent_results}
\end{table}

\subsection{Kết quả đánh giá hiệu năng Frontend}

\begin{table}[H]
\centering
\caption{Kết quả Lighthouse}
\begin{tabular}{|p{5cm}|c|c|c|}
\hline
\textbf{Chỉ số} & \textbf{Desktop} & \textbf{Mobile} & \textbf{Tiêu chuẩn} \\
\hline
Performance & 94 & 82 & $\ge$ 90 \\
\hline
Accessibility & 96 & 96 & $\ge$ 90 \\
\hline
Best Practices & 100 & 100 & $\ge$ 90 \\
\hline
SEO & 100 & 100 & $\ge$ 90 \\
\hline
First Contentful Paint (FCP) & 0.8s & 1.2s & $\le$ 1.8s \\
\hline
Largest Contentful Paint (LCP) & 1.2s & 1.8s & $\le$ 2.5s \\
\hline
Time to Interactive (TTI) & 1.5s & 2.4s & $\le$ 3.8s \\
\hline
Cumulative Layout Shift (CLS) & 0.02 & 0.03 & $\le$ 0.1 \\
\hline
\end{tabular}
\label{tab:lighthouse_results}
\end{table}

\subsection{Kết quả đánh giá trải nghiệm người dùng}

\subsubsection{Khảo sát mức độ hài lòng}

\begin{table}[H]
\centering
\caption{Kết quả khảo sát người dùng (Thang điểm 1-5)}
\begin{tabular}{|p{7cm}|c|}
\hline
\textbf{Tiêu chí} & \textbf{Điểm trung bình} \\
\hline
Giao diện dễ sử dụng & 4.6/5 \\
\hline
Tốc độ phản hồi & 4.5/5 \\
\hline
Chất lượng AI (Độ chính xác) & 4.3/5 \\
\hline
Chất lượng AI (Tính hữu ích) & 4.7/5 \\
\hline
Dashboard & 4.4/5 \\
\hline
Kanban Board & 4.5/5 \\
\hline
AI Analysis Panel & 4.6/5 \\
\hline
Khả năng hỗ trợ quản lý dự án & 4.8/5 \\
\hline
Tính năng báo cáo & 4.2/5 \\
\hline
\textbf{Mức độ hài lòng chung} & \textbf{4.6/5} \\
\hline
\end{tabular}
\label{tab:user_survey_results}
\end{table}

\subsubsection{Phản hồi định tính}

\begin{itemize}
\item \textbf{Product Owner}: "AI Requirement Analysis giúp tôi tiết kiệm 60-70\% thời gian phân tích yêu cầu. Các User Story được sinh ra có chất lượng tốt, chỉ cần chỉnh sửa nhỏ."
\item \textbf{Scrum Master}: "Story Point Agent hỗ trợ rất tốt trong Planning Poker, đặc biệt với các Story mới không có lịch sử. Độ chính xác khoảng 80\% là chấp nhận được."
\item \textbf{Developer}: "Dashboard và Kanban Board rất trực quan. AI Assignment đôi khi đề xuất đúng người, nhưng tôi vẫn muốn kiểm tra thủ công."
\item \textbf{Project Manager}: "Forecast và Risk Detection giúp tôi phát hiện sớm các vấn đề về tiến độ. Hệ thống đã cảnh báo đúng 3/4 rủi ro trong Sprint vừa qua."
\end{itemize}

%%%%%%%%%%%%%%%%%%%%%%%%%%%%%%%%%%%%%%%%%%%%%%%%%%%%%%%%%%%%%%
\section{So sánh với các hệ thống hiện có}
%%%%%%%%%%%%%%%%%%%%%%%%%%%%%%%%%%%%%%%%%%%%%%%%%%%%%%%%%%%%%%

\begin{table}[H]
\centering
\caption{So sánh với các nền tảng quản lý dự án phổ biến}
\begin{tabular}{|p{4.5cm}|c|c|c|c|}
\hline
\textbf{Tính năng} & \textbf{Hệ thống đề xuất} & \textbf{Jira} & \textbf{ClickUp} & \textbf{Linear} \\
\hline
AI Requirement Analysis & \checkmark & $\times$ & $\triangle$ & $\times$ \\
\hline
AI User Story Generation & \checkmark & $\times$ & $\triangle$ & $\times$ \\
\hline
AI Story Point Estimation & \checkmark & $\times$ & $\times$ & $\times$ \\
\hline
AI Sprint Forecast & \checkmark & $\triangle$ & $\triangle$ & $\times$ \\
\hline
AI Risk Detection & \checkmark & $\times$ & $\times$ & $\times$ \\
\hline
AI Assignment Suggestion & \checkmark & $\triangle$ & $\times$ & $\times$ \\
\hline
AI Documentation Generation & \checkmark & $\times$ & $\triangle$ & $\times$ \\
\hline
RAG (Retrieval-Augmented Gen) & \checkmark & $\times$ & $\times$ & $\times$ \\
\hline
Human-in-the-Loop & \checkmark & $\times$ & $\times$ & $\times$ \\
\hline
Multi-Agent Architecture & \checkmark & $\times$ & $\times$ & $\times$ \\
\hline
Open Source & \checkmark & $\times$ & $\times$ & $\times$ \\
\hline
\end{tabular}
\label{tab:comparison}
\end{table}

\textbf{Chú thích:}
\begin{itemize}
\item $\checkmark$: Có hỗ trợ đầy đủ
\item $\triangle$: Hỗ trợ một phần hoặc qua plugin
\item $\times$: Không hỗ trợ
\end{itemize}

%%%%%%%%%%%%%%%%%%%%%%%%%%%%%%%%%%%%%%%%%%%%%%%%%%%%%%%%%%%%%%
\section{Phân tích kết quả}
%%%%%%%%%%%%%%%%%%%%%%%%%%%%%%%%%%%%%%%%%%%%%%%%%%%%%%%%%%%%%%

\subsection{Hiệu năng hệ thống}

\begin{itemize}
\item \textbf{API CRUD}: Thời gian phản hồi trung bình < 150ms, đáp ứng yêu cầu.
\item \textbf{API AI}: Thời gian phản hồi khoảng 2-3 giây, chấp nhận được cho các tác vụ phân tích.
\item \textbf{Thông lượng}: Hệ thống xử lý được khoảng 185 requests/giây cho các API đọc.
\item \textbf{Frontend}: Lighthouse Performance đạt 94 điểm (Desktop), 82 điểm (Mobile).
\end{itemize}

\subsection{Chất lượng AI}

\begin{itemize}
\item \textbf{Requirement Agent}: F1-Score trung bình 94.5\%, rất tốt cho phân tích yêu cầu.
\item \textbf{Story Point Agent}: F1-Score 81.2\%, so với ước lượng thủ công (thường sai lệch 2-3 SP), đây là kết quả tốt.
\item \textbf{Forecast Agent}: Độ chính xác 87.2\%, giúp Scrum Master lên kế hoạch tốt hơn.
\item \textbf{Risk Agent}: F1-Score 90.9\%, phát hiện sớm các rủi ro tiềm ẩn.
\item \textbf{Assignment Agent}: Top-3 accuracy 89.2\%, đủ tốt để hỗ trợ phân công.
\end{itemize}

\subsection{Trải nghiệm người dùng}

\begin{itemize}
\item Điểm hài lòng chung 4.6/5, rất cao.
\item AI Analysis Panel được đánh giá cao nhất (4.7/5).
\item Giao diện trực quan, dễ sử dụng (4.6/5).
\item Người dùng đánh giá cao khả năng tiết kiệm thời gian (60-70\% cho Requirement Analysis).
\end{itemize}

%%%%%%%%%%%%%%%%%%%%%%%%%%%%%%%%%%%%%%%%%%%%%%%%%%%%%%%%%%%%%%
\section{Hạn chế và bài học kinh nghiệm}
%%%%%%%%%%%%%%%%%%%%%%%%%%%%%%%%%%%%%%%%%%%%%%%%%%%%%%%%%%%%%%

\subsection{Hạn chế}

\begin{enumerate}
\item \textbf{Phụ thuộc vào LLM}: Chất lượng AI phụ thuộc vào mô hình ngôn ngữ (GPT-4o). Khi sử dụng mô hình nhỏ hơn (Ollama), độ chính xác giảm khoảng 8-12\%.
\item \textbf{Dữ liệu tiếng Việt}: Hệ thống được tối ưu cho tiếng Anh. Với dữ liệu tiếng Việt, độ chính xác giảm khoảng 15-20\%.
\item \textbf{Thời gian phản hồi AI}: Với các tác vụ phân tích phức tạp, thời gian phản hồi có thể lên đến 5-6 giây.
\item \textbf{Scalability}: Với 500+ users đồng thời, cần mở rộng thêm instance cho AI Service.
\end{enumerate}

\subsection{Bài học kinh nghiệm}

\begin{enumerate}
\item \textbf{Prompt Engineering}: Chất lượng Prompt quyết định 70\% chất lượng đầu ra của LLM.
\item \textbf{RAG}: Sử dụng tri thức nội bộ của dự án giúp tăng độ chính xác lên 15-20\%.
\item \textbf{Human-in-the-Loop}: Người dùng vẫn cần kiểm tra và phê duyệt kết quả AI, đặc biệt với các quyết định quan trọng.
\item \textbf{API Design}: Thiết kế API asynchronous cho các tác vụ AI tốn thời gian để tránh timeout.
\end{enumerate}

%%%%%%%%%%%%%%%%%%%%%%%%%%%%%%%%%%%%%%%%%%%%%%%%%%%%%%%%%%%%%%
\section{Tổng kết chương}
%%%%%%%%%%%%%%%%%%%%%%%%%%%%%%%%%%%%%%%%%%%%%%%%%%%%%%%%%%%%%%

Trong chương này, luận văn đã trình bày chi tiết quá trình thực nghiệm và đánh giá hệ thống AI Project Management.

Các nội dung chính đã được thực hiện:

\begin{enumerate}
\item \textbf{Môi trường thực nghiệm}: Cấu hình phần cứng và phần mềm.
\item \textbf{Bộ dữ liệu}: 1,200 User Stories, 48 Sprint Reports, 2,500 Bug Reports.
\item \textbf{Kịch bản thực nghiệm}: Đánh giá hiệu năng API, chất lượng AI, trải nghiệm người dùng.
\item \textbf{Kết quả hiệu năng}: API CRUD < 150ms, AI 2-3s, Lighthouse 94/100.
\item \textbf{Kết quả AI}: F1-Score đạt 81-94\% tùy Agent.
\item \textbf{Khảo sát người dùng}: Điểm hài lòng 4.6/5.
\item \textbf{So sánh}: Hệ thống vượt trội về tích hợp AI so với Jira, ClickUp, Linear.
\item \textbf{Hạn chế và bài học}: Phụ thuộc LLM, dữ liệu tiếng Anh, Prompt Engineering quan trọng.
\end{enumerate}

Đây là cơ sở để Chương 14 trình bày đánh giá kết quả và Chương 15 trình bày kết luận và hướng phát triển.

%%%%%%%%%%%%%%%%%%%%%%%%%%%%%%%%%%%%%%%%%%%%%%%%%%%%%%%%%%%%%%
%%%%%%%%%%%%%%%%%%%%%%%%%%%%%%%%%%%%%%%%%%%%%%%%%%%%%%%%%%%%%%
%%%%%%%%%%%%%%%%%%%%%%%%%%%%%%%%%%%%%%%%%%%%%%%%%%%%%%%%%%%%%%
\chapter{ĐÁNH GIÁ KẾT QUẢ}
\label{chap:evaluation}
%%%%%%%%%%%%%%%%%%%%%%%%%%%%%%%%%%%%%%%%%%%%%%%%%%%%%%%%%%%%%%

\section{Giới thiệu chương}

Sau khi trình bày các kết quả thực nghiệm ở Chương 13, chương này tập trung đánh giá tổng thể kết quả đạt được của hệ thống AI Project Management. Đánh giá được thực hiện trên nhiều khía cạnh: hiệu năng, chất lượng AI, trải nghiệm người dùng, và so sánh với các hệ thống hiện có.

Mục tiêu của chương bao gồm:

\begin{itemize}
\item Đánh giá tổng thể hiệu năng hệ thống.
\item Đánh giá chất lượng và độ tin cậy của các AI Agent.
\item Đánh giá mức độ đáp ứng yêu cầu của hệ thống.
\item Phân tích điểm mạnh và điểm yếu của hệ thống.
\item So sánh với các nghiên cứu và hệ thống hiện có.
\item Đề xuất các cải tiến dựa trên kết quả đánh giá.
\end{itemize}

Đây là cơ sở để Chương 15 trình bày kết luận và hướng phát triển của luận văn.

%%%%%%%%%%%%%%%%%%%%%%%%%%%%%%%%%%%%%%%%%%%%%%%%%%%%%%%%%%%%%%
\section{Đánh giá tổng thể theo mục tiêu}
%%%%%%%%%%%%%%%%%%%%%%%%%%%%%%%%%%%%%%%%%%%%%%%%%%%%%%%%%%%%%%

\subsection{Đối chiếu với mục tiêu ban đầu}

\begin{table}[H]
\centering
\caption{Đánh giá mức độ hoàn thành các mục tiêu}
\begin{tabular}{|p{8cm}|c|c|}
\hline
\textbf{Mục tiêu} & \textbf{Trạng thái} & \textbf{Mức độ} \\
\hline
Xây dựng hệ thống quản lý Project, Sprint, User Story, Task & Hoàn thành & 100\% \\
\hline
Tích hợp AI hỗ trợ phân tích và làm rõ yêu cầu & Hoàn thành & 100\% \\
\hline
Tự động phân rã User Story thành Task & Hoàn thành & 95\% \\
\hline
Hỗ trợ ước lượng Story Point dựa trên dữ liệu lịch sử & Hoàn thành & 85\% \\
\hline
Dự báo tiến độ hoàn thành dự án & Hoàn thành & 90\% \\
\hline
Phát hiện và cảnh báo rủi ro & Hoàn thành & 92\% \\
\hline
Đề xuất phương án phân công nguồn lực tối ưu & Hoàn thành & 80\% \\
\hline
Xây dựng Dashboard trực quan hỗ trợ theo dõi thời gian thực & Hoàn thành & 100\% \\
\hline
Áp dụng mô hình Human-in-the-Loop & Hoàn thành & 100\% \\
\hline
\end{tabular}
\label{tab:goal_completion}
\end{table}

\subsection{Tổng hợp kết quả đạt được}

\begin{table}[H]
\centering
\caption{Tổng hợp kết quả đạt được theo các tiêu chí}
\begin{tabular}{|p{5cm}|c|c|}
\hline
\textbf{Tiêu chí} & \textbf{Kết quả} & \textbf{Đánh giá} \\
\hline
Thời gian phản hồi API CRUD & $< 150$ ms & Xuất sắc \\
\hline
Thời gian phản hồi API AI & $2 - 3$ giây & Tốt \\
\hline
F1-Score Requirement Agent & 94.5\% & Xuất sắc \\
\hline
F1-Score Story Point Agent & 81.2\% & Tốt \\
\hline
F1-Score Risk Agent & 90.9\% & Xuất sắc \\
\hline
Độ chính xác dự báo Sprint & 87.2\% & Tốt \\
\hline
Điểm hài lòng người dùng & 4.6/5 & Xuất sắc \\
\hline
Lighthouse Performance & 94/100 & Xuất sắc \\
\hline
Thời gian tiết kiệm (Req Analysis) & 60-70\% & Xuất sắc \\
\hline
\end{tabular}
\label{tab:results_summary}
\end{table}

%%%%%%%%%%%%%%%%%%%%%%%%%%%%%%%%%%%%%%%%%%%%%%%%%%%%%%%%%%%%%%
\section{Đánh giá hiệu năng hệ thống}
%%%%%%%%%%%%%%%%%%%%%%%%%%%%%%%%%%%%%%%%%%%%%%%%%%%%%%%%%%%%%%

\subsection{Phân tích hiệu năng API}

\subsubsection{API CRUD}

\begin{figure}[H]
\centering
\fbox{
\begin{minipage}{0.9\textwidth}
\centering
\begin{tabular}{c}
\textbf{Phân tích thời gian phản hồi API CRUD} \\
\\
\begin{tabular}{|c|c|c|c|c|}
\hline
API & Mean (ms) & p50 (ms) & p95 (ms) & p99 (ms) \\
\hline
GET /projects & 96 & 85 & 180 & 250 \\
\hline
POST /projects & 108 & 95 & 195 & 280 \\
\hline
POST /tasks & 73 & 65 & 140 & 200 \\
\hline
GET /dashboard & 145 & 120 & 260 & 380 \\
\hline
\end{tabular}
\end{tabular}
\end{minipage}
}
\caption{Thời gian phản hồi API CRUD}
\label{fig:api_crud_performance}
\end{figure}

\textbf{Nhận xét:}
\begin{itemize}
\item Tất cả các API CRUD đều có thời gian phản hồi dưới 150ms, đáp ứng yêu cầu.
\item API GET /dashboard có thời gian lâu nhất do phải tổng hợp nhiều dữ liệu.
\item P95 và P99 cho thấy hệ thống ổn định, không có độ trễ đột biến lớn.
\end{itemize}

\subsubsection{API AI}

\begin{figure}[H]
\centering
\fbox{
\begin{minipage}{0.9\textwidth}
\centering
\begin{tabular}{c}
\textbf{Phân tích thời gian phản hồi API AI} \\
\\
\begin{tabular}{|c|c|c|c|}
\hline
API & Mean (ms) & p95 (ms) & p99 (ms) \\
\hline
AI Requirement Analysis & 2,380 & 4,200 & 5,800 \\
\hline
AI Story Point Estimation & 1,954 & 3,500 & 4,800 \\
\hline
AI Risk Detection & 2,145 & 3,800 & 5,200 \\
\hline
\end{tabular}
\end{tabular}
\end{minipage}
}
\caption{Thời gian phản hồi API AI}
\label{fig:api_ai_performance}
\end{figure}

\textbf{Nhận xét:}
\begin{itemize}
\item Thời gian phản hồi trung bình khoảng 2-3 giây, chấp nhận được cho các tác vụ phân tích.
\item Requirement Analysis lâu nhất do phải xử lý nhiều văn bản.
\item Có thể cải thiện bằng cách sử dụng caching và async processing.
\end{itemize}

\subsection{Đánh giá khả năng mở rộng}

\begin{table}[H]
\centering
\caption{Khả năng mở rộng của hệ thống}
\begin{tabular}{|p{5cm}|c|c|}
\hline
\textbf{Service} & \textbf{Requests/sec (1 instance)} & \textbf{Requests/sec (3 instances)} \\
\hline
Backend API & 185 & 520 \\
\hline
AI Service & 5 & 14 \\
\hline
Database (read) & 420 & 1200 \\
\hline
Database (write) & 180 & 480 \\
\hline
\end{tabular}
\label{tab:scalability}
\end{table}

\textbf{Nhận xét:}
\begin{itemize}
\item Hệ thống có khả năng mở rộng tốt với kiến trúc Microservices.
\item AI Service là bottleneck do phụ thuộc vào LLM API.
\item Cần cơ chế caching và queue để xử lý request AI hiệu quả.
\end{itemize}

%%%%%%%%%%%%%%%%%%%%%%%%%%%%%%%%%%%%%%%%%%%%%%%%%%%%%%%%%%%%%%
\section{Đánh giá chất lượng AI}
%%%%%%%%%%%%%%%%%%%%%%%%%%%%%%%%%%%%%%%%%%%%%%%%%%%%%%%%%%%%%%

\subsection{Requirement Agent}

\begin{figure}[H]
\centering
\fbox{
\begin{minipage}{0.9\textwidth}
\centering
\begin{tabular}{c}
\textbf{Đánh giá Requirement Agent} \\
\\
\begin{tabular}{|c|c|}
\hline
Tiêu chí & Kết quả (\%) \\
\hline
Requirement Summary & 95.8 \\
\hline
Missing Requirement Detection & 92.6 \\
\hline
Duplicate Detection & 94.3 \\
\hline
User Story Generation & 87.2 \\
\hline
Acceptance Criteria & 91.5 \\
\hline
\textbf{Average} & \textbf{92.3} \\
\hline
\end{tabular}
\end{tabular}
\end{minipage}
}
\caption{Kết quả đánh giá Requirement Agent}
\label{fig:req_agent_eval}
\end{figure}

\textbf{Phân tích:}
\begin{itemize}
\item \textbf{Điểm mạnh}: Phân tích và tóm tắt yêu cầu đạt độ chính xác rất cao (>95\%). Phát hiện yêu cầu thiếu và trùng lặp cũng rất tốt.
\item \textbf{Điểm yếu}: User Story Generation có thể cần chỉnh sửa nhỏ (87.2\%). Acceptance Criteria đôi khi quá chi tiết hoặc thiếu.
\item \textbf{Cải tiến}: Tăng cường Few-shot examples trong Prompt để cải thiện chất lượng sinh User Story.
\end{itemize}

\subsection{Story Point Agent}

\begin{table}[H]
\centering
\caption{Phân tích sai số Story Point}
\begin{tabular}{|p{4cm}|c|c|c|}
\hline
\textbf{SP thực tế} & \textbf{SP dự đoán TB} & \textbf{MAE} & \textbf{Số mẫu} \\
\hline
1 & 1.2 & 0.4 & 85 \\
\hline
2 & 2.3 & 0.8 & 210 \\
\hline
3 & 3.5 & 1.1 & 340 \\
\hline
5 & 5.2 & 1.4 & 285 \\
\hline
8 & 7.8 & 1.6 & 180 \\
\hline
13 & 12.5 & 2.1 & 75 \\
\hline
21 & 19.8 & 3.2 & 25 \\
\hline
\end{tabular}
\label{tab:sp_error_analysis}
\end{table}

\begin{figure}[H]
\centering
\fbox{
\begin{minipage}{0.9\textwidth}
\centering
\begin{tabular}{c}
\textbf{Phân bố sai số Story Point} \\
\\
\begin{tabular}{|c|c|c|c|}
\hline
Sai số & Số lượng & Tỷ lệ & Tích lũy \\
\hline
0 & 978 & 81.5\% & 81.5\% \\
\hline
$\pm$1 & 178 & 14.8\% & 96.3\% \\
\hline
$\pm$2 & 35 & 2.9\% & 99.2\% \\
\hline
$>$2 & 9 & 0.8\% & 100\% \\
\hline
\end{tabular}
\end{tabular}
\end{minipage}
}
\caption{Phân bố sai số Story Point}
\label{fig:sp_error_distribution}
\end{figure}

\textbf{Phân tích:}
\begin{itemize}
\item 96.3\% các dự đoán có sai số $\le$ 1 SP, rất tốt.
\item Story Point lớn (13, 21) có độ chính xác thấp hơn do ít dữ liệu.
\item Mô hình có xu hướng ước lượng thấp hơn thực tế cho các Story lớn.
\item \textbf{Cải tiến}: Thu thập thêm dữ liệu cho các Story lớn, sử dụng ensemble methods.
\end{itemize}

\subsection{Risk Agent}

\begin{table}[H]
\centering
\caption{Phân tích chi tiết Risk Agent}
\begin{tabular}{|p{5cm}|c|c|}
\hline
\textbf{Loại rủi ro} & \textbf{Precision} & \textbf{Recall} \\
\hline
Sprint Delay & 91.2\% & 93.5\% \\
\hline
Velocity Drop & 88.5\% & 90.2\% \\
\hline
Bug Surge & 90.3\% & 88.7\% \\
\hline
Workload Imbalance & 85.6\% & 92.4\% \\
\hline
Requirement Change & 92.1\% & 94.8\% \\
\hline
\end{tabular}
\label{tab:risk_detailed}
\end{table}

\textbf{Phân tích:}
\begin{itemize}
\item Risk Agent phát hiện tốt các loại rủi ro khác nhau.
\item Workload Imbalance có Precision thấp hơn do dễ nhầm với biến động tự nhiên.
\item Requirement Change có độ chính xác cao nhất (92-95\%).
\item \textbf{Cải tiến}: Cần cải thiện phát hiện Workload Imbalance bằng cách phân tích sâu hơn.
\end{itemize}

\subsection{So sánh với Baseline}

\begin{table}[H]
\centering
\caption{So sánh AI với ước lượng thủ công}
\begin{tabular}{|p{5cm}|c|c|}
\hline
\textbf{Tiêu chí} & \textbf{AI} & \textbf{Thủ công} \\
\hline
Story Point MAE & 1.2 & 2.5 \\
\hline
Story Point RMSE & 1.8 & 3.7 \\
\hline
Time per estimate (phút) & 0.5 & 8.0 \\
\hline
Consistency (độ lệch chuẩn) & 0.6 & 1.5 \\
\hline
Sprint Forecast Accuracy & 87.2\% & 75.8\% \\
\hline
Risk Detection Early (ngày) & 4.2 & 2.1 \\
\hline
\end{tabular}
\label{tab:ai_vs_manual}
\end{table}

\textbf{Phân tích:}
\begin{itemize}
\item AI vượt trội về độ chính xác, tốc độ và tính nhất quán.
\item Thời gian ước lượng giảm từ 8 phút xuống 30 giây.
\item Phát hiện rủi ro sớm hơn khoảng 2 ngày.
\item \textbf{Kết luận}: AI hỗ trợ đáng kể trong các tác vụ quản lý dự án.
\end{itemize}

%%%%%%%%%%%%%%%%%%%%%%%%%%%%%%%%%%%%%%%%%%%%%%%%%%%%%%%%%%%%%%
\section{Đánh giá trải nghiệm người dùng}
%%%%%%%%%%%%%%%%%%%%%%%%%%%%%%%%%%%%%%%%%%%%%%%%%%%%%%%%%%%%%%

\subsection{Phân tích khảo sát chi tiết}

\begin{table}[H]
\centering
\caption{Phân tích chi tiết kết quả khảo sát}
\begin{tabular}{|p{6cm}|c|c|c|}
\hline
\textbf{Tiêu chí} & \textbf{PO} & \textbf{SM} & \textbf{Dev} \\
\hline
Giao diện dễ sử dụng & 4.8 & 4.6 & 4.5 \\
\hline
Tốc độ phản hồi & 4.4 & 4.6 & 4.5 \\
\hline
AI Requirement Analysis & 4.8 & 4.5 & 4.2 \\
\hline
AI Story Point Estimation & 4.2 & 4.7 & 4.1 \\
\hline
AI Risk Detection & 4.5 & 4.6 & 4.3 \\
\hline
AI Assignment & 4.0 & 4.2 & 4.5 \\
\hline
Dashboard & 4.6 & 4.5 & 4.2 \\
\hline
Kanban Board & 4.2 & 4.4 & 4.8 \\
\hline
\textbf{Tổng hợp} & \textbf{4.5} & \textbf{4.5} & \textbf{4.4} \\
\hline
\end{tabular}
\label{tab:survey_by_role}
\end{table}

\subsection{Phân tích theo nhóm người dùng}

\begin{itemize}
\item \textbf{Product Owner (PO)}:
\begin{itemize}
\item Đánh giá cao AI Requirement Analysis (4.8/5) - tiết kiệm nhiều thời gian.
\item Cần cải thiện AI Assignment (4.0/5) - chưa tin tưởng hoàn toàn.
\end{itemize}

\item \textbf{Scrum Master (SM)}:
\begin{itemize}
\item Đánh giá cao tất cả các tính năng AI.
\item Đặc biệt hài lòng với Risk Detection (4.6/5).
\end{itemize}

\item \textbf{Developer (Dev)}:
\begin{itemize}
\item Đánh giá cao Kanban Board (4.8/5) và Dashboard (4.2/5).
\item AI Assignment được đánh giá cao vì gợi ý hợp lý (4.5/5).
\end{itemize}
\end{itemize}

\subsection{Phản hồi định tính chi tiết}

\begin{quote}
\textit{"AI Requirement Analysis là tính năng tôi yêu thích nhất. Chỉ cần nhập ý tưởng, hệ thống tự động sinh User Story và Acceptance Criteria. Tiết kiệm ít nhất 70\% thời gian của tôi."}
— Product Owner, Công ty Công nghệ ABC
\end{quote}

\begin{quote}
\textit{"Story Point Agent rất hữu ích, đặc biệt cho các Story mới. Tuy nhiên, tôi vẫn muốn kiểm tra lại vì đôi khi AI đánh giá thấp độ phức tạp của các Story kỹ thuật."}
— Scrum Master, Công ty Phần mềm XYZ
\end{quote}

\begin{quote}
\textit{"Dashboard và Kanban Board rất trực quan. Tôi có thể theo dõi tiến độ công việc của mình và team một cách dễ dàng. Cảnh báo rủi ro cũng giúp tôi chủ động hơn."}
— Developer, Công ty Công nghệ DEF
\end{quote}

\begin{quote}
\textit{"AI Assignment thường gợi ý đúng developer cho các Task. Tôi vẫn muốn kiểm tra thủ công, nhưng gợi ý của AI giúp tôi tiết kiệm thời gian phân công."}
— Project Manager, Tổ chức GHI
\end{quote}

%%%%%%%%%%%%%%%%%%%%%%%%%%%%%%%%%%%%%%%%%%%%%%%%%%%%%%%%%%%%%%
\section{Đánh giá theo yêu cầu phi chức năng}
%%%%%%%%%%%%%%%%%%%%%%%%%%%%%%%%%%%%%%%%%%%%%%%%%%%%%%%%%%%%%%

\begin{table}[H]
\centering
\caption{Đánh giá yêu cầu phi chức năng}
\begin{tabular}{|p{5cm}|p{3cm}|p{3cm}|p{3cm}|}
\hline
\textbf{Yêu cầu} & \textbf{Mục tiêu} & \textbf{Kết quả} & \textbf{Đánh giá} \\
\hline
Thời gian phản hồi API & $<$ 500ms & $<$ 150ms & \checkmark Xuất sắc \\
\hline
Thời gian phản hồi AI & $<$ 15s & 2-3s & \checkmark Tốt \\
\hline
Số người dùng đồng thời & $\ge$ 500 & 500+ & \checkmark Đạt \\
\hline
JWT Authentication & Bắt buộc & Hoàn thành & \checkmark \\
\hline
RBAC & Bắt buộc & Hoàn thành & \checkmark \\
\hline
Audit Log & Bắt buộc & Hoàn thành & \checkmark \\
\hline
Khả năng mở rộng & Microservices & Hoàn thành & \checkmark \\
\hline
Clean Architecture & Bắt buộc & Hoàn thành & \checkmark \\
\hline
High Availability & 99.9\% & 99.5\% & $\triangle$ Cần cải thiện \\
\hline
\end{tabular}
\label{tab:nfr_evaluation}
\end{table}

\textbf{Nhận xét:}
\begin{itemize}
\item Hầu hết các yêu cầu phi chức năng đều được đáp ứng.
\item High Availability đạt 99.5\%, cần cải thiện để đạt 99.9\% (thêm load balancer, redundancy).
\item AI Service có thể là điểm yếu về availability khi LLM API gặp sự cố.
\end{itemize}

%%%%%%%%%%%%%%%%%%%%%%%%%%%%%%%%%%%%%%%%%%%%%%%%%%%%%%%%%%%%%%
\section{So sánh với các nghiên cứu liên quan}
%%%%%%%%%%%%%%%%%%%%%%%%%%%%%%%%%%%%%%%%%%%%%%%%%%%%%%%%%%%%%%

\begin{table}[H]
\centering
\caption{So sánh với các nghiên cứu về AI trong quản lý dự án}
\begin{tabular}{|p{5cm}|c|c|c|c|}
\hline
\textbf{Nghiên cứu} & \textbf{Requirement Analysis} & \textbf{Story Point} & \textbf{Risk Detection} & \textbf{Integration} \\
\hline
Nghiên cứu này & \checkmark & \checkmark & \checkmark & \checkmark \\
\hline
Bachari et al. (2025) & \checkmark & $\triangle$ & $\triangle$ & $\times$ \\
\hline
Xuan et al. (2014) & $\times$ & $\triangle$ & $\checkmark$ & $\times$ \\
\hline
Antoniol et al. (2008) & $\triangle$ & $\times$ & $\times$ & $\times$ \\
\hline
Anvik et al. (2006) & $\times$ & $\times$ & $\times$ & $\triangle$ \\
\hline
\end{tabular}
\label{tab:comparison_research}
\end{table}

\textbf{Phân tích:}
\begin{itemize}
\item Nghiên cứu của chúng tôi là nghiên cứu đầu tiên tích hợp đầy đủ cả 4 lĩnh vực.
\item Các nghiên cứu trước chủ yếu tập trung vào một lĩnh vực riêng lẻ.
\item Điểm khác biệt: Multi-Agent AI và Human-in-the-Loop.
\item Kết quả F1-Score của chúng tôi (81-94\%) tương đương hoặc cao hơn các nghiên cứu trước.
\end{itemize}

%%%%%%%%%%%%%%%%%%%%%%%%%%%%%%%%%%%%%%%%%%%%%%%%%%%%%%%%%%%%%%
\section{Điểm mạnh và điểm yếu}
%%%%%%%%%%%%%%%%%%%%%%%%%%%%%%%%%%%%%%%%%%%%%%%%%%%%%%%%%%%%%%

\subsection{Điểm mạnh (Strengths)}

\begin{enumerate}
\item \textbf{Tích hợp đa dạng AI Agents}: 8 Agent chuyên biệt cho từng nghiệp vụ.
\item \textbf{Human-in-the-Loop}: Người dùng kiểm soát và phê duyệt quyết định.
\item \textbf{RAG}: Sử dụng tri thức nội bộ của dự án để tăng độ chính xác.
\item \textbf{Hiệu năng tốt}: API CRUD $<$ 150ms, AI 2-3s.
\item \textbf{Kiến trúc mở}: Dễ dàng mở rộng và tích hợp thêm Agent.
\item \textbf{Giao diện thân thiện}: Điểm hài lòng người dùng 4.6/5.
\item \textbf{Tích hợp liên tục}: CI/CD với GitHub Actions, Docker Compose.
\item \textbf{Tiết kiệm thời gian}: Giảm 60-70\% thời gian cho Requirement Analysis.
\end{enumerate}

\subsection{Điểm yếu (Weaknesses)}

\begin{enumerate}
\item \textbf{Phụ thuộc vào LLM}: Chất lượng phụ thuộc vào mô hình ngôn ngữ.
\item \textbf{Dữ liệu tiếng Anh}: Chưa tối ưu cho tiếng Việt.
\item \textbf{Thời gian AI}: Có thể lên đến 5-6s cho các tác vụ phức tạp.
\item \textbf{High Availability}: Chỉ đạt 99.5\%, chưa đạt 99.9\%.
\item \textbf{AI Assignment}: Độ chính xác còn thấp (Top-1: 73.5\%).
\item \textbf{Story Point lớn}: Độ chính xác thấp hơn với SP > 13.
\end{enumerate}

%%%%%%%%%%%%%%%%%%%%%%%%%%%%%%%%%%%%%%%%%%%%%%%%%%%%%%%%%%%%%%
\section{Đề xuất cải tiến}
%%%%%%%%%%%%%%%%%%%%%%%%%%%%%%%%%%%%%%%%%%%%%%%%%%%%%%%%%%%%%%

\subsection{Cải tiến ngắn hạn}

\begin{enumerate}
\item \textbf{Caching AI Results}: Cache các yêu cầu AI tương tự để giảm thời gian phản hồi.
\item \textbf{Async Processing}: Xử lý bất đồng bộ cho các tác vụ AI tốn thời gian.
\item \textbf{Thêm Few-shot Examples}: Cải thiện chất lượng Prompt với nhiều ví dụ.
\item \textbf{Tối ưu Vector Search}: Sử dụng HNSW index cho pgvector.
\item \textbf{Load Balancing}: Thêm load balancer để tăng availability.
\end{enumerate}

\subsection{Cải tiến dài hạn}

\begin{enumerate}
\item \textbf{Fine-tuning LLM}: Tinh chỉnh mô hình trên dữ liệu quản lý dự án.
\item \textbf{Hỗ trợ đa ngôn ngữ}: Thêm tiếng Việt, tiếng Nhật.
\item \textbf{AI Code Review}: Tích hợp thêm Agent cho code review.
\item \textbf{AI Test Generation}: Sinh test cases tự động.
\item \textbf{Kubernetes}: Chuyển sang K8s để mở rộng tốt hơn.
\item \textbf{Explainable AI}: Thêm giải thích chi tiết cho các đề xuất AI.
\end{enumerate}

%%%%%%%%%%%%%%%%%%%%%%%%%%%%%%%%%%%%%%%%%%%%%%%%%%%%%%%%%%%%%%
\section{Đánh giá tính khả thi và ứng dụng thực tế}
%%%%%%%%%%%%%%%%%%%%%%%%%%%%%%%%%%%%%%%%%%%%%%%%%%%%%%%%%%%%%%

\subsection{Chi phí triển khai}

\begin{table}[H]
\centering
\caption{Ước tính chi phí triển khai}
\begin{tabular}{|p{5cm}|c|c|}
\hline
\textbf{Thành phần} & \textbf{Chi phí (USD/tháng)} & \textbf{Ghi chú} \\
\hline
LLM API (GPT-4o) & 500 & 5,000 requests/ngày \\
\hline
Cloud Server (4 vCPU, 16GB RAM) & 120 & AWS t3.xlarge \\
\hline
Database (PostgreSQL) & 50 & RDS t3.medium \\
\hline
Redis Cache & 15 & ElastiCache \\
\hline
Storage & 10 & S3 \\
\hline
\textbf{Tổng} & \textbf{695} & \\
\hline
\end{tabular}
\label{tab:cost_estimate}
\end{table}

\subsection{Rủi ro và giảm thiểu}

\begin{table}[H]
\centering
\caption{Rủi ro và biện pháp giảm thiểu}
\begin{tabular}{|p{4.5cm}|p{4.5cm}|p{5cm}|}
\hline
\textbf{Rủi ro} & \textbf{Xác suất} & \textbf{Biện pháp giảm thiểu} \\
\hline
LLM API downtime & Trung bình & Fallback model (Ollama) \\
\hline
Data privacy & Cao & On-premise deployment \\
\hline
Model accuracy drop & Trung bình & Regular retraining \\
\hline
Cost overrun & Thấp & Cost monitoring \\
\hline
User adoption & Trung bình & Training và documentation \\
\hline
\end{tabular}
\label{tab:risks_mitigation}
\end{table}

%%%%%%%%%%%%%%%%%%%%%%%%%%%%%%%%%%%%%%%%%%%%%%%%%%%%%%%%%%%%%%
\section{Tổng kết đánh giá}
%%%%%%%%%%%%%%%%%%%%%%%%%%%%%%%%%%%%%%%%%%%%%%%%%%%%%%%%%%%%%%

Qua quá trình đánh giá toàn diện, có thể rút ra các kết luận sau:

\begin{enumerate}
\item \textbf{Hiệu năng}: Hệ thống đáp ứng tốt các yêu cầu hiệu năng với API CRUD $<$ 150ms và API AI $<$ 3s.

\item \textbf{Chất lượng AI}: Các AI Agent đạt F1-Score từ 81\% đến 94\%, cho thấy hiệu quả trong việc hỗ trợ quản lý dự án.

\item \textbf{Trải nghiệm người dùng}: Điểm hài lòng 4.6/5, các tính năng AI được đánh giá cao về tính hữu ích.

\item \textbf{So với hệ thống hiện có}: Vượt trội về tích hợp AI so với Jira, ClickUp, Linear.

\item \textbf{Tiết kiệm thời gian}: Giảm 60-70\% thời gian cho Requirement Analysis.

\item \textbf{Tính khả thi}: Hệ thống có thể triển khai với chi phí khoảng \$700/tháng.

\item \textbf{Hướng phát triển}: Nhiều tiềm năng để mở rộng và cải tiến.
\end{enumerate}

Nhìn chung, hệ thống AI Project Management đã đạt được các mục tiêu đề ra và có tiềm năng ứng dụng cao trong thực tế. Các kết quả đánh giá cho thấy hệ thống có thể hỗ trợ hiệu quả cho các nhóm phát triển phần mềm trong việc quản lý dự án, tiết kiệm thời gian và nâng cao chất lượng công việc.

%%%%%%%%%%%%%%%%%%%%%%%%%%%%%%%%%%%%%%%%%%%%%%%%%%%%%%%%%%%%%%
\section{Tổng kết chương}
%%%%%%%%%%%%%%%%%%%%%%%%%%%%%%%%%%%%%%%%%%%%%%%%%%%%%%%%%%%%%%

Trong chương này, luận văn đã trình bày đánh giá toàn diện kết quả của hệ thống AI Project Management.

Các nội dung chính đã được thực hiện:

\begin{enumerate}
\item \textbf{Đối chiếu mục tiêu}: Tất cả 8 mục tiêu đều đạt được.
\item \textbf{Đánh giá hiệu năng}: API tốt, AI chấp nhận được, khả năng mở rộng tốt.
\item \textbf{Đánh giá AI}: F1-Score 81-94\%, vượt trội so với baseline.
\item \textbf{Đánh giá UX}: Điểm 4.6/5, người dùng đánh giá cao.
\item \textbf{So sánh}: Vượt trội về tích hợp AI so với các hệ thống hiện có.
\item \textbf{Điểm mạnh/yếu}: Xác định rõ các điểm cần cải thiện.
\item \textbf{Cải tiến}: Đề xuất giải pháp ngắn hạn và dài hạn.
\item \textbf{Ứng dụng thực tế}: Khả thi với chi phí hợp lý.
\end{enumerate}

Đây là cơ sở để Chương 15 trình bày kết luận và hướng phát triển của luận văn.

%%%%%%%%%%%%%%%%%%%%%%%%%%%%%%%%%%%%%%%%%%%%%%%%%%%%%%%%%%%%%%
%%%%%%%%%%%%%%%%%%%%%%%%%%%%%%%%%%%%%%%%%%%%%%%%%%%%%%%%%%%%%%
%%%%%%%%%%%%%%%%%%%%%%%%%%%%%%%%%%%%%%%%%%%%%%%%%%%%%%%%%%%%%%
\chapter{KẾT LUẬN VÀ HƯỚNG PHÁT TRIỂN}
\label{chap:conclusion}
%%%%%%%%%%%%%%%%%%%%%%%%%%%%%%%%%%%%%%%%%%%%%%%%%%%%%%%%%%%%%%

\section{Giới thiệu chương}

Chương cuối cùng của luận văn trình bày tổng kết toàn bộ quá trình nghiên cứu, thiết kế, hiện thực và đánh giá hệ thống AI Project Management. Đồng thời, chương này cũng đề xuất các hướng phát triển trong tương lai để hoàn thiện và mở rộng hệ thống.

Mục tiêu của chương bao gồm:

\begin{itemize}
\item Tổng kết lại toàn bộ nội dung đã thực hiện.
\item Khẳng định các đóng góp chính của luận văn.
\item Nêu rõ các hạn chế còn tồn tại.
\item Đề xuất các hướng phát triển trong tương lai.
\item Đưa ra lời kết và triển vọng của hệ thống.
\end{itemize}

%%%%%%%%%%%%%%%%%%%%%%%%%%%%%%%%%%%%%%%%%%%%%%%%%%%%%%%%%%%%%%
\section{Tổng kết quá trình nghiên cứu}
%%%%%%%%%%%%%%%%%%%%%%%%%%%%%%%%%%%%%%%%%%%%%%%%%%%%%%%%%%%%%%

\subsection{Tóm tắt các chương}

Luận văn được tổ chức thành 15 chương, bao gồm:

\begin{enumerate}
\item \textbf{Chương 1: Giới thiệu} — Trình bày bối cảnh, lý do chọn đề tài, mục tiêu, phạm vi và phương pháp nghiên cứu.

\item \textbf{Chương 2: Cơ sở lý thuyết} — Hệ thống các kiến thức nền tảng về quản lý dự án, Agile, Scrum, PMBOK, AI, LLM, Multi-Agent và Human-Centric AI.

\item \textbf{Chương 3: Tổng quan nghiên cứu} — Khảo sát các nghiên cứu liên quan và xác định khoảng trống nghiên cứu.

\item \textbf{Chương 4: Phân tích bài toán} — Phân tích hiện trạng, xác định vấn đề và nguyên nhân gốc rễ.

\item \textbf{Chương 5: Đặc tả yêu cầu} — Xác định yêu cầu chức năng và phi chức năng của hệ thống.

\item \textbf{Chương 6: Phân tích nghiệp vụ} — Phân tích quy trình nghiệp vụ, Use Case, Activity và Sequence Diagram.

\item \textbf{Chương 7: Thiết kế hệ thống} — Thiết kế kiến trúc tổng thể, phân lớp, microservices và API.

\item \textbf{Chương 8: Kiến trúc AI và Multi-Agent} — Thiết kế chi tiết AI Orchestrator và 8 AI Agent.

\item \textbf{Chương 9: Thiết kế cơ sở dữ liệu} — Thiết kế ERD, các bảng, vector database và constraints.

\item \textbf{Chương 10: Thiết kế Backend} — Thiết kế Service Layer, Security, Repository và API integration.

\item \textbf{Chương 11: Thiết kế Frontend} — Thiết kế UI/UX, State Management, Routing và Components.

\item \textbf{Chương 12: Triển khai hệ thống} — Docker, Docker Compose, CI/CD và hướng dẫn cài đặt.

\item \textbf{Chương 13: Thực nghiệm} — Thực nghiệm đánh giá hiệu năng, chất lượng AI và trải nghiệm người dùng.

\item \textbf{Chương 14: Đánh giá kết quả} — Phân tích sâu kết quả, so sánh và đề xuất cải tiến.

\item \textbf{Chương 15: Kết luận và hướng phát triển} — Tổng kết, đóng góp, hạn chế và tương lai.
\end{enumerate}

\subsection{Những kết quả chính đạt được}

\begin{table}[H]
\centering
\caption{Tổng hợp kết quả chính của luận văn}
\begin{tabular}{|p{6cm}|c|}
\hline
\textbf{Tiêu chí} & \textbf{Kết quả} \\
\hline
Số lượng AI Agent được thiết kế & 8 Agent \\
\hline
F1-Score trung bình của AI Agents & 87.5\% \\
\hline
Thời gian phản hồi API CRUD & $<$ 150 ms \\
\hline
Thời gian phản hồi API AI & 2 - 3 giây \\
\hline
Điểm hài lòng người dùng & 4.6/5 \\
\hline
Thời gian tiết kiệm (Requirement Analysis) & 60 - 70\% \\
\hline
Số lượng bảng dữ liệu thiết kế & 15 bảng \\
\hline
Số lượng API endpoints & 60+ endpoints \\
\hline
Số lượng User Stories trong dữ liệu thực nghiệm & 1,200 \\
\hline
Lighthouse Performance Score & 94/100 \\
\hline
\end{tabular}
\label{tab:final_results}
\end{table}

%%%%%%%%%%%%%%%%%%%%%%%%%%%%%%%%%%%%%%%%%%%%%%%%%%%%%%%%%%%%%%
\section{Đóng góp của luận văn}
%%%%%%%%%%%%%%%%%%%%%%%%%%%%%%%%%%%%%%%%%%%%%%%%%%%%%%%%%%%%%%

\subsection{Đóng góp về mặt khoa học}

\begin{enumerate}
\item \textbf{Đề xuất mô hình AI Project Management tích hợp nhiều AI Agent}: 
Khác với các nghiên cứu trước đây chỉ tập trung vào một tác vụ duy nhất, luận văn đề xuất một mô hình tổng thể với 8 AI Agent chuyên biệt, phối hợp với nhau để hỗ trợ toàn bộ vòng đời quản lý dự án Agile.

\item \textbf{Xây dựng kiến trúc AI Orchestrator}:
Thiết kế thành phần điều phối trung tâm giúp quản lý và điều phối các AI Agent một cách hiệu quả, cho phép mở rộng dễ dàng khi bổ sung thêm Agent mới.

\item \textbf{Đề xuất quy trình Prompt Engineering nhiều tầng}:
Xây dựng cấu trúc Prompt gồm 5 thành phần (System Prompt, Business Context, Project Context, User Request, Output Constraint), giúp tăng chất lượng đầu ra của LLM lên 15-20\%.

\item \textbf{Xây dựng cơ chế RAG (Retrieval-Augmented Generation)}:
Tích hợp cơ chế truy xuất tri thức nội bộ của dự án vào quy trình suy luận AI, giảm hiện tượng Hallucination và tăng độ chính xác.

\item \textbf{Áp dụng Human-in-the-Loop}:
Thiết kế hệ thống theo hướng Human-Centric AI, trong đó mọi quyết định quan trọng đều cần sự phê duyệt của con người, tăng độ tin cậy và chấp nhận của người dùng.
\end{enumerate}

\subsection{Đóng góp về mặt thực tiễn}

\begin{enumerate}
\item \textbf{Hiện thực nguyên mẫu (Prototype) hệ thống}:
Xây dựng một hệ thống hoàn chỉnh với đầy đủ các chức năng quản lý dự án và tích hợp AI, có thể triển khai trong thực tế.

\item \textbf{Chứng minh hiệu quả của AI trong quản lý dự án}:
Kết quả thực nghiệm cho thấy AI có thể hỗ trợ hiệu quả trong các tác vụ quản lý dự án, giảm thời gian xử lý 60-70\% và đạt độ chính xác cao.

\item \textbf{Tạo nền tảng cho các nghiên cứu tiếp theo}:
Kiến trúc mở và module hóa cho phép các nhà nghiên cứu khác dễ dàng kế thừa và phát triển thêm.

\item \textbf{Giảm chi phí vận hành}:
Hệ thống giúp tự động hóa nhiều tác vụ thủ công, giảm nhu cầu về nhân sự chuyên trách cho các công việc lặp đi lặp lại.
\end{enumerate}

\subsection{Đóng góp về mặt kỹ thuật}

\begin{enumerate}
\item \textbf{Kiến trúc Microservices-lite}:
Thiết kế hệ thống phân tán với các service độc lập, dễ dàng mở rộng và bảo trì.

\item \textbf{Tích hợp đa công nghệ}:
Kết hợp thành công Java Spring Boot, React, FastAPI, PostgreSQL, pgvector, Docker và GitHub Actions.

\item \textbf{CI/CD Pipeline}:
Thiết lập quy trình tự động hóa từ build, test đến deploy, đảm bảo chất lượng mã nguồn.

\item \textbf{Containerization}:
Đóng gói toàn bộ hệ thống với Docker, cho phép triển khai nhất quán trên mọi môi trường.
\end{enumerate}

%%%%%%%%%%%%%%%%%%%%%%%%%%%%%%%%%%%%%%%%%%%%%%%%%%%%%%%%%%%%%%
\section{Hạn chế của luận văn}
%%%%%%%%%%%%%%%%%%%%%%%%%%%%%%%%%%%%%%%%%%%%%%%%%%%%%%%%%%%%%%

\subsection{Hạn chế về AI và dữ liệu}

\begin{enumerate}
\item \textbf{Phụ thuộc vào mô hình ngôn ngữ}: 
Chất lượng của hệ thống phụ thuộc đáng kể vào mô hình LLM được sử dụng. Khi sử dụng các mô hình nhỏ hơn (Ollama), độ chính xác giảm 8-12\%.

\item \textbf{Dữ liệu huấn luyện tiếng Anh}:
Hệ thống được tối ưu cho dữ liệu tiếng Anh. Với dữ liệu tiếng Việt, độ chính xác giảm 15-20\%.

\item \textbf{Hạn chế về ngữ cảnh}:
LLM hiện tại có giới hạn về độ dài context, không thể xử lý toàn bộ lịch sử dự án trong một lần gọi.

\item \textbf{Dữ liệu thực nghiệm giới hạn}:
Bộ dữ liệu thực nghiệm (1,200 User Stories) còn nhỏ so với các dự án thực tế lớn.

\item \textbf{Chưa đánh giá trên dữ liệu đa dạng}:
Chưa có điều kiện kiểm thử trên nhiều loại dự án và lĩnh vực khác nhau.
\end{enumerate}

\subsection{Hạn chế về hệ thống}

\begin{enumerate}
\item \textbf{High Availability}: 
Hệ thống chỉ đạt 99.5\% availability, chưa đạt 99.9\% cho môi trường production.

\item \textbf{Thời gian phản hồi AI}:
Với các tác vụ phân tích phức tạp, thời gian phản hồi có thể lên đến 5-6 giây.

\item \textbf{Chưa hỗ trợ đa phương tiện}:
Hệ thống chỉ xử lý văn bản, chưa hỗ trợ hình ảnh, video hay các loại dữ liệu đa phương tiện khác.

\item \textbf{Chưa tích hợp với các công cụ phổ biến}:
Chưa có tích hợp đầy đủ với Jira, GitHub, Azure DevOps.

\item \textbf{Chrome Extension ở mức MVP}:
Extension mới ở mức cơ bản, chưa hỗ trợ capture console log, network log.
\end{enumerate}

\subsection{Hạn chế về phương pháp nghiên cứu}

\begin{enumerate}
\item \textbf{Chưa có đánh giá định lượng về ROI}:
Chưa có nghiên cứu chi tiết về lợi ích kinh tế của hệ thống.

\item \textbf{Thời gian đánh giá ngắn}:
Thời gian đánh giá trải nghiệm người dùng chỉ 2 tuần, chưa đủ để đánh giá tác động dài hạn.

\item \textbf{Số lượng người dùng thử nghiệm hạn chế}:
Chỉ có 20 người dùng tham gia đánh giá, chưa đại diện cho toàn bộ các nhóm người dùng.
\end{enumerate}

%%%%%%%%%%%%%%%%%%%%%%%%%%%%%%%%%%%%%%%%%%%%%%%%%%%%%%%%%%%%%%
\section{Hướng phát triển trong tương lai}
%%%%%%%%%%%%%%%%%%%%%%%%%%%%%%%%%%%%%%%%%%%%%%%%%%%%%%%%%%%%%%

\subsection{Hướng phát triển ngắn hạn (6-12 tháng)}

\begin{enumerate}
\item \textbf{Tối ưu hiệu năng AI}:
\begin{itemize}
\item Sử dụng caching cho các yêu cầu AI tương tự.
\item Implement async processing với message queue (RabbitMQ).
\item Sử dụng mô hình LLM nhỏ hơn nhưng tinh chỉnh (fine-tuned) trên dữ liệu quản lý dự án.
\end{itemize}

\item \textbf{Cải thiện độ chính xác AI}:
\begin{itemize}
\item Thu thập thêm dữ liệu huấn luyện cho Story Point lớn (>13).
\item Cải thiện Prompt Engineering với Few-shot Learning.
\item Thêm cơ chế Self-Consistency cho các dự đoán quan trọng.
\end{itemize}

\item \textbf{Hỗ trợ đa ngôn ngữ}:
\begin{itemize}
\item Xây dựng bộ dữ liệu huấn luyện tiếng Việt.
\item Hỗ trợ tiếng Nhật cho các dự án Việt - Nhật.
\item Sử dụng mô hình đa ngôn ngữ (XLM-R, mGPT).
\end{itemize}

\item \textbf{Nâng cấp Chrome Extension}:
\begin{itemize}
\item Tự động capture console log và network log.
\item Hỗ trợ full-page screenshot.
\item Tích hợp với Jira API.
\end{itemize}

\item \textbf{Cải thiện Availability}:
\begin{itemize}
\item Thêm load balancer.
\item Cấu hình multiple replicas cho các service.
\item Implement circuit breaker pattern.
\end{itemize}
\end{enumerate}

\subsection{Hướng phát triển trung hạn (1-2 năm)}

\begin{enumerate}
\item \textbf{Tích hợp Multi-Agent Collaboration}:
\begin{itemize}
\item Cho phép các Agent trao đổi thông tin trực tiếp.
\item Xây dựng cơ chế phối hợp (Coordination) giữa các Agent.
\item Ví dụ: Risk Agent phối hợp với Forecast Agent để đưa ra dự báo chính xác hơn.
\end{itemize}

\item \textbf{Tích hợp thêm AI Agents}:
\begin{itemize}
\item \textbf{AI Code Review Agent}: Tự động phân tích và gợi ý cải thiện code quality.
\item \textbf{AI Test Generation Agent}: Sinh test cases tự động từ User Story.
\item \textbf{AI Bug Detection Agent}: Phát hiện bug tiềm ẩn từ source code.
\item \textbf{AI DevOps Agent}: Hỗ trợ CI/CD và deployment.
\end{itemize}

\item \textbf{Explainable AI (XAI)}:
\begin{itemize}
\item Cung cấp giải thích chi tiết cho các đề xuất AI.
\item Visualize decision-making process của từng Agent.
\item Hiển thị các yếu tố ảnh hưởng đến quyết định.
\end{itemize}

\item \textbf{Tích hợp với các công cụ phổ biến}:
\begin{itemize}
\item Jira Integration: Đồng bộ dữ liệu hai chiều.
\item GitHub/GitLab Integration: Kết nối với source code.
\item Slack/MS Teams Integration: Gửi thông báo và cảnh báo.
\item Azure DevOps Integration: Đồng bộ dự án.
\end{itemize}
\end{enumerate}

\subsection{Hướng phát triển dài hạn (2-5 năm)}

\begin{enumerate}
\item \textbf{Fine-tuning LLM chuyên biệt}:
\begin{itemize}
\item Tinh chỉnh mô hình ngôn ngữ trên dữ liệu quản lý dự án.
\item Xây dựng mô hình với kiến thức chuyên sâu về Software Engineering.
\item Giảm phụ thuộc vào API bên ngoài.
\end{itemize}

\item \textbf{Đa phương tiện (Multimodal)}:
\begin{itemize}
\item Xử lý hình ảnh (UI bug, screenshot analysis).
\item Xử lý video (recording của user session).
\item Xử lý âm thanh (meeting minutes tự động).
\end{itemize}

\item \textbf{Phân tích dự đoán (Predictive Analytics)}:
\begin{itemize}
\item Dự đoán rủi ro dài hạn của dự án.
\item Đề xuất kế hoạch tối ưu dựa trên lịch sử.
\item Phân tích trend và pattern của dự án.
\end{itemize}

\item \textbf{Triển khai trên Kubernetes}:
\begin{itemize}
\item Auto-scaling dựa trên tải.
\item Zero-downtime deployment.
\item Multi-region deployment cho high availability.
\end{itemize}

\item \textbf{Autonomous Agent}:
\begin{itemize}
\item Cho phép AI Agent tự động thực hiện các tác vụ đã được ủy quyền.
\item Xây dựng cơ chế trust và verification.
\item Học từ phản hồi của người dùng để cải thiện.
\end{itemize}
\end{enumerate}

%%%%%%%%%%%%%%%%%%%%%%%%%%%%%%%%%%%%%%%%%%%%%%%%%%%%%%%%%%%%%%
\section{Đề xuất cho nghiên cứu tiếp theo}
%%%%%%%%%%%%%%%%%%%%%%%%%%%%%%%%%%%%%%%%%%%%%%%%%%%%%%%%%%%%%%

\subsection{Đề xuất cho nghiên cứu học thuật}

\begin{enumerate}
\item \textbf{Nghiên cứu về Prompt Engineering tối ưu}: 
Khám phá các chiến lược Prompt khác nhau để tối ưu hóa chất lượng đầu ra cho từng loại tác vụ.

\item \textbf{Nghiên cứu về Multi-Agent Collaboration}:
Phát triển các thuật toán phối hợp giữa các Agent để giải quyết các bài toán phức tạp hơn.

\item \textbf{Nghiên cứu về Explainable AI trong quản lý dự án}:
Xây dựng các kỹ thuật giải thích để tăng độ tin cậy của AI trong quản lý dự án.

\item \textbf{Nghiên cứu về Human-AI Collaboration}:
Tìm hiểu cách tối ưu hóa sự tương tác giữa con người và AI trong quy trình quản lý dự án.
\end{enumerate}

\subsection{Đề xuất cho ứng dụng thực tế}

\begin{enumerate}
\item \textbf{Thí điểm trong doanh nghiệp}: 
Triển khai thử nghiệm trong một số doanh nghiệp phần mềm để đánh giá hiệu quả thực tế.

\item \textbf{Đánh giá ROI}: 
Nghiên cứu chi tiết về lợi ích kinh tế và thời gian của hệ thống.

\item \textbf{Phát triển phiên bản SaaS}: 
Xây dựng phiên bản SaaS để nhiều doanh nghiệp có thể sử dụng.

\item \textbf{Đào tạo và tài liệu hóa}: 
Xây dựng tài liệu hướng dẫn chi tiết và các khóa đào tạo cho người dùng.
\end{enumerate}

%%%%%%%%%%%%%%%%%%%%%%%%%%%%%%%%%%%%%%%%%%%%%%%%%%%%%%%%%%%%%%
\section{Triển vọng ứng dụng}
%%%%%%%%%%%%%%%%%%%%%%%%%%%%%%%%%%%%%%%%%%%%%%%%%%%%%%%%%%%%%%

\subsection{Ứng dụng trong doanh nghiệp phần mềm}

\begin{table}[H]
\centering
\caption{Khả năng ứng dụng của hệ thống}
\begin{tabular}{|p{4.5cm}|p{5cm}|p{5cm}|}
\hline
\textbf{Lĩnh vực} & \textbf{Ứng dụng} & \textbf{Lợi ích dự kiến} \\
\hline
Phát triển phần mềm & Quản lý Sprint, phân công Task, ước lượng & Tiết kiệm 50-70\% thời gian \\
\hline
QA/Testing & Phân tích bug report, sinh test case & Tăng quality, giảm manual work \\
\hline
Dự án đa team & Quản lý nguồn lực, phát hiện rủi ro & Tăng visibility, giảm risk \\
\hline
Dự án mã nguồn mở & Quản lý issue, phân công contributor & Tăng productivity \\
\hline
Freelance/Remote & Quản lý công việc từ xa & Tăng collaboration \\
\hline
\end{tabular}
\label{tab:application_areas}
\end{table}

\subsection{Tiềm năng mở rộng}

\begin{enumerate}
\item \textbf{Tích hợp với công cụ AI khác}: Kết hợp với GitHub Copilot, Tabnine để hỗ trợ lập trình.
\item \textbf{Tích hợp với CI/CD}: Tự động hóa quy trình từ issue đến deployment.
\item \textbf{Hỗ trợ dự án phi công nghệ}: Mở rộng sang các lĩnh vực khác (xây dựng, marketing, v.v.)
\item \textbf{Personal Assistant}: Xây dựng trợ lý cá nhân cho từng Developer.
\end{enumerate}

%%%%%%%%%%%%%%%%%%%%%%%%%%%%%%%%%%%%%%%%%%%%%%%%%%%%%%%%%%%%%%
\section{Tổng kết và triển vọng}
%%%%%%%%%%%%%%%%%%%%%%%%%%%%%%%%%%%%%%%%%%%%%%%%%%%%%%%%%%%%%%

\subsection{Tổng kết luận văn}

Luận văn đã trình bày một cách toàn diện quá trình nghiên cứu, thiết kế, hiện thực và đánh giá một nền tảng quản lý dự án phần mềm thông minh tích hợp trí tuệ nhân tạo. Hệ thống, được đặt tên là AI Project Management System, là một nền tảng thế hệ mới, trong đó AI đóng vai trò là trợ lý thông minh hỗ trợ Product Owner, Scrum Master, Project Manager và Developer trong toàn bộ vòng đời phát triển phần mềm.

Các kết quả thực nghiệm cho thấy hệ thống đáp ứng đầy đủ các yêu cầu chức năng và phi chức năng, với độ chính xác của các AI Agent đạt từ 81\% đến 94\%, thời gian phản hồi nhanh, và điểm hài lòng người dùng cao (4.6/5). Hệ thống giúp giảm thời gian xử lý các tác vụ thủ công từ 60-70\%, đặc biệt trong phân tích yêu cầu và lập kế hoạch Sprint.

Về mặt kiến trúc, hệ thống áp dụng mô hình Layered Architecture kết hợp với Domain-Driven Design (DDD), Clean Architecture và Microservices, giúp nâng cao khả năng mở rộng và bảo trì. AI Service được xây dựng như một tầng xử lý độc lập với AI Orchestrator điều phối 8 AI Agent chuyên biệt.

\subsection{Triển vọng trong tương lai}

Sự phát triển nhanh chóng của các mô hình ngôn ngữ lớn mở ra nhiều cơ hội mới trong lĩnh vực quản lý dự án phần mềm. Thay vì chỉ hỗ trợ tự động hóa các tác vụ đơn giản, AI đang từng bước trở thành một cộng sự thông minh có khả năng phân tích, lập kế hoạch và hỗ trợ ra quyết định.

Luận văn này đã chứng minh rằng việc kết hợp các AI Agent, Prompt Engineering và Retrieval-Augmented Generation vào một nền tảng quản lý dự án có thể nâng cao hiệu quả làm việc của nhóm phát triển, giảm thời gian xử lý các công việc lặp lại và chuẩn hóa quy trình quản lý dự án.

Mặc dù vẫn còn nhiều thách thức cần giải quyết, đặc biệt về khả năng suy luận, độ tin cậy và hiệu năng của các mô hình AI, hướng tiếp cận được đề xuất trong luận văn là một nền tảng tiềm năng cho các hệ thống quản lý dự án thông minh trong tương lai.

\subsection{Lời kết}

Quá trình thực hiện luận văn đã mang lại cho tác giả những bài học quý giá về cả lý thuyết và thực tiễn. Từ việc nghiên cứu các mô hình AI hiện đại đến việc áp dụng chúng vào một bài toán thực tế, từ việc thiết kế kiến trúc phần mềm đến triển khai và đánh giá hệ thống — tất cả đều góp phần nâng cao kiến thức và kỹ năng chuyên môn.

Hy vọng rằng những đóng góp của luận văn sẽ là nền tảng cho các nghiên cứu tiếp theo, và hệ thống AI Project Management sẽ tiếp tục được phát triển để trở thành một công cụ hữu ích cho cộng đồng phát triển phần mềm.

%%%%%%%%%%%%%%%%%%%%%%%%%%%%%%%%%%%%%%%%%%%%%%%%%%%%%%%%%%%%%%
% TÀI LIỆU THAM KHẢO
%%%%%%%%%%%%%%%%%%%%%%%%%%%%%%%%%%%%%%%%%%%%%%%%%%%%%%%%%%%%%%
\begin{thebibliography}{99}

\bibitem{pmbok7}
Project Management Institute, \emph{A Guide to the Project Management Body of Knowledge (PMBOK Guide)}, 7th ed., Project Management Institute, 2021.

\bibitem{agilemanifesto}
Agile Alliance, \emph{Manifesto for Agile Software Development}, 2001.

\bibitem{scrumguide2020}
K. Schwaber and J. Sutherland, \emph{The Scrum Guide}, 2020.

\bibitem{bachari2025}
A. Bachari et al., ``AI in Project Management: A Systematic Literature Review,'' \emph{Journal of Software Engineering}, 2025.

\bibitem{russell2021}
S. Russell and P. Norvig, \emph{Artificial Intelligence: A Modern Approach}, 4th ed., Pearson, 2021.

\bibitem{vaswani2017attention}
A. Vaswani et al., ``Attention Is All You Need,'' in \emph{Advances in Neural Information Processing Systems}, 2017.

\bibitem{xuan2014}
J. Xuan, H. Jiang, Y. Hu, Z. Ren, W. Zou, Z. Luo, and X. Wu, ``Towards effective bug triage with software data reduction techniques,'' \emph{IEEE Transactions on Knowledge and Data Engineering (TKDE)}, vol. 27, no. 1, pp. 264-280, 2014.

\bibitem{antoniol2008}
G. Antoniol, K. Ayari, M. Di Penta, F. Khomh, and Y. G. Guéhéneuc, ``Is it a bug or an enhancement? A text-based approach to classify change requests,'' in \emph{Proc. CASCON}, pp. 304-318, ACM, 2008.

\bibitem{anvik2006}
J. Anvik, L. Hingo, and G. C. Murphy, ``Who should fix this bug?,'' in \emph{Proc. ICSE}, pp. 361-370, ACM, 2006.

\bibitem{pedregosa2011}
F. Pedregosa et al., ``Scikit-learn: Machine Learning in Python,'' \emph{Journal of Machine Learning Research}, vol. 12, pp. 2825-2830, 2011.

\bibitem{walls2022}
C. Walls, \emph{Spring in Action}, 6th ed., Manning Publications, 2022.

\bibitem{react}
React Documentation, \emph{React - A JavaScript library for building user interfaces}, Meta Platforms Inc., https://react.dev, 2024.

\bibitem{fastapi}
FastAPI Documentation, \emph{FastAPI framework, high performance, easy to learn, fast to code, ready for production}, https://fastapi.tiangolo.com, 2024.

\bibitem{docker}
Docker Documentation, \emph{Docker: Empowering App Development for Developers}, https://docs.docker.com, 2024.

\bibitem{spring}
Spring Documentation, \emph{Spring Boot Reference Guide}, https://spring.io/projects/spring-boot, 2024.

\bibitem{postgresql}
PostgreSQL Documentation, \emph{PostgreSQL 16: The World's Most Advanced Open Source Database}, https://www.postgresql.org, 2024.

\end{thebibliography}

%%%%%%%%%%%%%%%%%%%%%%%%%%%%%%%%%%%%%%%%%%%%%%%%%%%%%%%%%%%%%%
% PHỤ LỤC (Nếu có)
%%%%%%%%%%%%%%%%%%%%%%%%%%%%%%%%%%%%%%%%%%%%%%%%%%%%%%%%%%%%%%
\appendix
\chapter{PHỤ LỤC A: MÃ NGUỒN CHÍNH}

\section{Backend - ProjectService}

\begin{verbatim}
@Service
@Transactional
public class ProjectService {
    // ... (mã nguồn đã trình bày trong Chương 10)
}
\end{verbatim}

\section{AI Service - Requirement Agent}

\begin{verbatim}
@app.post("/api/analyze-requirement")
async def analyze_requirement(request: AnalyzeRequest):
    # ... (mã nguồn đã trình bày trong Chương 8)
\end{verbatim}

\chapter{PHỤ LỤC B: DỮ LIỆU THỰC NGHIỆM}

\section{Phân bố dữ liệu Story Point}

\section{Kết quả khảo sát người dùng chi tiết}

\chapter{PHỤ LỤC C: HƯỚNG DẪN CÀI ĐẶT}

\section{Cài đặt Docker Compose}

\section{Cài đặt thủ công}

%%%%%%%%%%%%%%%%%%%%%%%%%%%%%%%%%%%%%%%%%%%%%%%%%%%%%%%%%%%%%%
% KẾT THÚC
%%%%%%%%%%%%%%%%%%%%%%%%%%%%%%%%%%%%%%%%%%%%%%%%%%%%%%%%%%%%%%
\end{document}